\documentclass[11pt]{article}
\usepackage[T1]{fontenc}
\usepackage[utf8]{inputenc}
\usepackage{lmodern,amsmath,amssymb,amsthm,mathtools,microtype}
\usepackage[margin=1in]{geometry}
\usepackage[round,authoryear]{natbib}
\usepackage[hidelinks]{hyperref}
\usepackage{enumitem}
\setlist{itemsep=3pt,topsep=5pt}
\newcommand{\workbenchurl}{https://github.com/KokunoYumeto/schemes-and-moduli-workbench}
\newcommand{\recordurl}{https://github.com/KokunoYumeto/mathematics-workbenches}
\newcommand{\B}{\mathbb B}
\newcommand{\N}{\mathbb N}
\newcommand{\Z}{\mathbb Z}
\newcommand{\F}{\mathbb F}
\newcommand{\supp}{\operatorname{supp}}
\newcommand{\Sup}{\operatorname{Sup}}
\newcommand{\vc}{\operatorname{vc}}
\newcommand{\cmon}{\operatorname{C}_{\mathrm{mon}}}
\newcommand{\cbool}{\operatorname{C}_{\mathrm{bool}}}
\newcommand{\car}{\operatorname{C}_{\mathrm{arith}}}
\newcommand{\cone}{\operatorname{C}_{1}}
\newcommand{\cread}{\operatorname{C}_{\mathrm{read}\,1}}
\newcommand{\Fin}{\mathcal P_{\mathrm{fin}}}
\newtheorem{theorem}{Theorem}[section]
\newtheorem{proposition}[theorem]{Proposition}
\newtheorem{lemma}[theorem]{Lemma}
\newtheorem{corollary}[theorem]{Corollary}
\theoremstyle{definition}
\newtheorem{definition}[theorem]{Definition}
\newtheorem{example}[theorem]{Example}
\theoremstyle{remark}
\newtheorem{remark}[theorem]{Remark}
\title{The P = NP? edge-circuit conjecture at degree one,\protect\\and zero-adjoined rings after Faul, Goswami, Joubert and Manuell}
\author{}
\date{24 November 2026}
\begin{document}
\maketitle
\begin{abstract}
Faul, Goswami, Joubert and Manuell adjoined a new zero to a ring and identified the free commutative inverse semiring; Caviglia, Faul, Manuell and Mesiti describe inverse semirings by groups over an idempotent semiring. We use these results to keep the polynomial and Boolean function of one circuit together. For a graph without a four-cycle, the minimum nonnegative arithmetic complexity among quadratic polynomials with its edge support is exactly its number of edges plus its vertex-cover number minus one. By the classical lower-envelope argument, the same count holds when each edge has a squarefree witness in a Boolean circuit's formal polynomial. We extend the full lower bound to circuits missing at most three such witnesses, to four missing witnesses with maximal vertex-cover gain, and to arbitrarily many missing witnesses when extraction of a nonzero quadratic part deletes at most one addition. If extraction deletes at most two additions, every nontrivial component of the missing-edge graph is a single edge or a two-edge path, with no restriction on multiplication count. With exactly four missing edges, at least four deleted multiplications are required in that two-addition regime. The proof combines the equality case of the quadratic count, exchanges of star centers, and a normalization that retains the whole quadratic circuit while absorbing redundant higher terms. For a bipartite graph without a four-cycle, the exact count also holds for circuits of degree at most one in one side of the bipartition, by transposition and a disjunction-count lemma permitting arbitrary conjunctions. For four missing witnesses, the only remaining numerical saving possibility has an optimal quadratic remainder, three deleted products, three deleted additions and cover gain three. A circuit with at most four deleted products and two deleted additions has vertex-cover number at most two more than its original retained product count. The five-deletion, cover-gain-two case is excluded by bounds that retain the full insertion contexts and count the original quadratic products. The unrestricted Boolean problem remains open here. Finally we separate cancellation from repeated-variable identification, and specify the output condition under which zero-adjoined circuits have monotone complexity. The mathematics and text were produced and checked by GPT-6 Astra Pro; no human has checked them.
\end{abstract}

\section{The statements and their scope}
The purpose of this paper is to make precise which monotone-circuit questions can be transported through a zero-adjoined ring without changing their gate count.
For a commutative unital ring $R$, let
\[
 R_0=R\sqcup\{\tau\}.
\]
The old ring operations are unchanged; $\tau$ is an additive identity and a multiplicative absorber. We write $e=0_R$ for the old zero, so $e\ne\tau$. This is exactly \citet[Definition~3.6]{FGJM}. Their Proposition~5.2 identifies $\Z_0[x_1,\ldots,x_n]$ as the free commutative inverse semiring, and the paragraph following that proposition identifies its polynomial and finite-support projections. Their Proposition~6.1 proves that semimodules have inverse additive monoids. Our arguments concern algebras and circuits, and do not require a further semimodule classification.

Every circuit used in an exact count below is a finite acyclic list of binary gates. Inputs are variables and cost no gates; every addition, multiplication, conjunction or disjunction costs one, and arbitrary reuse of a wire is allowed. Constants are excluded from the exact arithmetic model. Write $\car(P)$ for the minimum size of a variable-input $+,\times$ circuit computing the formal polynomial $P$ over $\N$, and $\cmon(f)$ for the minimum size of an AND/OR circuit computing $f$ on all Boolean inputs. Later, $\cbool(f)$ permits NOT gates as well. The Boolean function of a simple graph $G=(V,E)$ is
\[
 f_G(x)=\bigvee_{\{u,v\}\in E}(x_u\wedge x_v),
 \qquad H_G(x)=\sum_{\{u,v\}\in E}x_ux_v.
\]
We assume $E\ne\varnothing$. Write $\vc(G)$ for the least size of a vertex cover. Isolated vertices can be deleted.

For a polynomial $P\in\N[x_v:v\in V]$, define $P\in\Sup(G)$ to mean that every monomial of $P$ contains the endpoints of an edge in its variable support, and every edge is the exact variable support of at least one monomial of $P$. Coefficients of occurring monomials are positive integers, but their exponents need not be one.

Following \citet[Definition~1, p.~6 of the author manuscript]{JuknaRead}, a Boolean circuit for $f_G$ is \emph{read-1} if its formal nonnegative polynomial contains $x_ux_v$ for every edge $uv$. Other monomials may have larger degree and repeated variables. This condition concerns witnesses for minimal true assignments; it does not restrict the number of circuit occurrences of a variable. Denote the minimum size of such a circuit by $\cread(f_G)$.

\begin{theorem}[The exact polynomial and the support class]\label{thm:main}
For every finite simple graph with at least one edge,
\[
 \cmon(f_G)=\min_{P\in\Sup(G)}\car(P).
\]
If $G$ has no four-cycle, then
\[
 \min_{c_e\in\N_{>0}}\car\!\left(\sum_{e=\{u,v\}\in E}c_ex_ux_v\right)
 =\car(H_G)=\cread(f_G)=|E|+\vc(G)-1.
\]
All complexities count binary gates in the respective variable-input models.
\end{theorem}
Thus the Boolean optimization problem is an optimization over all polynomial lifts with the prescribed minimal supports. The quadratic and read-1 cases allow shared wires and repeated coefficients. The next result allows two edges to be witnessed only by higher powers.

\begin{theorem}[At most two missing quadratic witnesses]\label{thm:twomissing}
Let $G$ be a nonempty finite simple graph without a four-cycle. Let a variable-input $+,\times$ circuit $C$ compute $P\in\Sup(G)$, and let $E_2$ be the set of edges whose monomial $x_ux_v$ occurs in $P$. If $|E\setminus E_2|\le2$, then
\[
 |C|\ge |E|+\vc(G)-1.
\]
Consequently the minimum over this class is exactly the displayed count, attained by the star circuit.
\end{theorem}
This first extension excludes savings with one or two edges witnessed solely through repeated variables. Proposition~\ref{prop:missing} gives quantitative bounds for any number of missing witnesses. The proof of Theorem~\ref{thm:twomissing} also shows why counting discarded gates alone is insufficient: an optimal quadratic remainder has additional rigidity. The next theorem extends the same full bound to three missing witnesses.

\begin{theorem}[At most three missing quadratic witnesses]\label{thm:threemissing}
Let $G$ be a finite simple graph with a nonempty edge set and no four-cycle. Let a variable-input circuit $C$ compute $P\in\Sup(G)$, and let $E_2$ be its quadratic edge support. If $|E\setminus E_2|\le3$, then
\[
 |C|\ge |E|+\vc(G)-1.
\]
The star circuit attains this bound, so it is the exact minimum over the stated class. There is no restriction on the degrees or multiplicities of the higher witnesses.
\end{theorem}
Thus a circuit improving the conjectured count must lack quadratic witnesses for at least four edges. For three missing edges, the preliminary count reduces a putative saving to five deleted gates and an optimal quadratic remainder. We exclude every distribution of those five gates; the maximal cover-number gain in that reduction is essential to one of the cases. Example~\ref{ex:missingmatching} shows why an induced missing matching alone would not suffice. The exact count also holds for every four-cycle-free graph with at most four edges, without a witness restriction. In any graph, an induced matching of four missing quadratic edges requires at least six discarded gates. Theorem~\ref{thm:maxfour} below also proves the star lower bound at maximal cover gain four, with arbitrary extracted multiplication and addition counts.

The next restriction allows arbitrarily many missing witnesses. To extract the smallest-degree part, replace each unequal-degree sum by its smaller-degree operand and prune gates no longer leading to the output. When the quadratic part is nonzero, this gives its quadratic circuit.

\begin{theorem}[At most one deleted addition]\label{thm:oneaddition}
Let $G$ be a nonempty finite simple graph without a four-cycle, and let a positive variable-input circuit $C$ compute $P\in\Sup(G)$ with nonzero quadratic part. If the canonical extraction just described deletes at most one original addition gate, then
\[
 |C|\ge |E(G)|+\vc(G)-1.
\]
The number of missing quadratic witnesses, the retained multiplication and addition counts, and the cover-number gain are unrestricted.
\end{theorem}
Thus a saving with nonempty quadratic support requires at least two deleted additions. Proposition~\ref{prop:oneadditionall} proves the more precise single-cut bound: $r$ missing edges must form an induced matching and require at least $2^r-2$ deleted products. For a graph with neither triangles nor four-cycles, Proposition~\ref{prop:commonoldcenter} further forces $r\le2$ in that case.

\begin{theorem}[Missing components with at most two deleted additions]\label{thm:twosumscomponents}
Let $G$ be a finite simple graph without a four-cycle, with triangles permitted. Let a positive variable-input circuit compute $P\in\Sup(G)$ with nonempty quadratic support $H$. If canonical extraction deletes at most two original additions, every nontrivial connected component of $G\setminus H$ is $K_2$ or $P_3$. Deleted products, retained counts, fanout and cover gain are unrestricted.
\end{theorem}
This condition does not require an induced missing target: even a three-edge missing path with an old diagonal is excluded. It places no bound on the number of allowed components and does not assert that every such union is realizable. The proof retains the actual low inputs of the two cuts; the missing-degree bound and the induced-path argument alone would not exclude the old diagonal.

\begin{theorem}[Four missing witnesses with maximal cover gain]\label{thm:maxfour}
Let $G$ be a nonempty finite simple graph without a four-cycle. Let a positive variable-input circuit $C$ compute $P\in\Sup(G)$ with nonempty quadratic support $H$. If
\[
 |E(G)\setminus E(H)|=4,\qquad \vc(G)-\vc(H)=4,
\]
then $|C|\ge |E(G)|+\vc(G)-1$. There is no restriction on the extracted multiplication or addition counts.
\end{theorem}
The proof combines the exact quadratic gate counts with an algebraic obstruction for three recursively dependent products. Corollary~\ref{cor:fourdeletedproducts} shows that four induced missing matching edges require four original deleted products, regardless of the number of deleted additions. At cover gain three, Proposition~\ref{prop:gainthreesix} requires six deleted gates whenever the retained multiplication count is minimal, without restricting retained additions. Proposition~\ref{prop:gainthreefiveprivate} also excludes the five-deletion saving with one extra retained product; its proof classifies every hidden-addition placement in that budget. Theorem~\ref{thm:twosumscomponents} requires three deleted additions whenever the missing graph has a component other than a single edge or a two-edge path, without restricting multiplication counts. Corollary~\ref{cor:hiddenunequalany} excludes the two-product, three-addition hidden unequal-degree branch even at cover gain two. Proposition~\ref{prop:threeproductstwosums} excludes all three-product, two-addition configurations with four missing witnesses, without retained-count or cover-gain assumptions. Propositions~\ref{prop:onevisibleexceptional} and~\ref{prop:twovisibleexceptional}, together with the hidden-input and all-visible bounds, exclude the three-product, three-addition profile on the two exceptional outside graphs at minimum retained multiplication count and cover gain at least three. Proposition~\ref{prop:dependentlinearglobal}, together with the corresponding nonancestral and visibility bounds, reduces the five-deletion, cover-gain-two case to one visible cut and two hidden additions; Proposition~\ref{prop:onevisibletwocover} excludes that final configuration. Proposition~\ref{prop:sevenfactorprofiles} also excludes the one-visible and first-high-all-four-products configurations in the four-product, two-addition case at cover gain at least three. Proposition~\ref{prop:unaffectedzerocross} also gives a cover bound when both high networks are multiplicative on old forms and the mixed coefficient is zero. Proposition~\ref{prop:twowitnesscover} extends that cover argument to a precise cubic-high configuration with a possibly positive mixed coefficient. Proposition~\ref{prop:onevisibleindependentsets} bounds every $H$-independent induced restriction by seven maximal independent sets in the one-visible profile, excluding that profile on both remaining outside shapes, with arbitrary retained counts and hidden-addition degrees. Proposition~\ref{prop:fourproductstwosumscover} bounds the vertex-cover number by $K+2$ whenever at most four products and two additions are deleted, where $K$ counts the original retained products. Corollary~\ref{cor:fourresidualfinal} identifies the remaining six-deletion case. Its six deleted gates must now be three products and three additions.

For a bipartite graph with sides $X$ and $Y$, mark the $X$-variables. The \emph{syntactic $X$-degree} of an input is one for a marked variable and zero for an unmarked variable. A disjunction takes the maximum of its operands' degrees and a conjunction takes their sum. Let $\cone(f_G)$ be the least size of a circuit whose output has $X$-degree at most one.

\begin{theorem}[One-sided degree one]\label{thm:onesided}
Let $G=(X\sqcup Y,E)$ be a finite bipartite graph without a four-cycle and with $E\ne\varnothing$. Then
\[
 \cone(f_G)=|E|+\vc(G)-1.
\]
In fact every such circuit has at least $|E|-1$ disjunction gates and at least $\vc(G)$ conjunction gates.
\end{theorem}
This is the mathematical formulation attributed to node L-meta-647 of the public project \emph{P = NP?} in \citet{Fleet}. Its proof below expands the degree-one argument in our unpublished project notes \citep{ProjectNotes} and proves the required addition-count lemma directly. The theorem settles that mathematical restriction of the edge-count question. No statement about acceptance by the public project's Lean checker is made. The unrestricted conjectures L-meta-608 and L-meta-612, and the degree-two bound discussed as L-meta-649, concern larger classes. Theorems~\ref{thm:main}--\ref{thm:onesided} do not settle them or the Boolean P versus NP problem.

\subsection{Credit and relation to earlier work}
\citet[Section~5]{FGJM} already distinguish a formal inverse-semiring polynomial from its ring value and its finite set of exponents. We exhibit the group-valued data for that construction and its monoid version. The general correspondence used to describe those data is \citet[Theorem~8.5]{CFMM}; their Remark~8.7 identifies the two-idempotent construction. The additional quotient used here records a monotone Boolean function together with its polynomial.

Monomial support is a classical source of monotone arithmetic lower bounds. \citet{Schnorr} counts additions in nonnegative computations. \citet[Sections~2--4]{JS} treat formal polynomial computation over semirings and explain explicitly why two formal Boolean polynomials can define the same function. \citet[Sections~3--4 and Theorem~9]{JuknaTropical} compares arithmetic, tropical and Boolean equivalence and isolates the role of homogeneity. The support criterion defining $\Sup(G)$ is the graph case of \citet[Lemma~1, p.~5]{JuknaRead}, stated there as folklore. His Lemma~2 and Theorem~1, p.~7, give the lower-envelope reduction and the equivalence between read-1 Boolean complexity and arithmetic complexity when positive coefficients are allowed to vary. We spell out the coefficient-independent edge count because its exact constant and shared wires matter here; we make no priority claim for it.

The distinction from unrestricted Boolean circuits is substantive. \citet[Theorem~2.1, p.~4 of the author manuscript]{JuknaSingle} disproves the general single-level conjecture for quadratic Boolean functions: allowing several levels of conjunction can give large savings. That result does not itself refute the four-cycle-free conjecture. It does rule out using a general, cost-preserving flattening as an unstated step.

A recent construction of Krapivin, Przybocki, and Subercaseaux (2026, Section~3.4, Theorem~9; full version Theorem~6 and Appendix~A.2) gives a monotone circuit for $f_{K_n}=T_2$ with $2n+\sqrt{2n}+O(n^{1/3})$ gates. Its proof uses a threshold-three subcircuit. Complete graphs contain four-cycles for $n\ge4$, so this construction lies outside the graph hypothesis used here.

For four-cycle-free graphs, \citet[Section~3, p.~11 of the author manuscript]{JuknaSingle} attributes the single-level lower bound $\cmon^{\mathrm{single}}(f_G)\ge |E|$ to \citet{AM}; this counts all binary gates but restricts the conjunction levels. The same manuscript distinguishes their total-gate separation for several designated outputs from its own single-output result. A different nearby result is \citet[Theorem~7.3, p.~18]{JuknaGraphs}: for a graph without triangles or four-cycles, the edge function has monotone formula length at least $|E|/2$ and monotone depth at least $\log_2|E|-1$. Formula length counts leaves of a tree and does not give the corresponding size bound for a circuit with shared wires. These prior results therefore do not supply the missing unrestricted step here.

Cancellation can independently yield exponential arithmetic savings. \citet[Theorem~1, p.~307]{Valiant} proves an exponential monotone lower bound for perfect-matching polynomials of triangular grids, while the Pfaffian construction computes those planar matching polynomials efficiently with subtraction. His Lemma~1, p.~305, also gives the homogeneity principle used in our elementary quadratic argument. These are formal-polynomial results. Our Boolean quotient instead identifies repeated variables and absorbs larger supports, so cancellation and Boolean identification must be distinguished.

A recent preprint by \citet[Theorem~1.1 and Section~1.3]{Anthropic} states lower bounds for the permanent and Hamiltonian-cycle polynomials over fields, with an accompanying Lean formalization. Its model requires a polynomial identity, whereas the optimization in Theorem~\ref{thm:main} permits every lift in $\Sup(G)$. That distinction prevents a direct inference of the Boolean edge bound from such algebraic results. None of the proofs below depends on that preprint.

\citet{OpenAICubic} proves a cubic lower bound for exact and border determinantal representations of the complex permanent; the accompanying formalization's scope states the explicit bound for matrix size at least $1408$. A determinantal representation is a restricted algebraic model, and coefficientwise limits still concern a formal polynomial. This result therefore provides neighboring algebraic context, without giving the gate bound for arbitrary Boolean lifts considered here.

\section{Proof overview}
The algebraic part separates two operations that ordinary ring notation combines: an exponent can be present, and its coefficient can be zero. Finite subsets of a monoid form the idempotent semiring; the coefficients over each subset form an abelian group. Passing from exponents to their variable supports, and then taking upward closure, gives the Boolean-function quotient.

The quadratic arithmetic count uses positivity to force each useful gate to have a single degree. Every multiplication forms a rectangle of edges, and the absence of a four-cycle makes that rectangle a star. Linear subcomputations can be shared among stars with the same center, but cannot be shared between distinct centers. Counting components in each center's linear network, and reads in the quadratic network, gives the exact lower bound even when coefficients repeat. Taking the smallest-degree part of a positive circuit costs no extra gates. For read-1 circuits that part contains all edges, so the same bound applies. One missing witness costs at least two discarded gates, and two cost at least three. The same structural analysis gives four deletions for at least three missing witnesses and five for at least four, with sharp examples at each threshold. In the only remaining case for two witnesses, a hypothetical saving would require an optimal quadratic remainder and exactly three discarded gates. The missing endpoints form an induced matching outside every minimum old cover. Classifying those three gates then forces a star whose center can be exchanged for two of the missing endpoints, contradicting the assumed cover number.

For three missing witnesses, a four-gate classification combines two kinds of arguments. Products of old-neighborhood linear forms restrict the possible missing supports, while an optimal quadratic remainder permits exchanges of assigned star centers. Exactly four deletions yield a missing-edge graph covered by two vertices; with an optimal remainder the cover gain is at most one. Integer gate counts leave only a five-deletion circuit with an optimal quadratic remainder and maximal cover-number gain. We first remove off-frame uses of full quadratic carriers, preserving the retained circuit rather than just its polynomial support. With at most two private products this yields a template that cannot compute a three-edge matching. The one-cut profiles reduce to fewer than eight linear clauses. The remaining three-product, two-cut profile is excluded by its actual insertion contexts: a small context would allow an assigned center to be exchanged for missing endpoints. These cases complete Theorem~\ref{thm:threemissing}.

With only one deleted addition, every useful deleted product lies before that cut. Its high input is a product of old linear forms, and the entire downstream linear context lies in one old neighborhood. Independent transversals of the missing matching require exponentially many distinct factors. Counting boundary nodes of the product DAG converts this into a bound on original deleted gates, sufficient for Theorem~\ref{thm:oneaddition}. When at least three edges are missing, overlapping transversals force all their endpoints to have one common old neighbor; this is impossible in a triangle-free graph.

For degree-one Boolean circuits, higher powers of the unmarked variables remain possible, so this homogeneity argument is unavailable. Instead we transpose the part of the circuit depending on the marked inputs. Its outputs become the neighborhood disjunctions. A component-counting argument shows that a family of such disjunctions with pairwise intersections of size at most one requires the sum of their sizes minus one disjunctions, even if conjunction gates are freely interspersed. Separately, first conjunctions on successful input paths supply a vertex cover. The two counts give the exact total.

\section{Finite-support polynomials over a zero-adjoined ring}
This section gives explicit formulas for the inverse-semiring structures used by the circuit translations.

Let $M$ be a commutative monoid, written multiplicatively. The semiring $\B[M]$ is the set $\Fin(M)$ of finite subsets of $M$, with union as addition and
\[
 ST=\{st:s\in S,\ t\in T\}
\]
as multiplication. Its zero is $\varnothing$ and its unit is $\{1_M\}$. Its additive order is inclusion. For $S\subseteq T$, extend a function $a:S\to R$ by zero to $T$.

\begin{proposition}\label{prop:monoid}
The semiring $R_0[M]$ consists of pairs $(S,a)$, where $S\in\Fin(M)$ and $a\in R^S$. Its operations are
\begin{align*}
 (S,a)+(T,b)&=(S\cup T,\widetilde a+\widetilde b),\\
 (S,a)(T,b)&=(ST,c),\qquad
 c(u)=\sum_{\substack{s\in S,\ t\in T\\st=u}}a(s)b(t).
\end{align*}
The additive inverse in the inverse-monoid sense is $(S,-a)$, and the additive idempotent associated to $(S,a)$ is $(S,0)$. Under the correspondence of \citet[Theorem~8.5]{CFMM}, the idempotent semiring is $\B[M]$, the functor is $S\mapsto R^S$ with extension-by-zero maps, and its lax multiplication is convolution.
\end{proposition}
\begin{proof}
By the definition of a formal monoid semiring, an element of $R_0[M]$ is a function $M\to R_0$ equal to $\tau$ outside a finite set. Let $S$ be precisely the set where the coefficient is different from $\tau$. Its remaining coefficients lie in $R$, including possibly $e$. Addition produces a coefficient different from $\tau$ exactly on the union: a sum of old ring elements is always an old ring element. Multiplication produces a coefficient different from $\tau$ exactly on $ST$, even when all contributions sum to $e$. These observations give the two formulas and their uniqueness.

Associativity of multiplication follows by regrouping the finite sum over triples $(s,t,u)$, and distributivity by distributing the finite sums. They remain valid if different pairs have the same product in $M$. The identity is $(\{1_M\},1_R)$ and the zero is $(\varnothing,0)$. In the additive monoid, the identities
\[
 (S,a)+(S,-a)+(S,a)=(S,a),\qquad
 (S,-a)+(S,a)+(S,-a)=(S,-a)
\]
hold. Conversely an additive inverse must have the same support, by these identities, and then has coefficient $-a$ in the group $R^S$. An additive idempotent has $a+a=a$, hence $a=0$. Their operations are therefore union and setwise product. Extension by zero is functorial. Convolution is bilinear, natural for inclusions, associative and unital, so these are precisely the lax monoidal data in the stated correspondence.
\end{proof}

The maps
\[
 R_0[M]\longrightarrow R[M],\quad (S,a)\longmapsto\sum_{m\in S}a(m)m,
 \qquad
 R_0[M]\longrightarrow\B[M],\quad (S,a)\longmapsto S
\]
are semiring homomorphisms. The first may erase coefficients; the second erases their values. A \emph{supported coefficient zero} is an $m\in S$ with $a(m)=0_R$.

\begin{example}
In $\Z_0[x]$, the expression $x-x$ is $(\{x\},0)$, whereas the zero polynomial is $(\varnothing,0)$. Multiplying the former by $x$ gives $(\{x^2\},0)$. Both have ring value zero and nonempty support.
\end{example}

\section{The inverse semiring of a Boolean function and a polynomial}
We next replace exponent support by the monotone predicate it generates.
Let $Q_n$ be the semiring of monotone functions $\{0,1\}^n\to\{0,1\}$ under pointwise disjunction and conjunction. For a monomial $m$, write $\supp(m)$ for its set of variables, and $1_A$ for the Boolean vector of a subset $A\subseteq\{1,\ldots,n\}$. Put
\[
 F(f)=\bigoplus_{\substack{m\text{ monomial}\\f(1_{\supp(m)})=1}}\Z m
 \ \subseteq\ \Z[x_1,\ldots,x_n].
\]
Every element is a finite sum, even though the group has infinitely many potential basis monomials.

\begin{proposition}\label{prop:booleanpair}
The assignments $f\mapsto F(f)$ and the inclusions for $f\le g$ form a lax monoidal functor from $Q_n$ to abelian groups. The associated inverse semiring $T_n$ has elements $(f,P)$ with $P\in F(f)$ and operations
\[
 (f,P)+(g,Q)=(f\vee g,P+Q),\qquad
 (f,P)(g,Q)=(f\wedge g,PQ).
\]
There is a surjective homomorphism
\[
 \Phi:\Z_0[x_1,\ldots,x_n]\longrightarrow T_n,
 \qquad (S,P)\longmapsto(f_S,P),
 \quad
 f_S(x)=\bigvee_{m\in S}\bigwedge_{i\in\supp(m)}x_i.
\]
An empty inner conjunction is true, and an empty outer disjunction is false.
\end{proposition}
\begin{proof}
If $f\le g$, every allowed monomial for $f$ is allowed for $g$, so the inclusion is defined and these inclusions compose. If $m$ is allowed for $f$ and $m'$ for $g$, then both $f$ and $g$ are true on $\supp(mm')=\supp(m)\cup\supp(m')$, by monotonicity. Thus polynomial multiplication gives a bilinear map $F(f)\otimes F(g)\to F(f\wedge g)$. These maps commute with the inclusions. Their associativity is polynomial associativity, and their unit is the integer-to-polynomial map $\Z\to F(1)$ sending $1$ to the constant monomial. This proves the lax monoidal assertion and the displayed formulas.

The relation $f_{S\cup T}=f_S\vee f_T$ is immediate. Also $f_{ST}=f_S\wedge f_T$: a Boolean assignment contains the support of some product $mm'$ if and only if it contains the supports of one monomial from each set. Proposition~\ref{prop:monoid} now proves multiplicativity and additivity of $\Phi$.

For surjectivity, take $(f,P)\in T_n$. Let $A_1,\ldots,A_k$ be the minimal true sets of $f$, and let $S$ be the union of the monomial support of $P$ with the squarefree monomials $\prod_{i\in A_j}x_i$. Give the added monomials coefficient zero where necessary. Every support in $S$ is a true set of $f$, while every true set contains one of the $A_j$. Therefore $f_S=f$. When $f=0$, the group is zero and the empty support gives its preimage.
\end{proof}

\begin{remark}
This quotient keeps the formal polynomial: in $T_n$, the pairs $(x_i,x_i)$ and $(x_i,x_i^2)$ differ. Only their Boolean coordinates agree. One must not impose $x_i^2=x_i$ on the polynomial coordinate when using its degree.
\end{remark}

\begin{proposition}[One gate list, two evaluations]\label{prop:lift}
Let $C$ have variable inputs and binary addition and multiplication gates, and let $P_C\in\N[x_1,\ldots,x_n]$ be its formal polynomial. Replacing the gates by disjunction and conjunction gives a Boolean function $f_C$, and evaluation in $T_n$ gives $(f_C,P_C)$. For $A\subseteq\{1,\ldots,n\}$, substitute $\tau$ for inputs outside $A$ and arbitrary integers $a_i\in\Z$ for inputs in $A$. The value in $\Z_0$ is $\tau$ if $f_C(1_A)=0$; otherwise it is the old ring element
\[
 \sum_{\supp(m)\subseteq A}[m]P_C\,m(a).
\]
If a set $X$ of variables is marked, the syntactic marked degree of $C$ is the total $X$-degree of $P_C$.
\end{proposition}
\begin{proof}
The first assertion follows by induction using Proposition~\ref{prop:booleanpair}. Equivalently, expanding the formal polynomial shows that the Boolean evaluation is true exactly when some occurring monomial uses only variables in $A$. In the zero-adjoined ring, a monomial using an input $\tau$ evaluates to $\tau$; every other monomial evaluates to an old ring element. A finite sum is $\tau$ exactly when every summand is $\tau$, which proves the evaluation formula, including the case where the integer sum is zero.

Each input polynomial is nonzero, and every gate has nonnegative coefficients, so no gate polynomial is zero and no monomial cancels. The $X$-degree of a sum is consequently the maximum of the degrees, and the $X$-degree of a product is their sum. Induction gives the last assertion.
\end{proof}

\begin{proof}[Proof of the first assertion of Theorem~\ref{thm:main}]
Suppose $C$ computes $f_G$. An occurring monomial with edge-free variable support would make $C$ true on an independent set, a contradiction. On the assignment whose true variables are the two endpoints of an edge, some monomial must be active. Its support contains an edge and is contained in those two endpoints, so it is exactly that edge. Thus $P_C\in\Sup(G)$. Conversely the two conditions defining $\Sup(G)$ make $f_C$ true exactly on assignments containing an edge. The replacement of operations preserves the complete gate list in both directions, so minimizing gives the equality.
\end{proof}

\section{The exact quadratic arithmetic count}
We now prove the second assertion of Theorem~\ref{thm:main}, retaining the constant in the gate count.

\begin{lemma}\label{lem:quadstructure}
Let $P=\sum_{uv\in E}c_{uv}x_ux_v$, with every $c_{uv}>0$. Delete gates that do not feed the output of a variable-input nonnegative arithmetic circuit for $P$. Every remaining gate computes a homogeneous polynomial of degree one or two. Every multiplication gate is the product of two linear forms with positive coefficients on disjoint variable sets. If $G$ has no four-cycle, the support of this product is a star.
\end{lemma}
\begin{proof}
Fix a gate $g$ and a directed path from it to the output. At each multiplication on that path choose a monomial in the other operand; at each addition follow the chosen path. Positivity implies that the resulting fixed monomial multiple of every monomial at $g$ occurs at the output. Two distinct degrees at $g$ would give two output degrees. Since no input is a constant, every gate has positive degree; hence its degree is one or two.

The two operands of any useful multiplication have positive degrees adding to two. They are therefore linear forms. Their variable sets $A,B$ are disjoint, because a common variable would give a square; every pair in $A\times B$ is an edge. If both sets had size at least two, four distinct endpoints would form a four-cycle. One set is therefore a singleton, so the support of the product is a star.
\end{proof}

\begin{proof}[Proof of the quadratic assertion of Theorem~\ref{thm:main}]
Apply Lemma~\ref{lem:quadstructure} to an arbitrary circuit for $P=\sum c_{uv}x_ux_v$. Because the required lower bound depends only on support, we may replace any degree-one gate supported on a single variable by that variable and then delete unused gates. This changes some positive coefficients but neither the output monomials nor its degree, and cannot increase the size. Every remaining linear gate contains at least two variables.

Let $K$ be the number of multiplication gates and $A_2$ the number of quadratic additions. Every gate following a quadratic gate is an addition. The quadratic network has $K+A_2$ nodes and $2A_2$ read occurrences. Each node except the output is read at least once, so $2A_2\ge K+A_2-1$, or $A_2\ge K-1$. This argument permits arbitrary sharing and repeated coefficients.

A product whose support has at least two edges has a unique star center. Any linear gate feeding its nontrivial factor can feed factors only at that center: otherwise two distinct centers would share two of its variables as neighbors, giving a four-cycle. Partition the linear gates according to this unique center. For center $v$, let $k_v$ be the number of product factors they feed, and let $d_v$ be the number of distinct variables occurring in the union of those factors. Let $L_v$ be the number of linear additions in this part of the circuit.

Its underlying undirected graph has $L_v+d_v$ vertices and at most $2L_v$ distinct edges, since each addition reads twice. Every connected component contains one of the $k_v$ designated factor outputs. A graph with $c$ components has at least its number of vertices minus $c$ edges, so
\[
 2L_v\ge L_v+d_v-k_v,\qquad L_v\ge d_v-k_v.
\]
Let $k_0$ count the one-edge products. The star supports cover $E$, so $\sum_v d_v+k_0\ge|E|$, while $\sum_vk_v+k_0=K$. Hence the number $L=\sum_vL_v$ of linear additions satisfies $L\ge|E|-K$. Choose one center for every product, either endpoint for a one-edge product. These centers cover $E$, giving $K\ge\vc(G)$. Altogether,
\[
 L+K+A_2\ge (|E|-K)+K+(K-1)\ge|E|+\vc(G)-1.
\]

For the converse, take a minimum vertex cover $U$. It is inclusion-minimal, so each $u\in U$ has an edge to a vertex outside $U$. Assign each edge to one of its endpoints in $U$; the preceding private edges ensure every center can receive an edge. For each $u$, form the sum of the variables at the other endpoints assigned to it, multiply by $x_u$, and sum the star outputs. There are $|E|-|U|$ inner additions, $|U|$ multiplications and $|U|-1$ outer additions. This is the claimed size.
\end{proof}

\begin{lemma}[Rigidity at the quadratic addition minimum]\label{lem:quadoptimal}
Under the quadratic hypotheses of Theorem~\ref{thm:main}, let a pruned circuit have $K$ multiplication gates. If its total size is $|E|+K-1$, its output is $H_G$: every edge coefficient is one. Its $K$ products are nonempty stars partitioning $E$. Every neighbor form is a binary sum tree of distinct variables, products with the same center have disjoint neighbor supports, and the products are joined by a binary sum tree. Every nonoutput gate is read exactly once; input variables may be shared. The centers need not be distinct.

In particular, a circuit of exactly $|E|+\vc(G)-1$ gates has $K=\vc(G)$, with distinct centers forming a minimum cover. No pruning qualification is needed in this optimal-total-size special case, since an unused gate would contradict the lower bound.
\end{lemma}
\begin{proof}
The preceding proof gives the fixed-product inequality $L+A_2\ge|E|-1$, where $L$ counts linear additions and $A_2$ quadratic additions. Replacing a singleton-support linear gate by its variable leaves all $K$ products useful on their unchanged quadratic output routes and preserves support. Thus there can be no such gate at equality; otherwise this replacement would give fewer than $|E|-1$ additions. Every center factor is therefore a literal variable.

Retain the notation $k_v,d_v,L_v,k_0$ from the quadratic proof, and let $c_v$ count the components in the linear operand graph assigned to center $v$. Its $k_v$ designated product-factor outputs meet every component. Hence
\[
 2L_v\ge L_v+d_v-c_v\ge L_v+d_v-k_v.
\]
The total equality forces
\[
 A_2=K-1,\qquad L_v=d_v-k_v,\qquad c_v=k_v,
 \qquad \sum_v d_v+k_0=|E|.
\]
Thus each center's operand multigraph is a forest with exactly one designated factor output in each component. Parallel or repeated operand incidences would give an extra edge and are impossible. Each component is a binary sum tree of distinct variable leaves, with coefficient one. Variable nodes were identified by name in this graph, so different components at the same center have disjoint leaf supports. In particular there is no shared nontrivial linear computation between two products at that center.

The equality in the support-incidence sum now says that all product-star supports, including the one-edge products, are pairwise disjoint. They cover $E$, so they partition it. The quadratic addition network is connected, has $K+A_2$ nodes and $2A_2$ operand incidences, and has one fewer edge than nodes. It too is a tree. Every product reaches the output once, proving that every output coefficient is one and establishing all read counts.

For the optimal-total-size special case, the fixed-product inequality and $K\ge\vc(G)$ force $K=\vc(G)$. The chosen centers cover $E$ and there are $K$ of them; they must therefore be distinct and form a minimum cover. Conversely, any partition of $E$ into $K$ nonempty stars gives $|E|-K$ leaf additions, $K$ products and $K-1$ outer additions, attaining $|E|+K-1$ gates. The general statement requires pruning: adding an unused multiplication increases both $K$ and the total size by one while preserving the numerical equality.
\end{proof}

\begin{lemma}[The lower envelope]\label{lem:lowerenvelope}
The nonzero terms of smallest total degree of a nonzero nonnegative variable-input circuit polynomial can be computed using no more gates than the original circuit.
\end{lemma}
\begin{proof}
For each gate polynomial $P_g$, let $d_g$ be its smallest monomial degree and $P_g^{\min}$ its homogeneous part of that degree. At a product gate, the desired part is the product of the two desired operand parts. At an addition gate with unequal operand degrees, it is the part from the operand of smaller degree, so that gate can be bypassed. With equal operand degrees, add the two parts. Processing the original gate list in order gives the assertion, using positivity to exclude cancellation. This is the lower-envelope argument of \citet[Lemma~2, p.~7]{JuknaRead} in the present variable-input model.
\end{proof}

\begin{proof}[Proof of the read-1 assertion of Theorem~\ref{thm:main}]
Let $C$ be read-1 for $f_G$. Every monomial of its polynomial contains an edge, so its smallest degree is at least two. Every edge also occurs as the squarefree monomial $x_ux_v$, by the read-1 hypothesis. The degree-two part is therefore $\sum_{uv\in E}c_{uv}x_ux_v$ with every coefficient positive. Lemma~\ref{lem:lowerenvelope} and the quadratic lower bound apply without increasing the gate count. The star circuit is read-1 and attains the bound.
\end{proof}

\begin{lemma}[Two gates lost with a missing edge]\label{lem:lostgates}
Let a variable-input nonnegative arithmetic circuit compute $P\in\Sup(G)$. Suppose its degree-two part is nonzero, and some edge $uv$ has no quadratic witness. The lower-envelope construction, followed by deletion of unused gates, removes at least two gates: an addition and a multiplication.
\end{lemma}
\begin{proof}
The polynomial has degree-two terms and also a monomial supported on $\{u,v\}$ of degree at least three. It is therefore not homogeneous. If every addition in the original circuit joined operands of equal minimum degree, induction from the homogeneous variable inputs would make every gate homogeneous. Thus an addition with unequal minimum degrees occurs, and the lower-envelope construction bypasses it.

Expand a parse tree for a monomial supported on $\{u,v\}$. At an addition choose one operand, and at a multiplication use both; after contracting the addition steps, the leaves are variables and the internal nodes are products. A product with no product below it has two variable leaves. At the corresponding original gate $g$, the selected monomial is one of $x_u^2,x_ux_v,x_v^2$, and its minimum degree is two, since each operand has positive minimum degree.

None of these three monomials occurs in the degree-two output: the squares have edge-free support, and $x_ux_v$ is missing by hypothesis. If $g$ remained useful in the lower-envelope circuit, no multiplication could follow it on a path to the degree-two output, because all factors have positive degree. Along an addition-only path its degree-two monomial would survive by positivity. This is a contradiction. Thus $g$ is removed as well, and it is distinct from the bypassed addition.
\end{proof}

\begin{lemma}[Two missing edges cost three discarded gates]\label{lem:threelost}
Under the hypotheses of Lemma~\ref{lem:lostgates}, assume in addition that $G$ has no four-cycle. If at least two edges lack quadratic witnesses, the lower-envelope construction followed by deletion of unused gates removes at least three original gates.
\end{lemma}
\begin{proof}
Suppose exactly two gates are removed. Lemma~\ref{lem:lostgates} identifies one unequal-minimum-degree addition $t$ and one multiplication $g$ whose minimum-degree part has a monomial forbidden at the quadratic output. These are all the removed gates. Every retained useful gate has minimum degree one or two.

Every path from $g$ to the output must be cut at an unequal-degree addition: the construction retains multiplications and equal-degree additions. The only such addition is $t$, so $g$ precedes $t$. Its forbidden quadratic monomial cannot occur at any retained useful gate, because any retained gate containing it has degree at least two, and an addition-only path to the quadratic output would preserve it. Thus the larger-degree operand of $t$ is itself $g$; an intermediate gate on that branch would also be removed. Since $g$ has minimum degree two, we have
\[
 t=a+bc,\qquad g=bc,
\]
where $a,b,c$ are nonzero linear forms. Indeed all three precede the only unequal-degree addition, so the homogeneous-input induction makes their original polynomials homogeneous; their minimum degrees are one. Their wires survive. The multiplication $g$ has no surviving use other than through $t$.

Treat the output of $t$ as a fresh degree-one variable $z$ in the retained quadratic circuit. Its output is homogeneous of degree two. A nonzero $z^2$ coefficient would give a square after substituting the nonzero linear form $a$, contrary to the squarefree quadratic output. Hence that output has the form $zD+Q$, where $D$ is a nonzero linear form and $Q$ is quadratic, both with nonnegative coefficients. The original and extracted polynomials are therefore
\[
 P=P_2+bcD,\qquad P_2=aD+Q.
\]
Every variable of $D$ is an old neighbor, in $G_2=(V,E_2)$, of any fixed variable of $a$.

Put $B=\supp(b)$, $C=\supp(c)$ and $D_0=\supp(D)$. Each of these sets with at least two elements lies in an old star neighborhood: this was just proved for $D_0$; for $B$ or $C$, its surviving linear wire reaches some first multiplication of the retained quadratic circuit, so every variable of the other factor is a common old neighbor. Choose one such center. No variable belongs to all three sets, since then $bcD$ would have a singleton-support monomial, forbidden by $P\in\Sup(G)$.

Consider a missing edge witnessed in $bcD$. One endpoint, say $u$, must occur in two of its three factors, and the other endpoint $v$ in the remaining factor, with support set $S$. If $|S|\ge2$, then $u$ is adjacent in $G$ to every vertex of $S$, because all monomials $u^2w$, $w\in S$, occur. The old center of $S$ has the same neighbors. Four-cycle-freeness forces that center to be $u$, which makes $uv$ an old edge, a contradiction. Thus the remaining factor is the singleton $\{v\}$.

If exactly one of $B,C,D_0$ is a singleton, every missing edge joins its vertex $v$ to the intersection of the other two sets. Two such edges would give $v$ two common neighbors with an old center of either nonsingleton set. That center cannot equal $v$, since the edges are missing, so again there would be a four-cycle. If two singleton factors have different vertices, there is at most one possible two-vertex support. If two singleton factors have the same vertex $u$, the third set either is a singleton or has at least two vertices; in the latter case its old center must be $u$ by the same four-cycle argument, and every resulting edge is old. These cases prove that at most one edge is missing, the required contradiction.
\end{proof}

\begin{proposition}[Edges without quadratic witnesses]\label{prop:missing}
Let $C$ be any monotone Boolean circuit for a four-cycle-free graph $G$. Let $E_2$ be the set of edges whose squarefree monomial occurs in its formal polynomial, put $G_2=(V,E_2)$ and $r=|E\setminus E_2|$. If $E_2\ne\varnothing$, then
\[
 |C|\ge
 \begin{cases}
 |E_2|+\vc(G_2)-1,&r=0,\\
 |E_2|+\vc(G_2)+1,&r=1,\\
 |E_2|+\vc(G_2)+2,&r\ge2.
 \end{cases}
\]
For $r>0$, in all cases,
\[
 |C|\ge\max\{1,\ |E|+\vc(G)-1-2r+2\}.
\]
If $r\ge2$ and $E_2\ne\varnothing$, the final $+2$ improves to $+3$. Consequently the full star lower bound holds when $r\le1$, and also when $r=2$, $E_2\ne\varnothing$ and $\vc(G)-\vc(G_2)\le1$. A saving of $\delta>0$ gates requires at least $\lceil(\delta+2)/2\rceil$ edges with no quadratic witness; if $E_2\ne\varnothing$, it requires at least $\lceil(\delta+3)/2\rceil$.
\end{proposition}
\begin{proof}
When $E_2$ is nonempty, the smallest-degree part is a positive quadratic polynomial with support exactly $E_2$. Apply the quadratic bound and Lemma~\ref{lem:lowerenvelope}; for $r=1$, Lemma~\ref{lem:lostgates} adds two gates, and for $r\ge2$, Lemma~\ref{lem:threelost} adds three. A cover of $G_2$, supplemented by one endpoint of each missing edge, covers $G$; hence $\vc(G_2)\ge\vc(G)-r$. Since $|E_2|=|E|-r$, the bounds in terms of $r$ follow. Every circuit for a nonempty edge function needs at least one gate, because a single input variable is true on a singleton assignment. If $E_2$ is empty, then $r=|E|$ and $\vc(G)\le|E|$, so the expression with $+2$ is at most one. This covers that case too. For $r=1$ the displayed bound is the full star bound; for $r=0$ Theorem~\ref{thm:main} applies. If $r=2$ and the cover number drops by at most one, the first displayed bound is at least $|E|+\vc(G)-1$. Rearranging proves the saving assertions; a saving already forces $r\ge2$.
\end{proof}

\subsection{Two missing witnesses and the equality case}
We give the additional argument needed when deleting two edges lowers the cover number by two.

\begin{lemma}[Four factors for a two-edge matching]\label{lem:fourfactors}
Let $T=\{u,v,w,t\}$, with matching edges $uv,wt$. Suppose a product of positive linear forms in these four variables has an exact-support witness for each matching edge and every monomial contains a matching edge. Then it has at least four distinct factor supports. If it has exactly four factors, their supports are the four cross pairs $\{u,w\},\{u,t\},\{v,w\},\{v,t\}$, each once.
\end{lemma}
\begin{proof}
Every factor meets both matching edges. There are four maximal independent subsets of $T$, namely the cross pairs. Each such set $J$ must be disjoint from a factor; otherwise choosing a $J$-variable in every factor gives a monomial with independent support. A factor disjoint from $J$ must contain the complementary cross pair, because it meets both edges, and it can contain nothing else in $T$. It is therefore exactly $T\setminus J$. One factor can reject only one of the four sets. This proves both assertions, including the requirement of distinct supports.
\end{proof}

\begin{proof}[Proof of Theorem~\ref{thm:twomissing}]
Write $H=(V,E_2)$ and $r=|E\setminus E_2|$. If $E_2$ is empty, then $G$ has at most two edges. A circuit depending on $n$ distinct inputs has at least $n-1$ binary gates: its connected ancestor multigraph has $n+k$ vertices and $2k$ input edges for $k$ gates. One edge, two incident edges and two disjoint edges therefore require at least one, two and three gates, respectively, exactly the asserted bounds.

Assume $E_2\ne\varnothing$. Proposition~\ref{prop:missing} settles $r\le1$, and also $r=2$ unless $\vc(G)=\vc(H)+2$. Suppose a counterexample exists in this last case. Put
\[
 B_2=|E_2|+\vc(H)-1.
\]
The quadratic lower bound and Lemma~\ref{lem:threelost} force its extracted circuit to have exactly $B_2$ gates, with exactly three original gates removed. The two missing edges are disjoint; write them $uv,wt$ and put $T=\{u,v,w,t\}$. Every minimum cover of $H$ avoids $T$: a cover containing an endpoint, supplemented by one endpoint of the other missing edge, would cover $G$ with only $\vc(H)+1$ vertices. Hence $H$ has no edge within $T$, and $G[T]$ is precisely the matching $uv,wt$.

\emph{The optimal quadratic remainder.} Lemma~\ref{lem:quadoptimal} forces exactly $\vc(H)$ products. Choose a center for each product star. The centers are distinct, since otherwise their set would be a smaller cover; they form a minimum cover $Z$, disjoint from $T$. There is no singleton-support linear addition: replacing such a gate by its variable preserves positive quadratic support and would give fewer than $B_2$ gates. Thus every center factor is its literal input variable. At most one product has any given center. Every retained quadratic wire has an addition-only path to the original output, obtained by restoring the bypassed additions along its extracted path. Its entire original polynomial therefore occurs at the output.

The last observation permits a useful simplification. A removed unequal-degree addition whose larger operand is a retained quadratic wire can be bypassed without changing the Boolean function or $E_2$. Any removed contribution contains a monomial of that wire, already available at the output along the independent addition-only path. In particular an exact two-variable witness cannot disappear. The extracted circuit is unchanged. Such a bypass would leave at most two removed gates, contrary to Lemma~\ref{lem:threelost}. We may therefore exclude it.

\emph{One removed addition and two removed products.} The addition is the only cut of a path in the minimum-degree extraction. Both removed products precede it and form a connected two-gate multiplication circuit in its larger operand. No retained intermediate wire can separate them from that operand, since its use there would retain the products. A retained quadratic input before the cut has only old edge monomials, and cannot supply a missing witness; it also vanishes when variables outside $T$ are set to zero. A contribution capable of supplying both missing edges is consequently either $UVW$ or $(UV)^2$, with $U,V,W$ retained linear forms.

If the addition keeps a degree-two input, its contribution reaches the output through additions. After restricting to $T$, Lemma~\ref{lem:fourfactors} excludes both alternatives: they have at most three distinct factor supports. If it keeps a linear input, a nonzero contribution on $T$ must enter a center factor, since all unchanged centers lie outside $T$. The kept input is therefore a literal center $h\in Z$, and only its one product is affected. Its other factor $L_h$ is unaffected by the replaced signal; otherwise substituting $h$ into both factors would give a square in the quadratic remainder. The new contribution is $UVW L_h$ or $(UV)^2L_h$.

The square alternative still has only three distinct supports. In the other alternative, the four restricted factors must be exactly the four cross pairs. Relabel so that $\supp(L_h)\cap T=\{v,t\}$. In fact $L_h$ has no additional variable $z$. If it did, the other three factors are all hit by $\{u,w\}$, giving a monomial of the original output supported within $\{u,w,z\}$. This set must contain an edge, so $zu$ or $zw$ is an edge. The first gives the four-cycle $u,v,h,z,u$, and the second gives $w,t,h,z,w$. The vertices are distinct: $h\notin T$, $z\notin T$, and $z\ne h$ because the quadratic star is squarefree. Both alternatives are impossible. Thus $L_h$ has exactly the two neighbors $v,t$.

Replace $h$ in $Z$ by $v,t$. All other product stars remain covered; the $h$-star is covered by its two leaves; and the two missing edges are covered as well. This is a cover of $G$ of size $\vc(H)+1$, a contradiction.

\emph{Two removed additions and one removed product.} A bottom multiplication in a parse tree for a missing witness shows that the removed product $g$ has minimum degree two, as in Lemma~\ref{lem:lostgates}. If both removed additions have unequal operand degrees, their larger operands must be $g$: all other possible larger operands are retained quadratic wires, already excluded by the bypass argument. Both additions follow $g$, whose inputs are therefore unchanged homogeneous linear forms $U,V$. Expanding the retained homogeneous quadratic circuit after the two substitutions gives
\[
 P=P_2+UVL+c(UV)^2,
\]
where $L$ is a nonnegative linear form and $c\ge0$. This follows by treating each replaced linear signal as a fresh degree-one variable; the output has total degree two in these variables and the ordinary inputs. Constant multiplicities and a zero $L$ are harmless.

Restrict to $T$. If $c=0$, the product $UVL$ would supply both matching-edge witnesses and contradict Lemma~\ref{lem:fourfactors}. If $c>0$, every monomial of $(UV)^2$ contains a matching edge. Squaring each monomial of $UV$ shows that every monomial of the quadratic product $UV$ must itself be a matching edge. A product of two linear forms can contain only one of the two disjoint matching edges without also containing a square or crosspair. Thus every perturbation monomial contains that one edge, and none can witness the other. Again there is a contradiction.

There remains one unequal addition and one private homogeneous addition $a$. If $a$ is quadratic and follows $g$, it adds a retained quadratic contribution to $g$, or repeats $g$. The retained contribution is already available at the output and repetition is Boolean-idempotent, so bypassing $a$ again contradicts the three-gate minimum. If it is quadratic and precedes $g$, all its terms contain old edges, and the removed product cannot create any missing witness. Therefore $a$ is linear, feeds $g$ directly, and $g=aV$ or $g=a^2$ for a retained linear $V$. The unequal addition keeps a linear input, and the output expansion is respectively
\[
 P=P_2+aVD\qquad\text{or}\qquad P=P_2+a^2D
\]
with a nonnegative linear context $D$. A squared occurrence of the replaced signal in the quadratic remainder is impossible because it would give a square after substitution. On $T$, each displayed product has at most three distinct factor supports. It cannot supply both matching edges by Lemma~\ref{lem:fourfactors}. These cases exhaust the three removed gates and prove the lower bound. The star circuit attains it.
\end{proof}

\subsection{Products of old neighborhood forms}
The support argument extends to any number of factors, although this does not yet handle an arbitrary circuit.

\begin{proposition}\label{prop:oldfactors}
Let $H\subseteq G$ be finite simple graphs on the same vertex set, with $G$ four-cycle-free. Let $F=\prod_{i=1}^k L_i$, where $k\ge1$ and the $L_i$ are nonzero positive linear forms. Suppose every nonsingleton support $S_i=\supp(L_i)$ is contained in $N_H(z_i)$ for some vertex $z_i$, and every monomial of $F$ contains a $G$-edge. Let $M$ be the edges of $G\setminus H$ that occur as exact two-variable supports of $F$, and put $r=|M|$. Then $M$ is an induced matching in $G$ and
\[
 k\ge2^r.
\]
If $r>0$ and $k=2^r$, the full supports $S_i$ are precisely the $2^r$ transversals choosing one endpoint of each edge of $M$, each once, with no extra vertices.
\end{proposition}
\begin{proof}
Let $A$ be the graph of exact two-variable supports of $F$. Then $A\subseteq G$, and $F$ has no singleton-support monomial. A pair $uv$ lies in $A$ exactly when every factor meets $\{u,v\}$: choosing such a variable in every factor produces a monomial on that pair, and singleton supports are forbidden.

A missing edge $uv\in M$ is isolated in $A$. If $uw\in A$ with $w\ne v$, choose a factor omitting $u$, which exists because $u^k$ is forbidden. It contains $v,w$ and has an $H$-center $z$ adjacent to both. The center cannot equal $u$, since $uv\notin H$; simplicity excludes equality with $v,w$. Thus $z,v,u,w,z$ is a four-cycle, a contradiction. In particular $M$ is a matching.

Take two of its edges $uv,wt$. Isolation gives $vw\notin A$, so some factor avoids $\{v,w\}$. That factor still meets both missing edges, and therefore contains $u,t$. Its $H$-center $z$ satisfies $zu,zt\in H$. If $uw\in G$, the vertices $z,u,w,t$ form a four-cycle. They are distinct, since $z=w$ would make $wt$ an old edge. Relabeling excludes every crossedge between the two missing pairs. Hence the matching is induced in $G$.

Write $W$ for its $2r$ endpoints. Each transversal $T$ of $M$ is independent in $G$, so some factor avoids $T$; otherwise $F$ has a monomial supported within $T$. Every factor meets every missing edge. A factor avoiding $T$ consequently has restriction $S_i\cap W=W\setminus T$, and can avoid at most this one transversal. The $2^r$ transversals require at least $2^r$ factors. For $r=0$ the inequality is simply $k\ge1$.

At equality with $r>0$, each transversal $T$ has exactly one rejecting factor $S_T$, whose restriction is $W\setminus T$. Suppose $z\in S_T\setminus W$. Every other factor meets $T$, so choosing $z$ in $S_T$ and a $T$-variable in every other factor gives a monomial supported within $T\cup\{z\}$. For $t\in T$, let $s$ be its missing-edge mate. The factor $S_T$ contains $s,z$ and has an $H$-center $a$ adjacent to both. An edge $zt\in G$ would give the four-cycle $a,s,t,z,a$; equality $a=t$ is excluded because $st\notin H$. Therefore $z$ has no neighbor in $T$. The set $T\cup\{z\}$ is independent, a contradiction. There are no extra vertices, proving the equality description.
\end{proof}
This exponential count concerns linear factors subject to the old-neighborhood hypothesis. It is not an exponential lower bound for arbitrary circuits: a circuit may use sums of products, share subcomputations, and have factors of higher degree. For three such factors there is at most one missing edge; four factors can supply two only in the displayed matching configuration.

\begin{lemma}[Common cubic factors force a missing-edge star]\label{lem:commoncubic}
Let $H\subseteq G$ be graphs on the same vertex set, with $G$ four-cycle-free. Let $U,V,D_1,\ldots,D_t$ be nonzero positive-coefficient linear forms, every nonsingleton support being contained in an $H$-neighborhood. Suppose every monomial of each $UVD_i$ contains a $G$-edge. Then all edges outside $H$ that occur as exact two-variable supports in these products share one vertex, unless there are none.
\end{lemma}
\begin{proof}
For a missing edge $ab$ witnessed by a cubic monomial, choose one variable from each factor. One endpoint, say $a$, is chosen twice. Let $W$ be the factor from which $b$ is chosen. For every $w\in\supp(W)$ the monomial $a^2w$ occurs, so $aw$ is an edge and $a\notin\supp(W)$. If $W$ were nonsingleton, take $w\ne b$ and an old center $z$ adjacent in $H$ to both. Since $ab\notin H$, we have $z\ne a$. The four distinct vertices $a,b,z,w$ form a four-cycle, a contradiction. Hence $W$ is the singleton $b$.

If $U$ or $V$ is a singleton, all missing edge supports already contain its vertex. Otherwise the singleton factor is always $D_i$. Orient each missing edge as $a_i b_i$, with $a_i\in\supp(U)\cap\supp(V)$ and $\supp(D_i)=\{b_i\}$. Fix an old center $z$ for $U$. If $a_i\ne a_j$, choose $a_j$ from both $U,V$ in the $i$th product. This forces $a_jb_i\in E(G)$. Also $za_i,za_j\in E(H)$, while $z\ne b_i$ because $a_ib_i$ is missing. The vertices $a_i,b_i,a_j,z$ are distinct and form a four-cycle. Thus every $a_i$ is the same vertex.
\end{proof}
This is a restriction on a family with two fixed linear factors. A general circuit need not have this form, so the lemma alone gives no unrestricted edge-circuit bound.

\begin{lemma}[Cubic grid]\label{lem:cubicgrid}
Let $H\subseteq G$ be graphs on the same vertex set, with $G$ four-cycle-free. Let $U_1,U_2,B,D_1,D_2$ be nonzero positive linear forms, each nonsingleton support lying in an $H$-neighborhood. Suppose every monomial of each $U_iBD_j$, $i,j\in\{1,2\}$, contains a $G$-edge. The four products together have at most three distinct missing edge supports.
\end{lemma}
\begin{proof}
In this proof we identify each linear form with its variable support. Each product has at most one missing edge by Proposition~\ref{prop:oldfactors}. If four distinct missing edges occurred, each cell would supply one. The singleton-factor argument in the proof of Lemma~\ref{lem:commoncubic} says that every cubic missing-edge witness repeats one endpoint in two factors and has a singleton third factor.

Suppose first that $B$ is nonsingleton. Each cell therefore has a singleton $U_i$ or $D_j$. If neither pair consisted entirely of singletons, a nonsingleton $U_i$ and a nonsingleton $D_j$ would meet in an impossible cell. By symmetry let $U_i=\{u_i\}$, with $u_1\ne u_2$ because all four edges are distinct. If $D_j=\{d_j\}$, both column edges are $u_1d_j,u_2d_j$; equality $d_j=u_i$ would leave the nonsingleton $B$ as the third factor of a repeated-endpoint witness, which is impossible. If $D_j$ is nonsingleton, both witnesses repeat their other endpoint in $B\cap D_j$. Every vertex of this intersection is adjacent to both $u_1,u_2$, and neither $u_i$ is in the intersection, since that would give a singleton-support monomial. Four-cycle-freeness makes the intersection a singleton $v_j$. Thus each column has the form $u_1v_j,u_2v_j$. Four distinct edges force $v_1\ne v_2$ and give a four-cycle.

Now let $B=\{b\}$. A singleton $U_i\ne\{b\}$ would fix one edge throughout its row. If $U_i=\{b\}$, both $D_j$ would have to be singletons, fixing one edge throughout each column. Either contradicts four distinct edges. By symmetry all $U_i,D_j$ are nonsingleton. The missing edge in cell $(i,j)$ is therefore $bw_{ij}$ with $w_{ij}\in\supp(U_i)\cap\supp(D_j)$. In particular $U_1$ contains two distinct new neighbors of $b$. Their old common center either is $b$, making both edges old, or forms a four-cycle with $b$. This final contradiction proves the bound.
\end{proof}

\begin{lemma}[Two extensions of a cubic]\label{lem:sharedquartics}
Let $H\subseteq G$ be graphs on the same vertex set, with $G$ four-cycle-free. Let $A,B,C,D,E$ be nonzero positive linear forms, each nonsingleton support contained in an $H$-neighborhood. Suppose every monomial of $ABCD$ and $ABCE$ contains a $G$-edge. Their combined exact two-variable supports outside $H$ comprise at most two edges.
\end{lemma}
\begin{proof}
Each quartic has at most two missing edges. If neither has two, the assertion follows. Otherwise let $ABCD$ have two. Proposition~\ref{prop:oldfactors} makes them an induced matching on four vertices and makes the four full factor supports its four cross pairs. The three supports of $A,B,C$ have empty common intersection. Every monomial of $ABC$ therefore uses at least two variables, all among those four. An exact two-variable support of $ABCE$ cannot contain an outside vertex, since its $ABC$ part would then have singleton support. Among the four matching vertices, the only $G$-edges are the two matching edges. Thus $ABCE$ adds no missing edge beyond them.
\end{proof}

\begin{lemma}[Three related cubic products]\label{lem:relatedcubics}
Let $H\subseteq G$ be graphs on the same vertex set, with $G$ four-cycle-free. Let $U,V,L,D,A,B$ be nonzero positive linear forms, each nonsingleton support contained in an $H$-neighborhood. Let $a$ be a nonzero positive linear form supported in $\supp(A)\cap\supp(B)$, with $aL$ supported on $H$-edges. Suppose every monomial of each of $UVL$, $UVD$ and $ABD$ contains a $G$-edge. The first two products together have at most one missing exact edge support; all three together have at most two.
\end{lemma}
\begin{proof}
An old-neighborhood support meeting two missing edges $xy,xz$ must contain $x$: otherwise it contains $y,z$, whose common old center gives a four-cycle with $x$, or makes the edges old. Also an old-neighborhood support consisting of $G$-neighbors of $x$ and containing a missing neighbor $y$ must be the singleton $y$, by the same four-cycle argument.

Each cubic has at most one missing edge. If the first two have different edges, Lemma~\ref{lem:commoncubic} says that they share a vertex, so write them as $xy,xz$ with $y\ne z$. Each of $U,V$ meets both pairs and therefore contains $x$. Validity of the two cubics makes every variable of $L,D$ a neighbor of $x$. Their respective missing witnesses and the preceding observation force $L=\{y\}$ and $D=\{z\}$ as supports.

Choose $w\in\supp(a)$. The old-edge condition gives $wy\in H$, excluding $w=y$ and also $w=x$, since $xy$ is missing. The monomial $w^2z$ occurs in $ABD$, so validity excludes $w=z$ and gives $wz\in G$. Thus $x,y,w,z$ are distinct and form a four-cycle, a contradiction. The first two products have at most one missing edge altogether, and $ABD$ adds at most one.
\end{proof}

\begin{lemma}[Three missing edges cost four discarded gates]\label{lem:fourlost}
Let $P\in\Sup(G)$ for a four-cycle-free graph $G$, and assume its quadratic part is nonzero. If at least three edges lack quadratic witnesses, minimum-degree extraction and pruning remove at least four gates. No optimality assumption on the quadratic remainder is required.
\end{lemma}
\begin{proof}
Lemma~\ref{lem:threelost} already gives three removals. Suppose there are exactly three. A retained quadratic wire has an addition-only path to the original output; every retained nonsingleton linear support lies in an old neighborhood by following it to its first retained product. As in the proof of Theorem~\ref{thm:twomissing}, a removed unequal addition whose larger input is a retained quadratic wire can be bypassed while preserving the Boolean function and the extracted quadratic support. This would leave at most two removals, contradicting Lemma~\ref{lem:threelost}. The same applies to a private quadratic sum of the removed product with a retained quadratic wire, or to a repeated sum of the removed product. These absorption arguments use no optimality of the remainder.

There is at least one removed product and one unequal addition. First suppose there are two products and one addition. The two products form a chain before the unique cut. A retained quadratic factor makes every resulting monomial contain an old edge and cannot contribute a missing witness. Otherwise the discarded input is $BCD$ or $(BC)^2$, with retained linear factors. If the kept input is quadratic, its output context is scalar. If it is linear, the context is linear, say $L$, and lies in an old neighborhood because its product with any variable of the kept linear signal occurs in the quadratic remainder. A squared occurrence of that signal is excluded by squarefreeness. In the nonsquared case there are at most four old-neighborhood factors, so Proposition~\ref{prop:oldfactors} allows at most two missing edges. In the squared case there are at most three distinct factor supports. The transversal-count proof of that proposition requires $2^r$ distinct restrictions, so repeated factors cannot evade the count; here there is at most one missing edge.

Now suppose there is one removed product $g$ and two unequal additions. The bottom-product argument in Lemma~\ref{lem:lostgates} gives $\deg_{\min}g=2$. Both larger inputs must be $g$, since retained quadratic larger inputs were excluded, and the other unequal addition has minimum degree one. Hence $g=UV$ for unchanged retained linear forms. Expanding the two insertions in the retained quadratic circuit gives
\[
 P=P_2+UV(L_1+L_2)+c(UV)^2.
\]
Each nonzero $L_i$ is a linear context in an old neighborhood: its product with any variable of the corresponding kept linear signal has support in $P_2$. This remains true for nested insertions; individual squared signal terms are forbidden by squarefreeness, and the displayed square comes from the cross term of the two insertions. If $c=0$, each product $UVL_i$ contributes at most one missing edge by Proposition~\ref{prop:oldfactors}, giving at most two in all.

If $c>0$, squaring each monomial of $UV$ shows that all of them have edge support. Thus $U,V$ have disjoint supports and their complete bipartite graph is a star. Suppose $U$ is the singleton $u$. If $V$ has at least two variables, its old center $h$ and $u$ have two common neighbors. Four-cycle-freeness forces $h=u$, so every edge of $UV$ is old. If $V$ is a singleton, the product has only one edge. All perturbation monomials are multiples of $UV$, so in either case at most one missing edge can result.

The last case has one removed product, one unequal addition, and one private homogeneous addition $a$. A quadratic $a$ after the product is eliminated by the absorption argument; a quadratic $a$ before it gives only old-edge multiples. Higher degree is impossible with only one removed product of minimum degree two. Thus $a=U+V$ is linear with retained linear operands, and directly feeds the removed product. If that product is $aB$, the new contribution is
\[
 (U+V)BL=UBL+VBL,
\]
where $B$ and the linear context $L$ have old-neighborhood supports. Each summand contributes at most one missing edge. If the product is $a^2$, the contribution is $a^2L$. The supports of $a,L$ are disjoint, since otherwise a pure cube occurs; every pair between them is an edge, since every $x^2y$ occurs. Four-cycle-freeness makes one support singleton. If $a$ is singleton, an old center of a nonsingleton $L$ forces all edges to be old; with two singletons there is only one edge. If $L$ is singleton, a nonsingleton $U$ or $V$ likewise contributes only old edges. Each singleton among $U,V$ contributes at most one new edge, so there are at most two in total. Every case contradicts three missing edges.
\end{proof}

\begin{lemma}[Two independent inserted products]\label{lem:independentproducts}
Let $H\subseteq G$ be graphs on the same vertex set, with $G$ four-cycle-free. Let $U,V,W,X,L,M$ be nonzero positive linear forms, each nonsingleton support contained in an $H$-neighborhood. Let $A,B$ be nonzero positive linear forms with
\[
 \supp(A)\subseteq\supp(M),\qquad
 \supp(B)\subseteq\supp(L).
\]
Suppose the quadratic products $AB$, $AL$ and $BM$ are supported on $H$-edges. Suppose also that every monomial of each of $UVL$, $WXM$ and $UVWX$ contains a $G$-edge. Then the union of exact two-variable supports outside $H$ in these three latter products has at most two edges.
\end{lemma}
\begin{proof}
Write $g=UV$, $h=WX$. Each cubic has at most one missing edge and the quartic has at most two. If the quartic has two, its factors are exactly the four cross pairs of an induced matching, by Proposition~\ref{prop:oldfactors}. If $U,V$ are disjoint pairs, $g$ has no pure square, so $gL$ cannot witness an edge with a vertex outside the matching; inside its four endpoints there are no other $G$-edges. Otherwise $U,V$ meet at one endpoint $u$, while $W,X$ meet at its matching mate $v$. Any further edge of $gL$ has form $uz$, with $z$ outside the matching and $z\in\supp(L)$. Choose $a\in\supp(A)$. Since $A$ is supported in $M$, validity of $hM$ forces $va\in G$ through the monomial $v^2a$, and $AL$ gives $az\in H$. Equality $a=u$ makes $uz$ old, $a=z$ gives an old square, and $a=v$ gives a pure cube. Otherwise $u,v,a,z$ are distinct and form a four-cycle. Thus this cubic adds no missing edge. The other cubic is symmetric.

It remains to exclude a quartic missing edge $e_0$ and two different additional cubic missing edges $e_1,e_2$. We show that $e_1=uz$ with $u^2$ occurring in $g$ and $\supp(L)=\{z\}$. If $U,V$ are nonsingleton, the cubic-witness argument in Lemma~\ref{lem:commoncubic} forces the singleton third factor to be $L$ and repeats $u$ in $U,V$, proving this assertion. Otherwise let $U$ be the singleton $u$. Both $e_0$ and $e_1$ contain $u$, say $e_0=ut$, $e_1=uz$. The support of $V$ meets both pairs. If it omitted $u$, it would contain $t,z$; its old center and the two missing edges would give a four-cycle, or make the edges old. Thus $u\in\supp(V)$ and $u^2$ occurs in $g$. Validity of $gL$ makes every variable of $L$ a neighbor of $u$ and excludes $u$ itself. Since $uz$ is missing and $z\in\supp(L)$, a second variable of $L$ would again give a four-cycle through its old center. Hence $L$ is the singleton $z$.

Similarly write $e_2=vw$, with $v^2$ in $h$ and $M$ the singleton $w$. The nonzero forms $A,B$ are supported on $w,z$, respectively, so $wz\in H$. Validity of $gh$ gives $uv\in G$ through $u^2v^2$, and excludes $u=v$. The four vertices $u,z,w,v$ are distinct: the two missing edges exclude $u=z$ and $v=w$; the old edge excludes $w=z$; $u=w$ or $v=z$ would make one of those missing edges old. They form the four-cycle $u,z,w,v,u$, a contradiction. Thus there are at most two missing edges.
\end{proof}

\begin{lemma}[Four missing edges cost five discarded gates]\label{lem:fivelost}
Let $P\in\Sup(G)$ for a four-cycle-free graph $G$, and assume its quadratic part is nonzero. If at least four edges lack quadratic witnesses, minimum-degree extraction and pruning remove at least five gates.
\end{lemma}
\begin{proof}
Put $H=G_2$. Call a nonzero linear form \emph{old} if its support is a singleton or is contained in an $H$-neighborhood. Every retained linear signal is old. We count missing exact two-variable supports; positive coefficients have no effect on that count. A product is called valid when every monomial contains a $G$-edge.

We use Proposition~\ref{prop:oldfactors} in the following forms. A valid product of at most three old forms has at most one missing edge; one with at most five distinct old factors has at most two. Repeated powers of a factor do not affect its Boolean function, and the transversal argument requires distinct factor supports. Lemma~\ref{lem:sharedquartics} bounds two quartics sharing three old factors by two missing edges altogether.

Suppose four gates are deleted; Lemma~\ref{lem:fourlost} excludes fewer. We may remove a Boolean-redundant private branch only when doing so preserves the retained extracted circuit itself. Merely decreasing the whole circuit size would not suffice, because a retained gate might otherwise become a newly deleted gate. Here the following reductions do preserve that circuit. A retained quadratic gate has an addition-only path to the original output, so its complete original polynomial $Q$ is already an additive summand there. Any further product branch divisible by $Q$ is absorbed by this summand. A private homogeneous sum $Q+F$ can be replaced by $F$: its whole output branch is pruned, since an equal-degree addition on the extracted circuit would be retained. A cut whose larger input is $Q$ can be bypassed through its unchanged smaller input. Repeated private sums $F+F$ can also be bypassed. Propagating a formally removed zero branch to the first surviving sum leaves a variable-input circuit with no constants. These operations leave all retained paths intact and remove a deleted gate, contradicting Lemma~\ref{lem:fourlost}. Thus none occurs in an essential four-deletion configuration. In what follows a \emph{cut} is an addition of unequal minimum degrees.

There is a deleted product and a cut. The first deleted product has minimum degree two: its boundary inputs are retained linear signals, since a retained quadratic factor was excluded. With four gates the product counts are one, two, or three. Any further private additions are linear, except a possible sum of two distinct private quadratic products. A higher-degree private sum would need two distinct higher-degree products, unavailable with at most two products and the required cut; a retained quadratic summand or a repeated sum was excluded.

\emph{Three products and one cut.} The high branch is a multiplication DAG with at most four distinct old boundary inputs. Its context is scalar if the cut keeps a quadratic signal, and old linear if it keeps a linear signal. In the latter case multiplying the context by any variable of the old low part gives quadratic output monomials. Thus there are at most five distinct old factors, allowing at most two missing edges.

\emph{Two products, one cut, and one private sum.} A private quadratic sum joins the two independent quadratic products. Including its context gives two old cubics, or two quadratics, with at most two missing edges. A private linear sum is $A+B$ with $A,B$ old. The products form a chain above it. If the sum occurs once, expansion gives two products of at most four old factors; the four-factor case has three common factors and is bounded by Lemma~\ref{lem:sharedquartics}. If the sum occurs repeatedly, Boolean idempotence leaves at most three distinct factors including the context, one being $A+B$. Expansion then gives two old cubics. In either case there are at most two missing edges.

\emph{One product and three additions.} No cut can precede the sole product $g$, because its high input would have to be retained quadratic. Private additions before $g$ are linear; those after it repeat $g$ or add a retained quadratic signal and were excluded. With three cuts and no private sum, $g=UV$ with old forms and the perturbation is
\[
 g(D_1+D_2+D_3)+c g^2,
\]
where the $D_i$ are old linear contexts, possibly zero, and $c\ge0$. If $c=0$, the three cubics give at most three missing edges. If $c>0$, squaring each monomial shows that $g$ itself is valid. Its support is then a complete bipartite star, and every later edge witness contains an edge already in $g$. Proposition~\ref{prop:oldfactors} allows at most one missing edge.

With two cuts and one private sum, $g=(A+B)W$ or $(A+B)^2$, and the perturbation is $g(D_1+D_2)+c g^2$. If $c>0$, $g$ is valid and its quadratic support is a star. Each old summand can contribute at most one new leaf: two neighbors of a different center in one old neighborhood would form a four-cycle. Hence at most two missing edges occur. If $c=0$ and $g$ is a square, Boolean reduction gives $(A+B)(D_1+D_2)$, whose valid complete bipartite support is a star; the same argument gives at most two. Otherwise expansion gives the four cubics $AWD_j,BWD_j$, and Lemma~\ref{lem:cubicgrid} gives at most three.

With one cut and two private linear additions, the used addition DAG has either two parallel sums or a chain $J=A+B$, $K=J+D$. The product uses $K$ with an old form, $K^2$, or $KJ$; other chain uses leave a private gate unused or repeat a sum. A three-old-form sum times an old form expands, after its old context is included, into three cubics. Its square reduces to a bipartite star with at most one new leaf per old summand. Two parallel sums give the cubic grid, with the context as common factor. Finally $J(J+D)$ is Boolean-equivalent to $J$, giving at most two missing edges after its old context. Scalar contexts only reduce these bounds. This exhausts the one-product case.

\emph{Two products and two cuts.} Write $g=UV$ for the first product in topological order, and $f$ for the second. If one cut is the high input of the other, its low degree is two and the outer low degree one. An inner retained quadratic low signal can be removed from this wholly pruned branch using its unaffected baseline path. Otherwise that low signal is $g$, and $f$ is a multiple of $g$; Boolean absorption $g+f=g$ removes the inner cut. In each case the outer cut keeps its same low input, so the extracted circuit is unchanged. Thus both cut high inputs are products directly.

If both insert $f$, neither cut precedes it. The second product is $ABC$ or $(AB)^2$, with old forms and possible repeated factors. The perturbation is $fD_1+fD_2+c f^2$, with scalar or old linear contexts. With $c=0$, Lemma~\ref{lem:sharedquartics} handles two quartics sharing three factors, or Boolean reduction gives two cubics. With $c>0$, $f$ itself is valid and every missing witness is already an edge of $f$, which has at most one. These cases give at most two missing edges.

It remains that one cut inserts each product. The ordering of the products does not order their insertion cuts. Interpret every old low part as its extracted homogeneous part. Temporarily replace the two high inputs by independent symbols $y_1,y_2$. Retained linear signals are positive linear expressions in the variables and degree-one insertion symbols; retained quadratic signals are sums of products of two such expressions, together with additive degree-two insertions. A squared symbol is impossible, since the same variable from its nonzero old low part would then produce a square in $H$. Consequently the frame output is
\[
 P_2+y_1D_1+y_2D_2+c y_1y_2.
\]
Each linear $D_i$ is old because its product with the corresponding old low part occurs in $P_2$. A scalar $D_2$ is possible only when $f$ has minimum degree above two, and then $c=0$. If either cut's low input depends on the other insertion, their two paths cannot meet at a retained product, again by the forbidden old square; thus $c=0$ in either temporal order. For $c>0$, the product supplying the cross term also supplies the product of the two old low parts and both mixed terms. These are mandatory, even with shared paths.

Substitute $y_1=g,y_2=f$. If $f=WX$ is independent of the first insertion, the perturbation is
\[
 UVD_1+WXD_2+c UVWX.
\]
For $c=0$ the two cubics give at most two missing edges. For $c>0$, the mandatory old and mixed contexts give every hypothesis of Lemma~\ref{lem:independentproducts}, with $L=D_1,M=D_2$ and $A,B$ the extracted old low parts. Thus this case has at most two missing edges as well.

If an input to $f$ is affected by the $g$ insertion, that cut necessarily precedes $f$. Such an input has form $A+\alpha g$, where $A$ is old and contains the nonzero old low part $a$ of the cut. Suppressing positive multiplicities, if exactly one input is affected then $f=(A+g)W$. For $c>0$, Boolean absorption reduces the perturbation to
\[
 UVD_1+AWD_2+UVW,
\]
These are valid old cubics. The mandatory cross context puts the old low part $a$ in $D_2$, while $a$ is supported in $A$ and $aD_1$ is old. If $D_1\ne0$, Lemma~\ref{lem:relatedcubics}, with $L=D_1$, $D=W$ and its last two factors $A,D_2$, gives at most two missing edges. If $D_1=0$, two cubics remain. For $c=0$ the parts are
\[
 UVD_1,\qquad AWD_2,\qquad UVWD_2.
\]
In fact if $D_1\ne0$, there are at most two missing edges. If $D_2=0$, only the first cubic remains. Otherwise find $k\in\supp(D_1)$ with $\supp(A)\subseteq N_H(k)$. If the affected input is the cut output itself, $A=a$ as supports and any $k\in\supp(D_1)$ works, since $aD_1$ is old. Otherwise the input is a retained linear gate: no private-addition budget remains. Follow its retained path to its first product, with opposite extracted linear form $B$. The quadratic output contains $AB$, and the direct insertion context $D_1$ contains $B$, because the affected signal contributes a positive multiple of $g$ along this addition path. Any $k\in\supp(B)$ works. Apply Lemma~\ref{lem:independentproducts} to the products $UV$ and $WD_2$, contexts $D_1,A$, and old low forms $a,\{k\}$. Its required old products are $aD_1$, $kA$ and $ak$, and $a$ is supported in $A$. This gives the stronger bound of two. If $D_1=0$, only one cubic and one quartic remain, giving at most three.

If both inputs are affected, Booleanly $f=(A+g)(B+g)=AB+g$. For $c=0$, the perturbation is the sum of the three old cubics $UVD_1,UVD_2,ABD_2$. The first old low part $a$ is supported in both $A,B$ and $aD_1$ is old, so Lemma~\ref{lem:relatedcubics} gives at most two missing edges; a zero context leaves fewer terms. For $c>0$, the perturbation reduces to $g+ABD_2$, again with at most two. A direct $g$ input leaves $f=gW$, $g(B+g)$, or $g^2$. The latter two are Boolean-equivalent to $g$. The first gives one cubic and one quartic if $c=0$, and two cubics if $c>0$. Scalar contexts cannot increase these bounds. All positive summands after these Boolean simplifications are valid, since the whole expression still computes $f_G$; these are support comparisons, distinct from the earlier frame-preserving circuit rewrites.

Every four-deletion configuration has at most three missing edges, a contradiction. Therefore at least five gates are deleted.
\end{proof}

\begin{corollary}[An affected input with a direct context]\label{cor:affectedcase}
Consider the two-product, two-cut configuration established in the proof of Lemma~\ref{lem:fivelost}. Let its products be $g=UV$ and $f=(A+\alpha g)W$, with $\alpha>0$, and suppose the quadratic cross coefficient of the two insertions is zero. If the direct linear output context $D_1$ of the $g$ insertion is nonzero, at most two edges lack quadratic witnesses. Optimality of the retained quadratic circuit is not required.
\end{corollary}
\begin{proof}
If $D_2=0$, there is only one valid old cubic. Otherwise the retained-path argument in that proof gives $k\in\supp(D_1)$ with $kA$ supported on old edges. Together with the old low part $a$, this supplies all the mandatory contexts for Lemma~\ref{lem:independentproducts}, applied to $UV$ and $WD_2$. The three valid contributions are $UVD_1$, $AWD_2$ and $UVWD_2$.
\end{proof}

\begin{corollary}[An affected input with zero direct context]\label{cor:zeroaffected}
In the configuration of Corollary~\ref{cor:affectedcase}, suppose instead that $D_1=0$, the retained quadratic circuit is optimal, and exactly three edges lack quadratic witnesses. Then
\[
 \vc(G)\le\vc(H)+1.
\]
Thus neither direct-context case permits a below-star circuit with three missing witnesses and cover-number increase two, when this one-affected-input configuration has zero cross coefficient.
\end{corollary}
\begin{proof}
The contributions outside the quadratic part are $AWD_2+\alpha UVWD_2$. Three missing edges force two in the quartic, so Proposition~\ref{prop:oldfactors} makes its factors the four full cross pairs of an induced matching. In particular $K=\supp(D_2)$ has two vertices. Every missing edge meets $K$, because every witness has the factor $D_2$.

Let $b$ be the extracted low part of the second cut. Since $bD_2$ is supported on $H$-edges, the two supports are disjoint. If $b$ had two variables, their complete bipartite product would contain a four-cycle. Thus $b$ is supported on a singleton $h$.

By Lemma~\ref{lem:quadoptimal}, the retained products are stars at the distinct vertices of a minimum cover $C$ of $H$, with their edge sets partitioning $H$. An insertion path entering a neighbor factor contributes only that star's center to $D_2$, hence a vertex of $C$. A contribution outside $C$ can arise only by entering the center factor of the star assigned to $h$. In that event $h\in C$ and its entire assigned leaf set is contained in $K$. If $K\subseteq C$, the old cover already covers all missing edges. Otherwise replace $h$ by $K$: its assigned edges are covered by $K$, the other assigned stars retain their centers, and every missing edge meets $K$. This gives a cover of size at most $|C|+1$. Fanout into additional stars merely adds their centers to $D_2$ and does not affect this argument.
\end{proof}

\begin{corollary}[A direct dependent product]\label{cor:directproduct}
Consider the two-product, two-cut configuration of Lemma~\ref{lem:fivelost}, with products $g=UV$ and $f=gW$, zero quadratic cross coefficient, and optimal extracted quadratic circuit for $H$. If exactly three edges of $G$ lack quadratic witnesses, then $\vc(G)\le\vc(H)+1$.
\end{corollary}
\begin{proof}
The output has the form
\[
 P=P_2+UVD_1+UVWD_2,
\]
with positive coefficients absorbed into the linear forms. Three missing edges force the quartic to supply two and the cubic to supply a third; a zero or scalar context cannot do this. Proposition~\ref{prop:oldfactors} therefore makes the four quartic factors the four full cross pairs of an induced matching. If $U,V$ were disjoint, any two-variable support meeting both would be contained in those four matching vertices, so the cubic could add no third edge. After renaming vertices, their supports are consequently
\[
 U=\{x,z\},\quad V=\{x,t\},\quad
 W=\{y,z\},\quad K:=\supp(D_2)=\{y,t\},
\]
and the matching is $xy,zt$. The third edge must be $xw$ for a vertex $w$ outside these four vertices: an outside vertex in an exact two-variable support leaves $x$ as the only possible common choice from $U,V$.

In fact $D_1$ has singleton support $\{w\}$. Validity of the terms obtained by selecting $x$ in both $U,V$ excludes $x$ from $D_1$ and requires $xd\in E(G)$ for every $d\in\supp(D_1)$. If $d\ne w$ also belonged to this support, its old neighborhood center $a$ would give old edges $aw,ad$. The center cannot equal $x$, since $xw$ is missing from $H$, and $x,w,a,d$ would form a four-cycle. Thus no such $d$ exists. Also $wt\notin E(H)$: an old center $b$ for $V=\{x,t\}$ would otherwise give the four-cycle $b,x,w,t$, and $b=w$ is impossible because $xw$ is not old.

Let $C$ be the distinct assigned star centers from Lemma~\ref{lem:quadoptimal}; they form a minimum cover of $H$. We use the context observation proved in Corollary~\ref{cor:zeroaffected}: if a linear insertion context contains a vertex outside $C$, the insertion enters the literal center of one assigned star, whose entire assigned leaf set is contained in that context. Other insertion paths contribute only centers in $C$.

If $K\subseteq C$, adjoining $w$ covers all three missing edges. Otherwise the second cut enters a center $h_2\in C$ whose entire assigned leaf set lies in $K$, and $h_2K$ consists of old edges. If $w\in C$, then $h_2\ne w$ because $wt$ is not old. Replacing $h_2$ by $K$ preserves $w$, covers the old assigned stars and the three new edges, and increases the cover size by at most one.

Finally suppose $w\notin C$. The first cut then enters a center $h_1\in C$ whose nonempty assigned leaf set is $\{w\}$. Since $w\notin K$, this star differs from the star at $h_2$. Replace both centers by $\{w\}\cup K$. Their assigned edges are covered by their respective leaves; every other assigned star keeps its center; and the three missing edges are covered. The resulting cover has at most $|C|-2+3=|C|+1$ vertices. The context observation allows fanout and either order of the insertion cuts, so these cases exhaust the configuration.
\end{proof}

\begin{corollary}\label{cor:fourlostbound}
If $r\ge3$ and $E_2\ne\varnothing$, the bound in Proposition~\ref{prop:missing} improves to
\[
 |C|\ge |E_2|+\vc(G_2)+3
 \ge |E|+\vc(G)-1-2r+4.
\]
If $r\ge4$, it improves further to
\[
 |C|\ge |E_2|+\vc(G_2)+4
 \ge |E|+\vc(G)-1-2r+5.
\]
Any saving of $\delta>0$ gates below the star count requires at least three missing quadratic witnesses. When $E_2\ne\varnothing$, it requires
\[
 r\ge\max\{3,\lceil(\delta+4)/2\rceil\}.
\]
\end{corollary}
\begin{proof}
Add the four or five discarded gates to the quadratic lower bound, and use $\vc(G_2)\ge\vc(G)-r$. Theorem~\ref{thm:twomissing} excludes a saving for $r\le2$; rearranging the first inequality gives the stated uniform saving condition.
\end{proof}

\begin{example}[The first four discarded-gate estimates are sharp]
For the path with edges $xu,uy,yv$,
\[
 (x+y)(u+yv)=xu+yu+xyv+y^2v
\]
has four gates and computes the edge function. Its lower envelope $u(x+y)$ has two gates. Exactly one edge has no quadratic witness, and exactly one addition and one multiplication disappear. For the star with edges $ad,dx,dy$,
\[
 d\bigl(a+(x+y)^2\bigr)=ad+dx^2+2dxy+dy^2
\]
also has four gates. Reusing the wire $x+y$ in the square is allowed. Its lower envelope $ad$ has one gate, so two missing edges can cost exactly three discarded gates. This second circuit is one gate larger than the optimal three-gate star circuit; it is no counterexample to the conjecture.

More generally, build $x_1+\cdots+x_t$ with $t-1$ additions and compute
\[
 d\bigl(a+(x_1+\cdots+x_t)^2\bigr).
\]
Its Boolean graph is the star with edges $da,dx_1,\ldots,dx_t$, while its quadratic part is $da$. The circuit has $t+2$ gates and extraction retains one, so exactly $t+1$ disappear. In particular $t=3$ attains Lemma~\ref{lem:fourlost}, and $t=4$ attains the five-gate bound in Lemma~\ref{lem:fivelost}. This family gives examples, not a universal lower bound of $t+1$.
\end{example}

\begin{example}
For the five-cycle with consecutive vertices $1,\ldots,5$, assign edges to the cover $\{1,3,5\}$. The expression
\[
 x_1(x_2+x_5)+x_3(x_2+x_4)+x_4x_5
\]
has two inner additions, three multiplications and two outer additions, hence seven gates. Theorem~\ref{thm:main} proves this arithmetic count is optimal. It does not infer the unrestricted Boolean lower bound from this calculation.
\end{example}

\subsection{Three missing witnesses and cover exchanges}

For the rest of this subsection let $J$ denote the graph consisting of the missing edges $E\setminus E_2$. Recall that an \emph{old} linear form is one whose nonsingleton support has a common neighbor in $H$. We first complete the one-product part of the four-gate classification. The argument separates statements about $J$ from exchanges in a minimum cover of $H$.

\begin{lemma}[One deleted product]\label{lem:oneproductcover}
Suppose $H$ is nonempty, $J$ has three edges, and minimum-degree extraction and pruning deletes exactly four gates, precisely one of them a product. Then $\vc(J)\le2$. If the retained quadratic circuit has minimum size $|E_2|+\vc(H)-1$, then $\vc(G)\le\vc(H)+1$.
\end{lemma}
\begin{proof}
We use the frame-preserving reduction from the proof of Lemma~\ref{lem:fivelost}. It applies here because Lemma~\ref{lem:fourlost} already requires four deletions when three witnesses are missing. A permitted simplification preserving the entire retained frame and removing a deleted gate would contradict that earlier lemma. We therefore have the three one-product forms classified there, according to whether one, two or three additions are insertion cuts.

We record two support observations. If an old support $S$ contains $y$ and $xy$ is missing, then
\[
 S\cap N_G(x)\subseteq\{y\}.
\]
Indeed, two such neighbors together with a common old neighbor of $S$ would form a four-cycle; that old neighbor cannot be $x$, since $xy$ is missing. In particular, an old support meeting two incident missing edges $xy,xz$ must contain $x$. Also, in a valid old cubic, a missing witness $x^2y$ forces the factor supplying only $y$ to be the singleton $\{y\}$: varying that factor makes every one of its variables adjacent to $x$, and the preceding observation applies.

When the retained quadratic circuit is optimal, fix the assigned minimum cover $Z$ given by Lemma~\ref{lem:quadoptimal}; its assigned stars partition $E_2$. Let $D$ be a nonzero direct linear context of an insertion. If $z\in\supp(D)\setminus Z$, its contribution cannot come from an insertion into a star's leaf factor, since that contributes the opposite center in $Z$. It comes from an insertion into a center factor, whose extracted value is a literal input $h\in Z$. The full assigned leaf set $L_h$ is then contained in $\supp(D)$. This conclusion allows fanout: further affected stars add their centers to $D$. No cut enters both factors of one retained product, since its old low part would then produce a square.

It follows that a context of at most two vertices can be included in a cover of $H$ of size at most $|Z|+1$: if necessary replace $h$ by $\supp(D)$. Furthermore, all singleton contexts can simultaneously be included in a minimum cover of $H$. For each singleton $\{z\}$ outside $Z$, the associated star has $L_h=\{z\}$; replace $h$ by $z$. Distinct targets have distinct centers, and the targets lie outside the original cover. Every assigned star of a removed center is covered by its target, and every other assigned star keeps its center. These are cover exchanges, not circuit rewrites.

\emph{One cut.} The perturbation is $FGD$, where $F,G$ are arbitrary positive linear forms and $D$ is old. Every missing edge meets $\supp(D)$. If this support has at most two vertices, it covers $J$, and the preceding exchange proves the optimal-frame conclusion. We show that if $S=\supp(D)$ has at least three vertices, the missing edges form a star.

Put $I=S\cap\supp(F)$ and $K=S\cap\supp(G)$. A missing cubic witness repeated in $F,G$ would force $D$ to be singleton. Every remaining missing witness repeats an endpoint in $F,D$ or in $G,D$. If a witness of the first kind is $ux$, then every element of $I$ is adjacent to $x$, by the monomials $i^2x$. The old-support observation gives $I=\{u\}$. All missing edges of that kind therefore have center $u$. The other kind similarly has a common center $v$ and $K=\{v\}$. If both kinds occur, validity excludes $u=v$ and gives $uv\in E$.

Suppose their union is not a star. Choose missing edges $ux,vy$ with $x\ne v$ and $y\ne u$. We have $x,y\notin S$, since $K=\{v\}$ and $I=\{u\}$, and $x\ne y$ by the old-support observation applied to the two neighbors $u,v\in S$. Choose $z\in S\setminus\{u,v\}$. The missing edges $ux,vy$ force $xz,yz\notin E$. The monomial $yxz$ occurs in $FGD$, so validity forces $xy\in E$. Now $u,x,y,v$ form a four-cycle. Thus $J$ is a star; adding its center to $Z$ proves the required cover bound.

\emph{Three cuts.} Here $g=UV$ with old $U,V$, and the perturbation is
\[
 g(D_1+D_2+D_3)+c g^2,\qquad c\ge0.
\]
If $c=0$, Lemma~\ref{lem:commoncubic} makes $J$ a star. If $c>0$, the square of every monomial of $g$ is valid, so $g$ itself is a valid quadratic product. Its two disjoint supports span a complete bipartite subgraph, which is a star by four-cycle-freeness. Every exact two-variable witness in a multiple of $g$ is already an edge of that star. Both conclusions follow by adding its center. The same reasoning handles every positive square cross term in the two-cut case.

\emph{Two cuts with zero cross term.} The private product is $(A+B)W$ or $(A+B)^2$, with all displayed forms old, and the perturbation is $g(D_1+D_2)$. For a square $S^2$, validity forces disjoint complete bipartite adjacency between $\supp(S)$ and the union of the contexts. Four-cycle-freeness makes that graph a star, proving both conclusions.

In the remaining case consider the four valid old cubics $AWD_i,BWD_i$. Each has at most one missing edge. If their union consists of three distinct edges, three cells in an L-shaped part of this two-by-two array have distinct labels: if all four cells are occupied, place the repeated label in a row, column or diagonal and choose the other three cells. Relabel so these cells are $AWD_1$, $AWD_2$, $BWD_1$. Lemma~\ref{lem:commoncubic} makes the first edge incident with each of the others. If all three have a common endpoint, $J$ is a star. Otherwise write their edges as
\[
 uv,\qquad ux,\qquad vy,
 \qquad x\ne v,\quad y\ne u.
\]
The old supports $A,W$ meet the first two edges, so contain $u$. The old supports $D_1,W$ meet the first and third, so contain $v$. The singleton-factor observation applied to $u^2v$ gives $D_1=\{v\}$, and applied to $u^2x$ in the second cubic gives $D_2=\{x\}$. Thus the two singleton contexts cover $J$. In the optimal frame they can simultaneously be included in a minimum cover of $H$, so in this case $\vc(G)\le\vc(H)$. A zero context leaves at most two missing edges and cannot occur under the hypotheses. All one-product forms have been covered.
\end{proof}

\begin{proposition}[Four deletions with three missing edges]\label{prop:threefour}
Suppose $H$ is nonempty, exactly three edges are missing, and exactly four gates disappear under minimum-degree extraction and pruning. Then $\vc(J)\le2$. If the retained quadratic circuit is optimal, then $\vc(G)\le\vc(H)+1$.
\end{proposition}
\begin{proof}
The frame-preserving essential classification from Lemma~\ref{lem:fivelost} applies by Lemma~\ref{lem:fourlost}, as in the preceding proof. We inspect every distribution of the four deleted gates. All missing supports considered below are exact two-variable supports of valid positive summands; Boolean support simplifications do not assert equality of formal coefficients.

Three products and one cut give a product of at most five distinct old factors, including its scalar or linear context. Proposition~\ref{prop:oldfactors} would require at least eight distinct factors for three missing edges, so this case is impossible. Two products, one cut and a private addition give either two old cubics or two quartics sharing three factors, with repeated private factors reduced Booleanly as in the classification. Lemma~\ref{lem:sharedquartics} bounds their missing union by two. One product and three additions is precisely Lemma~\ref{lem:oneproductcover}.

It remains to treat two products and two cuts. If both cuts insert the second product $f$, then $f=ABC$ or $(AB)^2$ on old forms, and the perturbation is $fD_1+fD_2+c f^2$. A positive $c$ makes $f$ itself valid, with at most one missing edge, and every missing edge of a multiple of $f$ is already an edge of $f$. For $c=0$, Lemma~\ref{lem:sharedquartics} bounds two quartics, or Boolean reduction gives two cubics. A scalar context makes $f$ itself valid and only reduces the possibilities. Hence this configuration cannot have three missing edges.

Suppose one cut inserts each product. Write $g=UV$ for the first product in product order, keeping the separate cut-order convention of the classification. The independent-symbol frame is
\[
 P_2+y_1D_1+y_2D_2+c y_1y_2.
\]
If the second product $f$ is independent, $c=0$ gives two old cubics, while $c>0$ gives the mandatory old cross contexts of Lemma~\ref{lem:independentproducts}. Both give at most two missing edges.

If $f=(A+\alpha g)W$ has one affected input and $c=0$, a nonzero $D_1$ is excluded by Corollary~\ref{cor:affectedcase}. If $D_1=0$, the perturbation is $AWD_2+UVWD_2$. Three missing edges force the quartic to supply two. Equality in Proposition~\ref{prop:oldfactors} makes $K=\supp(D_2)$ a two-element transversal of that matching. Every missing edge meets $K$, proving $\vc(J)\le2$, and Corollary~\ref{cor:zeroaffected} proves the stronger cover exchange when the frame is optimal. A zero or scalar second context cannot supply three missing edges. For $c>0$, the perturbation reduces Booleanly to $UVD_1+AWD_2+UVW$. If $D_1\ne0$, the first cut's old low part $a$ is supported in $A\cap D_2$ and $aD_1$ is old; Lemma~\ref{lem:relatedcubics} bounds the missing union by two. If $D_1=0$, only two cubics remain.

If both inputs are affected, $f=(A+\alpha g)(B+\beta g)$ is Boolean-equivalent to $AB+g$. With $c=0$, the missing supports come from $UVD_1+UVD_2+ABD_2$. The first old low part $a$ is supported in $A\cap B$ and $aD_1$ is old. Lemma~\ref{lem:relatedcubics} gives at most two missing edges when both contexts are nonzero; a zero context leaves at most two cubics. With $c>0$, the expression reduces to $g+ABD_2$, again with at most two missing edges.

If $f=gW$ reads the first product directly, positive $c$ leaves two old cubics. For $c=0$, the perturbation is $UVD_1+UVWD_2$. A scalar context gives at most two cubics, so the three-edge case has a linear second context and a quartic with two missing edges. Its four factors are two-element transversals of a matching by Proposition~\ref{prop:oldfactors}. Any third missing edge in the cubic meets $\supp(U)$, hence meets an endpoint of that matching. This endpoint and one endpoint of the other matching edge cover all three missing edges. Thus $\vc(J)\le2$; when the frame is optimal, Corollary~\ref{cor:directproduct} gives $\vc(G)\le\vc(H)+1$.

Finally, $f=g(B+\beta g)$ or $g^2$ is Boolean-equivalent to $g$. The perturbation reduces to $g(D_1+D_2)+c g^2$. Two linear contexts give at most two old cubics. A scalar context or positive square term makes $g$ itself valid and leaves at most one missing edge. These cases are impossible. Every essential configuration has now been excluded or has the two asserted cover bounds.
\end{proof}

\begin{lemma}[At most three graph edges]\label{lem:smalledge}
For every nonempty finite simple graph with at most three edges,
\[
 \cmon(f_G)=|E|+\vc(G)-1.
\]
There is no restriction on the degrees of its formal witnesses.
\end{lemma}
\begin{proof}
Delete isolated variables by setting them to zero and simplifying constants, without increasing circuit size. Every remaining vertex is essential: set one neighbor to one and all other variables to zero. If a trimmed binary circuit has $k$ gates, its $2k$ operand slots include at least $k-1$ uses of nonoutput gates. Thus it has at most $k+1$ variable slots and at most $k+1$ distinct essential inputs. For $n$ essential variables, $k\ge n-1$; equality makes every nonoutput gate occur once and every variable occupy one slot, so the circuit is a read-once formula.

The input count proves the required bounds for one edge, two incident edges, two disjoint edges, a three-edge star, a two-edge path plus a disjoint edge, and three disjoint edges. The only remaining graphs are the four-vertex path and the triangle, each with target count four.

A three-gate circuit for the four-vertex path would be a read-once formula. At an OR root, the two children use disjoint variable sets and every edge lies within one set, contrary to connectivity. At an AND root, every output minterm is the disjoint union of one minterm from each child. Since all output minterms have size two, both children have only singleton minterms and are disjunctions of their variables. The graph is then complete bipartite, whereas the four-vertex path is not. It therefore needs four gates.

For the triangle, write $f=xy\vee xz\vee yz$. This function has no singleton true assignment and is not bounded above by any input variable. If a circuit of at most three gates computed it, neither operand at the output could be an input variable. Duplicate gate operands can be removed, so its two children must be distinct gates $u,v$.

If these gates are independent, each is a single binary operation on variables. At an OR output both must imply $f$, so each is an AND of distinct variables; their union covers at most two of the three edge assignments. At an AND output each must be implied by $f$, so each is an OR of two distinct variables. Their supports intersect among the three variables, giving a forbidden singleton true assignment. Both alternatives fail.

If $v$ depends on $u$, it has the form $u\wedge x$ or $u\vee x$ for an input $x$, or reduces to $u$. At an OR output, absorption gives $u$ or $u\vee x$; the latter has a singleton true assignment and the former is a one-gate function. At an AND output it gives $u$ or $u\wedge x$; the latter misses the opposite edge assignment. No case is the triangle function. Thus four gates are necessary. The star construction attains every stated lower bound.
\end{proof}

\begin{proof}[Preliminary reduction for Theorem~\ref{thm:threemissing}]
Assume exactly three witnesses are missing; Theorem~\ref{thm:twomissing} handles the smaller cases. Put $H=(V,E_2)$ and $d=\vc(G)-\vc(H)$. If $E_2=\varnothing$, the graph has exactly three edges. Theorem~\ref{thm:main} translates its arithmetic circuit to a Boolean circuit of the same size, and Lemma~\ref{lem:smalledge} applies. Otherwise put $q=|E_2|+\vc(H)-1$, let $D$ be the number of deleted gates, and let $e\ge0$ be the retained quadratic circuit's excess over $q$. Then
\[
 |C|=q+e+D,\qquad D\ge4
\]
by Lemma~\ref{lem:fourlost}. If $d\ge2$, Proposition~\ref{prop:threefour} excludes $e=0,D=4$. Hence $e+D\ge5$. Since the star count is $q+3+d$, this gives the full count for $d=2$ and one below it for $d=3$. For $d=0,1$, the elementary $e+D\ge4$ already gives the full count.

If $d=3$ and $|C|$ is below the star count, integrality forces $e+D=5$. The only possibilities are $(e,D)=(1,4)$ or $(0,5)$. In the first case Proposition~\ref{prop:threefour} gives $\vc(J)\le2$, and adjoining a cover of $J$ to one of $H$ gives $d\le2$, a contradiction. Thus a possible saving has exactly five deleted gates and an optimal retained quadratic circuit. We exclude this remaining configuration below.
\end{proof}

\subsection{A two-product obstruction on the remaining matching}

The following algebraic lemma has no old-neighborhood hypothesis. We then prove a circuit normalization that applies when at most two deleted gates are products. We then exclude the other product distributions in the five-deletion configuration from the preliminary reduction.

\begin{lemma}[Two-product matching template]\label{lem:twoproductmatching}
On six variables let $\mathcal M=x_1y_1\vee x_2y_2\vee x_3y_3$. Let $U,V,W,X,A,B$ be nonnegative homogeneous linear forms, possibly zero, and put $g=UV$. Let $f$ have one of the forms
\[
 WX,\quad (A+ag)W,\quad (A+ag)(B+bg),\quad
 gW,\quad g(B+bg),\quad g^2,
 \qquad a,b>0.
\]
For nonnegative affine linear forms $L,M$ and constants $\alpha,\beta,\gamma\ge0$, the polynomial
\[
 F=gL+fM+\alpha g^2+\beta gf+\gamma f^2
\]
does not have Boolean function $\mathcal M$.
\end{lemma}
\begin{proof}
All calculations in this proof are Boolean support identities. A positive constant makes an affine clause identically one; a zero factor removes its summand. Call a function valid if it vanishes on every independent set of the three-edge matching. If a positive sum computes $\mathcal M$, every summand is valid and each matching edge witnesses some summand.

We first give three elementary support facts. A valid product of $k$ nonzero linear clauses, allowing constant-one factors, with $r$ different matching-edge witnesses requires $k\ge2^r$. Restrict to those edges. Every clause meets each edge, while each of the $2^r$ independent transversals must make some clause false. A clause can reject at most one such transversal, and a constant clause rejects none. If $k=4,r=2$, equality forces the four distinct two-element transversals as the four factor supports on those endpoints. There are no outside variables: setting an outside variable of one factor to one leaves the other three unable to reject all four transversals, producing an independent satisfying assignment. Thus a product of at most three affine clauses has at most one edge witness, and a four-clause product with two witnesses is confined to their four endpoints.

Second, suppose a valid cubic $H_1H_2K$ witnesses an edge $ab$, and $H_1,H_2$ are both true on another matching edge $xy$. Then, after interchanging endpoints,
\[
 K=a,\qquad H_1=b\vee x,\qquad H_2=b\vee y.
\]
Indeed, if $K=1$, the two-factor product has two witnesses, contradicting the preceding count. If $H_1=1$, validity with witness $ab$ makes $H_2,K$ the opposite singleton endpoints of that edge, contradicting the $xy$ hypothesis; the other case is symmetric. Hence all factors are nonconstant. The cubic cannot witness $xy$, so $K$ misses that pair. It cannot contain both $a,b$: the other two factors could not reject all four transversals of the two pairs. Write $K\cap\{a,b\}=\{a\}$. On the two transversals containing $a$, validity forces both $H_1,H_2$ to miss $a$, contain $b$, and select opposite singleton endpoints of $xy$. An outside variable $z$ in $K$ makes the independent set $\{b,z\}$ true; one in $H_1$ makes $\{a,y,z\}$ true, and similarly for $H_2$. Thus the stated supports are exact.

Third, for affine clauses, if
\[
 UVK,\qquad WZJ,\qquad UVWZ
\]
are valid, their union has at most two different matching-edge witnesses. The cubics each have at most one. If the quartic has two, its four factors are confined to those two edges, so neither cubic can witness the outside third edge. If the quartic has none, the assertion is immediate. In the only remaining putative case it witnesses the third edge, while the cubics witness the first and second. Applying the second fact to each cubic gives
\[
 UV=b_1\vee x_3y_3,\qquad WZ=b_2\vee x_3y_3,
\]
where $b_i$ is an endpoint of edge $i$. Their product is true on the independent pair $\{b_1,b_2\}$, a contradiction. The same bound holds when the quartic summand is absent.

We now reduce the displayed expression. Boolean idempotence absorbs $\alpha g^2$ into $gL$ by replacing $L$ with $L\vee[\alpha>0]$, and similarly replaces $M$ with $M\vee[\gamma>0]$. These are still affine clauses. Put $\delta=[\beta>0]$. The six choices of $f$ become, respectively,
\[
 WX,\qquad AW\vee gW,\qquad AB\vee g,\qquad gW,\qquad g,\qquad g.
\]
For $f=WX$, the result is $UVL\vee WXM\vee\delta UVWX$, covered by the third support fact. For $f=AW\vee gW$, distributivity and absorption give
\[
 UV(L\vee\delta W)\vee AWM\vee UVWM.
\]
Apply that fact with pair factors $U,V$ and $W,M$, and outer clauses $L\vee\delta W$ and $A$. For $f=AB\vee g$, the result is
\[
 UV(L\vee M\vee\delta)\vee ABM,
\]
a union of two valid products with at most three clauses each, so at most two witnesses.

For $f=gW$, the result is $UV(L\vee\delta W)\vee UVWM$. Three witnesses would force two in the quartic and an outside third in the cubic. The quartic equality case confines $U,V$ to the first two edges, making that third witness impossible. Finally $f=g$ gives $UV(L\vee M\vee\delta)$, with at most one witness. Zero forms only delete terms, and constant clauses were covered throughout. No allowed case computes the three-edge matching.
\end{proof}

\begin{lemma}[Normalization with two deleted products]\label{lem:twoproductnormal}
Let a positive variable-input circuit have minimum output degree two and squarefree quadratic part $P_2$. Perform minimum-degree extraction and pruning, and suppose that at most two deleted gates are products. By deleting or bypassing Boolean-absorbed branches outside the retained quadratic circuit, one can preserve that entire circuit and obtain a Boolean-equivalent circuit whose polynomial is
\[
 P_2+gL+fM+\alpha g^2+\beta gf+\gamma f^2.
\]
Here $g=UV$, the form of $f$ is one of those in Lemma~\ref{lem:twoproductmatching}, $L,M$ are nonnegative affine linear forms, and $\alpha,\beta,\gamma\ge0$. An absent private product may be represented by zero. There is no bound on the number of deleted additions, and no graph or optimality assumption.
\end{lemma}
\begin{proof}
Let $\mathcal R$ consist of the original nodes of minimum degree two that have an addition-only path to the output, including the output. Every node on such a path again has minimum degree two, so the path stays in $\mathcal R$. Each retained quadratic node belongs to $\mathcal R$: a selected path from it cannot contain a product, which would raise its degree. Restoring the bypassed additions leaves an addition-only path. The other nodes of $\mathcal R$ are precisely the restored additions on these quadratic paths whose higher-degree input extraction omitted. We call their full current values quadratic carriers, without identifying them with their quadratic parts.

In topological order, replace every operand occurrence from a node $Q\in\mathcal R$ to a node outside $\mathcal R$ by a formal zero. Keep all edges within $\mathcal R$. At each step the full current value of $Q$ still reaches the output on its unchanged addition-only path inside $\mathcal R$. Every contribution lost by the replacement is a multiple of that current value, and is therefore Boolean-absorbed by it. Earlier replacements may have changed higher terms of $Q$; the argument uses the whole value at the current step and remains valid.

No degree-two output term is lost. A removed route leaves $\mathcal R$ and must encounter a product before reaching the output; otherwise its immediate successor would also have an addition-only output path and would belong to $\mathcal R$. That product raises the degree of the contribution from at least two to at least three. All paths of the extracted circuit remain intact: retained products have degree-one inputs, and retained quadratic additions use edges within $\mathcal R$. Propagate the formal zeros through private products and bypass a private sum when just one input remains. These operations introduce no gate or final constant and remove no retained gate, since each still has its unchanged extracted path. Thus both $P_2$ and the entire retained quadratic circuit are preserved.

There is now no edge from $\mathcal R$ to its complement. All remaining products outside $\mathcal R$ are deleted products. Order them as $g$ and $f$, omitting either if absent. Before $g$, every signal outside $\mathcal R$ is a homogeneous linear form, so $g=UV$. Before $f$, every such signal has the form $A+ag$ with $a\ge0$, by induction through the addition gates. Multiplying two such signals, and separating the possibilities that a linear part or a coefficient vanishes, gives exactly the six forms in Lemma~\ref{lem:twoproductmatching}, up to exchanging inputs and absorbing positive scalar multiples. After $f$, every outside signal has the form $A+ag+bf$.

The actual inputs to a retained product therefore have the form
\[
 A_i+a_i g+b_i f,\qquad B_i+c_i g+d_i f.
\]
Their homogeneous linear parts $A_i,B_i$ are their unchanged extracted inputs. Expanding their product gives
\begin{align*}
 A_iB_i&+g(a_iB_i+c_iA_i)+f(b_iB_i+d_iA_i)\\
 &+a_ic_i g^2+(a_id_i+b_ic_i)gf+b_id_i f^2.
\end{align*}
The output inside $\mathcal R$ is an addition network of these retained products and higher private signals. The latter are linear combinations of $g,f$; an original variable cannot reach the output through additions alone because the minimum output degree is two. Summing the expansions, with their nonnegative multiplicities, gives the claimed expression. The sum of the $A_iB_i$ terms is exactly $P_2$. A private product of degree two cannot enter the output additively without having been retained, so the direct private terms add no quadratic coefficient. Such direct occurrences explain the allowed constant terms of $L,M$. This is an identity for the polynomial of the modified circuit, whose Boolean output is unchanged.
\end{proof}

\begin{lemma}[The two one-cut profiles]\label{lem:onecutprofiles}
Suppose $G$ is four-cycle-free, $H$ is the nonempty quadratic support of a circuit for $f_G$, exactly three edges are missing, and $\vc(G)-\vc(H)=3$. Suppose extraction deletes exactly five gates and retains an optimal quadratic circuit. The deleted gates cannot be four products and one addition, or three products and two additions of which just one has unequal operand minimum degrees.
\end{lemma}
\begin{proof}
The three missing edges form a matching, and every minimum cover of $H$ avoids their six endpoints $T$: otherwise the missing edges could be covered by at most two additional vertices. Hence $H[T]$ is empty and $G[T]$ is precisely the three-edge matching. Restrict to $T$ by setting other variables to zero. The output must compute that matching, while $P_2$ vanishes.

There is no unused original gate. Deleting one would preserve the extracted circuit and leave at most four deletions, contradicting Lemma~\ref{lem:fourlost} or Proposition~\ref{prop:threefour}, since the missing-edge matching has cover number three. Write the unique unequal-degree addition as $t=a+F$, with $F$ its higher-degree input. Every other deleted gate lies in the ancestry of $F$. Indeed every path from a deleted useful gate to the output must cross a bypassed input edge, and the high-input edge of $t$ is the only such edge. All retained boundary signals of this private ancestry precede $t$. They are homogeneous, since no unequal-degree addition precedes them, and have degree one or two. Their quadratic monomials occur in $P_2$ along retained addition-only paths, so every quadratic boundary vanishes on $T$; the remaining boundaries are linear clauses, possibly zero.

The low degree of $t$ is one or two. If it is two, the output after the cut is $P_2+cF$ for a positive scalar $c$. If it is one, temporarily replace $t$ by a degree-one symbol $z$. The quadratic frame has form $Q+zD+cz^2$ with $D$ linear. Its $z^2$ coefficient vanishes: substituting the nonzero old linear input $a$ would otherwise give a square in $P_2$. Thus the output is $P_2+FD$. Fanout only adds to $D$. After restriction, the high computation times a scalar or a single linear context must therefore compute the matching.

With four products and no private addition, that computation is a multiplication DAG with at most five distinct linear boundary inputs. For a connected DAG of $p$ binary products and $b$ distinct boundary nodes, its $2p$ operand edges give $p+b-1\le2p$, hence $b\le p+1$. Boolean idempotence removes repeated factors. Including the context leaves at most six clauses, whereas the matching-clause bound in Lemma~\ref{lem:twoproductmatching} requires eight.

For the second profile, let $S=F_1+F_2$ be its sole private equal-degree addition. If a restricted input vanishes, bypass $S$ and use the product-only bound with at most three products. If $S$ is linear, regard it as another linear boundary input; again at most five clauses remain including the context. We may thus assume that both inputs are nonzero homogeneous signals of degree at least two. Each is rooted at a private product, because retained quadratic boundaries vanished.

Outside $S$, the private computation uses products alone. If one operand of $S$ is a multiplicative ancestor of the other, their sum is Booleanly absorbed by that ancestor. If either whole operand is also a multiplicative factor of the output outside $S$, it absorbs $S$ there. In either case the high computation becomes a product-only expression with at most three product gates, again too few clauses. Otherwise the two operand root products are distinct, neither depends multiplicatively on the other, and neither reaches the high output along a path avoiding $S$.

Resolve $S$ once by choosing $F_1$, and once by choosing $F_2$. Repeated occurrences of $S$ are removed by Boolean idempotence. In each resolution the root product of the unselected operand disappears: a surviving occurrence would give precisely the excluded path avoiding $S$. Each resolved summand therefore uses at most two of the three product gates and is a product of at most three linear clauses before its context, or four after it. Each is valid separately, since their positive union computes the matching.

If the context is scalar, the two summands witness at most one matching edge each. If it is linear and their union witnesses all three edges, one summand must witness two. It has exactly four clauses, so the equality case in Lemma~\ref{lem:twoproductmatching} confines all its clauses, including the common context, to those two edges. The other summand contains the same context and cannot witness the disjoint third edge. This final contradiction excludes both profiles.
\end{proof}

\begin{lemma}[Three products and two cuts]\label{lem:threeproductstwocuts}
Under the hypotheses of Lemma~\ref{lem:onecutprofiles}, the five deleted gates cannot consist of three products and two unequal-degree additions.
\end{lemma}
\begin{proof}
Fix the assigned minimum cover $Z$ of the optimal quadratic circuit, and put $T=V\setminus Z$. As before, $G[T]$ is the three-edge missing matching together with possibly isolated vertices. Set the variables in $Z$ to zero, but retain those isolated variables for the cover argument. Call a positive summand safe if it is false on every independent set of this restricted graph.

We first record a constraint on an actual direct insertion context $D$. A quadratic insertion has scalar context. A degree-one insertion contributes on $T$ only when it perturbs the literal center input $h$ of a star; insertions solely into leaf factors have opposite centers in $Z$ and vanish. Its context contains the whole set $N_H(h)\cap T$, since every edge from $h$ to $T$ is assigned to that star. Thus a nonzero linear context cannot be supported on a transversal of the three missing edges, with no other variables. Otherwise replace $h$ in $Z$ by $N=N_H(h)\cap T$, then add an endpoint of each missing edge not met by $N$. The resulting cover of $G$ has at most
\[
 |Z|-1+|N|+(3-|N|)=|Z|+2
\]
vertices. Here $N\ne\varnothing$, because $Z$ is minimum. This contradicts the cover gain three. Keeping all of $T$, including its isolated vertices, is necessary for this inference.

The matching facts in Lemma~\ref{lem:twoproductmatching} apply with those isolated variables present. In particular a safe product of at most three linear factors and an actual direct context witnesses at most one missing edge: two would require a quartic equality case and a forbidden transversal context. Cubic rigidity says that if $UVD$ witnesses one edge and $U,V$ are both true on another, then $D$ is a singleton on the first edge. When $D$ is an actual context this is forbidden. If instead the actual context occupies one of the other two positions, rigidity makes it a two-element transversal and again forbids it.

We also need a short quintic observation. A safe product of at most five affine clauses with two edge witnesses cannot share three factor occurrences with a product witnessing a third edge. Each shared factor is either already true from a constant term or meets that third edge. Hence one of its endpoints makes at least two of the three shared factors true. Set that endpoint to one. At most three factors remain to reject the four independent transversals of the first two edges, and each rejects at most one. A surviving transversal gives an independent satisfying assignment, a contradiction.

Any permitted frame-preserving deletion of a private gate would leave at most four deletions, impossible by Proposition~\ref{prop:threefour}. The full-carrier erasure in Lemma~\ref{lem:twoproductnormal} applies before its final two-product parametrization and does not depend on the number of private products. Hence no essential off-frame use of a full quadratic carrier remains. Every private product must feed a cut's higher input. Call a cut visible when its own gate belongs to the selection subgraph, obtained by following both inputs at products and equal-degree additions and only the low input at a cut. Reuse of the low-input node elsewhere does not suffice. The last cut is visible: any private signal after it would need a later cut to enter the frame.

Suppose only one cut is visible. Its high input $F$ is computed by the three products and the other cut, with old linear boundary inputs. This connected four-gate DAG has at most five distinct boundary nodes, by the operand-slot count used in Lemma~\ref{lem:onecutprofiles}. With only one addition, its Boolean function is
\[
 F=C_0(A_0\vee B_0),
\]
where the three displayed factors are products of linear clauses. Factor common clauses into $C_0$ and use idempotence to remove repetitions; altogether there are at most five distinct clauses. Its visible output context $D$ is scalar or linear, and the restricted output is $FD$. If $D$ is scalar, distributive conversion to conjunctive normal form has at most six clauses: the number is $ab+c$ with $a+b+c\le5$. If $D$ is linear, absorb it into $C_0$; now $a+b+c\le6$ and $c\ge1$, giving at most seven clauses. Zero arms or constant-one arms only reduce the bound. The matching requires eight, so this case is impossible. This argument includes a private nested cut with a quadratic low branch and makes no frame-changing degree reduction.

Both cuts are therefore visible. Their high inputs are private-product outputs: an old linear signal cannot be the higher operand, while a full quadratic carrier would be absorbed. In cut order call the high inputs $g,f$. Replacing them temporarily by independent insertion symbols in the retained frame gives the restricted output
\[
 F=gD_1+fD_2+c\,gf,\qquad c\ge0.
\]
The $D_i$ are the actual scalar or linear direct contexts. An individual squared insertion is impossible because its nonzero old low part would produce a square in $P_2$. Dependence of $f$ on the earlier cut is substituted only after this frame identity is formed. All Boolean summands below are safe. Zero and scalar contexts are allowed throughout.

Let $k$ be the number of private products ancestral to $g$. Before the first cut they form a multiplication-only network on old linear forms, so $g$ uses at most $k+1$ distinct clauses. Every other private product is ancestral to $f$.

If $k=3$, the signal $f$ is an ancestor of $g$ or is $g$ itself. Write $g=fA$ Booleanly. The output becomes $f(AD_1\vee D_2\vee cA)$. For $c>0$, the $AD_1$ term is absorbed; for $c=0$, conjunctive expansion adds one clause $D_2\vee a$ for each factor $a$ of $A$, and the clause $D_2\vee D_1$. Factors already present in $f$ absorb their repetitions. Since $g$ uses at most four distinct clauses, there are at most five in all, a contradiction.

Suppose $k=2$. Write $u=UV$ for the first product; Booleanly $g=u$ or $uW$, with $W$ linear. The last product $f$ can read old linear forms plus a multiple of $g$, or the direct earlier signals $u,g$. An uninserted $u$ cannot enter a retained linear sum without another private addition. If $f$ has a direct $u$ input, it is Booleanly $uA$, $uA\vee g$, $u$ or $g$. Factoring out $u$ gives, among their specializations, brackets
\[
 WD_1\vee AD_2\vee cWA,
 \qquad W(D_1\vee D_2\vee c)\vee AD_2.
\]
Each has a conjunctive normal form of at most four linear clauses; with $U,V$ there are at most six. The cases $g=u$ are obtained by taking $W=1$.

Otherwise the possible values of $f$ are $AB$, $AB\vee gB$, and $AB\vee g$, allowing zero old factors. For $f=AB$, put $P=gD_1$, $Q=ABD_2$ and $R=cgAB$. The actual-context rule gives at most one witness each to $P,Q$. If $R$ has two, an outside witness of $P$ contradicts the quintic observation using the three $g$ factors, or quartic confinement if $g$ has only two. An outside witness of $Q$ contradicts cubic rigidity, since $A,B$ are true on an $R$ edge and its context $D_2$ would be a forbidden singleton. If $R$ has just one witness, the same rigidity prevents $Q$ from having a different one. Thus the union cannot witness three edges; if $R=0$, two summands already suffice for this conclusion.

For $f=AB\vee gB$ and $c>0$, absorption gives $g(D_1\vee B)\vee ABD_2$. Three witnesses would require two in the first quartic. Its equality case confines $B$ to those two edges, excluding an outside third witness in the second term. For $c=0$, use $P=gD_1$, $Q=ABD_2$, $R=gBD_2$. The first two have at most one witness each. Two witnesses in $R$ exclude an outside $P$ witness by the same quintic or quartic argument. An outside $Q$ witness makes $B,D_2$ true on another edge, so cubic rigidity forces the actual context $D_2$ to be a forbidden two-element transversal. The same argument excludes a $Q$ witness different from a sole $R$ witness. Again the union has at most two.

For $f=AB\vee g$, the expression is $g(D_1\vee D_2\vee c)\vee ABD_2$. The second term is cubic. If the first has two witnesses, it is a quartic whose equality case confines $D_2$ to those edges, preventing an outside third witness in the second term. A positive constant in that context would leave only three factors and could not yield two witnesses. This completes $k=2$, including direct $g$ factors and zero old parts.

Finally let $k=1$, so $g=UV$. The remaining two-product network makes $f$ a product of at most three factors of the form $A_i+a_i g$, including pure $g$ factors. If none is affected, $f$ is a product of at most three old forms; interchange the two inserted signals in the preceding independent-product case, whose Boolean argument did not use cut order. Otherwise group unaffected factors into $H_0$ and old parts of affected factors into $K_0$. Boolean absorption gives
\[
 f=H_0K_0\vee gH_0,
\]
where $H_0$ has $h\le2$ factors and $H_0K_0$ at most three. Let $P=gD_1$ and $Q=H_0K_0D_2$; each has at most one witness by the context rule.

For $c=0$, the remaining term is $R=gH_0D_2$, with at most five factors. If it has two witnesses, an outside $P$ witness contradicts cubic rigidity. An outside $Q$ witness shares $H_0,D_2$ with $R$: for $h=2$ the quintic observation applies, for $h=1$ quartic confinement applies, and for $h=0$ the cubic $R$ cannot have two witnesses. If $R$ has one witness, cubic rigidity prevents $P$ from having a different one, leaving at most two in total. For $c>0$, the term $gH_0$ absorbs $gH_0D_2$, so take $R=gH_0$. Two witnesses require $h=2$ and a quartic equality case, which confines both $g$ and $H_0$ and prevents an outside witness of either $P$ or $Q$. With only one $R$ witness, cubic rigidity again prevents a different $P$ witness. Every case contradicts the required three witnesses, proving the lemma.
\end{proof}

\begin{example}[An induced missing matching does not suffice]\label{ex:missingmatching}
Five deletions and an optimal quadratic remainder can coexist with an induced missing matching of size three. Take four centers $h_{ac},h_{ad},h_{bc},h_{bd}$ with respective leaf pairs $\{a,c\},\{a,d\},\{b,c\},\{b,d\}$, and a disjoint center $z$ with leaves $x,y$. The old graph $H$ has polynomial
\[
 P_H=h_{ac}(a+c)+h_{ad}(a+d)+h_{bc}(b+c)+h_{bd}(b+d)+z(x+y).
\]
Add edges $ab,cd,xy$ to obtain $G$, and put
\[
 P=P_H+(a+c)(a+d)(b+c)(b+d)+xy(x+y).
\]
The quartic Boolean function is $ab\vee cd$: its first two clauses give $a\vee cd$ and its last two give $b\vee cd$. The last term gives $xy$. Thus $P\in\Sup(G)$, its quadratic part is exactly $P_H$, and the six missing endpoints induce a three-edge matching.

The first component of $H$ is an eight-cycle and has cover number four. Adding $ab,cd$ leaves that cover number unchanged, since all four endpoints still cover every edge. It creates no four-cycle: two centers in such a cycle would need the same leaf pair; one center would require consecutive edges in the endpoint graph, which is a matching. The other component changes from a two-edge star to a triangle, with cover numbers one and two. Therefore $H,G$ are four-cycle-free,
\[
 |E(H)|=10,\quad\vc(H)=5,\qquad |E(G)|=13,\quad\vc(G)=6.
\]
The five-star quadratic circuit has five leaf additions, five products and four output additions, totaling fourteen gates, the minimum for $H$. Reuse its leaf sums to form $(a+c)(a+d)(b+c)$ with two extra products and insert it into the center input $h_{bd}$. Form $xy$ with one product and insert it into $z$. These three products and two additions give a nineteen-gate circuit for $P$. Extraction deletes exactly those five gates and retains the fourteen-gate frame. The star count for $G$ is eighteen, so this circuit is above that bound. The example disproves a stronger exclusion based only on the induced missing matching; its cover gain is one, and it does not meet the gain-three hypothesis of the preceding lemma.
\end{example}

\begin{corollary}[Three missing matching edges require three deleted products]\label{cor:threedeletedproducts}
Let $G$ be any finite simple graph, and let a circuit computing $P\in\Sup(G)$ have nonzero quadratic support $H$. If six vertices induce a matching of three edges in $G$ and no edge in $H$, minimum-degree extraction and pruning delete at least three multiplication gates. This bound is sharp. No four-cycle, cover-number, optimality or deleted-addition hypothesis is required.
\end{corollary}
\begin{proof}
The output has minimum degree two, since $P\in\Sup(G)$ and its quadratic part is nonzero. If at most two products were deleted, Lemma~\ref{lem:twoproductnormal} would give its exact-frame-preserving template. Set variables outside the six vertices to zero. The quadratic part vanishes, while all template factors remain nonnegative linear forms, possibly zero, and its contexts remain affine. The Boolean output is precisely the induced three-edge matching, contradicting Lemma~\ref{lem:twoproductmatching}. Example~\ref{ex:missingmatching} attains three deleted products. Its complete circuit need not be optimal; the asserted sharpness concerns only the number of deleted multiplications.
\end{proof}

\begin{proof}[Completion of Theorem~\ref{thm:threemissing}]
Theorem~\ref{thm:twomissing} handles at most two missing witnesses. With exactly three, Lemma~\ref{lem:smalledge} handles empty quadratic support. The preliminary reduction settled all cover gains at most two and showed that any remaining saving has cover gain three, exactly five deleted gates and an optimal quadratic remainder. Its three missing edges form an induced matching on six vertices with no old edge among them.

Let $k$ be the number of deleted products. There must be an unequal-degree addition: if none existed, induction from the variables would make every useful gate homogeneous, and extraction would retain the whole circuit. Thus $k\le4$. If $k\le2$, Corollary~\ref{cor:threedeletedproducts} gives a contradiction. If $k=4$, the only addition is a cut and Lemma~\ref{lem:onecutprofiles} excludes it. If $k=3$, the two additions are either one cut and one private equal-degree addition, excluded by that same lemma, or two cuts, excluded by Lemma~\ref{lem:threeproductstwocuts}. These exhaust all five-gate distributions. No saving is possible. The star circuit has the asserted size and all its witnesses are quadratic, so it also proves attainability in the stated class.
\end{proof}

\subsection{Four missing witnesses: boundary cases and remaining gate counts}

\begin{lemma}[Gate loss under zero restriction]\label{lem:zerorestriction}
Trim a constant-free binary monotone Boolean circuit to the ancestors of its output. Set $k$ distinct occurring variables to zero, and suppose that the restricted output is nonzero. Zero propagation, bypassing disjunctions with one zero input, and pruning remove at least $k$ gates.
\end{lemma}
\begin{proof}
Merge input vertices carrying the same variable, retaining parallel wires. Let $X$ be the $k$ zeroed inputs, $Z$ the gates whose restricted functions are zero, and $B$ the nonzero disjunction gates with a zero input. Zero propagation detects every identically zero restricted subcircuit: a monotone function is nonzero exactly when it is true on the all-one assignment of its remaining variables. Each gate in $B$ has exactly one zero input.

Form an undirected multigraph on $X\cup Z\cup B$, keeping the original wires from $X\cup Z$ to $Z\cup B$. Every vertex in $X\cup Z$ has a directed route to $B$: follow an original output path to its first nonzero gate. That gate cannot be a conjunction and is a member of $B$. Thus every connected component meets $B$. If $z=|Z|$, $b=|B|$ and $c$ is the component count, then $b\ge c$. The graph has $k+z+b$ vertices, at least $k+z+b-c$ edges, and at most $2z+b$ edges, by counting incoming wires. Hence $z+b\ge k+b-c\ge k$. All $Z$ gates are removed and all $B$ gates bypassed. Further pruning only increases the loss. Repeated operands and fanout are included by the multigraph count.
\end{proof}

\begin{proposition}[The exact count through four edges]\label{prop:fouredgeboundary}
For every nonempty four-cycle-free simple graph with at most four edges, the monotone Boolean edge-circuit complexity is $|E|+\vc(G)-1$. Consequently every positive variable-input circuit computing any $P\in\Sup(G)$ has at least that many gates, including when its quadratic part is zero.
\end{proposition}
\begin{proof}
The cases with at most three edges are Lemma~\ref{lem:smalledge}. We give the full four-edge classification. First, a circuit depending on $n$ variables and having $s$ ancestor gates has at least $n+s$ vertices and exactly $2s$ operand wires in its connected ancestor multigraph. Thus $s\ge n-1$. Equality forces exactly those $n$ input vertices and makes the graph a tree, so the circuit is a read-once formula.

A connected edge graph represented by such a formula must be complete bipartite. At an OR root, the children use disjoint variable sets and no edge crosses them, contradicting connectedness. At an AND root, every minimal true set is a union of a nonempty minimal true set from each child. All these unions have size two, so both children have only singleton minimal true sets and are disjunctions of their variables. Every cross pair is then an edge. A four-cycle-free complete bipartite graph is a star. Therefore a connected four-cycle-free nonstar needs at least $n$ gates.

If a four-edge graph has a cycle, it is a triangle with either a pendant edge or a disjoint edge. In the pendant case set its outside vertex to zero. The remaining triangle needs four gates and Lemma~\ref{lem:zerorestriction} charges one more, giving five, equal to $4+2-1$. For a disjoint extra edge, zero its two endpoints; the same lemma gives six, equal to $4+3-1$.

Otherwise the graph is a forest. If connected, it is a tree on five vertices. A star has cover number one and the input count gives four gates. A nonstar has cover number two: a smaller bipartition class has at most two vertices, while a one-vertex cover would make it a star. The read-once exclusion gives five gates. For disconnected forests, partition the four edges among components. In the partition $3+1$, the three-edge tree is a star or a four-vertex path. The former has six relevant inputs and total cover number two, giving five gates. For the latter, zero the separate edge's endpoints. The path needs four gates and restriction charges two, giving six, its star count. The partitions $2+2$, $2+1+1$, and $1+1+1+1$ have respectively $(n,\vc)=(6,2),(7,3),(8,4)$; the input lower bound gives their respective star counts $5,6,7$. These exhaust all cases.

The assigned-star construction attains each bound. Booleanization preserves the gate count of a positive variable-input arithmetic circuit and computes $f_G$ exactly when its polynomial lies in $\Sup(G)$. The lower bound therefore applies to every such polynomial; the attainment assertion concerns the Boolean function or the ordinary all-one edge polynomial, not arbitrary prescribed coefficients.
\end{proof}

\begin{lemma}[Four induced missing edges require six deletions]\label{lem:fourmatchinglost}
Let $G$ be any finite simple graph, and let a positive variable-input binary circuit compute $P\in\Sup(G)$ with nonzero quadratic support $H$. Suppose eight vertices induce exactly four disjoint edges in $G$ and no edge in $H$. Minimum-degree extraction and pruning delete at least six gates. No four-cycle, cover-number, quadratic-optimality or deleted-addition hypothesis is required.
\end{lemma}
\begin{proof}
Suppose at most five gates are deleted. The output has minimum degree two and squarefree quadratic part. Apply the full-carrier erasure at the beginning of Lemma~\ref{lem:twoproductnormal}: protect every minimum-degree-two node with an addition-only output path and erase its off-carrier uses. At each step the whole current carrier still reaches the output additively and absorbs its lost multiples. This argument precedes that lemma's two-product parametrization and is independent of the number of deleted products. It preserves the exact retained quadratic circuit, introduces no gate, and leaves the Boolean function unchanged. Propagate zero branches and trim. The remaining circuit still has at most five deletions and still belongs to $\Sup(G)$.

By Corollary~\ref{cor:threedeletedproducts}, applied to any three of the four matching edges, at least three deleted products remain. At least one deleted addition has unequal operand minimum degrees; call such an addition a cut. Without a cut all useful gates would be homogeneous, and a nonzero quadratic output would have no missing quadratic witness. The only possible profiles are therefore four products and one cut; three products and one cut; or three products and two additions, at least one a cut.

Define the selection subgraph by keeping both input edges at products and equal-degree additions and only the lower-degree edge at a cut, starting from the output. A cut is \emph{visible} when its own gate lies on this selected subgraph; reuse of its low-input node elsewhere does not suffice. Every private product feeds some cut's higher input. The last cut is visible, because a private signal after it would need another cut to enter the quadratic frame. Off-frame uses of full quadratic carriers have already been erased.

The direct context of a visible insertion is scalar if its low degree is two, and linear if its low degree is one. An individual insertion cannot enter both factors of a retained product: its nonzero old linear low part would then contribute a square to $P_2$. With two visible cuts, give their higher inputs independent symbols $g,f$, in cut order. Before substituting any dependence of $f$ on the earlier cut, the exact frame polynomial is
\[
 P_2+gL+fM+c\,gf,\qquad c\ge0,
\]
with nonnegative scalar or linear direct contexts $L,M$. Allowing affine contexts only enlarges the class considered below. Fanout changes coefficients and does not alter this form.

Set all variables outside the eight specified vertices to zero. This is a semantic restriction, not a claimed frame-preserving rewrite. The quadratic part vanishes and the Boolean output is the four-edge matching. Any conjunction of positive affine-linear clauses computing that matching requires at least sixteen clauses. Indeed each clause must be true on all four edge assignments. A constant-true clause rejects none of the sixteen independent endpoint transversals; every other clause meets each edge and rejects at most one such transversal. All sixteen transversals must be rejected. We shall produce fewer than sixteen clauses in each remaining profile. All subsequent idempotent simplifications are Boolean identities for this counting argument.

If there is one visible cut and no other private addition, its high input is a multiplication DAG with at most four products. Such a connected $p$-product DAG has at most $p+1$ distinct boundary nodes, because its $2p$ operand slots include at least $p-1$ internal uses. After carrier erasure its boundaries are old linear forms. Including the visible context gives at most six clauses.

With three products and two private additions but only one visible cut, the other addition lies in the high-input computation. This includes a hidden unequal-degree cut and a private equal-degree sum. A hidden cut cannot feed a retained linear node on a selected output path, since it would then be visible. Thus the high computation has old linear boundaries and three products plus one addition. It has at most five distinct boundary forms. With one addition its Boolean value is $C_0(A_0\vee B_0)$, where the three displayed terms are products of these forms. Remove repetitions and move shared factors into $C_0$. If their factor counts are $c_0,a_0,b_0$, their total is at most five. A scalar context yields at most $a_0b_0+c_0\le6$ clauses. A linear context adds one common factor, so the total is at most six with common-factor count at least one; the clause bound is then seven. Zero or constant-one arms only simplify the expression. Both bounds contradict sixteen.

It remains to consider three products and two visible cuts. Both high inputs are private-product outputs. An old linear signal cannot be a higher operand, and a full quadratic carrier cannot have a private use after erasure. A preceding visible cut either has degree one, too low to be a later high input, or is a protected quadratic carrier.

Let $k$ be the number of private products ancestral to the first high signal $g$. Before that first cut they form a multiplication-only network on old linear forms. Every other private product feeds $f$. Thus $k=1,2,3$.

\emph{Case $k=3$.} Then $f$ is $g$ or one of its product ancestors. Write $g=fA$ Booleanly, with at most four distinct linear factors in their union. The output is
\[
 f(AL\vee M\vee cA).
\]
If $c>0$, the term $AL$ is absorbed. If $c=0$, conjunctive expansion adds $M\vee a$ for each factor $a$ of $A$, and $M\vee L$. Repeated factors already in $f$ are absorbed. There are at most five clauses.

\emph{Case $k=2$.} Write $u=UV$ for the first product. Booleanly $g=u$ or $uW$; in particular $g$ has at most $b=3$ factors. The last product reads old linear forms plus multiples of $g$, or direct earlier signals $u,g$. An uninserted $u$ cannot enter a retained linear sum without an additional cut or private addition. If the last product has a direct $u$ factor, it is $uA$, $uA\vee g$, $u$, or $g$. Factoring out $u$ leaves a specialization of
\[
 WL\vee AM\vee cWA,
 \qquad W(L\vee M\vee c)\vee AM.
\]
The first bracket has at most four clauses: for $c=0$ use
$(W\vee A)(W\vee M)(L\vee A)(L\vee M)$;
for $c>0$ use $(W\vee A)(W\vee M)(A\vee L)$.
The second bracket is an OR of two products with two clauses each and also has at most four. Including $U,V$ gives at most six. Taking $W=1$ covers $g=u$.

Otherwise $f$ is $AB$, $AB\vee gB$, or $AB\vee g$, including zero old parts and direct $g$ factors. For $f=AB$ and $c=0$, pairwise distribution of $gL\vee ABM$ gives $3(b+1)\le12$ clauses. For $c>0$ its CNF is
\[
 (g\vee AB)(g\vee M)(AB\vee L),
\]
with $3b+2\le11$ linear clauses after distribution. For $f=AB\vee gB$ and $c=0$, use
\[
 (g\vee A)(g\vee B)(g\vee M)(L\vee B)(L\vee M),
\]
with $3b+2\le11$ clauses; for $c>0$ omit the final factor, leaving at most ten. For $f=AB\vee g$ and $c=0$, use
\[
 (g\vee A)(g\vee B)(g\vee M)(L\vee M),
\]
with $3b+1\le10$ clauses; for $c>0$ omit the final factor, leaving at most nine. Each identity follows by setting $g$ to zero and one. A factor $g\vee X$ expands over the at most $b$ factors of $g$; in $g\vee AB$ there are $2b$ such clauses. Every bound is below sixteen.

\emph{Case $k=1$.} Now $g=UV$, with at most $b=2$ factors. The remaining two-product network makes $f$ a product of at most $t=3$ factors, each an old linear form $A_i$ or $A_i\vee g$; a direct $g$ factor is the case $A_i=0$. If none is affected, write $f=K$. For $c=0$, distributing $gL\vee KM$ gives at most $(b+1)(t+1)\le12$ clauses. For $c>0$, use
\[
 (g\vee K)(g\vee M)(K\vee L),
\]
with at most $bt+b+t\le11$ clauses.

If a factor is affected, group the unaffected factors into $H_0$ and the old parts of the affected factors into $K_0$. Then $H_0$ has $h\le2$ factors, $H_0K_0$ has at most three, and
\[
 f=H_0K_0\vee gH_0.
\]
For $c=0$ the output identity is
\[
 gL\vee H_0K_0M\vee gH_0M
   =(g\vee H_0K_0M)(L\vee H_0M).
\]
It uses at most $b(t+1)+(h+1)\le11$ linear clauses. For $c>0$, the cross term absorbs $gH_0M$ and the identity becomes
\[
 gL\vee H_0K_0M\vee gH_0
   =(g\vee H_0K_0M)(L\vee H_0),
\]
using at most $b(t+1)+h\le10$ clauses. These identities also follow by evaluating $g=0,1$. Empty products, zero old parts and affine constants can only reduce the counts.

All profiles contradict the sixteen-clause lower bound. Erasure never increased the original deletion count, so the original circuit deletes at least six gates.
\end{proof}

\begin{corollary}[Residual cases with four missing witnesses]\label{cor:fourmissingledger}
Let a circuit $C$ compute $P\in\Sup(G)$ for a four-cycle-free graph and have exactly four missing quadratic witnesses. If it has fewer than $|E(G)|+\vc(G)-1$ gates, its quadratic support $H$ is nonempty. Put
\[
 q=|E(H)|+\vc(H)-1,\quad
 e=|C_2|-q,\quad D=|C|-|C_2|,\quad
 d=\vc(G)-\vc(H),
\]
where $C_2$ is the canonically extracted and pruned circuit. Then $e\ge0$, and the only pairs not excluded by the preceding results are
\[
\begin{array}{c|l}
 d & (D,e)\\ \hline
 2 &(5,0)\\
 3 &(5,0),(5,1),(6,0)\\
 4 &(6,0),(6,1),(7,0).
\end{array}
\]
These are necessary conditions on a hypothetical saving; no existence assertion is made for any listed pair.
\end{corollary}
\begin{proof}
If $H$ is empty, $G$ has four edges and Proposition~\ref{prop:fouredgeboundary} gives the full star count. Otherwise the quadratic bound gives $e\ge0$, and $0\le d\le4$. A strict saving requires $e+D\le3+d$, while Lemma~\ref{lem:fivelost} gives $D\ge5$. Thus $d\le1$ is excluded and the rows for $d=2,3$ follow by integer enumeration.

If $d=4$, the four missing edges must form a matching, since otherwise three vertices cover them. Every minimum cover of $H$ avoids all their endpoints, since otherwise adjoining three vertices would cover $G$. Their eight endpoints consequently contain no edge of $H$. Lemma~\ref{lem:fourmatchinglost} gives $D\ge6$, yielding the final row. The excess $e$ and the deletion count $D$ remain separate throughout; neither is silently set to zero.
\end{proof}


\begin{lemma}[A full context has no cover vertex]\label{lem:contextsupport}
Let $H\subseteq G$ be finite simple graphs with $G$ four-cycle-free, and let $C$ cover $H$. Suppose $G[V\setminus C]$ is a matching containing distinct edges $a_1b_1,a_2b_2$. Let $A,L$ be nonzero positive linear forms, with $AL$ supported on $H$-edges and $\supp(L)\setminus C=\{a_1,a_2\}$. If a positive polynomial $g$ has a monomial with exact variable support $\{b_1,b_2\}$ and every monomial of $gL$ contains a $G$-edge, then
\[
 \supp(L)=\{a_1,a_2\}.
\]
\end{lemma}
\begin{proof}
Choose any $h\in\supp(A)$. The old-edge condition gives $ha_1,ha_2\in E(H)$, so $h\in C$. If $u\in\supp(L)\cap C$, then $hu\in E(H)$ and $u\ne h$. Multiplying the stipulated monomial of $g$ by $u$ gives support $\{u,b_1,b_2\}$. Since $b_1b_2$ is not an edge, safety forces $ub_i\in E(G)$ for some $i$. The four distinct vertices $h,a_i,b_i,u$ form a four-cycle. This contradiction proves the assertion. No degree restriction on $g$ or singleton assumption on $A$ is needed.
\end{proof}

\begin{lemma}[Four products among six deletions]\label{lem:fourproductssix}
Let $G$ be four-cycle-free, let $P\in\Sup(G)$ have nonempty quadratic support $H$, and suppose exactly four edges are missing. Assume $\vc(G)-\vc(H)=4$. Suppose the pruned extracted frame has $K\le\vc(H)+1$ products and $|E(H)|-1$ additions. Exactly six deleted gates cannot consist of four products and two additions. This includes an optimal-total-size frame, the frame with one extra product and minimum additions, one cut with a private equal-degree addition, and two-cut configurations.
\end{lemma}
\begin{proof}
Fix any minimum cover $Z$ of $H$. The missing edges form a matching, and $Z$ avoids all eight endpoints: otherwise three additional vertices would cover $G$. On $T=V\setminus Z$, the graph is precisely that matching together with isolated vertices. Retain those isolated vertices in the restriction. The quadratic output vanishes. By Lemma~\ref{lem:quadoptimal}, the retained products partition $H$ into stars with literal center inputs, leaf sum trees and an outer addition tree. Their centers may repeat and need not form the chosen minimum cover.

Apply full-carrier erasure, preserving the exact frame. Any lost private gate would leave at most five deletions, contrary to Lemma~\ref{lem:fourmatchinglost}. Thus the six-gate profile remains essential. Use the selected-occurrence visibility and independent-symbol frame of that lemma. If only one cut is visible, its high input has four products and one addition, with at most six old linear boundary forms. The Boolean expression $C_0(A_0\vee B_0)$ has at most six distinct factors after extracting common factors. A scalar context gives at most $ab+c_0\le9$ clauses; a linear context gives total factor count at most seven with at least one common factor, hence at most ten clauses. This includes a hidden cut and a private equal-degree sum. The matching needs sixteen clauses, so both cuts are visible.

Write the restricted output as
\[
 F=gL\vee fM\vee cgf,
\]
where $L,M$ are the actual old contexts and $c\ge0$. Dependence of $f$ on the first cut is substituted after this exact-frame expression is formed. Let $k$ count the four private products ancestral to the first high signal $g$. Every remaining private product feeds $f$.

If $k=4$, the second high signal is an ancestor of $g$ or equals it. Write $g=fA$ Booleanly, with at most five distinct old factors. As in Lemma~\ref{lem:fourmatchinglost}, the expression $f(AL\vee M\vee cA)$ has at most six clauses.

If $k=3$, then $g$ has at most $b=4$ linear factors and only one private product remains for $f$. Without a direct proper private-ancestor input, $f$ is $AB$, $AB\vee gB$, or $AB\vee g$, including zero old parts. The explicit identities in Lemma~\ref{lem:fourmatchinglost} give respectively the following clause counts:
\[
\begin{array}{c|cc}
 f&c=0&c>0\\ \hline
 AB&3(b+1)\le15&3b+2\le14\\
 AB\vee gB&3b+2\le14&3b+1\le13\\
 AB\vee g&3b+1\le13&3b\le12.
\end{array}
\]
All are below sixteen. If both inputs of $f$ are private ancestors of $g$, its Boolean factors are among those of $g$ and the ancestor-factor argument applies. Otherwise let its direct ancestor input be $u$. Such an uninserted signal cannot enter a retained linear sum without another private addition. Write $g=uW$, and $f=uA$ or $uA\vee g$. If $u$ has $r\ge1$ distinct factors and $W$ has $t$ remaining factors, then $r+t\le4$. Factoring $u$ leaves a bracket
\[
 WL\vee AM\vee cWA,
 \qquad W(L\vee M\vee c)\vee AM.
\]
The first has at most $2(t+1)$ clauses when $c=0$, and at most $2t+1$ when $c>0$, by the identity $(W\vee A)(W\vee M)(A\vee L)$. The second has at most $2(t+1)$. Including $u$ gives at most $r+2(t+1)\le9$ clauses. Repeated-input squares can make $r=1$ and are included.

Suppose $k=2$. Write $u=UV$ and $g=uW$ Booleanly, allowing $W=1$ for a repeated product. The two-product network for $f$ has at most three boundary forms. If it reads $u$ directly, factor $u$ from both $g$ and $f$. Under that common factor, each $g$-affected boundary $A\vee g$ becomes the linear form $A\vee W$, and a direct $g$ becomes $W$. At most two other factors remain, so $f=uK$ with $K$ a product of at most two linear clauses. The bracket in
\[
 F=u(WL\vee KM\vee cWK)
\]
has at most six clauses for $c=0$, and at most five for $c>0$ using $(W\vee K)(W\vee M)(K\vee L)$. Including $U,V$ gives at most eight.

Otherwise $f$ is a product of at most $t=3$ old or $g$-affected forms, and $g$ has at most $b=3$ factors. If none is affected, write $f=K$. The clause counts are $(b+1)(t+1)\le16$ for $c=0$ and $bt+b+t\le15$ for $c>0$. If at least one is affected, write
\[
 f=H_0K_0\vee gH_0,
\]
with $h\le2$ factors in $H_0$ and $t\le3$ in $H_0K_0$. The identities proved in Lemma~\ref{lem:fourmatchinglost} give
\[
 F=(g\vee H_0K_0M)(L\vee H_0M)\quad(c=0),
\]
and the second factor becomes $L\vee H_0$ when $c>0$. The clause counts are at most $b(t+1)+h+1\le15$ and $b(t+1)+h\le14$, respectively.

The only count that can reach sixteen is therefore $c=0$, with independent $g,f$ each a product of three old forms. Then $F=gL\vee fM$. Each safe summand has at most four clauses and witnesses at most two matching edges. To witness all four, each must witness two, on disjoint pairs. Relabel so the first pair consists of $a_1b_1,a_2b_2$ and $\supp(L)\cap T=\{a_1,a_2\}$; similarly the second pair gives $\supp(M)\cap T=\{a_3,a_4\}$. Quartic equality gives the other three factors of $g$ on $T$ as $\{a_1,b_2\},\{b_1,a_2\},\{b_1,b_2\}$. Thus $g$ has a monomial with exact support $\{b_1,b_2\}$ before restriction as well.

Let $A$ be the first cut's selected old low form. It is linear and nonzero, and $AL$ is supported in $H$. Lemma~\ref{lem:contextsupport} therefore proves that the full context has support $\{a_1,a_2\}$, including before restriction. The same holds for $M$ and $\{a_3,a_4\}$. The rectangle $AL\subseteq H$ and $|\supp(L)|=2$ force $A$ to have singleton support; the coefficient-one frame makes it a literal input. Likewise the second low is literal.

Every star affected by the first cut is covered by $\supp(L)$. If the literal cut enters its center wire, its whole leaf support occurs in $L$. If it enters a leaf wire, the center itself occurs in $L$. The same statement holds for the second cut and $\supp(M)$. At least one product is affected by each cut. A product affected in opposite factors would give a positive $gf$ coefficient, contradicting $c=0$; a product affected twice in the same factor would put its other nonzero factor into both contexts, contradicting their disjoint supports. Thus two distinct assigned stars can be covered by the four vertices in $\supp(L)\cup\supp(M)$.

Choose centers for the other $K-2$ products and add those four vertices. Every old edge belongs to a covered assigned star, and every missing edge meets the four chosen endpoints. This is a cover of size at most
\[
 K-2+4\le\vc(H)+3,
\]
contrary to the cover gain. Edges between centers are covered by their assigned stars in exactly the same way; repeated centers only decrease the size of this cover.

Finally, if $k=1$, then $g=UV$ has $b\le2$ factors and the other three products give at most $t=4$ boundary factors for $f$. If none is affected, the preceding counts are at most $(b+1)(t+1)\le15$ and $bt+b+t\le14$. Otherwise the same $H_0,K_0$ decomposition has $h\le3$ and gives at most $b(t+1)+h+1\le14$ or $b(t+1)+h\le13$. Zero factors, constants and repeated inputs only simplify these Boolean formulas.

Every profile has been excluded. The minimum-addition star partition and the cover gain exclude the sole possible sixteen-clause case. The argument includes one extra retained product; it does not exclude the three-product/three-addition profile.
\end{proof}

\begin{lemma}[One deleted addition and a missing matching]\label{lem:oneadditionmatching}
Let $G$ be any finite simple graph and let $P\in\Sup(G)$ have nonempty quadratic support $H$. Suppose $2r$ vertices, with $r\ge1$, induce an $r$-edge matching in $G$ and no edge in $H$. If canonical extraction and pruning delete exactly one addition and $p$ multiplications, then
\[
 p\ge 2^r-2,\qquad p+1\ge2^r-1.
\]
There is no four-cycle, cover-number or optimality hypothesis. The estimate for $r=1$ is deliberately weak; no sharpness assertion is made.
\end{lemma}
\begin{proof}
The unique deleted addition must be an unequal-degree cut. Otherwise every useful gate would be homogeneous, and the nonzero quadratic output would contain all its edge witnesses. Apply full-carrier erasure from Lemma~\ref{lem:twoproductnormal}. The matching witnesses remain, so the unique cut cannot disappear; the operation preserves the exact frame and can only decrease the deleted product count. Every surviving private product is ancestral to the cut's high input $F$, since a private signal must enter the selected frame through a cut. Before the sole cut, all remaining off-carrier boundaries are old homogeneous linear forms. No full quadratic carrier is used privately.

If there are $p'\le p$ private products in the high-input ancestry, their connected multiplication DAG has at most $p'+1$ distinct boundary forms, by the operand-slot count. Repetitions disappear under Boolean idempotence. The low degree of the cut is one or two. In degree two its output context is scalar. In degree one, treating the cut as an independent symbol in the quadratic frame gives a linear context: a squared occurrence would put a positive square of the nonzero old low form into $P_2$, impossible for $\Sup(G)$. Thus the full output has form $P_2+FD$, with $D$ scalar or homogeneous linear, including all fanout multiplicities.

Set variables outside the specified matching endpoints to zero. The quadratic part vanishes, leaving the nonzero matching function. The product $FD$ is a conjunction of at most $p'+2$ positive linear clauses, including the context if nonconstant. Each clause must meet every matching edge and rejects at most one of its $2^r$ independent endpoint transversals. Constants reject none. Hence $2^r\le p'+2\le p+2$, giving both claims. The restriction is a Boolean analysis and is not asserted to preserve formal degrees or the deletion count.
\end{proof}

\begin{proposition}[Deleted products at maximal cover gain]\label{prop:maxgainproducts}
Let $G$ be four-cycle-free and let a positive variable-input circuit compute $P\in\Sup(G)$ with nonempty quadratic support $H$. Put $r=|E(G)\setminus E(H)|>0$. Suppose
\[
 \vc(G)-\vc(H)=r,
 \qquad \#\{\text{multiplications in }C_2\}=\vc(H),
\]
where $C_2$ is the canonically extracted quadratic circuit. At least $r$ deleted gates are multiplications. Its retained addition count is unrestricted. There is no bound on deleted additions or total deletions, and arbitrary fanout, repeated inputs and quadratic-carrier feedback are permitted. An optimal total-size quadratic circuit has the required product count by Lemma~\ref{lem:quadoptimal}.
\end{proposition}
\begin{proof}
Every retained multiplication is the product of two homogeneous linear forms with disjoint supports, and its cross pairs are edges of $H$. Four-cycle-freeness makes one support a singleton. Choose its vertex as center, choosing either endpoint for a one-edge product. The chosen centers cover $H$. Their number is at most the product count $\vc(H)$ and at least $\vc(H)$, so they are distinct and form a minimum cover $Z$, with exactly one retained product per center. This argument requires neither optimal additions nor coefficient one; edges between centers may occur repeatedly.

The $r$ missing edges form a matching, and every minimum cover of $H$ avoids all their endpoints: otherwise at most $r-1$ additional vertices would cover $G$. On $T=V\setminus Z$ the graph is that matching together with isolated vertices. Keep all those isolated vertices throughout.

For each center $h\in Z$, put $N_h=N_H(h)\cap T$. This is nonempty because $Z$ is minimum. Its vertices are precisely the support on $T$ of $h$'s assigned old leaf form: every old edge to a vertex outside $Z$ must be assigned to $h$. Moreover $N_h$ cannot lie in a partial transversal of the missing matching, meaning a set with at most one endpoint from each pair and no isolated vertex. If it did, replacing $h$ by $N_h$ and adding an endpoint of each missing edge not met by $N_h$ would give a cover with at most
\[
 |Z|-1+|N_h|+(r-|N_h|)=|Z|+r-1
\]
vertices, contrary to the cover gain.

We use a support observation. Let $M$ be a matching function with possible isolated variables, and let a nonempty vertex set $N$ not lie in any partial transversal. If a monotone function $A$ satisfies
\[
 A\wedge\bigvee_{v\in N}x_v\ \le\ M,
\]
then $A\le M$. Otherwise $A$ is true on an independent set $S$. For every $v\in N$, monotonicity makes $A$ true on $S\cup\{v\}$, so this enlarged set must contain a matching edge. Thus $v$ is the missing mate of an endpoint in $S$. All such mates form a partial transversal, containing no isolated vertex. They cannot contain $N$, a contradiction.

Apply full-carrier erasure from Lemma~\ref{lem:twoproductnormal}, preserving the entire extracted frame. This step does not assume a bound on private products. It introduces no multiplication and leaves no use of a full quadratic carrier outside its protected addition-only output routes. Let $p$ denote the original number of deleted products; normalization leaves at most $p$ private products.

Now set every variable in $Z$ to zero, so the output is exactly the matching function $M$ on $T$. Consider the retained product centered at $h$. Let $A_h,B_h$ be its full restricted center-input and leaf-input functions. Positivity and the unchanged extracted paths give
\[
 B_h\ge\bigvee_{v\in N_h}x_v.
\]
The product still has an addition-only output route, so
\[
 A_h\wedge\bigvee_{v\in N_h}x_v
 \ \le\ A_h\wedge B_h\ \le\ M.
\]
The support observation proves $A_h\le M$.

Replace every retained product by its center-input wire $A_h$. This is a Boolean modification after restriction; it need not preserve formal degrees. It introduces no gate and can only increase the output, since $A_h\wedge B_h\le A_h$. Yet every added output disjunct $A_h$ implies $M$, so the output cannot exceed $M$. More explicitly, after carrier erasure the retained products and unchanged private high inputs feed an addition-only carrier. Its center and leaf inputs have no retained-product ancestor. The latter functions therefore stay fixed under all replacements, while the other carrier inputs were already output disjuncts and imply $M$. The modified output is exactly $M$.

Every retained multiplication has now been bypassed. The resulting monotone circuit computes the $r$-edge matching with at most $p$ conjunction gates, and arbitrary disjunctions, constants and fanout. Such a circuit needs at least $r$ conjunctions. Here is a proof that includes sharing. Resolve each of its $p$ conjunctions globally by choosing one operand, leaving disjunctions unchanged. Each of the $2^p$ choices gives an OR-only circuit, hence a positive linear clause or a constant. The original function is at most every resolution. Conversely, on an input where it is false, choose a false operand at each originally false conjunction. Induction through the DAG shows that every originally false gate remains false in that resolution. Thus the original function is the conjunction of the at most $2^p$ resolved clauses.

The matching requires at least $2^r$ such clauses: each must be true on every matching edge and can reject at most one independent endpoint transversal, while all $2^r$ transversals must be rejected. Isolated variables are set to zero on those assignments; a constant-true clause rejects none, and a false clause is impossible for the nonzero matching function. Hence $2^r\le2^p$ and $p\ge r$.
\end{proof}

\begin{corollary}[The refined four-witness residual]\label{cor:refinedfourledger}
In Corollary~\ref{cor:fourmissingledger}, the maximal-gain pair $(d;D,e)=(4;6,0)$ is impossible. The remaining necessary pairs are
\[
\begin{array}{c|l}
 d&(D,e)\\ \hline
 2&(5,0)\\
 3&(5,0),(5,1),(6,0)\\
 4&(6,1),(7,0).
\end{array}
\]
Within $(d;D,e)=(4;7,0)$, only four deleted products with three additions, or five deleted products with two additions, remain unexcluded by these results. No existence assertion is made for these pairs or profiles.
\end{corollary}
\begin{proof}
For $d=4,e=0$, Proposition~\ref{prop:maxgainproducts} requires at least four deleted products. At least one deleted addition is a cut, since otherwise the useful output would be homogeneous quadratic. If $D=6$, the possibilities are four products and two additions, excluded by Lemma~\ref{lem:fourproductssix}, or five products and one addition. The latter is excluded by Lemma~\ref{lem:oneadditionmatching}, which for the induced four-edge matching requires at least fourteen products when there is only one deleted addition.

For $D=7,e=0$, the same product bound leaves four, five or six products. Six products and one addition are again excluded by that lemma. This gives exactly the two stated profiles. The positive-excess pair $(6,1)$ is outside the optimal-frame hypothesis and is not removed. The lower-gain rows are unchanged from Corollary~\ref{cor:fourmissingledger}.
\end{proof}

\begin{lemma}[Two additions with at most five products]\label{lem:fiveproductsseven}
Under the hypotheses of Proposition~\ref{prop:maxgainproducts}, suppose $r=4$. It is impossible to have at most five deleted multiplications and at most two deleted additions. The retained addition count is unrestricted. In particular, this excludes five products and two additions in an optimal-total-size quadratic frame.
\end{lemma}
\begin{proof}
Restrict the assigned minimum cover to zero, keeping isolated outside vertices. The safe-center replacement in Proposition~\ref{prop:maxgainproducts} removes every retained multiplication and preserves the four-edge matching function. The selected center input at $h$ is a positive multiple $c_hx_h$. Its retained linear addition network may have arbitrary sharing and duplication, but positivity forces every selected variable in it to be $x_h$. Restoring its cuts gives that old multiple plus a positive linear combination of cut-high signals. After restriction, the old term vanishes. The direct high inputs of the quadratic addition carrier are also output disjuncts. Thus the matching is the OR of a nonempty subset of the contributing high signals, each individually safe. Extra retained additions only join existing signals; they create no new cut. Leaf-only cuts need not be safe and are omitted by the center replacement.

Full-carrier erasure cannot increase the at most five private products or two private additions. If only one addition remained, Lemma~\ref{lem:oneadditionmatching} would require fourteen products. Thus there are two additions and at least one cut. If only one cut is visible, its high computation has at most five products and one private addition, on at most seven old linear boundaries. This includes a hidden cut and a private equal-degree sum. As before, its Boolean value is $C_0(A_0\vee B_0)$, with total distinct factor count at most seven after extracting common factors. It has at most $ab+c\le12$ clauses. That sole contributing high must itself be the matching, which needs sixteen clauses, a contradiction.

Suppose both cuts are visible, with high signals $g,f$ in cut order. Let $k$ be the number of private products ancestral to $g$. This first high signal is a product of at most $k+1\le6$ old linear clauses, so it cannot be the sole contributing signal. If only $f$ contributes, collapse the private ancestry of $g$ to boundary nodes of the remaining network for $f$. There are at most $n\le5-k$ products in this network and at most $n+1$ boundary signals. Each boundary is an old linear form, a form $A\vee g$ affected by the first cut, or a direct product ancestor of $g$, including $g$ itself. Each has a CNF with at most $k+1$ clauses. Uninserted ancestors cannot enter retained linear sums without another cut, and quadratic carriers were erased. Multiplying these boundaries gives at most
\[
 (k+1)(n+1)\le(k+1)(6-k)\le12
\]
clauses, including the ancestor case $n=0$. Thus $f$ alone is impossible.

It remains that both $g$ and $f$ contribute, so the output is $g\vee f$ and both are safe. The first cut's old low part vanishes under restriction: it lies in a selected linear center network supported on $\{h\}$ or directly in the quadratic carrier. Hence its restricted full value is $g$. Because $g$ is already an output disjunct, replace that cut value and direct uses of $g$ inside $f$ by zero, preserving the computation of $g$. The identity
\[
 g\vee f(g)=g\vee f(0)
\]
holds by evaluating $g=0,1$, including arbitrary fanout. Write $f_0=f(0)$. All affected retained linear boundaries in $f_0$ are now old linear forms. The two signals $g,f_0$ are outputs of a multiplication-only DAG on such forms, using at most five products in their combined ancestry. Proper ancestors of $g$ may still be shared. Both outputs are safe and their union is the matching.

Each output has at most six distinct boundary clauses, so witnesses at most two matching edges: three witnesses would require eight clauses. Therefore each witnesses exactly two, on disjoint pairs. Let $A,B$ be their sets of distinct boundary nodes, and put
\[
 u=|A\setminus B|,\qquad v=|B\setminus A|,\qquad w=|A\cap B|.
\]
The union of their product ancestries has at most two connected components. With $p\le5$ product gates, counting its $2p$ operand wires gives at most $p+2$ distinct boundary nodes. Thus $u+v+w\le7$.

On the sixteen independent endpoint transversals, with isolated variables zero, a clause private to $A$ meets both edges witnessed by $g$ and can reject at most four transversals. A shared clause meets all four witnessed edges and rejects at most one. Safety of $g$ requires $4u+w\ge16$; similarly safety of $f_0$ requires $4v+w\ge16$. Therefore
\[
 32\le4u+4v+2w\le4(u+v+w)\le28,
\]
a contradiction. Constants and repeated inputs cannot increase these rejection counts. All cases are excluded. No step used an optimal retained addition count; the required hypothesis is the minimum retained multiplication count.
\end{proof}


\begin{proposition}[Additions when the product bound is attained]\label{prop:equalproductadditions}
Under the hypotheses of Proposition~\ref{prop:maxgainproducts}, if exactly $r$ deleted gates are multiplications, at least $r$ deleted gates are additions. The addition count refers to original deleted gates. There is no restriction on retained additions, fanout, repeated operands or nested cuts. In particular the assertion applies to an optimal-total-size quadratic frame.
\end{proposition}
\begin{proof}
We first establish the needed equality case of the matching conjunction bound. Suppose a monotone circuit computes the $r$-edge matching $M$ with exactly $r$ binary conjunctions. Constant-one inputs can be propagated first: removing a conjunction would contradict the lower bound. Our application introduces only constant zero. Resolve the $r$ conjunctions globally as in Proposition~\ref{prop:maxgainproducts}. There are $2^r$ resolutions and $2^r$ independent endpoint transversals to reject. Each resolved clause rejects at most one. Hence every resolution rejects exactly one transversal and no two resolutions reject the same one.

Fix a conjunction $g$. It must be reachable from the output after every resolution of the other conjunctions; otherwise flipping its own choice gives the same clause twice. This implies that its output is at most $M$. Indeed force $g$ to one and all other input leaves to zero. If the output were false, choose a false operand at every other false conjunction. The false output's resolved OR cone would avoid $g$, a contradiction. Thus forcing $g$ to one already forces the output to one, and monotonicity proves $g\le M$ on every actual input.

Order the conjunctions topologically. Each matching edge has a first conjunction that is true on its two-endpoint assignment. Otherwise an OR-only output contribution would also accept a singleton. For a fixed conjunction, set all preceding conjunction outputs to zero in its two operands. Their remaining parts are positive linear clauses or constants. Their conjunction is at most $M$. If it witnesses an edge, its two factors must be precisely the singleton literals at that edge's endpoints: any overlap gives a singleton witness, and any additional cross pair is not an edge of the matching. Thus one conjunction can be first for at most one edge. There are $r$ edges and $r$ gates, so every gate is first for exactly one edge. On that edge every earlier conjunction is false, and each operand of the new gate has a nonzero raw-variable contribution. In particular, no conjunction can directly read another conjunction, or an OR-only expression in earlier conjunction outputs without a raw-variable contribution.

Return to the arithmetic circuit. Apply exact-frame full-carrier erasure, then restrict the chosen minimum cover to zero and replace retained products by their center inputs as in Proposition~\ref{prop:maxgainproducts}. Keep every deleted addition on a surviving path, including a cut whose old input has become zero; suppressing such OR-with-zero gates would destroy the count of original additions. The resulting matching circuit has at most $r$ private conjunctions. Its conjunction lower bound forces all $r$ to survive, and the preceding equality argument applies.

Every useful private product must enter a marked deleted addition before another product or the retained frame. Direct entry into another private product contradicts the equality argument. Direct entry into a selected retained degree-one position is impossible, since a product's minimum degree is at least two. An equal-degree entry into the retained quadratic addition carrier would select that private multiplication source into the extracted frame; a higher-degree entry makes the receiving addition a deleted cut. These are statements about the original selection, before Boolean simplification.

Extra retained center additions do not change this conclusion. Their selected polynomial has support $\{h\}$. A component made only of private products and private additions has minimum degree at least two, so it cannot supply the degree-one center wire without a cut having an old degree-one input. That input is a boundary node even when restriction turns its value into zero. This also handles new output-carrier edges created by bypassing retained products.

Form an undirected operand multigraph with the $r$ private product outputs, the $a'$ surviving marked deleted additions, and all other immediate input nodes of those additions. The last kind are boundary vertices. Include both operand incidences of every marked addition, including zero and repeated operands; there are exactly $2a'$ edges. Every private product has a surviving incidence by the preceding argument.

Each connected component has a boundary vertex. Otherwise its additions are OR-only expressions in preceding private products. Such a component cannot feed a later private conjunction, since its operand would be false on that conjunction's own matching edge. It cannot enter a retained linear position by degree. An equal-degree export into the quadratic carrier or original output would select a private multiplication source into the frame; a higher-degree export is another marked cut in the same component. A new export caused by a center replacement has an old degree-one boundary as just shown. Edges to another marked addition stay within the component. These exhaust all useful exits, so a component without a boundary would contain no useful private product; a component consisting only of additions and no boundary is impossible by acyclicity.

Let $b$ be the number of boundary vertices and $c$ the component count. Then $b\ge c$. Connectedness of the multigraph gives
\[
 2a'\ge r+a'+b-c,\qquad a'\ge r+b-c\ge r.
\]
No new addition was introduced, so the original number of deleted additions is at least $a'$, proving the claim. This counts original marked additions, not arbitrary OR gates in the transformed matching circuit.
\end{proof}

\begin{corollary}[Eight deletions at maximal four-edge gain]\label{cor:fourminproduct}
Under the hypotheses of Proposition~\ref{prop:maxgainproducts}, if $r=4$ then at least eight gates are deleted, and the original circuit has at least $|E(G)|+\vc(G)-1$ gates. The retained addition count is unrestricted.

Consequently, for the general four-missing-witness question of Corollary~\ref{cor:fourmissingledger}, a circuit below the star count can have only the following necessary pairs:
\[
\begin{array}{c|l}
 d&(D,e)\\ \hline
 2&(5,0)\\
 3&(5,0),(5,1),(6,0)\\
 4&(6,1).
\end{array}
\]
In the last row the extracted frame must have exactly $\vc(H)+1$ multiplications. Its additions are minimal for that product count, so Lemma~\ref{lem:quadoptimal} describes it as a partition into star trees with an outer sum tree; centers may repeat. Its six deleted gates must be three products and three additions. The restored higher terms in that remaining profile are not controlled by these results. No existence assertion is made for these remaining cases.
\end{corollary}
\begin{proof}
Write $p,a$ for the deleted product and addition counts. Proposition~\ref{prop:maxgainproducts} gives $p\ge4$, and a missing witness requires a cut, so $a\ge1$. If $p+a\le7$ and $p=4$, Proposition~\ref{prop:equalproductadditions} gives $a\ge4$, a contradiction. If $p=5$, then $a\le2$, contradicting Lemma~\ref{lem:fiveproductsseven}. If $p=6$, then $a=1$, contradicting the bound $p\ge14$ in Lemma~\ref{lem:oneadditionmatching}. If $p\ge7$, the required addition already makes the total at least eight. Thus $D=p+a\ge8$.

The retained circuit has at least $|E(H)|+\vc(H)-1$ gates. Both edge count and cover number increase by four, so adding eight gives the asserted star lower bound. Apply this result to the necessary pairs in Corollary~\ref{cor:refinedfourledger}. The optimal-frame row $(4;7,0)$ is removed. In the remaining row $(4;6,1)$, the product count cannot equal $\vc(H)$. The fixed-product inequality in the quadratic proof is $|C_2|\ge|E(H)|+K-1$, where $K$ is its multiplication count. Since its excess is one and $K\ge\vc(H)$, necessarily $K=\vc(H)+1$. Corollary~\ref{cor:threedeletedproducts} requires at least three deleted products. There is at least one cut; five products with one addition are excluded by Lemma~\ref{lem:oneadditionmatching}, and four products with two additions by the strengthened Lemma~\ref{lem:fourproductssix}. Hence only three products and three additions remain. The lower-gain rows are unchanged.
\end{proof}


\begin{lemma}[Three independent products and a matching]\label{lem:threeindependentmatching}
Let $M$ be a finite matching, with isolated variables allowed. For $i=1,2,3$, let $g_i=U_iV_i$, where $U_i,V_i$ are nonnegative homogeneous linear forms, possibly zero. Let $L_i$ be nonnegative affine-linear forms and $c_{ij}\ge0$. Suppose every monomial of
\[
 F=\sum_{i=1}^3g_iL_i+\sum_{i<j}c_{ij}g_ig_j
\]
contains an $M$-edge. Then at most three distinct $M$-edges occur as exact two-variable supports of monomials of $F$.
\end{lemma}
\begin{proof}
Call an edge available to $g_i$ when $g_i$ is true on its two-endpoint assignment. Each nonzero summand is safe by positivity. For a safe summand, truth on an edge assignment is equivalent to witnessing that exact edge. The matching factor count in Lemma~\ref{lem:twoproductmatching} shows that a direct summand $g_iL_i$ witnesses at most one edge, and a mixed summand $g_ig_j$ at most two. A nonzero constant term in an affine form makes its Boolean clause true and can only reduce these counts.

If a mixed summand witnesses two edges, quartic equality confines all four of its linear factors to their endpoints. Neither of its two signals is available on any other matching edge. Thus every term involving those signals witnesses only the two edges, and the sole remaining direct term supplies at most one more.

Now suppose every mixed term has at most one witness. If a direct term $g_iL_i$ witnesses $d$ while an incident mixed term $g_ig_j$ witnesses a different edge $e$, the cubic rigidity argument in the proof of Lemma~\ref{lem:twoproductmatching} gives
\[
 g_i=b_d\vee e
\]
Booleanly, where $b_d$ is one endpoint of $d$ and $e$ denotes the conjunction of its two endpoints. A positive constant term in $L_i$ would instead make $g_i$ a safe two-factor product, incapable of being available on both edges, so that boundary case cannot arise. In particular, $g_i$ is available only on $d,e$.

The other signal $g_j$ is also unavailable on a third matching edge $f$. Otherwise choose an endpoint $u$ of $e$ in $U_j$ and an endpoint $v$ of $f$ in $V_j$, possible because $g_j$ is available on both edges. The independent set $\{b_d,u,v\}$ makes $g_i$ true through $b_d$ and $g_j$ true through its two factors, contradicting safety of their mixed term. Consequently every term involving either signal witnesses only $d,e$, and the remaining direct term again supplies at most one further edge.

If neither situation occurs, let $m$ be the number of direct terms having a witness. Any mixed term incident with one of these indices can only repeat that direct witness. Only mixed terms between the other $3-m$ indices can add edges. The total number is at most
\[
 m+\binom{3-m}{2}\le3,\qquad 0\le m\le3.
\]
This exhausts the cases. Zero forms and repeated factors cause no difficulty. A positive term $g_i^2$ may also be included: its Boolean function is $g_i$, and safety of the square is equivalent to safety of $g_i$, since an unsafe monomial would remain unsafe when squared. It can therefore be folded into a positive constant term of $L_i$ for this support statement.
\end{proof}

This lemma concerns the displayed independent-product expression. It supplies no normalization of three dependent private multiplications and does not settle the remaining three-product/three-addition circuit profile.

\begin{proposition}[One deleted product and cover gain]\label{prop:oneproductcover}
Let $G$ be four-cycle-free, let a positive variable-input circuit compute $P\in\Sup(G)$ with nonempty quadratic support $H$, and let $K$ be the number of multiplication gates in its canonically extracted circuit. If at most one deleted gate is a multiplication, then
\[
 \vc(G)-\vc(H)\le\min\{2,\ K-\vc(H)+1\}.
\]
In particular, if $K=\vc(H)$ the cover gain is at most one. The retained and deleted addition counts are unrestricted, as is the number of missing quadratic witnesses. The universal bound two is sharp.
\end{proposition}
\begin{proof}
Apply the exact-frame carrier erasure of Lemma~\ref{lem:twoproductnormal}. If no private product remains, there is no missing edge. Otherwise the one-product specialization gives
\[
 P'=P_H+gD+\alpha g^2,\qquad g=AB,\quad \alpha\ge0,
\]
where $A,B$ are positive homogeneous linear forms, $D$ is nonnegative affine-linear, and the Boolean function and the extracted frame are unchanged. Arbitrarily many private additions may form $A,B$. If $\alpha>0$, every monomial of $g$ is safe, because its square has the same variable support. The same holds if $D$ has a positive constant term. Then the two supports of $g$ are disjoint and their complete bipartite cross support lies in $G$. Four-cycle-freeness makes this support a star, and every exact edge witness in a multiple of $g$ is already one of its edges. Adding its center to a minimum $H$-cover proves the result. We may therefore assume $\alpha=0$ and $D$ is nonzero homogeneous linear; $D=0$ supplies no missing edge.

We classify the exact two-variable witness graph $E$ of the cubic $ABD$. Every missing edge of $G$ lies in $E$, and every edge of $E$ lies in $G$. Write its three factor supports as $A_1,A_2,A_3$, with $A_3=\supp(D)$. Their triple intersection is empty, since a variable in all three would yield a forbidden pure cube. Let $X_i$ consist of variables in the other two supports but not $A_i$, and let $Y_i$ consist of variables in $A_i$ alone. The exact edges of $E$ are precisely
\[
 X_i\times X_j\quad(i\ne j),\qquad X_i\times Y_i\quad(i=1,2,3).
\]
This follows because two variables witness an edge exactly when their membership sets together meet all three factors.

If at most one $X_i$ is nonempty, $E$ is empty or complete bipartite, hence a star. If exactly two $X$-classes are nonempty and one has at least two vertices, the other must be a singleton and the larger class's attached $Y$-class must be empty: two vertices of that class otherwise have two common neighbors. Again $E$ is a star. If both are singletons, $E$ is a double star; the only non-star case has nonempty pendant leaf sets at both core vertices.

If all three $X$-classes are nonempty, each is a singleton $x_i$, since two vertices of one class and one from each other class would form a four-cycle. They form a triangle, with $Y_i$ attached at $x_i$. At least one $Y_i$ is empty. Otherwise choose $y_i\in Y_i$ for each $i$. The monomial $y_1y_2y_3$ occurs, while any possible edge $y_iy_j$ would complete the four-cycle $y_i,x_i,x_j,y_j$. This contradicts validity of that monomial. Thus $E$ is a star, a proper double star, or a triangle with pendant leaves at at most two core vertices. In particular, it has a cover of size at most two. Adjoining this cover to a minimum $H$-cover proves the universal gain bound two, without any retained-product hypothesis. Example~\ref{ex:oneproductnonminimal} attains it.

The quadratic product bound gives $K\ge\vc(H)$. If $K\ge\vc(H)+1$, the bound two already implies $\vc(G)\le K+1$. It remains to prove the gain bound one when $K=\vc(H)$.

Choose the distinct minimum-cover centers $C$ from the minimum-product argument in Proposition~\ref{prop:maxgainproducts}. There is exactly one retained product per center $h\in C$, with old center form a positive multiple of $h$ and old leaf form $L_h$. Its support outside $C$ is exactly $N_H(h)\setminus C$, and is nonempty by minimality of $C$. Expanding the actual retained products gives the following context property:
\begin{equation}\label{eq:oneproductcontext}
 w\in\supp(D)\setminus C\quad\Longrightarrow\quad
 w\in\supp(L_h)\subseteq\supp(D)\text{ for some }h\in C.
\end{equation}
Indeed, perturbing a leaf input contributes its center, which lies in $C$. An outside context variable must instead come from perturbing a center input, which contributes the whole old leaf form. All coefficients and output multiplicities are nonnegative, so no contribution disappears. This uses neither addition optimality nor coefficient one.

The star case is settled by adding its center to $C$. Suppose for contradiction that $\vc(G)>|C|+1$. Every two-element cover $W$ of $E$ must then be disjoint from $C$, since otherwise adding its other vertex would suffice. It is enough to find $h\in C$ with
\[
 N_H(h)\setminus C\subseteq W:
\]
then $(C\setminus\{h\})\cup W$ covers $H$ and all of $E$, hence $G$, with $|C|+1$ vertices. Two consequences of~\eqref{eq:oneproductcontext} will supply such an $h$.

First, if $\varnothing\ne\supp(D)\setminus C\subseteq W$, choose any outside context variable and its center in~\eqref{eq:oneproductcontext}. Second, suppose $W=\{v,w\}$, $w\in\supp(D)$, and $\supp(D)\subseteq N_G(v)$. The selected center $h$ for $w$ has no other outside neighbor $z$: otherwise $h,w,v,z$ form a four-cycle. These vertices are distinct because $h\in C$, $v,w,z\notin C$, and $v$ is not its own neighbor. This again gives the required center.

For the triangle case, $\supp(D)=\{x_1,x_2\}\cup Y_3$. If $Y_3$ is empty, take $W=\{x_1,x_2\}$ and use the first consequence. If $Y_3$ is nonempty, at least one of $Y_1,Y_2$ is empty; say $Y_2=\varnothing$. Then $W=\{x_3,x_1\}$ covers $E$, and the whole support of $D$ lies in $N_G(x_3)$. The second consequence applies with $v=x_3,w=x_1$. The other case is symmetric.

For a proper double star, let its nonempty core classes be $X_i=\{u\},X_j=\{v\}$, with nonempty leaf sets $Y_i,Y_j$, and let $k$ be the remaining index. Take $W=\{u,v\}$. If $k\ne3$, one core is $X_3$; call it $u$. Then $\supp(D)=\{v\}\cup Y_3\subseteq N_G(u)$, so the second consequence applies.

It remains $k=3$. Here $\supp(D)=\{u,v\}\cup Y_3$, and $Y_3$ has no $E$-neighbor. Neither $Y_i$ nor $Y_j$ can lie entirely in $C$: if, for example, $Y_i\subseteq C$, adding $v$ would cover all of $E$. Choose $y_i\in Y_i\setminus C$ and $y_j\in Y_j\setminus C$. If $z\in Y_3\setminus C$, the monomial $y_iy_jz$ occurs. The pair $y_iy_j$ cannot be a $G$-edge, since it would close the four-cycle $y_i,u,v,y_j$. Neither pair $y_iz,y_jz$ is an $E$-edge. Any such $G$-edge would therefore be old, but an $H$-edge cannot have both endpoints outside $C$. This contradicts validity. Hence $Y_3\subseteq C$, and $\supp(D)\setminus C=\{u,v\}$; the first consequence finishes the proof.
\end{proof}

The gain-one conclusion uses the minimum retained multiplication count; the universal gain-two conclusion permits arbitrary retained products. In particular, cover gain three requires at least two deleted multiplications, regardless of the quadratic excess. These conclusions do not close the remaining three-product/three-addition profile.

\begin{example}[Extra retained products]\label{ex:oneproductnonminimal}
Take eight distinct vertices $h_1,\ell_1,u,v,h_2,\ell_2,a,b$ and put
\[
 E(H)=\{h_1u,h_1\ell_1,h_2v,h_2\ell_2\},\qquad
 E(G)=E(H)\cup\{uv,ub,va\}.
\]
The graph $G$ is the path $\ell_1,h_1,u,v,h_2,\ell_2$ with additional leaves $b$ at $u$ and $a$ at $v$, hence a tree. With the shared product $g=(u+a)(v+b)$, compute
\[
 P=(h_1+g)u+h_1\ell_1+(h_2+g)v+h_2\ell_2.
\]
Booleanly $gu=uv\vee ub$ and $gv=uv\vee va$, so $P\in\Sup(G)$, and its quadratic support is exactly $H$. The cover $\{h_1,h_2\}$ and matching $\{h_1\ell_1,h_2\ell_2\}$ prove $\vc(H)=2$. The cover $\{h_1,h_2,u,v\}$ and matching $\{h_1\ell_1,h_2\ell_2,ub,va\}$ prove $\vc(G)=4$.

The circuit has five products and seven additions, hence twelve gates. Extraction bypasses the two cuts and removes the three-gate computation of $g$, leaving four products and three outer sums. Thus exactly one product and four additions are deleted, yet the cover gain is two. The retained product count is four, exceeding $\vc(H)$ by two. The full star count is $7+4-1=10$, so this example lies above it and does not refute the unrestricted conjecture.
\end{example}

\begin{lemma}[Nonancestral deleted products]\label{lem:nonancestralnormal}
Let a positive variable-input circuit have minimum output degree two and nonzero squarefree quadratic part $P_2$. Suppose its deleted multiplication gates are pairwise nonancestral: no directed path in the original circuit joins two of them, including paths through retained gates. There is a Boolean-equivalent positive normalization preserving the exact extracted quadratic frame and introducing no multiplication, whose polynomial is
\[
 P_2+\sum_i g_iL_i+\sum_{i\le j}c_{ij}g_ig_j,
 \qquad g_i=U_iV_i,
\]
where $U_i,V_i$ are nonnegative homogeneous linear forms, $L_i$ are nonnegative affine-linear forms, and $c_{ij}\ge0$. Erased products may be represented by zero. No algebraic independence of these forms is assumed, and the addition counts are unrestricted.
\end{lemma}
\begin{proof}
Apply the full-carrier erasure of Lemma~\ref{lem:twoproductnormal}, before its two-product parametrization. Protect all minimum-degree-two nodes having an addition-only output path. In topological order erase their off-carrier operand occurrences. The full current value of each protected node still occurs additively at the output and absorbs every removed multiple Booleanly. The exact extracted frame remains: retained linear inputs are unchanged, and retained quadratic addition paths stay inside the protected carrier. No new multiplication or ancestral relation is introduced.

Every retained product belongs to that carrier and has no remaining use feeding another product. A surviving private product $g_i$ has no private product ancestor, by the original nonancestry hypothesis, and no retained product ancestor, by erasure. Its operands are therefore addition-only expressions in variables, proving $g_i=U_iV_i$. More generally, every signal outside the protected carrier is an ordinary homogeneous linear form plus a nonnegative scalar linear combination of the $g_i$.

Each retained product consequently has inputs
\[
 A_h+\sum_i a_{hi}g_i,\qquad B_h+\sum_i b_{hi}g_i,
\]
where $A_h,B_h$ are its unchanged extracted linear inputs. Expansion gives $A_hB_h$, terms $g_i$ times old linear forms, and quadratic terms in the $g_i$, including their squares. The protected output carrier is an addition network. Any outside signal entering it additively contributes at most a linear combination of the $g_i$; an ordinary variable contribution would contradict minimum output degree two. Sum the expansions, with their nonnegative output multiplicities. The component without any $g_i$ is the unchanged $P_2$; additive contributions are permitted by the affine $L_i$. This is the stated exact identity for the normalized polynomial. Zero branches and repeated factors are included. The original nonancestry condition is essential to the operand argument; a condition only on the subgraph induced by private gates would not suffice.
\end{proof}

\begin{corollary}\label{cor:nonancestralmatching}
Suppose $P\in\Sup(G)$ has nonempty quadratic support $H$, and eight vertices induce four disjoint $G$-edges and no $H$-edge. Its circuit cannot have at most three deleted multiplications that are pairwise nonancestral. Consequently the three deleted products in the remaining maximal-gain profile of Corollary~\ref{cor:fourminproduct} must contain an ancestral pair.
\end{corollary}
\begin{proof}
Apply Lemma~\ref{lem:nonancestralnormal} and set every variable outside the eight vertices to zero. The quadratic remainder vanishes, and the Boolean function is the four-edge matching. The restricted expression has at most three products $g_i$ in precisely the form of Lemma~\ref{lem:threeindependentmatching}, including zero forms, affine contexts and diagonal squares. That lemma allows at most three matching witnesses, a contradiction. This argument does not normalize dependent private products and does not settle the remaining profile.
\end{proof}

\begin{lemma}[Literal multiplicative operands]\label{lem:literalproducts}
Let $M$ be a finite matching, with isolated variables allowed. Form at most three signals $g_1,g_2,g_3$ in acyclic order, each by multiplying two operands that are homogeneous nonnegative linear forms or positive scalar multiples of earlier signals. For nonnegative affine-linear $L_i$ and $c_{ij}\ge0$, suppose every monomial of
\[
 F=\sum_i g_iL_i+\sum_{i\le j}c_{ij}g_ig_j
\]
contains an $M$-edge. Then at most three matching edges have exact two-variable witnesses. If the three products form a literal chain
\[
 g_1=UV,\quad g_2=g_1W\text{ or }g_1^2,\quad
 g_3=g_2X,\ g_1g_2\text{ or }g_2^2,
\]
the bound is two. Both bounds are attained by abstract expressions of the stated type. Linear forms may coincide or vanish. An operand such as $A+g_i$ is outside this hypothesis unless it already equals one of the permitted operands.
\end{lemma}
\begin{proof}
Work with Boolean support functions, so positive coefficients can be replaced by one and squares folded into affine contexts. A safe conjunction of $t$ nontrivial affine-linear clauses witnessing $r$ matching edges needs $t\ge2^r$: each clause meets every witnessed edge and can reject at most one of their independent endpoint transversals. Constants and zero expressions only reduce the counts. We use the identity
\begin{equation}\label{eq:fourclauseidentity}
 xp\vee zq\vee xzr=(x\vee z)(x\vee q)(z\vee p)(p\vee q\vee r),
\end{equation}
which follows by evaluating the four pairs of truth values of $x,z$.

An antichain of products is Lemma~\ref{lem:threeindependentmatching}. For a chain, write $g_2=UVW_0$ and $g_3=UVW_0X_0$, where $W_0,X_0$ are the relevant linear forms or one. Nested products absorb into the later signal. Thus
\[
 F=UV(A\vee W_0B\vee W_0X_0C)
   =UV(A\vee W_0)(A\vee B\vee X_0)(A\vee B\vee C).
\]
There are at most five clauses, so at most two matching witnesses. The remaining strict ancestry types are one comparable pair, a fork, and a join.

For one comparable pair, write $x=g_1=UV$, $g_2=xW$, $z=g_3=XY$, allowing $W=1$ for a square. Folding nested products gives
\[
 F=x(A\vee WB)\vee zC\vee xz(\delta\vee\eta W),
 \qquad \delta,\eta\in\{0,1\}.
\]
Use~\eqref{eq:fourclauseidentity} with $p=A\vee WB$, $q=C$, $r=\delta\vee\eta W$. Its four factors expand into at most $4,2,4,2$ linear clauses, respectively. For the third count use $p=(A\vee W)(A\vee B)$. The last factor is one if $\delta=1$, a single clause $A\vee C\vee W$ if $\delta=0,\eta=1$, and otherwise $(A\vee W\vee C)(A\vee B\vee C)$. Thus there are at most twelve clauses.

For a fork, write $g_1=UV$, $g_2=g_1W$, $g_3=g_1X$, again allowing one for repeated products. The output has the form
\[
 F=UV(A\vee WB\vee XC\vee\delta WX).
\]
Apply~\eqref{eq:fourclauseidentity} to the last three terms and distribute the disjunction with $A$. There are at most four clauses inside the parentheses, hence at most six overall. For a join, $x=g_1=UV$, $z=g_2=WX$, $g_3=xz$, and
\[
 F=xA\vee zB\vee xzC=(x\vee z)(x\vee B)(z\vee A)(A\vee B\vee C).
\]
The respective clause counts are $4,2,2,1$, totaling nine. Each dependent type therefore has fewer than sixteen clauses and cannot witness four matching edges. This exhausts all three-element ancestry orders, including repeated signals and zero forms.

For sharpness of the chain bound, choose $U=a+c,V=a+d,W=b+c,X=b+d$ and output $g_3$; the four transversal clauses compute $ab\vee cd$. For the general bound, take
\[
 g_1=(a+c)(a+d),\qquad g_2=g_1(b+c),\qquad g_3=ef,
 \qquad F=g_2(b+d)+g_3.
\]
Its Boolean function is $ab\vee cd\vee ef$. These are sharpness examples for the algebraic expressions, not asserted realizations of a specified deletion budget.
\end{proof}

After exact full-carrier erasure, if every private-product operand has the literal form in the lemma, all other outside-carrier signals are linear forms plus nonnegative linear combinations of the $g_i$. Expanding the retained products gives the displayed expression after restriction to a missing matching. Thus the remaining four-matching profile with three deleted products must have a nonlinear-sum operand. This does not prove that forming that sum needs an additional gate beyond its possible role as an insertion cut, and does not settle the remaining three-addition configurations.

\begin{proposition}[Three products at cover gain three]\label{prop:gainthreeproducts}
Let $G$ be four-cycle-free and $P\in\Sup(G)$ have nonempty quadratic support $H$ missing exactly four $G$-edges. Suppose the extracted circuit has exactly $\vc(H)$ multiplication gates and $\vc(G)-\vc(H)\ge3$. Then at least three deleted gates are multiplications. Both retained and deleted addition counts are unrestricted.
\end{proposition}
\begin{proof}
Choose the distinct minimum-cover centers $C$ of the retained products and put $T=V(G)\setminus C$, $K=G[T]$. Since $C$ together with a $K$-cover covers $G$, we have $\vc(K)\ge3$. The graph $K$ has at most four edges, all missing from $H$. If it contains an induced three-edge matching, Corollary~\ref{cor:threedeletedproducts} proves the result.

Otherwise, apart from isolated vertices, $K$ is $P_4\sqcup K_2$ or $K_3\sqcup K_2$. Here $P_4$ is the path on four vertices. To check the classification, a connected four-edge tree has a bipartition class of size at most two; a triangle with a pendant edge likewise has a two-cover. The four-cycle is excluded. A three-edge component has cover two only when it is a path or triangle, and then a further isolated edge gives the two listed graphs. Three components with at most two edges each contain an induced three-edge matching, while two two-edge components have cover two. These exhaust the edge partitions. Both listed graphs have four edges and cover three. They therefore contain \emph{all} missing edges, so no missing edge meets $C$.

Every maximal independent set of either graph, including all isolated vertices, has complement a minimum three-vertex cover. For $P_4=a-b-c-d$ the maximal independent sets are $\{a,c\},\{a,d\},\{b,d\}$; for a triangle they are its single vertices, and the disjoint edge supplies one endpoint choice. Each listed graph has exactly six maximal independent sets.

Apply full-carrier erasure, preserving the exact frame, and restrict $C$ to zero while retaining all of $T$. Write $A_h,B_h$ for the full restricted center and leaf inputs of the product centered at $h\in C$. The old outside neighborhood $N_h=N_H(h)\setminus C$ is exactly the outside support of its old leaf form, because there is only one product at $h$. Its protected addition-only output path gives
\[
 A_h\wedge\bigvee_{v\in N_h}x_v\le f_K.
\]
We claim $A_h\le f_K$. Otherwise it is true on an independent set $I$. Every $v\in N_h$ must then have a $K$-neighbor in $I$, by evaluating the product on $I\cup\{v\}$. Extend $I$ to a maximal independent set $I'$. The cover $W=T\setminus I'$ has size three and contains $N_h$. Now $(C\setminus\{h\})\cup W$ covers all old edges and all missing edges, the latter being entirely in $K$. Its size $|C|+2$ contradicts the cover gain. This establishes the claim, including isolated outside variables.

Replace every retained product by its center input. After carrier erasure these inputs have no retained-product ancestors and remain fixed under the replacements. The output can only increase, and each newly inserted disjunct is at most $f_K$, so it stays equal to $f_K$. The resulting circuit has at most the original number $p$ of deleted conjunctions. The global conjunction-resolution argument in Proposition~\ref{prop:maxgainproducts} expresses it as a conjunction of at most $2^p$ positive clauses.

At least six clauses are necessary. Each of the six maximal independent sets must falsify a clause, and one valid clause cannot be false on two distinct such sets: it would also be false on their union, which contains an edge by maximality. Thus $2^p\ge6$, so $p\ge3$. The argument does not assert anything about a frame with an extra retained product.
\end{proof}

\begin{proposition}[One deleted addition]\label{prop:oneadditionall}
Let $G$ be four-cycle-free, and let a positive variable-input circuit compute $P\in\Sup(G)$ with nonempty quadratic support $H$. Suppose exactly one gate deleted by canonical minimum-degree extraction and pruning is an addition, and let $p$ be the number of deleted multiplication gates. If $r=|E(G)\setminus E(H)|\ge1$, those $r$ edges form an induced matching in $G$, and
\[
 p\ge2^r-2.
\]
If the unique unequal-degree addition retains an operand of minimum degree two, then $p\ge2^r-1$. Neither the number of retained gates nor the cover-number gain is restricted.
\end{proposition}
\begin{proof}
Unused original gates may first be removed; this can only decrease the number of deleted products. A missing witness requires a useful unequal-degree addition, because otherwise every useful gate would be homogeneous and the quadratic extraction would retain it. Every unequal-degree addition is deleted. Thus the sole deleted addition is the unique cut, say $a+g$, where the smaller minimum degree $d$ belongs to $a$. Before this cut all useful signals are homogeneous. Its selected low branch reaches the quadratic frame, so $d$ is one or two.

Every useful deleted product is ancestral to the high operand $g$. A useful product downstream of the cut with minimum degree two is retained if it reaches the output through additions. A larger minimum degree, or any later multiplication on such a route, would require another unequal-degree addition to reach the minimum-degree-two output. A private branch unrelated to the high operand would likewise need another cut. These possibilities are excluded.

The ancestor network of $g$, stopped at retained gates or variables, therefore contains only private multiplication gates. A retained quadratic boundary would make every monomial of $g$ contain an old $H$-edge. Since there is no private addition to provide another branch, every later multiple of $g$ would then contain an old edge and could not supply a missing-edge witness. This contradicts $r\ge1$. Hence every boundary is a nonzero old homogeneous linear form, with variables allowed as singleton forms. Every nonsingleton such support lies in an $H$-neighborhood: follow a retained linear route to a first quadratic product and choose any variable in its other nonempty operand. Positivity makes all the resulting pairs $H$-edges. No optimality hypothesis is used.

Let $p_g\le p$ be the number of products in this connected ancestor network and let $b$ be the number of its distinct boundary gates. Counting operand incidences gives $2p_g$ edges in a connected multigraph with $p_g+b$ vertices. Thus
\[
 b\le p_g+1\le p+1.
\]
Repeated input occurrences and shared subproducts are included in this count. The computed high signal is a product of positive powers of these boundary forms. In Boolean semantics repeated factors may be kept once; the resulting product has at most $b$ old-neighborhood factors and the same exact edge witnesses.

If $d=2$, every downstream path is addition-only, so the full output is $P=P_2+c g$ for some positive integer $c$. The product $g$ is safe, and its missing-edge witnesses are exactly the $r$ missing witnesses of $P$. Proposition~\ref{prop:oldfactors} makes them an induced matching and requires at least $2^r$ distinct factor supports. Hence $2^r\le b\le p+1$.

Suppose $d=1$. Temporarily replace the high operand by a symbol $z$ of minimum degree one. A retained linear signal affected by the cut has the form $A+c z$, where its old part $A$ contains $c a$ coefficientwise. Both operands of a retained product cannot have positive $z$-coefficients: selecting their low parts would give a positive multiple of $a^2$ in $P_2$, including a forbidden square. Thus no $z^2$ term occurs. All other downstream operations are additions, and expansion gives the exact identity
\[
 P=P_2+gD
\]
with a nonzero nonnegative homogeneous linear form $D$. Each contribution to $D$ is a positive multiple of the opposite, unmodified linear operand of a retained product. Fix any $h\in\supp(a)$. The old modified operand contains $h$, so every variable in each opposite operand is an $H$-neighbor of $h$. Consequently the \emph{entire} context satisfies
\[
 \supp(D)\subseteq N_H(h),
\]
even when the cut fans out to many products. This common center is the reason that the context counts as one old-neighborhood factor.

The product $gD$ is safe and supplies exactly all missing-edge witnesses. Remove repeated Boolean boundary factors from $g$ and apply Proposition~\ref{prop:oldfactors} to this product together with $D$. Boolean removal preserves safety and the exact missing witnesses, since the function is false on every independent assignment and on every singleton. The proposition again gives an induced matching, now with $2^r\le b+1\le p+2$. All bounds count original deleted gates, not occurrences in an expanded formula.
\end{proof}

\begin{proof}[Proof of Theorem~\ref{thm:oneaddition}]
Write $r=|E(G)\setminus E(H)|$ and $d=\vc(G)-\vc(H)$. Adding one endpoint of every missing edge to a minimum $H$-cover gives $d\le r$. The retained circuit has at least $|E(H)|+\vc(H)-1$ gates. If no addition is deleted, no useful higher branch can be discarded, so $r=0$ and this bound suffices. For $r\le2$, Theorem~\ref{thm:threemissing} already proves the assertion. For $r\ge3$, Proposition~\ref{prop:oneadditionall} gives at least $2^r-2$ deleted products and one deleted addition. Since $2^r-1\ge2r$ for every $r\ge3$, the original gate count is at least
\[
 |E(H)|+\vc(H)-1+2r
 \ \ge\ |E(G)|+\vc(G)-1.
\]
The exponential inequality holds at $r=3$ and persists by induction. The retained product count was not assumed minimal.
\end{proof}

\begin{proposition}[A common old neighbor]\label{prop:commonoldcenter}
In the setting of Proposition~\ref{prop:oldfactors}, suppose the safe product has at least three missing edge witnesses. There is one vertex $h$ adjacent in $H$ to every endpoint of all those missing edges. Consequently, if $G$ is triangle-free as well as four-cycle-free, every such product has at most two missing witnesses. A circuit with nonempty quadratic support and exactly one deleted addition then also has at most two missing edges, irrespective of its product count.
\end{proposition}
\begin{proof}
The missing witnesses form an induced matching $M$ by Proposition~\ref{prop:oldfactors}; write $r=|M|\ge3$ and let $W$ be its endpoint set. Every endpoint transversal $T$ is independent, so some factor support $S_T$ avoids it. That factor must meet every matching edge, because every edge is witnessed by the product. Therefore
\[
 S_T\cap W=W\setminus T.
\]
It contains $r\ge3$ vertices and hence has an old center $h_T$ with $S_T\subseteq N_H(h_T)$. This uses neither equality in the factor bound nor absence of extra support vertices.

If two transversals differ at exactly one pair, their complementary transversals have at least two distinct common vertices. Both old centers are adjacent to those vertices. Distinct centers would form a four-cycle with the two common neighbors; simplicity makes all four vertices distinct. Thus the centers coincide. The graph of endpoint transversals joined by one changed choice is a connected cube, so every $h_T$ is one vertex $h$. Every matching endpoint belongs to a complementary transversal, proving the claimed common neighborhood.

Simplicity ensures $h\notin W$. Each missing edge $uv$, together with the old edges $hu,hv$, therefore gives a triangle. This excludes three or more missing witnesses when $G$ is triangle-free. For the circuit assertion, Proposition~\ref{prop:oneadditionall} gives the safe product $g$ or $gD$ of old-neighborhood forms supplying exactly all missing edges; apply the product result.
\end{proof}

\begin{corollary}[The four-witness one-cut profiles]\label{prop:onecutgainthree}
With exactly four missing quadratic witnesses and one deleted addition, at least fourteen deleted gates are multiplications. In particular the earlier lower bound of six at cover gain at least three holds, with arbitrary retained multiplication and addition counts.
\end{corollary}
\begin{proof}
Apply Proposition~\ref{prop:oneadditionall} with $r=4$. No assumption about cover gain is needed for the stronger bound fourteen.
\end{proof}

\begin{lemma}[One shared nonlinear sum]\label{lem:sharednonlinearsum}
Let $M$ be a matching, with isolates allowed, let $Q=A+g_1$, and consider one of the exact templates
\[
\begin{array}{c|ccc}
 &g_1&g_2&g_3\\ \hline
 \text{fork}&UV&QW&QX\\
 \text{chain}&UV&QW&g_2X\\
 \text{join}&UV&WX&Qg_2.
\end{array}
\]
Here $A,U,V,W,X$ are nonnegative homogeneous linear forms, possibly zero. For affine-linear nonnegative $L_i$ and $c_{ij}\ge0$, if
\[
 F=\sum_i g_iL_i+\sum_{i\le j}c_{ij}g_ig_j
\]
is safe on $M$, it has exact witnesses for at most three matching edges. The one signal $Q$ may be shared between private operands and an insertion into the old frame; no additional gate is charged for that sharing.
\end{lemma}
\begin{proof}
Work Booleanly, put $x=UV$ and $Q=A\vee x$, and fold squares into affine contexts. For arbitrary $P,Z$,
\[
 xP\vee QZ=Q(x\vee Z)(P\vee Z).
\]
If two expressions have respectively $a,b$ linear clauses, their disjunction has at most $ab$ clauses by distribution.

In the fork, $xQ=x$, so the mixed terms involving $x$ enter its affine context. The output becomes
\[
 F=xP\vee QZ,\qquad Z=WB\vee XC\vee\delta WX,
\]
with $P,B,C$ affine-linear and $\delta\in\{0,1\}$. Identity~\eqref{eq:fourclauseidentity} gives at most four clauses for $Z$. The displayed factorization therefore gives at most $2+2\cdot4+4=14$ clauses for $F$.

In the chain, mixed products reduce to $xW,xWX,QWX$. Hence
\[
 F=xP\vee QZ,\quad P=P_0\vee\delta W\vee\eta WX,
 \quad Z=W(B\vee XC).
\]
The expression $P$ has at most two clauses: a $W$ term absorbs $WX$, and otherwise distribute $P_0\vee WX$. The expression $Z$ has the three clauses $W,B\vee X,B\vee C$. Thus $F$ has at most $2+2\cdot3+2\cdot3=14$ clauses.

For the join put $y=WX$. After folding absorbed terms, write
\[
 F=xP\vee yR\vee xyS,\qquad R=B\vee AC,\qquad S=C\vee\delta.
\]
Identity~\eqref{eq:fourclauseidentity} gives factors $x\vee y$, $x\vee R$, $y\vee P$, and $P\vee R\vee S$. These have at most $4,4,2,1$ linear clauses: use $R=(B\vee A)(B\vee C)$ and the absorption $P\vee R\vee S=P\vee B\vee C\vee\delta$. The total is eleven.

Each case has fewer than sixteen clauses, whereas four matching witnesses require sixteen distinct clauses by the independent-transversal bound. Zero forms, constants, shared variables and repeated factors only remove clauses or witnesses. This proves the assertion for the exact templates displayed.
\end{proof}

The lemma does not normalize every connected three-product ancestry into one of these forms. A mixed pair with an independent third product, two distinct nonlinear sums, and operands containing sums of earlier nonlinear signals still require separate arguments. In particular, neither repeated use of $Q$ nor its use as an insertion is counted as a separate addition here.

\begin{lemma}[A mixed pair and an independent product]\label{lem:mixedpairfull}
Let $M$ be a matching, with isolates allowed, and put
\[
 g=UV,\qquad h=(A+g)W,\qquad z=XY,
\]
where $U,V,A,W,X,Y$ are nonnegative homogeneous linear forms. Let $P,B,C$ be nonnegative affine-linear forms and $\delta,\eta\ge0$. If
\[
 F=gP+hB+zC+\delta gz+\eta hz
\]
is safe on $M$, then it witnesses at most three matching edges. Both cross coefficients are allowed, and the abstract bound three is attained. Zero forms, shared variables and repeated inputs are permitted. The sum $A+g$ is one shared signal; no separate gate is charged for its use in several roles.
\end{lemma}
\begin{proof}
Work Booleanly, discard zero terms and replace positive coefficients by one. Write
\[
 T=gWB,\qquad S=gP\vee AWB\vee T,\qquad
 R=\delta gz\vee\eta AWz\vee\eta gWz.
\]
Then $F=S\vee zC\vee R$, and every displayed summand is safe. Lemma~\ref{lem:twoproductmatching}, applied directly to the mixed pair $g,(A+g)W$, gives at most two witnesses for $S$. It also gives at most two for $AWB\vee zC\vee\eta AWz$.

We recall two support consequences of that lemma's proof. Call a product \emph{potential} on an edge when it is true on that edge's endpoint assignment, without assuming that the product is safe. If a two-clause product $IJ$ is potential on two disjoint edges $e,f$ and $IJD$ is safe with affine $D$, then every variable in $D$ belongs to their four endpoints, and $D$ has no positive constant. Indeed, two clauses cannot reject all four independent endpoint transversals of $e,f$. An outside context variable, or a constant, would turn a surviving transversal into an independent true assignment. Thus every witness of $IJD$ belongs to $e$ or $f$.

The cubic rigidity proved in Lemma~\ref{lem:twoproductmatching} is stronger: if $IJD$ safely witnesses $a$ and $IJ$ is potential on a different edge $t$, then, after naming the endpoints,
\[
 I=a_0\vee t_0,\qquad J=a_0\vee t_1,\qquad D=a_1,
 \qquad IJ=a_0\vee t.
\]
In particular this product is potential only on $a,t$. The same statement permits affine factors: a constant factor would reduce the safe product to two clauses, which cannot have the stipulated distinct witness and potential edge. We also use the four-factor equality case: two witnesses confine all four factors to their four endpoint transversals.

First, $R\vee zC$ has at most two witnesses. The possible nonzero forms of $R$ are
\[
 UVXY,\qquad XYW(A\vee U)(A\vee V),\qquad
 XY(U\vee A)(U\vee W)(V\vee A)(V\vee W).
\]
The last uses absorption when both coefficients are positive. These have respectively four, five and six clauses, so none witnesses three edges. If $R$ witnesses two, $z$ is potential on both and the preceding support observation confines every witness of $zC$ to them. If $R$ has at most one witness, the cubic $zC$ adds at most one.

\emph{The term $T$ has two witnesses.}
Its four-factor equality confines $U,V,W,B$ to those endpoints. Every term of $F$ except $zC$ contains $g$ or $W$, so has no outside witness. The cubic $zC$ adds at most one, proving the bound in this case.

\emph{The term $T$ has one witness $t$.}
Suppose first that $gP$ witnesses $a\ne t$. Cubic rigidity gives $g=a_0\vee t$ and $P=a_1$. Every term containing $g$ witnesses only $a,t$. If four witnesses existed, the remaining expression $AWB\vee zC\vee\eta AWz$ would have to supply two edges $c,d$ outside them. The term $AWB$ supplies neither, because $S$ already has $a,t$ and has at most two witnesses. If $AWz$ supplies both $c,d$, its four-factor equality confines $W$ to their endpoints, although $W$ is true on $t$, a contradiction. Otherwise $AWz$ witnesses, say, $c$ and $zC$ witnesses $d$. Then $\eta>0$ and rigidity gives $z=d_0\vee c$. The safe summand $gWz$ contains $a_0d_0W$. Hence every variable of $W$ must be $a_1$ or $d_1$; any other choice gives an independent monomial with $a_0,d_0$. This contradicts $W$ being true on $t$.

Next suppose $AWB$ witnesses $b\ne t$. Apply rigidity to the pair $W,B$, potential on $t$, with outer context $A$. It gives $WB=b_0\vee t$ and $A=b_1$. Safety of $gWB$ implies safety of $gb_0$. If a term containing $g$ witnessed an edge $c$ outside $b,t$, then $g$ would be potential on both $t,c$, contradicting the support observation with the outside context variable $b_0$. Thus all $g$-terms witness only $b,t$. The terms $AWB,AWz$ witness only $b$, because $A=b_1$; $zC$ adds at most one. If neither $gP$ nor $AWB$ has a witness different from $t$, then $S$ has only $t$, and the already proved bound on $R\vee zC$ finishes this case.

\emph{The term $T$ has no witness.}
If $gP\vee AWB$ witnesses at most one edge, the bound on $R\vee zC$ again suffices. Otherwise let their distinct witnesses be $a,b$. Four total witnesses require $R\vee zC$ to supply two further edges $c,d$.

Suppose $R$ witnesses both. Split it as
\[
 R=z g(\delta\vee\eta W)\vee\eta zAW.
\]
Rigidity applied to $gP$ shows that $g$ cannot be potential on both $c,d$: either potential edge already forces $g=a_0\vee c$ or $a_0\vee d$. Likewise $AW$ cannot be potential on both. Thus, after exchanging $c,d$, the first summand supplies $c$ and the second supplies $d$. In particular $\eta>0$, and rigidity gives
\[
 g=a_0\vee c,\qquad AW=b_0\vee d,
\]
with $b_0$ in $W$. There is an independent endpoint transversal of $c,d$ meeting both $X,Y$, because two clauses cannot reject all four such transversals. Add $a_0,b_0$ to it. This is still independent, but makes the safe summand $gWz$ true, a contradiction.

It remains that $R$ has one witness $c$ and $zC$ supplies $d$. Rigidity gives $z=d_0\vee c$. If $c$ comes from $\eta zAW$, then $AW=b_0\vee c$, and the safe term $AWz$ is true on the independent pair $\{b_0,d_0\}$, impossible. Thus $c$ comes from $zg(\delta\vee\eta W)$ and $g=a_0\vee c$. A positive $\delta$ makes $gz$ true on $\{a_0,d_0\}$, again impossible. Otherwise $\eta>0$, and safety of $gWz$ forces $\supp(W)\subseteq\{a_1,d_1\}$. But $AWB$ witnesses the different edge $b$, so $W$ must meet it. This final contradiction exhausts the cases.

For sharpness, take matching edges $ab,cd,ef$ and set
\[
 A=P=0,\quad U=a+c,\quad V=a+d,\quad W=b+c,
 \quad B=b+d,\quad z=ef,\quad C=1,\quad\delta=\eta=0.
\]
The resulting expression is the four transversal clauses for $ab\vee cd$, together with $ef$, and witnesses exactly those three edges. This is sharpness of the algebraic template, without any specified circuit-deletion budget.
\end{proof}

This lemma handles the displayed mixed pair with an independent third product. It does not normalize every three-product circuit to that form. Two distinct nonlinear sums, or sums of earlier nonlinear signals, remain outside the stated hypotheses.

\begin{proposition}[The exceptional outside graphs]\label{prop:exceptionaloutside}
Assume the hypotheses of Proposition~\ref{prop:gainthreeproducts}: exactly four edges are missing, the cover gain is at least three, and the extracted frame has exactly $\vc(H)$ products, with arbitrary retained additions. Let $C$ be its distinct minimum-cover centers. If $K=G[V(G)\setminus C]$, apart from isolates, is $P_4\sqcup K_2$ or $K_3\sqcup K_2$, then extraction deletes at least three original addition gates. The number of deleted multiplications is unrestricted.
\end{proposition}
\begin{proof}
The proof of Proposition~\ref{prop:gainthreeproducts} shows that all four missing edges lie outside $C$. After full-carrier erasure and restriction $C=0$, every full center input is safe on $K$, and replacing each retained product by that input preserves $f_K$. Keep outside isolated vertices. Every selected center form is a positive multiple of its center variable, which now vanishes. Its actual linear-input network is an addition network of old terms and the high signals inserted at private unequal-degree cuts. Consequently the outside output is the OR of a nonempty subset of the cut-high signals, including highs entering the quadratic output carrier directly. Call a cut \emph{contributing} when its high occurrence reaches this final disjunction. Reuse of its low variable alone does not make it contributing.

Suppose at most two original additions are deleted. We show that the outside output is a disjunction of at most two products of old linear forms. A product of restricted old forms still satisfies the neighborhood hypothesis in the original pair $H\subseteq G$; restriction only shrinks its supports. Retained quadratic carriers have no private outgoing uses after erasure. All changes below are Boolean support identities, without recounting gates in a new extraction.

With no contributing cut the outside output would be zero. A sole useful cut, with no other private addition in its high ancestry, has a multiplication-only high computation on old linear boundaries, and so gives one product. Zero branches or irrelevant additions only simplify the following two-addition cases.

\emph{One unequal cut and a private homogeneous addition.}
The private addition must lie upstream in the useful cut's high operand. After that cut, a useful degree-two addition on an output route would be retained, while a higher-degree route would need another unequal cut. Before the cut there is no earlier inhomogeneity. Thus the high computation is multiplicative on old forms except for one sum $U+V$, whose operands are products of old forms. Repeated uses of that sum collapse Booleanly, giving a high of the form $R(U\vee V)$, where $R$ is also an old-form product. After the center replacements only that high can contribute. Expansion gives the two old-form products $RU,RV$.

\emph{Two unequal cuts.}
Order their high operands as $g,f$; incomparable cuts have independent old-product highs and are immediate. Before the first cut there is no private addition, so $g$ is a product of old linear forms. There is no private addition between the two cuts. Each retained linear carrier affected by the first has form $D_i+\lambda_i g$, where $D_i$ is its unchanged old linear form. Booleanly a positive coefficient gives $D_i\vee g$. A private product depending on this cut cannot reenter a retained degree-one network without another unequal cut; the only remaining private addition is the second cut receiving $f$.

If the first cut's low part is itself private inside the high ancestry of the second, it is a product $B$ of old forms. A low degree above one cannot enter a retained linear network, so the cut $B\vee g$ has only multiplicative uses there. If its low part is retained and quadratic, the full cut is a protected quadratic carrier and its private uses were erased; direct uses of $g$ may remain. These cases include low branches not selected in the final frame.

It follows that the high computation for $f$ is multiplicative on old forms, proper product ancestors of $g$, possible direct copies of $g$, and factors $D_i\vee g$. Here the $D_i$ are old linear forms, except that a private low branch may give the old-form product $B$. Proper ancestors of $g$ are themselves products of old forms. If a direct $g$ factor occurs, it absorbs all affected factors and $f=gR$ for an old-form product $R$. Otherwise, for a nonempty affected list,
\[
 f=R\prod_i(D_i\vee g)
   =R\left(\prod_iD_i\right)\vee Rg.
\]
This exact identity allows different $D_i$, arbitrary fanout, repeated inputs and shared subproducts. If no affected boundary occurs, $f$ is one old-form product. Empty factor products mean one in this Boolean analysis, and a constant true output disjunct is impossible for $f_K$.

After the second cut, no useful private signal can return to the quadratic frame without a further unequal cut. Retained center networks therefore contribute only combinations of the two highs $g,f$. If only $f$ contributes, the displayed expansion already has at most two old-form products. If $g$ also contributes, its disjunct absorbs $Rg$, leaving
\[
 g\vee R\prod_i D_i.
\]
The case $f=gR$ is even simpler. A hidden addition after a contributing cut cannot affect that earlier high and has no further output contribution unless its own high contributes. Thus every temporal and contribution pattern yields at most two products of old forms. No multiplication bound was used.

Each nonzero resulting product is safe because it is a positive disjunct of $f_K$. Every $K$-edge is an exact witness of one of them, by evaluation on its two endpoints and falsity on singletons. These products use only outside variables, so safety on $K$ is safety on $G$ for their monomials. Their witnessed edges are all outside $H$. Proposition~\ref{prop:oldfactors} therefore makes each witness set an induced matching in $G$.

An induced matching contains at most one of the three edges of a triangle. The same holds for the three-edge path: adjacent edges meet, and its two end edges are joined by the middle edge. Two such witness sets cannot cover that component of $K$. This contradicts the representation of $f_K$ and proves that at least three additions were deleted. The minimum retained product count and the exceptional outside graphs are still required for the initial safe-center replacement.
\end{proof}

\begin{corollary}\label{cor:exceptionalbounded}
Under the hypotheses of Proposition~\ref{prop:exceptionaloutside}, it is impossible to delete at most three products and at most two additions.
\end{corollary}
\begin{proof}
The proposition excludes two additions for every product count. In particular, it excludes the stated bounds on both kinds of gate.
\end{proof}

In the remaining cover-gain-three profile with five deleted gates and the minimum retained product count, the outside graph is consequently, apart from isolates, either $3K_2$ or $P_3\sqcup2K_2$. To verify this restriction, its cover number is at least three and it has at most four edges. The component enumeration in Proposition~\ref{prop:gainthreeproducts} leaves the two exceptional graphs, these two graphs, and $4K_2$. The proposition just proved removes the exceptional graphs, while Lemma~\ref{lem:fourmatchinglost} removes $4K_2$ with only five deletions. In the $3K_2$ case the fourth missing edge meets $C$; in the $P_3\sqcup2K_2$ case all four missing edges lie outside $C$. This restriction also permits extra retained additions. Proposition~\ref{prop:paththreeproducts} below excludes the path shape. Propositions~\ref{prop:adjacentfourth} and~\ref{prop:disjointfourth} then exclude every placement of the fourth missing edge relative to the three-edge matching; Proposition~\ref{prop:gainthreesix} assembles the resulting six-deletion bound.


\begin{lemma}[A two-sided first dependency]\label{lem:twosidedfirst}
Let $\mathcal M$ be a matching, with isolates allowed, and put
\[
 g_1=UV,\qquad g_2=(A+g_1)(B+g_1),
\]
\[
 g_3=(C+c_1g_1+c_2g_2)(D+d_1g_1+d_2g_2),
\]
where $U,V,A,B,C,D$ are nonnegative homogeneous linear forms, zero forms are allowed, and the four scalar coefficients are nonnegative. Let $L_i$ be nonnegative affine linear forms and $b_{ij}\ge0$. If
\[
 F=\sum_{i=1}^3g_iL_i+\sum_{1\le i\le j\le3}b_{ij}g_ig_j
\]
is safe on $\mathcal M$, it witnesses at most three matching edges. The algebraic bound three is attained. No old-neighborhood or gate-budget hypothesis is required.
\end{lemma}
\begin{proof}
Work Booleanly. Write $x=g_1=UV$, $y=AB$ and $z=g_3$. Absorption gives $g_2=x\vee y$. Folding diagonal terms into contexts, and using $x(x\vee y)=x$, puts the output in the form
\[
 F=xP\vee yQ\vee zR\vee\alpha xz\vee\beta yz,
\]
with affine $P,Q,R$ and nonnegative flags $\alpha,\beta$. Allowing the two flags independently enlarges the class. Each operand of $z$ contains its old linear form and either no earlier signal, $x$, or $x\vee y$.

If neither operand contains $x$, then $z=CD$ and Lemma~\ref{lem:threeindependentmatching} applies directly. If exactly one contains $x$, write
\[
 z=(C\vee x\vee\epsilon y)D.
\]
For $\epsilon=0$, the identity $xz=xD$ absorbs its coefficient into $P$; Lemma~\ref{lem:mixedpairfull} applies to the mixed pair $x,z$ and independent product $y$. For $\epsilon=1$, both $xz=xD$ and $yz=yD$. Absorb their coefficients into $P',Q'$ and set $k=DR$. The output is
\[
 xP'\vee yQ'\vee kC\vee xk\vee yk.
\]
If $R$ has no positive constant, $k$ is a product of homogeneous linear forms, possibly zero, so Lemma~\ref{lem:threeindependentmatching} applies. If $R$ has a positive constant, it is Booleanly one and $k=D$. If $C$ or $D$ vanishes, the expression is a union of at most two safe cubics, each witnessing at most one matching edge. Otherwise safety of $CD$ forces $C,D$ to be opposite singleton endpoint forms of one matching edge: every cross pair must be an edge, and each matching vertex has only one neighbor. Every witness involving $D$ is then that edge. The remaining cubics $xP',yQ'$ add at most one each. This handles the affine boundary without treating $D$ as a quadratic polynomial.

If both operands contain $x$, write
\[
 z=(C\vee x\vee\epsilon y)(D\vee x\vee\eta y)=x\vee t,
 \qquad t=(C\vee\epsilon y)(D\vee\eta y).
\]
The output becomes
\[
 x(P\vee R\vee\alpha)\vee yQ\vee tR
       \vee\beta xy\vee\beta yt.
\]
If $\epsilon=\eta=0$, then $t=CD$ and the three-independent-product lemma applies. If exactly one flag is positive, say $t=(C\vee y)D$, then $yt=yD$ is absorbed into $Q$, and Lemma~\ref{lem:mixedpairfull} applies to $y,t$ and the independent product $x$. If both flags are positive, then $t=y\vee CD$ and $yt=y$; the expression reduces to three independent products $x,y,CD$, with affine contexts and the possible cross term $\beta xy$. Again there are at most three witnesses. These are Boolean equalities, so safety and exact witnesses are preserved, including zero forms, constants and repeated factors.

For attainment, take matching edges $ab,cd,ef$, let $g_1=ab$, $A=c$, $B=d$, and $g_3=ef$ with its earlier-signal coefficients zero. The output $g_2+g_3$ Booleanizes to those three edges. This is algebraic sharpness, with no prescribed private-gate budget.
\end{proof}

The second signal in this lemma has the displayed dependency in both operands. The lemma does not replace a classification of the first one-sided dependency or a general circuit normalization.

\begin{lemma}[An independent pair on a path and two edges]\label{lem:pathindependentpair}
Let $H\subseteq G$, with $G$ simple and four-cycle-free. Suppose $H[T]$ is empty and $G[T]$, retaining isolates, has precisely the edges
\[
 vx,\quad vy,\quad ab,\quad cd
\]
on seven distinct named vertices. Call a homogeneous linear form on $T$ old if every nonsingleton support lies in one $H$-neighborhood; singleton and zero forms are allowed. Let $g$ be a product of at most three old forms and $f$ a product of at most two. Let $L,M$ be nonnegative affine forms with old homogeneous parts. A safe expression
\[
 F=gL+fM+qgf,\qquad q\ge0,
\]
cannot witness all four edges of $G[T]$.
\end{lemma}
\begin{proof}
Write $E_A,E_B,E_C$ for the exact edge witnesses of $gL,fM,qgf$. Each nonzero summand is safe. A positive constant in a context makes its underlying high product safe, and every edge witnessed by a multiple of a safe product is already witnessed by that product. Proposition~\ref{prop:oldfactors} therefore gives
\[
 |E_A|\le2,\qquad |E_B|\le1,\qquad |E_C|\le2.
\]
If $gL$ witnesses two edges, its four-factor equality case applies: the three factors of $g$ and $L$ are exactly the four two-element transversals of that induced matching. The three supports in $g$ have empty common intersection and lie in its four endpoints. An edge outside those two cannot meet all three supports: an outside endpoint leaves only one endpoint in them, while the induced graph on their four vertices contains only the matching edges. Hence $E_C\subseteq E_A$, and the union has at most three edges.

If all four edges were witnessed, it would follow that $E_A,E_B,E_C$ have respectively one, one and two elements, pairwise disjoint. Write $E_C=\{e_1,e_2\}$, $E_B=\{e_4\}$ and $E_A=\{e_3\}$. Every factor of $f$ meets the two disjoint edges $e_1,e_2$, so none is singleton. The context $M$ has no positive constant: otherwise $f$ itself is safe and its two-factor equality case cannot be potential on both edges. Fewer than two factors in $f$ are equally impossible. Thus $f=XY$ and the cubic singleton-factor lemma makes $M$ the singleton $\{w\}$, where $e_4=tw$.

Both $X,Y$ contain $t$ and omit $w$. Indeed, if a witness used $w$ in either factor, then together with $M$ it would repeat $w$, forcing the other factor to be singleton $\{t\}$. The same argument excludes any other occurrence of $w$. Safety forces $w$ to have a neighbor among the endpoints of $e_1,e_2$: otherwise choose an endpoint of $e_1$ in $X$ and of $e_2$ in $Y$; these two vertices and $w$ form an independent true set. In the stipulated graph, $e_4$ must consequently be one path arm and one of $e_1,e_2$ the other. Relabel so that
\[
 e_4=vy,\qquad\{e_1,e_2\}=\{vx,ab\},\qquad e_3=cd.
\]
If $w=v$, both $X,Y$ contain $y$ and omit $v$. Potential on $vx$ forces both to contain $x$. An old form containing $x,y$ has an $H$-neighbor $h$ outside $T$, and $h,x,v,y$ form a four-cycle. Hence $w=y$, so $t=v$ and $f$ contains $v^2$.

Set $x=y=0$ in $g$, giving $g'$. Every monomial $m$ of $g'$ yields the safe monomial $mv^2$ in $gf$. Since $v$ has no remaining neighbor, $m$ must itself contain $ab$ or $cd$. All isolated variables remain and cannot repair an independent support. The product $g'$ is true on both $ab$ and $cd$, by the $gf$- and $gL$-witnesses. Its safety therefore gives exact witnesses for both. But it has at most three old factors, whereas Proposition~\ref{prop:oldfactors} requires four. This contradiction proves the assertion. Omitted factors, zero terms and affine constants only give the cases already considered.
\end{proof}

This is a restriction on the displayed independent-factor template. It makes no assertion that all three-product circuits have that form.


\begin{proposition}[The path and two-edge circuit boundary]\label{prop:paththreeproducts}
Let $G$ be four-cycle-free and let a positive variable-input circuit compute $P\in\Sup(G)$ with nonempty quadratic support $H$. Suppose $H[T]$ is empty and $G[T]$, apart from isolates, has edges $vx,vy,ab,cd$ on seven distinct vertices. Canonical extraction cannot delete exactly three products and two additions. No cover-gain or retained-count hypothesis is required.
\end{proposition}
\begin{proof}
Write $K=G[T]$ and keep its isolates. Perform exact full-carrier erasure before restricting variables. Lemma~\ref{lem:fivelost} requires at least five deletions for four missing witnesses, so none of the five private gates can disappear. Restrict outside $T$ to zero; the output is $f_K$ and the quadratic part vanishes. This is a semantic restriction, not a new claim about extraction of the restricted circuit. Restricted old forms remain old in the original pair $H\subseteq G$.

The function $f_K$ needs eight positive linear CNF clauses. Its minimal vertex covers choose either $\{v\}$ or $\{x,y\}$, and one endpoint of each isolated edge. For each such cover $C$, the assignment false on $C$ and true on its complement must falsify a clause. That clause is itself a cover contained in $C$, so minimality makes it exactly $C$. Thus all eight clauses are necessary; constants do not help.

Apply the selected-occurrence classification in Lemma~\ref{lem:fourmatchinglost}. One visible cut gives at most seven clauses. When both cuts are visible, give their highs independent symbols $g,f$ before substituting their dependence. The output is
\[
 F=gL\vee fM\vee qgf,\qquad q\ge0,
\]
with affine old contexts. Let $k$ count private products ancestral to the first high $g$. The $k=3$ case has at most five clauses, and the $k=2$ direct-private-ancestor cases have at most six. These are impossible. We treat the remaining cases using the same classification, without changing its hypotheses.

A safe old product of at most three factors witnesses at most one missing edge; four or five factors witness at most two. A two-witness four-factor product has precisely the four transversal supports. For this particular $K$, any edge outside such a pair that meets its endpoint set must be the other path arm. There is at most one such edge; another remaining edge is disjoint. A product containing three of the four transversal factors cannot witness even that other arm, since their common intersection is empty. A positive constant in a safe context makes the underlying product safe; its multiples then have no additional edge witnesses.

\emph{The remaining $k=2$ cases.}
The high $g$ has at most three old factors, and $f$ is $AB$, $AB\vee gB$, or $AB\vee g$, allowing zero old parts and direct-$g$ degenerations. The independent case is Lemma~\ref{lem:pathindependentpair}. If $f=AB\vee g$, then
\[
 F=gL\vee ABM\vee gM\vee qg.
\]
A positive $q$ or constant in either context makes $g$ safe, leaving at most one witness from its multiples and one from the cubic $ABM$. Otherwise $gL,gM$ each have at most four factors. If either witnesses two edges, its three $g$-factors confine every other $g$-witness to the same pair. If neither does, the three summands each have at most one witness. Both alternatives give at most three.

If $f=AB\vee gB$, then
\[
 F=gL\vee ABM\vee gBM\vee qgB.
\]
For $q=0$ and a context $M$ without a positive constant, regroup as $gL\vee(BM)A\vee g(BM)$ and apply Lemma~\ref{lem:pathindependentpair}. If $q>0$ or $M$ has a positive constant, $gB$ is safe and its multiples add no witnesses. The terms $gL,gB$ share all factors of $g$, so the same four-factor confinement argument, with the cubic $ABM$, gives at most three. This exhausts $k=2$.

\emph{The $k=1$ cases.}
Now $g=UV$. An unaffected second high is an old product of at most three factors, so the independent-pair lemma applies with the roles exchanged. Otherwise the admitted classification gives
\[
 f=DE\vee gD,
\]
where $D$ has at most two old factors, $DE$ at most three, and $D$ is a subproduct of $DE$. Empty products and zero old parts are allowed. Expansion and absorption give
\begin{equation}\label{eq:pathmixed}
 F=gL\vee DEM\vee gDM\vee qgD.
\end{equation}
If $D=1$, a positive $q$ or constant in $L$ or $M$ makes $g$ safe, so all $g$-terms have at most one witness while $EM$ has at most two. Otherwise $q=0$ and the two $g$-based cubics each have at most one. Four total witnesses would force $EM$ to attain four-factor equality. Its context $M$ is then a nonsingleton transversal. The outside witness of $gM$ can only be the other path arm, while that of $gL$ is the disjoint remaining isolated edge. Both factors $U,V$ meet these two disjoint edges, so neither is singleton. The safe cubic $UVM$ would have an exact missing-edge witness with all three factors nonsingleton, contrary to the cubic singleton lemma.

Suppose $D$ is nonempty. If it has two factors, $q=0$, and $M$ has no positive constant, regroup~\eqref{eq:pathmixed} as
\[
 (DM)E\vee gL\vee g(DM).
\]
Here $DM$ and $g$ have at most three and two factors, and $E$ is one old factor or a unit/zero degeneration. Lemma~\ref{lem:pathindependentpair} applies.

In every remaining case put $A_0=gL$, $B_0=DEM$, and let $C_0$ combine the last two terms in~\eqref{eq:pathmixed}. The first has at most one witness and the second at most two. The witnesses of $C_0$ are those of a safe product of at most four old factors: it is $gD$ if $q>0$ or $M$ has a positive constant, and otherwise it is $gDM$ with $D$ a single factor.

If $C_0$ witnesses two edges, four-factor equality confines all factors of $g$ and $D$ to their four endpoints. Both $A_0$ and $B_0$ contain such confined factors. Neither can witness an outside edge except possibly the same other path arm. Thus their whole union has at most three witnesses. Consequently four witnesses would require $B_0$ to witness two, and $A_0,C_0$ to witness the two distinct remaining edges.

Equality for $B_0$ makes its $DE$ and $M$ factors the four transversals. Since $D$ is nonempty, the outside witness of $C_0$ must be the other path arm; the $A_0$ witness is the remaining isolated edge. Relabel so that $B_0$ witnesses $vx,ab$, $C_0$ witnesses $vy$, and $A_0$ witnesses $cd$. The two factors of $g$ meet both $vy$ and $cd$, so neither is singleton. The context $L$ cannot have a positive constant, since that would make the two-factor product $g$ safe and potential on both edges. The cubic singleton lemma applied to $gL$ therefore makes $L$ one endpoint of $cd$ and puts the other endpoint $t$ in both $g$-factors. Thus $g$ contains $t^2$.

Every factor of $D$ is confined to $\{v,x,a,b\}$ and must contain $v$ in order for $C_0$ to witness $vy$. If $C_0$ uses $M$, that transversal must also contain $v$; otherwise the term $qgD$ is already present. Selecting $t$ in both $g$-factors and $v$ in the required remaining factors gives a true monomial supported on the independent pair $\{t,v\}$. This contradiction exhausts $k=1$ and proves the proposition.
\end{proof}

Consequently, with the minimum retained product count, cover gain three and five deletions, the outside graph can only be $3K_2$, with the fourth missing edge meeting the old cover. The preceding outside enumeration and Proposition~\ref{prop:exceptionaloutside} remove the other shapes, and the present proposition removes $P_3\sqcup2K_2$. The remaining shape is not excluded. Because the present proposition does not assume a retained-count bound, its path-shape exclusion also applies to a frame with an extra retained product.

\begin{lemma}[Two different low forms]\label{lem:differentlowfork}
Let $\mathcal M$ be a matching with isolates, and set
\[
 x=UV,\qquad y=(A+x)B,\qquad z=(C+x)D,
\]
where all six forms are nonnegative homogeneous linear, possibly zero. For nonnegative affine linear $P,Q,R$ and $\delta\ge0$, a safe expression
\[
 F=xP+yQ+zR+\delta yz
\]
has at most three exact matching-edge witnesses. The same conclusion permits all scalar quadratic terms in $x,y,z$, since diagonal terms and $xy,xz$ can be absorbed into affine contexts Booleanly. The low forms $A,C$ are unrelated.
\end{lemma}
\begin{proof}
Work Booleanly. If $\delta=0$, put $v=BQ$, $w=DR$. Then
\[
 F=xP\vee vA\vee wC\vee xv\vee xw,
\]
and Lemma~\ref{lem:threeindependentmatching} applies. If a context has a positive constant, its corresponding product is Booleanly a linear form; replacing that form by its square fits the homogeneous-factor statement. This is a support identity, not an added circuit gate.

Suppose $\delta>0$. Expansion gives $F=F_0\vee F_1$, where
\[
 F_0=ABQ\vee CDR\vee ABCD,\qquad
 F_1=x(P\vee BQ\vee DR\vee BD).
\]
All summands are safe. Lemma~\ref{lem:twoproductmatching}, applied to independent products $AB,CD$, gives at most two witnesses for $F_0$. The identity
\[
 BQ\vee DR\vee BD=(B\vee D)(B\vee R)(D\vee Q)
\]
makes $F_1$ a product of at most five affine clauses, after adjoining $U,V$ and distributing $P$. Hence it too has at most two matching witnesses. Four witnesses would force disjoint two-edge sets $E,J$ for $F_0,F_1$.

Index $BQ,DR,BD$ by $1,2,3$. The corresponding $F_0$ components are $ABQ,CDR,ABCD$; the $F_1$ components are $xBQ,xDR,xBD$, with $xP$ additionally indexed by zero. The first two $F_0$ components are cubic. If the third supplied both $E$-edges, four-factor equality would confine $A,B,C,D$ to their endpoints. Every $F_1$ component except $xP$ would then have no $J$-witness, and that cubic supplies at most one. Thus $E$ uses at least two distinct indices in $\{1,2,3\}$, with no one component supplying both.

No component $xBQ,xDR,xBD$ can supply both $J$-edges either. Four-factor equality for the first confines $B,Q$, blocking $F_0$ indices 1 and 3 and leaving one cubic; the second is symmetric, while the third confines $B,D$ and blocks every $F_0$ component. Nor can $xP$ supply one $J$-edge. If $xKL$ supplied the other, the safe union
\[
 xP\vee xKL=UV(P\vee K)(P\vee L)
\]
would attain four-factor equality. It confines $K,L$ to the $J$-endpoints. For indices 1,2,3 this confines respectively $B$, $D$, or both, leaving at most one $F_0$ witness. Therefore $J$ also uses at least two distinct indices in $\{1,2,3\}$, without index zero.

Indices 1 and 2 cannot be shared by the two sets. For index 1, suppose $ABQ$ witnesses $e\in E$ and $xBQ$ witnesses an edge in $J$. Cubic rigidity makes $BQ$ contain a pure monomial at one endpoint $a$ of $e$: the remaining factor $A$ is its opposite singleton. A positive constant in $Q$ would instead force a safe two-factor product and contradict potential on the disjoint $J$-edge. Safety of $xBQ$ now implies safety of $ax$. But $U,V$ are each nonzero on both $J$-edges. Choose $U$ on one and $V$ on the other; together with $a$ this is an independent true support, a contradiction. The argument for index 2is identical.

Two subsets of $\{1,2,3\}$ of size at least two must meet. Since their intersection avoids 1 and 2, the only possibility, up to symmetry, is that $E$ uses $\{1,3\}$ and $J$ uses $\{2,3\}$. The safe union of the two $J$ components is
\[
 xDR\vee xBD=UVD(R\vee B).
\]
Its two witnesses force four-factor equality and confine $D$ to the $J$-endpoints. But the $E$ component at index 3contains $D$, so cannot witness an edge disjoint from those endpoints. The symmetric case confines $B$. This final contradiction proves the bound.

Zero forms and affine constants only remove terms or clauses in the preceding cases. For the quadratic-output extension, use $x^2=x$, $y^2=y$, $z^2=z$, $xy=xB$ and $xz=xD$; diagonal coefficients add constants to contexts, and the last two crosses add linear forms to $P$. Only $yz$ remains as a separate cross term.
\end{proof}

\begin{theorem}[Three recursive products on a matching]\label{thm:threerecursive}
Let $\mathcal M$ be a matching with isolates. For nonnegative homogeneous linear forms, allowing zero, define
\[
 x=UV,\qquad y=(A+ax)(B+bx),
\]
\[
 z=(C+cx+dy)(D+ex+fy),
\]
with all six scalar coefficients nonnegative. For nonnegative affine linear forms $P,Q,R$ and arbitrary nonnegative scalar coefficients, every safe expression $xP+yQ+zR+\sum_{i\le j}c_{ij}g_ig_j$, with $(g_1,g_2,g_3)=(x,y,z)$, witnesses at most three matching edges. The bound three is attained. The assertion concerns this algebraic expression, without a gate budget for its linear forms or output expansion.
\end{theorem}
\begin{proof}
Use Boolean semantics and suppress positive coefficient magnitudes. Absorb diagonals into affine contexts, so the output is
\[
 F=xP\vee yQ\vee zR\vee\alpha xy\vee\beta xz\vee\gamma yz.
\]
For constant boundaries, a linear signal $D$ can be represented by $D^2$. Also, for $x=UV$,
\[
 (C\vee x)(C\vee U)=C\vee x.
\]
These are homogeneous-factor representations of Boolean functions only; no circuit gate is being charged or deleted.

First suppose $y=AB$ is independent of $x$. Each operand of $z$ contains its linear part and one of $0,x,y,x\vee y$. If some signal, say $x$, occurs in both, write
\[
 z=x\vee t,\qquad t=(C\vee\epsilon y)(D\vee\eta y).
\]
Substitution gives
\[
 x(P\vee R\vee\beta)\vee yQ\vee tR
       \vee(\alpha\vee\gamma)xy\vee\gamma yt.
\]
If $t=CD$, use Lemma~\ref{lem:threeindependentmatching}; if $t$ has one affected operand, use Lemma~\ref{lem:mixedpairfull} with mixed pair $y,t$ and independent $x$; if both are affected, $t=y\vee CD$ reduces to three independent products. The same applies with $x,y$ exchanged.

The remaining earlier-signal choices in the two operands are disjoint. No signal gives three independent products, and one signal gives the mixed-pair lemma. If both occur in the same operand, then $z=(C\vee x\vee y)D$, $xz=xD$ and $yz=yD$. Absorb the crosses and set $k=DR$ to obtain
\[
 xP'\vee yQ'\vee kC\vee xk\vee yk\vee\alpha xy,
\]
a three-independent-product expression. A positive constant in $R$ permits $k=D^2$. If instead $z=(C\vee x)(D\vee y)$, use $xz=xD\vee xy$ and $yz=yC\vee xy$ to absorb the crosses. Set $k=(C\vee x)R$; then $zR=kD\vee ky$. The mixed-pair lemma applies to $x,k$ and independent $y$, with both allowed crosses. If $R$ is Booleanly one, use the homogeneous replacement for $C\vee x$ displayed above. This exhausts independent first products.

If both $a,b$ are positive, Lemma~\ref{lem:twosidedfirst} applies directly. Otherwise, up to symmetry, the remaining dependency is
\[
 y=(A\vee x)B,
\]
and $xy=xB$ can be absorbed into $P$. Label a third-product operand by $0,X,Y,W$ according as its earlier-signal part is $0,x,y,x\vee y$. Exchanging operands leaves ten unordered pairs. In the following cases, every primed context is affine and simply includes the indicated absorbed terms.

For $00$, $z=CD$ and the mixed-pair lemma applies. For $0X$, $z=(C\vee x)D$ and $xz=xD$; after absorption use Lemma~\ref{lem:differentlowfork}. For $0Y$, $z=(C\vee y)D$ and $yz=yD$. If $\beta=0$, set $k=DR$ and use $zR=kC\vee ky$, giving the mixed-pair lemma. If $\beta>0$, set $k=(R\vee x)D$; then $z(R\vee x)=kC\vee ky$, giving the different-low fork lemma. A positive constant in $R$ makes $k=D$, represented as $(D\vee x)D=D$ in the latter application. For $0W$, both crosses are absorbed as $xz=xD$, $yz=yD$; with $k=DR$ the output is $xP'\vee yQ'\vee kC\vee xk\vee yk$, again the mixed-pair lemma.

For $XX$, write $z=x\vee t$, $t=CD$. Since $xz=x$, $yz=xB\vee yt$, the output is
\[
 x(P\vee R\vee\beta\vee\gamma B)\vee yQ\vee tR\vee\gamma yt,
\]
a mixed-pair expression. For $YY$, use $z=y\vee t$, $xz=xB\vee xt$, $yz=y$, giving
\[
 x(P\vee\beta B)\vee y(Q\vee R\vee\gamma)\vee tR\vee\beta xt.
\]
For $WW$, $z=x\vee y\vee t$ and both crosses absorb, leaving three terms with mixed pair $x,y$ and independent $t$.

For $XY$, the identities
\[
 xz=x(D\vee B),\qquad yz=yC\vee xB
\]
absorb all crosses into $P',Q'$. Set $k=(C\vee x)R$; the remaining output is $xP'\vee yQ'\vee kD\vee yk$, the different-low fork lemma. A positive constant in $R$ uses the homogeneous replacement for $C\vee x$.

For $XW$, $z=x\vee CD\vee Cy$, $xz=x$, and $yz=xB\vee yC$. After absorption and $k=CR$, the output is $xP'\vee yQ'\vee kD\vee ky$, a mixed-pair expression. For $YW$, write $z=y\vee t$, $t=(D\vee x)C$; then $xz=xB\vee xC$, $yz=y$, and the output is $xP'\vee yQ'\vee tR$. This is the different-low fork lemma with zero later cross coefficient. These ten cases exhaust all flags.

Every step is a Boolean identity preserving safety and exact edge witnesses. Zero forms, repeated factors, affine constants and all isolates are included. Each resulting template has at most three witnesses. Three are attained by the independent products $ab,cd,ef$ with constant contexts one and zero recursive coefficients.
\end{proof}


\begin{proposition}[Recursive private-product normalization]\label{prop:recursivenormal}
Let a positive variable-input binary circuit have minimum output degree two and nonzero squarefree quadratic part $P_2$. Suppose canonical extraction and pruning delete $p$ multiplication gates, with arbitrary deleted additions. There is a Boolean-equivalent circuit obtained by erasing absorbed off-frame branches, preserving the exact extracted quadratic frame and introducing no gate, whose surviving private products can be ordered as
\[
 g_i=\left(A_i+\sum_{j<i}a_{ij}g_j\right)
     \left(B_i+\sum_{j<i}b_{ij}g_j\right),\qquad 1\le i\le s\le p.
\]
Here $A_i,B_i$ are nonnegative homogeneous linear forms and all coefficients are nonnegative. Its output polynomial is exactly
\[
 P_2+\sum_{i=1}^s g_iL_i+\sum_{1\le i\le j\le s}c_{ij}g_ig_j,
\]
where the $L_i$ are nonnegative affine linear forms. Zero forms, squares, arbitrary fanout and all dependencies are permitted. No old-neighborhood restriction on $A_i,B_i$ is asserted. The representation is a polynomial identity for the modified circuit, not a gate-free construction of its expanded expression.
\end{proposition}
\begin{proof}
Use the full-carrier erasure proved at the beginning of Lemma~\ref{lem:twoproductnormal}, whose erasure step is independent of the number of private products. For completeness, let $R$ be the original minimum-degree-two nodes with addition-only output paths. Every node on such a path has minimum degree two, since otherwise the output would have degree one. In topological order erase each operand occurrence leaving $R$. The full current polynomial $Q$ at its source retains its untouched addition-only output path; each removed contribution is a multiple of $Q$ and is Booleanly absorbed by that output copy. This uses the full carrier, including its higher terms.

A departing route must meet a product before returning to the output, or its successor would itself lie in $R$. Thus every lost contribution has degree at least three. No selected frame occurrence is erased: selected linear nodes have degree-one ancestors, retained products use those linear inputs, and retained quadratic additions keep their selected degree-two occurrences in $R$. Propagating zero branches introduces no gate and preserves the exact frame. Every retained product lies in $R$, and every product in $R$ is retained, since its minimum degree is two and its addition-only output route selects it. Products in $R$ feed only additions within $R$.

Order the surviving private products topologically. Outside $R$, every remaining node other than these products is a variable or addition. Stop an additive expansion at each private-product output. Each operand of $g_i$ is therefore an ordinary homogeneous linear form plus a nonnegative scalar combination of earlier $g_j$, proving the recursion. Private additions may occur inside these linear forms; no old-neighborhood property or free realization of them is claimed.

The same expansion makes each actual input of a retained product a homogeneous linear form plus a scalar combination of the $g_i$. Multiplying its two inputs gives an old-old quadratic term, terms $g_i$ times linear contexts, and scalar multiples of $g_ig_j$. The output carrier adds these products and any outside signals entering directly. Such an outside signal is again a homogeneous linear form plus a scalar combination of private products; positivity and output minimum degree two force its ordinary linear part to vanish. The resulting output has the displayed form with some ordinary quadratic polynomial $Q_2$ in place of $P_2$.

A constant coefficient of $g_i$ in this expansion means an actual addition-only path from that private product to the output. Its original minimum degree cannot be two, since then it would belong to $R$ and be retained. It is at least three, and erasure cannot lower it. Every other single-private-product term has a context of degree at least one, and every private product has minimum degree at least two. A term $g_ig_j$ has minimum degree at least four. Consequently all terms involving private products have minimum degree at least three. The exact quadratic part is $Q_2$, and exact frame preservation gives $Q_2=P_2$.

Omit erased or zero private products, or pad the list by zero signals. Their number never exceeds the original $p$. All statements are polynomial identities for the normalized circuit; subsequent variable restriction is permitted and can make forms vanish.
\end{proof}

\begin{corollary}[Four missing matching edges require four private products]\label{cor:fourdeletedproducts}
Let $G$ be any finite simple graph and let a positive variable-input circuit compute $P\in\Sup(G)$ with nonzero quadratic part $P_2$, of support $H$. Suppose eight vertices $T$ induce exactly four disjoint $G$-edges and no $H$-edge. Canonical extraction deletes at least four original multiplication gates. No four-cycle, cover-number, retained-count or deleted-addition assumption is required.
\end{corollary}
\begin{proof}
Apply Proposition~\ref{prop:recursivenormal} before restricting variables. Its modification preserves the exact quadratic frame and Boolean output and introduces no products. If at most three original products were deleted, pad the surviving private signals with zeros to obtain the recursion in Theorem~\ref{thm:threerecursive}. Set every variable outside $T$ to zero. The quadratic remainder vanishes because $H[T]$ is empty; all remaining homogeneous forms and affine contexts have the permitted form. The restricted Boolean output is exactly the four-edge matching function: positivity excludes false supports, and all four exact edge witnesses survive. This contradicts the theorem's three-witness bound.
\end{proof}

\begin{proof}[Proof of Theorem~\ref{thm:maxfour}]
Fix a minimum vertex cover $C$ of $H$. If any missing edge met $C$, adding at most one endpoint of each of the other three missing edges would cover $G$ with at most $\vc(H)+3$ vertices. If two missing edges shared an endpoint, their four-edge graph would itself have a cover with at most three vertices. Both possibilities contradict the stipulated cover gain. Thus the four missing edges form a matching outside $C$. There is no $H$-edge among their eight endpoints, since $C$ covers $H$; every other $G$-edge would have to be one of the four missing edges. They therefore induce exactly a four-edge matching, and Corollary~\ref{cor:fourdeletedproducts} gives at least four deleted products.

Put $m=|E(H)|$ and $\tau=\vc(H)$. Write the retained multiplication and addition counts as
\[
 K=\tau+s,\qquad A_2=m-1+t,\qquad s,t\ge0,
\]
using the separate quadratic lower bounds. Let $D=p+a$ count the original deleted products and additions. The original size is at least
\[
 m+\tau-1+s+t+D,
\]
whereas the required star count is $m+\tau-1+8$. Lemma~\ref{lem:fourmatchinglost} gives $D\ge6$. If $s+t\ge2$, this proves the result. If $s=0$, Corollary~\ref{cor:fourminproduct} gives $D\ge8$, without a restriction on retained additions. A putative saving consequently requires precisely
\[
 s=1,\qquad t=0,\qquad D=6.
\]
The retained frame then has $K=\tau+1$ products and exactly $m-1$ additions. At least one deleted addition is a cut: without a cut, the useful output paths stay homogeneous quadratic and no edge loses its quadratic witness. Since $p\ge4$, the only profiles are $(p,a)=(4,2)$ and $(5,1)$. The first is excluded by Lemma~\ref{lem:fourproductssix}, whose strengthened hypotheses permit exactly this extra retained product. The second is excluded by Proposition~\ref{prop:oneadditionall}, which requires $p\ge2^4-2=14$. These exhaust the possible savings.
\end{proof}

\begin{proposition}[Where the fourth missing edge can attach]\label{prop:fourthattachment}
Suppose $G$ has no four-cycle, its quadratic support $H$ is nonempty, exactly four edges are missing from $H$, and $\vc(G)-\vc(H)=3$. Suppose the extracted frame has exactly $\vc(H)$ products, and exactly three products and two additions were deleted. Let $C$ be the distinct assigned center cover of that frame. If $G[V\setminus C]$ is a three-edge matching with isolates, the fourth missing edge is $hz$ with $h\in C$ and $z\notin C$. If $z$ is a matching endpoint with mate $z'$, then $hz'\in E(H)$. In particular, when $G$ has girth at least five, $z$ must be isolated in $G[V\setminus C]$.
\end{proposition}
\begin{proof}
Write $T=V\setminus C$, let $M=G[T]$ be the three-edge matching, and let $e_0$ be the fourth missing edge. It either lies inside $C$ or crosses from $C$ to $T$. Assume it is internal, or that $e_0=hz$ has $z$ on $M$ and $hz'\notin H$ for its mate $z'$.

For each center $k\in C$ put $N_k=N_H(k)\cap T$. These are the complete old outside leaf supports: there is exactly one product per distinct center, so every old edge from $k$ to $T$ belongs to that star. Each $N_k$ is nonempty, since otherwise removing $k$ from the minimum cover would still cover $H$. We claim that $N_k$ is not contained in a partial endpoint transversal of $M$, meaning a set with at most one endpoint from each matching edge and no isolates.

Otherwise complete $N_k$ to a three-vertex cover $W$ of $M$. The set $(C\setminus\{k\})\cup W$ covers $H$ and $M$. If $e_0$ is internal to $C$, one endpoint remains. If $e_0=hz$ and $k\ne h$, $h$ remains. If $k=h$, neither $z$ nor $z'$ lies in $N_h$, because $hz$ is missing and $hz'$ is assumed not old. Thus $W$ can also be chosen to contain $z$. In every case this gives a $G$-cover of size $|C|+2$, a contradiction.

Perform exact full-carrier erasure and set $C$ to zero, retaining all of $T$, including isolates. The Boolean output is $M$. For the full center and leaf inputs $A_k,B_k$ of each retained product,
\[
 A_k\bigvee N_k\ \le\ A_kB_k\ \le\ M.
\]
The center input $A_k$ is itself safe for $M$. Indeed, if it were true on an independent support $S$, every $n\in N_k$ would complete a matching edge with $S$. Thus $N_k$ would consist of missing mates of endpoints in $S$, a partial transversal, contradicting the claim. Replacing each retained product by its full center input therefore preserves $M$ and removes all retained conjunctions. Carrier erasure ensures those inputs have no retained-product ancestors, so these replacements do not alter one another's inputs.

At most the original three private conjunctions remain. The matching conjunction bound forces all three to survive. Its equality argument shows that each is safe and is the first true conjunction on its own distinct matching edge; both operands have a nonzero raw-variable contribution on that edge. We must count original deleted additions, rather than the retained disjunctions that are free in this reduction.

Keep all surviving original deleted additions marked, including cuts whose low input becomes zero. Every useful private product must reach a marked addition before reaching another private product or the selected frame. Direct entry into a later private conjunction contradicts its own-edge property. Entry into a selected linear position is impossible by original minimum degree. An equal-degree direct export into the quadratic carrier would select a private multiplication source into the frame; a higher-degree export is a marked cut. The retained center networks have selected support $\{k\}$, so a component of private products and additions, of minimum degree at least two, cannot enter one without an old degree-one boundary. This boundary is counted even when restriction makes it zero.

Form the operand multigraph on the three private product outputs, the $a'$ surviving marked additions, and their other immediate boundary inputs. Retain both operand incidences at each addition. Every component has a boundary: a boundary-free component cannot feed a later conjunction on its own edge, cannot enter a retained linear position by degree, and cannot export at equal degree to the carrier without selecting a private multiplication source. A higher-degree export is another marked cut in that component. A new export caused by center replacement has the degree-one boundary just described. An addition-only component without any boundary is also impossible by acyclicity. If there are $b$ boundary vertices and $c$ components, then $b\ge c$, and connectedness gives
\[
 2a'\ge3+a'+b-c,\qquad a'\ge3.
\]
This contradicts the two original deleted additions. We have rechecked the original-gate charge under the present three-edge restriction; no substitution of the total missing count into Proposition~\ref{prop:equalproductadditions} is used.

Therefore $e_0$ has the asserted crossing form and old mate edge. If $z$ is a matching endpoint, $h,z,z'$ form a triangle, consisting of $hz$, $zz'$ and $hz'$. Girth at least five excludes this possibility, leaving an outside isolated endpoint. Proposition~\ref{prop:disjointfourth} below treats the isolated attachment.
\end{proof}

\begin{proposition}[An essential variable absent from the quadratic support]\label{prop:absentvariable}
Let $G$ be a finite simple graph without a four-cycle and let a positive variable-input circuit compute $P\in\Sup(G)$ with nonempty quadratic support $H$. Suppose six vertices $T$ induce exactly three disjoint $G$-edges and no $H$-edge. Let $z\notin T$ have positive $G$-degree and zero $H$-degree. Canonical extraction cannot delete at most three multiplication gates and at most two addition gates. All retained gate counts and private-product dependencies are unrestricted.
\end{proposition}
\begin{proof}
Every retained node of the quadratic frame lies on a selected path to the quadratic output. Its selected minimum-degree polynomial is $z$-free: positivity would otherwise carry a $z$-containing monomial along that path to a quadratic output monomial involving $z$, contrary to its zero $H$-degree. This also excludes a selected $z^2$ term, since the quadratic part is squarefree.

Set $z=0$, propagate zeros, bypass additions with a zero operand, and prune. The result is a positive variable-input circuit for $P'=P|_{z=0}\in\Sup(G-z)$. Its quadratic part is exactly the same nonzero $P_2$. We verify preservation of the entire canonical frame. Inducting through its selected nodes, every nonzero selected minimum part and its degree are unchanged. At a cut, the low operand therefore keeps its nonzero minimum part, whereas restriction can only raise the high operand's minimum degree or annihilate it. The high branch cannot become selected. At a retained equal-degree addition both selected parts survive, and at a retained product both selected operands survive. Thus no retained gate is erased or bypassed, no selected output route becomes useless, and no original private gate becomes newly selected. The extracted frame after restriction is exactly the original frame.

The variable $z$ is essential to $f_G$, since it has a neighbor. It occurs on a useful input path, and setting it to zero removes at least one original gate: an addition immediately reading it is bypassed, while a multiplication reading it becomes zero. Every removed gate is private by the preceding frame argument. If $p',a'$ are the private-product and private-addition counts after restriction, the assumed bounds give
\[
 p'\le3,\qquad a'\le2,\qquad p'+a'\le4.
\]
The three-edge matching on $T$ is unchanged and still has empty quadratic restriction. Corollary~\ref{cor:threedeletedproducts} forces $p'\ge3$, hence $p'=3$ and $a'\le1$. If $a'=0$, there is no useful cut; induction from the variables makes every useful gate homogeneous, and the nonzero quadratic output has no missing quadratic witness. If $a'=1$, Proposition~\ref{prop:oneadditionall} applies to the at least three remaining missing edges and requires $p'\ge2^3-2=6$. Both alternatives are impossible.
\end{proof}

In the isolated-endpoint attachment left by Proposition~\ref{prop:fourthattachment}, the outside endpoint must consequently have an old $H$-neighbor. That neighbor lies in the assigned cover $C$, because the endpoint is isolated in $G[V\setminus C]$, and it differs from its missing neighbor $h$. This gives a separate restriction on the endpoint. Proposition~\ref{prop:disjointfourth} below also handles old neighbors in the present three-product, two-addition profile.

\begin{lemma}[An independent pair with an attached missing edge]\label{lem:attachedpair}
Let $H\subseteq G$ be finite simple graphs, with $G$ containing no four-cycle. Suppose distinct vertices $h,z,z',a,b,c,d$ satisfy
\[
 G[\{z,z',a,b,c,d\}]=\{zz',ab,cd\},\qquad
 E(H[\{z,z',a,b,c,d\}])=\varnothing,
 \qquad hz\in E(G)\setminus E(H).
\]
The edge $hz'$ and other edges incident with $h$ are unrestricted. Let $g$ be a product of at most three old homogeneous linear forms, and $f$ a product of at most two. Let $L,M$ be nonnegative affine linear forms whose homogeneous parts are old, and let $q\geq0$. If every monomial of
\[
 F=gL+fM+qgf
\]
contains a $G$-edge, then $F$ cannot witness all four missing edges $hz,zz',ab,cd$. Zero factors and absent terms are permitted.
\end{lemma}
\begin{proof}
Write $J=\{hz,zz',ab,cd\}$, and let $A,B,C$ denote the subsets of $J$ witnessed by $gL,fM,qgf$, respectively. Positivity makes each summand safe. Proposition~\ref{prop:oldfactors} gives
\[
 |A|\leq2,\qquad |B|\leq1,\qquad |C|\leq2.
\]
A positive affine constant makes the underlying product safe; any exact edge witnessed by its multiple was already witnessed by that product. Thus such constants only improve these bounds.

If $|A|=2$, the four-factor equality case confines $g$ to three of the four two-element transversals of an induced two-edge matching. Their common intersection is empty. A pair meeting all three supports either lies within the four endpoints or cannot meet them all; within those endpoints, every $G$-edge is one of the two matching edges. Every $C$-edge meets all three $g$ factors, so $C\subseteq A$ and the union has size at most three.

Four witnesses would therefore require $|A|=|B|=1$ and $|C|=2$, with these sets pairwise disjoint. Write $C=\{e_1,e_2\}$ and $B=\{e_4\}$. The two $C$-edges form an induced matching in the full graph $G$. Each factor of $f$ meets both edges and is consequently nonsingleton. The context $M$ cannot have a positive constant, since a safe product of at most two old factors cannot witness both edges. Nor can $f$ have fewer than two factors: a safe product of two old forms witnessing one missing edge has the two singleton supports of that edge, whereas its $f$ factor must meet both $e_1,e_2$. Thus $f=XY$ with $X,Y$ nonsingleton and $M$ homogeneous.

Cubic missing-edge rigidity forces $M$ to be a singleton, say $\{w\}$, and $e_4=tw$, with $t$ in both $X$ and $Y$. Neither $X$ nor $Y$ contains $w$: otherwise the missing-edge witness repeating $w$ in that factor and $M$ forces the other factor to be the singleton $t$. If an old nonsingleton factor containing $t$ also contained another neighbor $u$ of $w$, its old center $v$ would give the four-cycle $v,t,w,u,v$. These vertices are distinct; in particular $v=w$ would make $tw$ old. Hence
\[
 N_G(w)\cap\bigl(\supp X\cup\supp Y\bigr)\subseteq\{t\}.
\]
Neither $e_1$ nor $e_2$ can contain $w$. If both avoided $t$, choose an endpoint of $e_1$ in $X$ and an endpoint of $e_2$ in $Y$. The two chosen endpoints are nonadjacent because $e_1,e_2$ form an induced matching; neither is adjacent to $w$. They produce an independent-support monomial of $wXY$, a contradiction. Thus one $C$-edge shares $t$ with $e_4$.

The only incident pair in $J$ is $hz,zz'$. Therefore $t=z$, the $B$-edge is one arm, and a $C$-edge is the other. Relabeling $ab,cd$, the other $C$-edge is $ab$ and the $A$-edge is $cd$. In particular, $f$ contains $z^2$, while $g$ is true on each of $ab,cd$. Restrict $g$ to $a,b,c,d$, obtaining $g'$. Every monomial $m$ of $g'$ gives $mz^2$ in the safe product $gf$. On $\{z,a,b,c,d\}$ the only edges are $ab,cd$, so $g'$ is safe for that two-edge matching and witnesses both edges. It has at most three old factors, contradicting the four-factor lower bound. This argument never restricts the other edges incident with $h$.
\end{proof}

\begin{proposition}[A fourth missing edge adjacent to a missing matching]\label{prop:adjacentfourth}
Let $G$ be finite and four-cycle-free, and let a positive variable-input circuit compute $P\in\Sup(G)$ with nonempty quadratic support $H$. Suppose six vertices $T=\{z,z',a,b,c,d\}$ induce exactly the three $G$-edges $zz',ab,cd$, and $H[T]$ is empty. If $h\notin T$ and $hz\in E(G)\setminus E(H)$, canonical extraction cannot delete at most three products and at most two additions. The edge $hz'$ and all other edges incident with $h$ are unrestricted. No cover-gain, retained-count, or total-missing-edge hypothesis is required.
\end{proposition}
\begin{proof}
Put $J=\{hz,zz',ab,cd\}$. Corollary~\ref{cor:threedeletedproducts} requires at least three private products. With at most three products, Proposition~\ref{prop:oneadditionall} and the absence-of-cuts observation exclude at most one private addition. Thus the assumed bounds force exactly three private products and two private additions. Exact full-carrier erasure preserves the entire quadratic frame and Boolean output. It cannot remove any private gate, since the same two bounds would fail for the resulting circuit. Hence all five private gates remain essential.

Restricting outside $T$ to zero kills the quadratic part and leaves the three-edge matching, which requires eight positive linear CNF clauses, one for each minimal vertex cover. Apply the selected-occurrence classification in Lemma~\ref{lem:fourmatchinglost}. One visible cut gives at most seven clauses. With two visible cuts, let $k$ count the private products ancestral to the first high $g$. The $k=3$ case has at most five clauses, and the $k=2$ direct-private-ancestor cases at most six. These cases are excluded on $T$.

For all remaining cases return to the full graph. Before substituting the dependence of the highs, the output is
\[
 P_2+F,\qquad F=gL+fM+qgf,
\]
with old affine contexts $L,M$. Every positive summand is safe for $G$, and $F$ must witness all four edges of $J$, since none occurs in $P_2$. We retain all old edges involving $h$ throughout this argument. Any two missing witnesses of one safe old product form an induced matching in the full graph. Of the other edges of $J$, at most one meets the endpoints of such a matching pair: it can only be the other arm of $hz,zz'$. Four-factor equality confines every factor to those four endpoints. Three of the four distinct transversal supports have empty common intersection, so a product containing them cannot witness a further $G$-edge.

\emph{The remaining $k=2$ cases.}
Here $g$ has at most three old factors and $f$ is $AB$, $AB\vee gB$, or $AB\vee g$, allowing zero old parts. The independent case is Lemma~\ref{lem:attachedpair}.

If $f=AB\vee g$, expansion and Boolean absorption give
\[
 F=gL\vee ABM\vee gM\vee qg.
\]
A positive $q$ or a positive constant in either context makes $g$ itself safe. All its multiples then have only its own missing witnesses, at most one, while $ABM$ has at most one. Otherwise $gL,gM$ each have at most two missing witnesses and $ABM$ at most one. If either of the first two witnesses two, equality confines the three $g$ factors to that pair, and the other $g$ term adds none. If neither does, the three terms each contribute at most one. Either way there are at most three targets.

If $f=AB\vee gB$, expansion gives
\[
 F=gL\vee ABM\vee gBM\vee qgB.
\]
For $q=0$ and homogeneous $M$, regroup as $gL\vee(BM)A\vee g(BM)$ and apply Lemma~\ref{lem:attachedpair}. Otherwise $gB$ is safe and absorbs its multiples. The two terms $gL,gB$ have at most two missing witnesses together: if either attains two, the same four-factor confinement controls every $g$-based target, and otherwise each has at most one. The term $ABM$ adds at most one. Positive constants and zero degenerations only reduce these counts.

\emph{The $k=1$ cases.}
Write $g=UV$. An independent $f$ has at most three old factors, so Lemma~\ref{lem:attachedpair} applies with the roles exchanged. Otherwise the classification gives
\[
 f=DE\vee gD,\qquad
 F=gL\vee DEM\vee gDM\vee qgD,                 \tag{*}
\]
where $D$ has at most two old factors, $DE$ at most three, and $D$ is literally a subproduct of $DE$. Empty products and zero parts are permitted.

First suppose $D=1$. If $q>0$ or either context has a positive constant, $g$ is safe and contributes at most one missing witness through all its multiples, while $EM$ contributes at most two. Otherwise $gL,gM$ are cubics with at most one each. Four targets would force $EM$ to have two at four-factor equality, and the two cubics to supply the distinct remaining edges. Its factor $M$ is a nonsingleton transversal confined to the $EM$ endpoint set. If the $EM$ pair were $ab,cd$, $gM$ could witness neither remaining arm. Otherwise its outside target must be the other path arm and the target of $gL$ the remaining isolated edge. Those two targets are disjoint, so both $U,V$ are nonsingleton. The safe cubic $UVM$ would then witness a missing edge with three nonsingleton factors, contradicting cubic rigidity.

Now suppose $D$ is nonempty. If it has two factors, $q=0$, and $M$ is homogeneous, regroup $(\ast)$ as
\[
 (DM)E\vee gL\vee g(DM).
\]
The two high products have at most three and two old factors; $E$ is one old factor or a unit/zero degeneration. Lemma~\ref{lem:attachedpair} applies.

In every remaining case put $A_0=gL$ and $B_0=DEM$. If $q>0$ or $M$ has a positive constant, the remaining terms have exactly the witnesses of the safe product $C_0=gD$. Otherwise $D$ has one factor and put $C_0=gDM$. Thus the respective missing-witness bounds are $1,2,2$. If $C_0$ has two targets, equality confines its $g$ and nonempty $D$ factors to their endpoints. Both $A_0$ and $B_0$ retain a confined factor, so their only possible further target is the same other path arm. The total is at most three.

Four targets would therefore force $B_0$ to have two and $A_0,C_0$ to supply the distinct remaining edges. Equality confines every $DE$ and $M$ factor to the $B_0$ endpoints. Since $D$ is nonempty, the outside target of $C_0$ must be the other arm of $hz,zz'$, and that of $A_0$ the remaining isolated edge. Relabel so that $B_0$ witnesses one arm and $ab$, $C_0$ the other arm, and $A_0$ the edge $cd$.

Each of $U,V$ meets both $cd$ and the $C_0$ arm, so neither is singleton. The context $L$ cannot have a positive constant, since $g$ would then be safe while potential on two disjoint missing edges. Cubic rigidity consequently makes $L$ a singleton endpoint of $cd$, with the other endpoint $t$ in both $U,V$. Thus $g$ contains $t^2$. Every factor of $D$ is confined to the $B_0$ endpoints and meets the other path arm; their only common endpoint is $z$. All these factors contain $z$, as does $M$ if it is required in $C_0$. Choosing $t$ twice and $z$ in the remaining factors gives a monomial of $C_0$ supported on $\{t,z\}$. This pair is independent in the prescribed graph on $T$, contradicting safety. All cases of the classification have now been excluded.
\end{proof}

\begin{lemma}[Four disjoint targets in an independent pair]\label{lem:disjointpair}
Let $H\subseteq G$ with $G$ finite, simple, and four-cycle-free. Let $g,f$ be products of at most three and two old homogeneous linear forms, respectively, and let $L,M$ be nonnegative affine linear forms with old homogeneous parts. For $q\geq0$, suppose every monomial of $F=gL+fM+qgf$ contains a $G$-edge. Then $F$ cannot witness four pairwise vertex-disjoint edges of $E(G)\setminus E(H)$. Their eight endpoints need not induce a matching: other old or new edges are permitted.
\end{lemma}
\begin{proof}
Let $J$ be four such targets, and let $A,B,C$ be their witness sets in the three summands. The old-factor bound gives $|A|\leq2$, $|B|\leq1$, and $|C|\leq2$, with positive context constants handled by safety of the underlying product. If $|A|=2$, four-factor equality confines the three $g$ factors to three distinct transversals of that induced matching. Their common intersection is empty, so every $G$-edge potential on $g$ is one of the same two edges. Thus $C\subseteq A$, leaving at most three targets overall.

Four targets therefore force the disjoint witness sets to have sizes $1,1,2$. Write $C=\{e_1,e_2\}$ and $B=\{tw\}$. The two $C$-edges form an induced matching in the full graph, by Proposition~\ref{prop:oldfactors}. Each factor of $f$ meets both and is nonsingleton. A positive constant in $M$ would make $f$ safe and witness both edges with at most two factors. Fewer than two factors in $f$ would force a singleton $f$ factor at the two-factor equality case for its missing $B$-edge. Both alternatives are impossible. Hence $f=XY$, $M$ is homogeneous, and cubic rigidity gives $M=\{w\}$ with $t$ in both nonsingleton factors $X,Y$.

As in the full old-center argument in Lemma~\ref{lem:attachedpair}, neither factor contains $w$, and every other element of either factor avoids $N_G(w)$. Explicitly, $w$ in one factor would force the other to be a singleton; an old factor containing $t$ and another neighbor $u$ of $w$ would give a four-cycle through its old center, $t,w,u$. Since the four targets are disjoint, $e_1,e_2$ avoid $t,w$. Choose an endpoint $x$ of $e_1$ in $X$ and an endpoint $y$ of $e_2$ in $Y$. The pair $e_1,e_2$ is induced, so $xy\notin E(G)$; also $wx,wy\notin E(G)$. The monomial $wxy$ in $XYM$ has independent support, contradicting safety. Only the two cross-term targets needed inducedness, and the old-factor theorem supplied it.
\end{proof}

\begin{proposition}[A disjoint fourth missing edge]\label{prop:disjointfourth}
Let $G$ be finite and four-cycle-free, and let a positive variable-input circuit compute $P\in\Sup(G)$ with nonempty quadratic support $H$. Suppose six vertices $T$ induce three disjoint $G$-edges, all absent from $H$, and that some missing edge $hw$ is disjoint from $T$. Canonical extraction cannot delete at most three products and at most two additions. Other edges between $\{h,w\}$ and $T$, retained gate counts, and cover gain are unrestricted.
\end{proposition}
\begin{proof}
Let $J$ consist of the three matching edges and $hw$. Corollary~\ref{cor:threedeletedproducts}, Proposition~\ref{prop:oneadditionall}, and the absence of inhomogeneous cuts when there is no deleted addition force exactly three private products and two private additions. Exact full-carrier erasure cannot remove any of these gates without violating the same bounds. Apply the selected-occurrence classification in Lemma~\ref{lem:fourmatchinglost}. Restriction outside $T$ kills the quadratic part and leaves the three-edge matching, which requires eight linear CNF clauses. As in Proposition~\ref{prop:adjacentfourth}, this excludes the one-visible-cut case, the $k=3$ case, and the $k=2$ direct-private-ancestor cases.

Keep the entire graph for the remaining cases. The safe perturbation is $F=gL+fM+qgf$ and must witness every edge of $J$. If $q=0=M$, only $gL$ remains and has at most two targets, so this boundary is already excluded. If a four-factor old product witnesses two targets, equality confines every factor to their four endpoints. Every other edge of $J$ is disjoint from those endpoints. Consequently a term containing any confined factor cannot witness another target. Three distinct transversal factors further confine every $G$-edge potential on their product to the same two edges.

For $k=2$, the product $g$ has at most three old factors. Independent $f=AB$ is Lemma~\ref{lem:disjointpair}. If $f=AB\vee g$, a positive $q$ or context constant makes $g$ safe; all its multiples then have at most one target and $ABM$ at most one. Otherwise the three terms $gL,gM,ABM$ have bounds $2,2,1$. Four targets require a two-witness $g$ quartic, whose three confined factors control both $g$ terms; the total is at most three. If $f=AB\vee gB$, then for $q=0$ and homogeneous $M$ regroup as $gL\vee(BM)A\vee g(BM)$ and apply Lemma~\ref{lem:disjointpair}. Otherwise $gB$ is safe and absorbs its multiples. The two $g$ terms again have at most two targets together by confinement, and $ABM$ adds at most one.

For $k=1$, write $g=UV$. An independent $f$ has at most three old factors, so the lemma applies with the roles exchanged. Otherwise
\[
 f=DE\vee gD,\qquad F=gL\vee DEM\vee gDM\vee qgD,
\]
where $D$ has at most two factors, $DE$ at most three, and $D$ is a subproduct of $DE$. If $D=1$, positive $q$ or a context constant makes $g$ safe, giving at most one target in its multiples and at most two in $EM$. Otherwise $gL,gM$ are cubics with at most one target each. Four targets would force $EM$ to have two at four-factor equality, while both cubics supply new ones. But its confined $M$ factor is disjoint from both remaining target edges, so $gM$ supplies neither.

If $D$ has two factors, $q=0$, and $M$ is homogeneous, regroup as $(DM)E\vee gL\vee g(DM)$ and apply the lemma. In all other nonempty-$D$ cases use the three safe products $A_0=gL$, $B_0=DEM$, and $C_0=gD$ when $q>0$ or $M$ has a positive constant, otherwise $C_0=gDM$. They have target bounds $1,2,2$. If $C_0$ has two, its $g$ and $D$ factors are confined, so neither $A_0$ nor $B_0$ can add an outside target. If $C_0$ has at most one, four targets require two from $B_0$ and a new one from each other term. Yet the nonempty confined $D$ subproduct remains in $C_0$, preventing its new target. This exhausts the classification and proves the claim.
\end{proof}

Combining Propositions~\ref{prop:adjacentfourth} and~\ref{prop:disjointfourth}, an induced missing three-edge matching together with \emph{any} further missing edge excludes three private products and two private additions. A further edge cannot have both endpoints in the six matching vertices; it either meets them once or is disjoint from them.

\begin{proposition}[Cover gain three with the minimum retained product count]\label{prop:gainthreesix}
Suppose $G$ is finite and four-cycle-free, $P\in\Sup(G)$ has nonempty quadratic support $H$, exactly four edges are missing from $H$, and $\vc(G)-\vc(H)=3$. If canonical extraction retains exactly $\vc(H)$ multiplication gates, it deletes at least six gates in total. The retained addition count is unrestricted.
\end{proposition}
\begin{proof}
Let $C$ be the minimum $H$-cover of assigned centers. The outside graph $N=G[V\setminus C]$ has at most four edges, all missing, and $\vc(N)\geq3$, because adjoining an $N$-cover to $C$ covers $G$. Keep its isolates in all restrictions.

Lemma~\ref{lem:fivelost} already gives at least five deletions. If there were exactly five, Proposition~\ref{prop:gainthreeproducts} would force at least three private products, while Proposition~\ref{prop:oneadditionall} would exclude at most one private addition. Hence there would be exactly three products and two additions.

Up to isolates, a four-cycle-free graph with at most four edges and cover number at least three is one of
\[
 3K_2,\quad4K_2,\quad P_3\sqcup2K_2,\quad
 P_4\sqcup K_2,\quad K_3\sqcup K_2.
\]
Indeed every connected tree with at most four edges has cover number at most two; splitting components gives precisely these possibilities, and the only possible cycle with at most four edges is a triangle. Proposition~\ref{prop:exceptionaloutside} excludes the last two graphs with two private additions. Proposition~\ref{prop:paththreeproducts} excludes $P_3\sqcup2K_2$, and Corollary~\ref{cor:fourdeletedproducts} excludes $4K_2$ with three private products. In the remaining $3K_2$ case, choose its six endpoints as $T$. A fourth missing edge exists globally and cannot have both endpoints in $T$. Propositions~\ref{prop:adjacentfourth} and~\ref{prop:disjointfourth} exclude respectively its meeting $T$ once or avoiding $T$. Thus five deletions are impossible.
\end{proof}

\begin{proposition}[A missing induced path or triangle requires three additions]\label{prop:twosumspathtriangle}
Let $G$ be a finite simple four-cycle-free graph, and let a positive variable-input circuit compute $P\in\Sup(G)$ with nonempty quadratic support $H$. If $H[T]$ is empty and $G[T]$ is an induced $P_4$ or $K_3$, canonical extraction deletes at least three addition gates. The number of deleted products, the retained counts, cover gain, all other edges and all other vertices are unrestricted.
\end{proposition}
\begin{proof}
An \emph{old product} means a product of old homogeneous linear forms. Empty products may occur as algebraic contexts. Every safe old product witnesses at most one of the three target edges, by Proposition~\ref{prop:oldfactors}: its missing witnesses form an induced matching, whereas no two of these target edges do. An old form cannot contain both $u,w$ for a missing two-edge path $u-v-w$, since its old center gives a four-cycle or makes the missing edges old. Thus an old product potential on both $uv,vw$ is true on the singleton $v$. For an induced path $a-b-c-d$, an old form meeting both end pairs has restriction either $\{a,d\}$ or $\{b,c\}$, since $\{a,c\}$ and $\{b,d\}$ are forbidden. An old product potential on both ends is consequently true on each of the nonedges $\{a,c\},\{b,d\}$.

We first prove a three-term obstruction. Let $A,L$ be nonzero old linear forms with $AL$ supported on $H$. Let $g,U$ be old products, and let $F$ be a nonempty product of old forms each containing $\supp(A)$. Suppose
\[
 gL,\qquad UF,\qquad Ug
\]
are safe. They cannot witness all three target edges. Otherwise each witnesses a different edge, say $e_1,e_2,e_3$. Since $L$ meets a target edge and $AL\subseteq H$, the support of $A$ avoids the target vertices. Also $A$ Booleanly implies $F$, so $UA$ is safe.

If $e_2=uv$ and $e_3=uw$ meet, the product $U$ is true on $u$. Safety of $UF$ forces $F$ false on $u$, while its $uv$ witness makes some factor $F_i$ omit $u$ and contain $v$. Choose $h\in\supp(A)$, outside the target set. Safety of $UA$ gives $hu\in E(G)$. The old form $F_i$ contains $h,v$ and hence has an old center $z$. This center differs from $h,v$ by simplicity and from $u$ because $uv$ is missing. The vertices $z,v,u,h$ form a four-cycle. If instead $e_2,e_3$ are the two path ends $ab,cd$, the remaining edge is $bc$. Then $g$ is true on $c$, and $U$ is true on the nonedge $\{a,c\}$ by the preceding end-pair observation. Thus $Ug$ is true on that nonedge. Both cases are impossible, proving the obstruction.

Suppose now at most two additions are deleted. Apply exact full-current-quadratic-carrier erasure, keeping the complete original selected frame and its Boolean output. Every surviving boundary outside the carrier is an old linear signal or a private multiplication expression. A useful private computation must enter the selected frame at a visible unequal-degree cut, where visibility refers to the cut's own selected occurrence. Without a cut no missing witness remains.

With one visible cut, at most one other private addition lies in its high computation. Before and after that addition there are only products, on old homogeneous boundaries. Booleanly its high therefore has the form $Z(X\vee Y)$, where $Z,X,Y$ are old products; with no internal addition there is only one old product. Shared multiplication paths merely repeat factors or the same sum. A hidden cut cannot affect a retained linear boundary, since that would make its own occurrence visible. The direct context is old linear or scalar, and an individual insertion cannot be squared in the frame: its nonzero selected low would contribute a forbidden square to $P_2$. Expansion gives at most two safe old products, insufficient for the three targets.

It remains that both deleted additions are visible. Name their highs $g,f$ in cut order. Before substituting any dependence of $f$ on the first cut, exact frame expansion is
\[
 P_2\vee gL\vee fM\vee qgf,
\]
where $q$ records a nonnegative coefficient and the actual direct contexts are scalar or old homogeneous linear. There are no individual square terms for the same old-low reason. If the earlier cut retains degree two, $g$ itself reaches the output with scalar context. No retained linear boundary is affected; every dependent later high has $g$ as a multiplicative factor and is absorbed, while an independent later high is another old product. Again at most two relevant old products remain.

Otherwise the earlier low is a nonzero linear form $A$, with $AL\subseteq H$, and $g$ is an old product. Every retained linear boundary affected before the second cut has the form $A_i\vee g$ with $\supp(A)\subseteq\supp(A_i)$: the same positive addition-only route carries the original low and its replacement. All other boundaries are old forms, uninserted private-product ancestors (also old products), or the raw high $g$. No other private product can enter a retained linear sum without a further private addition. Full quadratic-carrier uses were erased. The remaining high network contains only products, so its possibilities are
\[
 f=U,\qquad f=Ug,\qquad f=U(F\vee g),
\]
where $U$ is an old product and, in the last genuinely affected case, $F$ is a nonempty product of old forms each containing $\supp(A)$. The identity
\[
 \prod_i(A_i\vee g)=\Bigl(\prod_i A_i\Bigr)\vee g
\]
handles arbitrary multiplication count, repeated inputs and shared ancestors. A raw $g$ factor absorbs all affected factors. A high which vanishes is omitted.

If $q>0$, both selected lows have degree one and $M$ contains $\supp(A)$. Indeed a retained product producing $gf$ reads the old contribution of the first low in its direct $f$ context. This is an identity of the full original frame, not an inference after target restriction.

For $f=U$, absence of the cross term leaves two safe old products. Its presence gives $gL,UM,Ug$, so the obstruction applies with $F=M$. For $f=Ug$, the output beyond $P_2$ is $gL\vee UgM$ if $q=0$, and $gL\vee Ug$ if $q>0$; there are again only two old products.

For $f=U(F\vee g)$ and $q=0$, the output terms are
\[
 gL,\qquad UFM,\qquad UgM.
\]
Apply the obstruction with $U'=UM$ and the same $F$. If $q>0$, absorption instead leaves
\[
 gL,\qquad UFM,\qquad Ug.
\]
Apply it with $F'=FM$; every factor of $F$ contains $A$, and the actual mixed-term identity gives the same containment for $M$. Scalar identity factors may be removed. If a term is zero or witnesses none of the targets, the other at most two safe old products cannot supply all three. If an empty product makes one of the highs a unit, the resulting terms absorb to at most two, or violate safety directly. These observations include every zero or scalar degeneration.

All cases contradict the three target witnesses. The erasure never increased the original deleted-addition count, and no arbitrary Boolean re-extraction was used. Thus at least three original additions are deleted.
\end{proof}

\begin{proposition}[Two additions and missing degree]\label{prop:twosumsdegree}
Let $G$ be finite and four-cycle-free, and let a positive variable-input circuit compute $P\in\Sup(G)$ with nonempty quadratic support $H$. If extraction deletes at most two original additions, the missing-edge graph $J=G\setminus H$ has maximum degree at most two. Deleted products, retained counts, fanout and cover gain are unrestricted.
\end{proposition}
\begin{proof}
Suppose three missing edges $va,vb,vc$ have the same endpoint. An old linear form meeting both $\{v,a\}$ and $\{v,b\}$ must contain $v$: otherwise its old center together with $a,v,b$ gives a four-cycle, or makes those edges old. Hence an old product potential on two such edges is true on singleton $v$. Every safe old product therefore witnesses at most one of the three edges.

Three safe old products of the form $gL,UC,Ug$ cannot witness all three. Otherwise each witnesses a different edge; the shared-endpoint observation makes both $g$ and $U$ true on singleton $v$, contradicting safety of $Ug$. Empty products satisfy the same implication, and zero terms leave fewer witnesses. No actual-low condition is needed for this particular obstruction.

Apply the full-carrier and arbitrary-product two-addition normalization proved in Proposition~\ref{prop:twosumspathtriangle}. With one visible cut there is at most one private sum in its high network; Boolean idempotence and distributivity expand its high and old context into at most two safe old products. A first cut retaining degree two similarly leaves at most two after absorption. With two visible cuts and first linear low, the exact independent-symbol frame is
\[
 P_2\vee gL\vee fM\vee qgf.
\]
Here $g$ is an old product, and the second high is $U$, $Ug$, or $U(C\vee g)$, where $U,C$ are old products. The last form follows from $\prod_i(A_i\vee g)=\prod_i A_i\vee g$; raw $g$ factors absorb the affected factors. Contexts are old linear forms or scalars. For independent $f=U$, a zero cross coefficient leaves two terms, while a positive one gives $gL,UM,Ug$. For $f=Ug$, absorption leaves at most two terms. For $f=U(C\vee g)$ and $q=0$, the terms are $gL,(UM)C,(UM)g$; for $q>0$ they reduce to $gL,U(CM),Ug$. Each three-term case has the forbidden form. All terms remain safe positive summands. Zero, unit, repeated-factor and arbitrary-fanout cases are covered by the same identities. This excludes three incident missing edges.
\end{proof}

\begin{proof}[Proof of Theorem~\ref{thm:twosumscomponents}]
Proposition~\ref{prop:twosumsdegree} gives maximum missing degree two. A missing triangle is induced and has no old edge, so Proposition~\ref{prop:twosumspathtriangle} excludes it. Suppose three successive missing edges are $ab,bc,cd$, on distinct vertices. The edge $ad$ would complete a four-cycle. Both diagonals $ac,bd$ would also give a four-cycle. If neither diagonal occurs, Proposition~\ref{prop:twosumspathtriangle} excludes the induced missing path. Otherwise reverse the path so $ac\in G$. This diagonal lies in $H$, since otherwise $c$ has three missing neighbors.

No old form can contain both endpoints of a length-two $G$-path having at least one missing edge: its old center differs from the middle vertex and gives a four-cycle. Apply this to the cross pairs between $\{a,b\}$ and $\{c,d\}$: the relevant paths for $ac,ad,bc,bd$ are $a-b-c$, $a-c-d$, $b-a-c$, $b-c-d$. Thus no nonempty old product is potential on both end edges $ab,cd$. A safe old product witnesses at most one of the three path edges; adjacent witnesses are excluded by the shared-endpoint argument in the preceding proof.

Use the same complete two-addition normalization. All cases with at most two safe old products fail. The remaining triples can be written
\[
 gL,\qquad VC,\qquad Vg,
\]
with an additional actual-context property: the first nonzero old low $A$ satisfies $AL\subseteq H$, and every factor of $C$ contains $\supp(A)$. For an independent second high and positive cross term, $C=M$ and the mandatory mixed context gives $A\subseteq M$. For an affected high $U(F\vee g)$ with zero cross coefficient, use $V=UM$, $C=F$; every affected old factor in $F$ contains $A$. With a positive cross coefficient, use $V=U$, $C=FM$, again using $A\subseteq M$. These are properties of the original cut occurrences, not assumptions about an arbitrary triple.

Each term would have to witness a distinct path edge. Unit cases cannot evade this: if $V=1$, the safe term $Vg=g$ absorbs the witnesses of $gL$; if $g=1$, it similarly absorbs those of $VC$. Thus both shared factors are nonempty. If $Vg$ witnesses an end edge, one of its shared factors is potential on both end edges, which was just excluded. Hence it witnesses $bc$.

Suppose $gL$ witnesses $ab$ and $VC$ witnesses $cd$. Then $g$ is true on singleton $b$, and $V$ on singleton $c$. Safety forces the context $L$ to be exactly the singleton $a$. It cannot be scalar; for a linear context, every support vertex is a neighbor of $b$, and any second vertex together with its old center and the missing neighbor $a$ would give a four-cycle. Choose $w\in\supp(A)$. The old product $AL$ gives $aw\in H$. Every factor of $C$ contains $w$, so $C$ is true on $w$; safety of $VC$ gives $cw\in G$. The vertices $a,b,c,w$ are distinct: old $aw$ excludes $w=a,b$, and safety excludes $w=c$. They form the four-cycle $a-b-c-w-a$. If the end terms are reversed, the same argument gives $d-c-b-w-d$. Both alternatives are impossible.

There is therefore no three-edge path in $J$. A finite graph of maximum degree two consists of paths, cycles and isolates. All cycles and all paths of at least three edges have been excluded. The nontrivial components are exactly of the asserted possible types $K_2$ and $P_3$. No realization claim is made.
\end{proof}

\begin{proposition}[The induced exceptional graphs with an arbitrary frame]\label{prop:exceptionalarbitrary}
Let $G$ be finite and four-cycle-free, and let a positive variable-input circuit compute $P\in\Sup(G)$ with nonempty quadratic support $H$. Suppose $H[T]$ is empty and, apart from isolates, $G[T]$ is $P_4\sqcup K_2$ or $K_3\sqcup K_2$. Canonical extraction cannot delete at most three products and at most two additions. Retained counts, cover gain, other vertices, and other missing edges are unrestricted.
\end{proposition}
\begin{proof}
The target contains an induced missing $P_4$ or $K_3$. Proposition~\ref{prop:twosumspathtriangle} requires three deleted additions, independently of the product count.
\end{proof}

\begin{proposition}[Two nonancestral products and three visible cuts]\label{prop:twooldthreecuts}
Let $G$ be finite and four-cycle-free, and let a positive variable-input circuit compute $P\in\Sup(G)$ with nonempty quadratic support $H$. Suppose $H[T]$ is empty and, apart from isolates, $G[T]$ is $P_4\sqcup K_2$ or $K_3\sqcup K_2$. At most two deleted products and at most three deleted additions are impossible if the products are pairwise nonancestral in the original circuit and every deleted addition is a visible unequal-degree cut. Visibility means that the cut's own node lies in the selected extraction subgraph. Retained counts, cover gain, other vertices and other missing edges are unrestricted.
\end{proposition}
\begin{proof}
Write $J$ for the three path or triangle edges, and $uv$ for the separate edge. Exact full-carrier erasure preserves the frame. If only one useful product remains, its three old-context cubics have missing witnesses forming a star, or a square contribution makes the product itself safe; neither supplies the cover-three target. With two useful nonancestral products, every operand is an unperturbed old linear signal, since any earlier inserted product reaching it would violate original nonancestry. Thus write
\[
 g=UV,\qquad f=WX.
\]
Every visible high is one of these products. All cuts retain degree one, since both highs have degree two. There are exactly three cuts by Proposition~\ref{prop:twosumspathtriangle}.

Treating their highs as independent insertion symbols gives three direct old linear contexts and possible pairwise cross terms. A cut with low $A_i$ and context $L_i$ has $A_iL_i\subseteq H$. No individual symbol is squared, since its nonzero low would then give a square in $P_2$. Distinct cuts carrying the same high may, however, give $g^2$ or $f^2$. If $g^2$ is present, $g$ is safe and all its multiples have only its own missing witnesses, at most one by the old-factor bound. The remaining at most two $f$-context cubics each add at most one. Thus fewer than four targets occur; the same holds for $f^2$. Without a mixed $gf$ term the three cubics are insufficient. We are left with
\[
 gL_1,\quad gL_2,\quad fM,\quad gf,
\]
after possibly exchanging the highs. Each displayed term is safe globally. The cubics have at most one missing witness each, and the quartic has at most two. Its target witnesses form an induced matching.

Suppose first that the quartic supplies exactly one target. All three cubics must supply three distinct others. The two $g$-cubic edges share a vertex $x$ by Lemma~\ref{lem:commoncubic}; they are adjacent edges $xy,xz$ of $J$. Every old factor of $g$ meets both and hence contains $x$, so $g$ has a pure power of $x$.

If the quartic target is $uv$, the $f$-cubic witnesses the remaining $J$-edge. Both factors of $f$ also meet $uv$, so neither can be the singleton factor in the cubic witness. Its context is therefore a singleton, and both factors of $f$ contain the other endpoint $s$ of that $J$-edge. The product $gf$ then has a monomial on $\{x,s\}$. For a path or triangle this is an earlier target, a nonedge, or a singleton, so the quartic either has a second target or violates safety. If instead the quartic supplies the third $J$-edge, the $f$-cubic witnesses $uv$. The same singleton-factor argument gives a pure power of $u$ or $v$ in $f$, and its product with the pure $x$ is a nonedge between the components. Both cases are impossible. No quartic target is plainly insufficient.

It remains that $gf$ has two targets, necessarily $uv$ and one $J$-edge $ab$. Equality in Proposition~\ref{prop:oldfactors} confines its four full factor supports to the four transversals of these two edges. Each is nonsingleton. A pair of these factors has either empty intersection or a singleton endpoint; in the latter case the complementary pair meets at its mate. An additional target cubic therefore has a singleton context and its repeated endpoint is that common endpoint. Intersections at $u,v$ give no additional $J$-edge. An end edge of a path is also impossible, because its opposite end has neither endpoint in the quartic's target set. Thus $ab$ is the path middle edge or a triangle edge, and the highs have pure powers of $a,b$. The two further $J$-edges are $az,bw$, where $z=w$ in the triangle, and require distinct high-cut occurrences with respective full contexts $\{z\},\{w\}$.

Any retained product supplying $gf$ uses a $g$-cut and an $f$-cut. If its $f$-cut is also the $bw$ witness occurrence, choose $h$ in that cut's nonzero old low. Its direct context gives $hw\in H$. The same retained product contributes $g$ times that old low, so safety and the pure power of $a$ give $ah\in G$. The four vertices $h,w,b,a$ are distinct, since $H[T]$ is empty, and form a four-cycle. Therefore the mixed-term $f$-cut differs from the $f$-witness cut. The symmetric argument separates the two $g$-cuts. Four cut occurrences would be necessary, contradicting three. This uses actual full contexts and remains valid with selected-path dependencies and arbitrary fanout.
\end{proof}

\begin{proposition}[Two dependent products and three visible cuts]\label{prop:twodependentthreecuts}
Under the graph and target hypotheses of Proposition~\ref{prop:twooldthreecuts}, at most two deleted products and at most three deleted additions are also impossible when the products are ancestrally related in the original circuit and every deleted addition is a visible unequal-degree cut. Retained counts, cover gain, other vertices and other missing edges are unrestricted.
\end{proposition}
\begin{proof}
The target has four edges and cover number three. Full-carrier erasure cannot remove a private product, by Proposition~\ref{prop:oneproductcover}, or a private addition, by Proposition~\ref{prop:twosumspathtriangle}. Thus the exact useful profile remains two products and three cuts. Original ancestry survives: if an erased occurrence broke a path between the products, its first subsequent private use would either erase the later product or remove a visible high entry, contradicting those same two necessities. With products ordered $g,f$, the first is $g=UV$ on old linear forms. The only other nonlinear private signal preceding $f$ is $g$; every retained affected linear operand is $A+\alpha g$, with old $A$. Boolean absorption leaves
\[
 f=gW,\qquad f=g,\qquad f=AB\vee g,
 \qquad f=AW\vee gW.
\]
Raw $g$ operands and squares are included. Every cut inserts $g$ or $f$, and at least one inserts $f$. If none inserts $g$, only the first two forms can occur. The frame gives one actual scalar or old linear context per cut and pairwise cross terms. An individual insertion symbol is not squared, because its nonzero old low would give a quadratic square; distinct cuts carrying the same high can produce $g^2$ or $f^2$. Counts refer throughout to original cut occurrences.

For $f=g$, the missing witnesses form one common-cubic star; a scalar or square term makes $g$ itself safe and gives the same conclusion. For $f=AB\vee g$, the witnesses lie in two common-cubic stars, based on $UV$ and $AB$. Cross and square terms are absorbed by these same signals. Their union has a two-vertex cover, so cannot contain the target.

Suppose $f=gW$. A positive $gf$ or $f^2$ term makes $gW$ itself safe and absorbs its context multiples. Its witnesses and those of the $g$-context cubics then form a common-cubic star. A $g^2$ term similarly makes $g$ safe and absorbs all its multiples. Otherwise only $g$-context cubics and $f$-context quartics remain. If no quartic has two missing witnesses, the three direct-context terms have at most three in total. If a quartic has two, equality in Proposition~\ref{prop:oldfactors} confines its three shared factors $U,V,W$ and all other quartic witnesses to that same pair. With zero or one $g$-cut, at most one further edge can occur. With two $g$-cuts there is one quartic: any additional cubic edge meets the common endpoint of the two transversals supporting $U,V$, if that endpoint exists. That endpoint and an endpoint of the other quartic edge cover the whole missing graph. If the transversals are disjoint, no additional edge is possible. This explicitly includes zero $g$-cuts and three $f$-cuts, and scalar contexts only reduce the bounds.

It remains that $f=AW\vee gW$ comes from an affected operand. Both products have minimum degree two, so their cuts retain degree one. There are either two $g$-cuts and one $f$-cut, or one $g$-cut and two $f$-cuts. A $g^2$ term makes $g$ safe, leaving a safe $g$-star and an $AW$ common-cubic star. An $f^2$ term requires two $f$-cuts and leaves one $g$-cut; safe $AW$ and safe $gW$ each have at most one missing witness, and the remaining $g$-context cubic adds at most one. A $gf$ term is Booleanly $gW$; it absorbs every $gW$ context multiple, leaving the two common-cubic families based on $UV$ and $AW$. None can contain the target. We may assume all cross coefficients vanish, giving
\[
 \sum_i gD_i+\sum_j AWM_j+\sum_j gWM_j.
\]
All displayed positive products are safe on the full graph.

Choose an actual $g$-cut occurrence affecting the operand $A+\alpha g$, with full direct context $D_0$ and extracted low $a$. Then $a\ne0$, $\supp(a)\subseteq\supp(A)$, and $aD_0$ is supported on $H$. Some $k\in\supp(D_0)$ also satisfies $kA\subseteq H$. If the product reads this cut directly, choose any context variable. Otherwise follow the affected retained linear signal along its selected path to its first retained product. Its opposite old form $B$ satisfies $AB\subseteq H$, and the chosen insertion contributes a positive multiple of $B$ to $D_0$. This remains true through other selected cuts and with arbitrary fanout: treat each cut occurrence as an independent symbol when taking its direct coefficient. Choose $k\in\supp(B)$.

Lemma~\ref{lem:independentproducts}, applied to $UV$ and $WM_j$ with contexts $D_0,A$ and lows $a,\{k\}$, now gives at most two missing witnesses in
\[
 gD_0,\qquad AWM_j,\qquad gWM_j.
\]
Its old products are $aD_0,kA,ak$ and its support containments are exactly the ones just proved. With two $g$-cuts and one $f$-cut, the sole remaining $g$-context cubic adds at most one witness.

With one $g$-cut and two $f$-cuts, put $R_j=AWM_j$ and $Q_j=gWM_j$. If one $Q_j$ has two missing witnesses, the shared-quartic equality confines both $Q$ terms to this pair. The triple bound confines $R_j$ and $gD_0$ to the same pair, and the other $R$ adds at most one. Otherwise each $Q_j$ has at most one. Four distinct targets would require four distinct edges from $R_1,R_2,Q_1,Q_2$, and no missing edge from $gD_0$, by the two triple bounds. The two $R$ edges share a vertex, say $xy,xz$. Old supports $A,W$ meeting both contain $x$, and the singleton-factor argument forces $M_1=\{y\}$ and $M_2=\{z\}$. If the distinct $Q_1$ edge is $yv$ with $v\ne x$, then $W$ contains $x$ and meets $\{y,v\}$. It cannot also contain $v$: its old center would form a four-cycle through the missing edge $xy$, or make $xy$ old. Hence $y\in\supp(W)$. Similarly $z\in\supp(W)$. The old form $W$ would contain two distinct missing neighbors of $x$, again impossible. Every case has been excluded.
\end{proof}

\begin{proposition}[One visible cut and independent sets]\label{prop:onevisibleclauses}
Let a positive variable-input circuit compute $P\in\Sup(G)$ for a finite simple graph $G$, with nonzero squarefree quadratic part supported on $H$. Suppose exactly one unequal-degree addition is visible in the selected quadratic frame, and let $p$ be the number of deleted multiplication gates. If $H[T]$ is empty and $G[T]$ has an edge, the number of maximal independent sets of $G[T]$ is at most $2^p+1$. If the visible cut retains degree two, the bound is $2^p$. No four-cycle hypothesis or bound on hidden additions or retained gates is required. In particular, an induced $H$-empty $P_4\sqcup K_2$ or $K_3\sqcup K_2$ is impossible when $p\le2$.
\end{proposition}
\begin{proof}
Erase only off-frame uses of full current quadratic carriers, preserving the exact selected frame and Boolean output as in Proposition~\ref{prop:recursivenormal}. Write $F$ for the high input of the sole visible cut and $A$ for its nonzero extracted low. Every useful private branch must eventually enter the selected frame through a visible cut. Thus every useful private gate is ancestral to $F$; a later private branch would require another visible entry. A retained linear boundary in this ancestral computation is unaffected, since an earlier insertion into its selected ancestry would be a second visible cut. Retained quadratic carrier boundaries have been erased. Equivalently, before erasure every selected boundary preceding this sole cut is homogeneous, and its quadratic boundaries vanish after the restriction to $T$. Hence the restricted high computation has at most $p$ conjunctions, arbitrary disjunctions, and old linear boundary forms, which themselves use disjunctions only. No retained conjunction has been counted as a free nonlinear boundary.

Temporarily replace the high insertion by an independent symbol. If $A$ has degree two, its downstream selected context is scalar. If $A$ has degree one, its context is an old linear form $L$. The insertion cannot reach both factors of a retained product: choosing the same variable of $A$ in both factors would put a square in the quadratic output. Positivity forbids cancellation of that square. Arbitrary fanout merely adds the linear contexts. Therefore, after Booleanization and setting variables outside $T$ to zero, the exact target is
\[
 f_{G[T]}=F_T L_T\quad(\deg A=1),\qquad
 f_{G[T]}=F_T\quad(\deg A=2).
\]
The nonzero target excludes a zero necessary factor. All gate counts here refer to the original extraction, not a new extraction of the simplified expression.

A monotone Boolean circuit with at most $p$ conjunctions has a positive CNF with at most $2^p$ clauses: choose one input at each conjunction, obtaining an OR-only resolution. The conjunction of these resolutions equals the circuit. Indeed, on a false assignment choose a false input at every false conjunction; this preserves every false wire, while every resolution is an upper bound for the original circuit. This argument permits sharing and constants. Apply it to $F_T$. The linear context adds at most one clause, while a scalar context adds none.

Every maximal independent set $I$ of $G[T]$ must falsify some clause of any positive CNF for its edge function. The support of that clause is exactly $T\setminus I$: it is disjoint from $I$, and omitting a vertex $v\notin I$ would also falsify the clause on $I\cup\{v\}$, which contains an edge by maximality. Distinct maximal independent sets therefore require distinct clauses. This proves both bounds, including isolated vertices. The path $P_4$ and the triangle each have three maximal independent sets, and a disjoint edge has two. Their respective disjoint unions have six, exceeding $2^2+1=5$.
\end{proof}

\begin{example}[A hidden cut realizes the exceptional target]\label{ex:hiddencutscope}
The visibility hypothesis in Propositions~\ref{prop:twooldthreecuts} and~\ref{prop:twodependentthreecuts} cannot be dropped using the induced target alone. Take distinct variables and put
\[
 U=u+a,\quad V=v+a,\quad W=u+b,\quad M=v+b,\quad g=UV.
\]
Compute these four sums and $g$ once, with reuse, and set
\[
 \begin{split}
 P={}&h_U(U+\ell_U)+h_V(V+\ell_V)+h_W(W+\ell_W)\\
    &+\bigl(h_M+(c+g)W\bigr)M+(h_D+g)d+h_Dz.
 \end{split}
\]
Let $H$ be its quadratic support: the stars with centers $h_U,h_V,h_W$ have leaves $\{u,a,\ell_U\}$, $\{v,a,\ell_V\}$, $\{u,b,\ell_W\}$ respectively; the remaining old edges are $h_Mv,h_Mb,h_Dd,h_Dz$. Thus $|E(H)|=13$. Let $G$ add the four edges $uv,ab,ad,bc$.

Booleanly $UV=a\vee uv$, $WM=b\vee uv$, and $UVWM=ab\vee uv$. Consequently
\[
 P=P_2+cWM+UVWM+UVd
\]
has Boolean function $f_H\vee uv\vee ab\vee bc\vee ad$. The extra terms have degrees three, four and three. Their exact witnesses include $u^2v^2,a^2b^2,b^2c,a^2d$, so $P\in\Sup(G)$ and its exact quadratic support is $H$. On $\{u,v,a,b,c,d\}$ the missing induced graph is $d-a-b-c$ together with $uv$.

The cyclic core is $K_4$ on $u,v,a,b$ with opposite edges $uv,ab$ unsubdivided and each other edge subdivided once by its own $h$ variable. Triangles of this $K_4$ become five-cycles and four-cycles become six- or eight-cycles. All remaining edges attach trees. Hence $G$ has girth five.

The five centers cover $H$, and the matching
\[
 h_U\ell_U,\quad h_V\ell_V,\quad h_W\ell_W,
 \quad h_Mv,\quad h_Dz
\]
proves $\vc(H)=5$. The set $\{v,a,b,h_U,h_V,h_W,h_D\}$ covers $G$, while
\[
 h_U\ell_U,\ h_V\ell_V,\ h_W\ell_W,\ h_Dz,
 \ uv,\ ad,\ bc
\]
is a seven-edge matching. Thus $\vc(G)=7$.

The quadratic frame uses seven linear additions, six products and five output additions: eighteen gates. The deleted gates are precisely the two products $g$ and $(c+g)W$ and the three additions $c+g$, $h_M+(c+g)W$, $h_D+g$. The last two cuts are visible, whereas $c+g$ lies only in a pruned high branch. The total is twenty-three gates, exactly $|E(G)|+\vc(G)-1$. The retained frame has one extra product and the minimum $|E(H)|-1=12$ additions; the cover gain is two, not three.

Omitting $\ell_U,\ell_V,\ell_W,z$ and the summand $h_Dz$ gives a smaller example. Its old graph has nine edges, its full graph thirteen, and both cover numbers are five: the cover $\{u,v,a,b,d\}$ and the old matching $h_Ua,h_Vv,h_Wu,h_Mb,h_Dd$ certify these values. Its thirteen-gate quadratic frame is optimal, and the same five deletions give eighteen total gates, one above its star count. Thus even an optimal frame can realize this target once a hidden cut is permitted, but neither example gives a saving or cover gain three. Both exact shared circuits are accompanied by separate symbolic and finite computational checks in the workbench.
\end{example}

\begin{proposition}[A hidden dependent sum and the cover number]\label{prop:hiddendependentcover}
Let $G$ be finite and four-cycle-free, and let a positive variable-input circuit compute $P\in\Sup(G)$ with nonempty quadratic support $H$ and exactly four missing edges. Let $K$ be the number of products in its pruned quadratic frame. Suppose that, after full-carrier erasure preserving that frame, the useful private computation consists of
\[
 g=UV,\qquad t=C+g,\qquad f=tW,
\]
and two visible unequal-degree cuts inserting $g$ and $f$. Here $U,V,W,C$ are nonzero old homogeneous linear forms, positive coefficients are allowed, and the displayed hidden sum and the two visible cuts are the only private additions. Then
\[
 \vc(G)\le K+1.
\]
There is no optimality, coefficient-one, or distinct-center assumption on the retained frame. In particular, $K\le\vc(H)+1$ implies cover gain at most two.
\end{proposition}
\begin{proof}
Both highs have minimum degree two, so both cuts retain nonzero linear low parts. Treating the high inputs as independent symbols, expansion of the exact quadratic frame gives the Boolean output
\[
 P_2\vee gD\vee fM\vee qgf,
\]
where $D,M$ are the full actual nonzero homogeneous linear contexts and $q\ge0$. An individual insertion cannot be squared: substituting its nonzero low part would put a square in $P_2$. Each context is old, since its product with the corresponding low part is supported in $H$. This includes every fanout contribution and output multiplicity.

If $q>0$, substitution of $f=(C\vee g)W$ and absorption reduce the perturbation to
\[
 gD\vee CWM\vee gW.
\]
These are three safe old cubics, each with at most one missing witness by Proposition~\ref{prop:oldfactors}. They cannot supply four missing edges. Thus $q=0$, and the three relevant safe terms are
\[
 UVD,\qquad CWM,\qquad UVWM.
\]
Their missing-witness bounds are respectively one, one, and two. Equality is therefore forced throughout, with disjoint witness sets.

The quartic equality case of Proposition~\ref{prop:oldfactors} confines the full supports of $U,V,W,M$ to the four distinct two-element transversals of its two missing edges, which form an induced matching in the full graph $G$. Since $U,V$ are nonsingleton, the missing cubic $UVD$ forces $D$ to be a singleton and $U,V$ to meet at its opposite endpoint. Similarly, $C$ is a singleton and $W,M$ meet at the opposite endpoint of the other cubic witness. Relabeling gives
\[
 \begin{gathered}
 \supp U=\{a,u\},\quad \supp V=\{a,v\},\quad \supp W=\{b,u\},\\
 \supp M=\{b,v\},\quad \supp D=\{d\},\quad \supp C=\{c\}.
 \end{gathered}
\]
The four missing edges are $uv,ab,ad,bc$. The vertices $c,d$ lie outside $\{a,b,u,v\}$, since those four vertices induce exactly $uv,ab$; $c=d$ is allowed. Thus $R=\{d,b,v\}$ covers every missing edge, and the two full context supports $\supp D$ and $\supp M$ are disjoint.

Mark a retained product if one of its linear inputs is affected by the first insertion, and mark it separately if affected by the second. Every first-marked product has its entire opposite old support contained in $D=\{d\}$: its positive contribution occurs in the full actual context. Consequently $d$ covers every old edge supplied by that product. Likewise $\{b,v\}$ covers every old edge supplied by a second-marked product.

The two marked sets are disjoint. If both insertions enter opposite operands of one retained product, it contributes a positive mixed coefficient $q$, already excluded. If they enter the same operand, the nonzero opposite old support occurs in both $D$ and $M$, contradicting their disjointness. An individual insertion cannot enter both operands by the same squarefree argument used above. Both marked sets are nonempty because the cuts are visible and the quadratic frame is pruned.

Every other retained product is a star. Choose one center for each of these at most $K-2$ products and adjoin $R$. This covers all of $H$, product by product, including old edges between chosen centers and all edges outside the missing endpoints. It covers the missing edges as well. Repeated centers only decrease its cardinality. Hence
\[
 \vc(G)\le (K-2)+3=K+1,
\]
as required. No Boolean re-extraction or alteration of the original frame count was used.
\end{proof}
The extended circuit in Example~\ref{ex:hiddencutscope} has $K=6$ and $\vc(G)=7$, so this bound is attained. The proposition excludes precisely the displayed hidden-dependent network from the cover-gain-three extra-product row. It does not classify other placements of a private linear sum or a hidden addition.

\begin{proposition}[A hidden unequal-degree addition]\label{prop:hiddenunequal}
Let $G$ be finite and four-cycle-free, and let a positive variable-input circuit compute $P\in\Sup(G)$ with nonempty quadratic support $H$ and exactly four missing edges. Suppose canonical extraction deletes two products and three additions, exactly two of the additions being visible cuts in the selected frame. If the remaining hidden addition has unequal original input minimum degrees, and the retained product count satisfies $K\le\vc(H)+1$, then
\[
 \vc(G)\le\vc(H)+2.
\]
Retained addition counts, fanout, repeated inputs, and outside vertices are unrestricted.
\end{proposition}
\begin{proof}
Suppose instead that the cover gain is at least three. Choose a minimum $H$-cover $C_0$ and put $N=G[V\setminus C_0]$, retaining isolates. It has at most four edges and cover number at least three. The small-graph enumeration in Proposition~\ref{prop:gainthreesix} gives an induced missing three-edge matching or one of $P_4\sqcup K_2,K_3\sqcup K_2$. Corollary~\ref{cor:threedeletedproducts} excludes the matching with only two private products. Thus $N$ is one of the exceptional graphs, contains all four missing edges, and has a separate edge component.

Every private gate is essential under a frame-preserving Boolean bypass: deleting one would leave at most four discarded gates, contrary to Lemma~\ref{lem:fivelost}. Apply full-current-quadratic-carrier erasure. A retained quadratic carrier, including all its higher terms, is already an output disjunct; its additional multiples and off-frame summands are absorbed. Such an erasure cannot make a private gate dispensable. Order the private products as $g,f$ and denote the hidden unequal addition by $h$.

The node $h$ cannot precede $g$: after carrier erasure all its available inputs there are old linear forms of equal degree. It cannot follow $f$ with minimum degree one, since no private product remains to take a hidden linear signal back to the selected frame. If both operands of $h$ have degree at least two, they must be $g,f$; the first product has degree two, and the larger-degree second product must contain $g$ as a factor. Then $f$ implies $g$ Booleanly and $h=g\vee f=g$, a forbidden bypass. A retained carrier operand was already excluded.

Therefore $h$ lies between the products, $g=UV$ on old forms, and $h$ has minimum degree one. Its high is $g$, and its low is an old linear form $A$, possibly already carrying $g$ through a visible cut. In the latter case that added high is absorbed by the low and $h$ is dispensable. Hence $h=A+g$ with old nonzero $A$. It feeds $f$ directly: an intervening selected linear addition would make $h$ visible, while another private addition is unavailable. The other operand of $f$ is an old form $W$ possibly carrying $g$, the raw $g$, or $h$ itself. Booleanly the complete list is
\[
 f=(A\vee g)W,\qquad f=AW\vee g,\qquad f=g,\qquad f=A\vee g.
\]
The last form comes from $h^2$. These are identities on the original network, not new circuits obtained by re-extraction.

Each visible high is $g$ or $f$; a carrier high is absorbed, and the degree-one $h$ cannot be the high of an unequal cut. At least one cut inserts $f$, so the allocations are $g,f$ or $f,f$. Their exact frame expansions are respectively
\[
 P_2\vee gL\vee fM\vee qgf,
 \qquad
 P_2\vee fM_1\vee fM_2\vee qf^2,
\]
with full old linear contexts, except that a higher-degree raw $f$ may have a scalar context when its cut retains degree two. Such a scalar makes $f$ safe and only simplifies the following bounds. Individual insertion squares are forbidden by the nonzero low square in $P_2$.

For the $g,f$ allocation, $f=g$ leaves a common-cubic missing star, or safe $g$. The forms $f=A\vee g$ and $f=AW\vee g$ give at most three safe old products with zero cross; a positive cross makes $g$ safe and absorbs its multiples, leaving at most two missing edges. For $f=(A\vee g)W$, a positive cross leaves the three cubics $gL,AWM,gW$, again insufficient. With zero cross the full actual contexts and factorization are precisely those in Proposition~\ref{prop:hiddendependentcover}; its product-by-product argument gives $\vc(G)\le K+1\le\vc(H)+2$.

For the $f,f$ allocation, a scalar context or $q>0$ makes $f$ safe. Its displayed expansion then has at most two missing edges, except that $A\vee g$ cannot be safe because of its nonzero linear term. Otherwise both contexts are old linear and $q=0$. The forms $g$, $AW\vee g$, and $A\vee g$ give respectively one common missing star, two common missing stars, or the safe complete bipartite star $A(M_1\vee M_2)$ together with a common-cubic star. These unions have cover number at most two and cannot contain $N$.

The remaining form gives four safe terms
\[
 AWM_1,\qquad AWM_2,\qquad UVWM_1,\qquad UVWM_2.
\]
The two cubics have at most one missing edge each; the shared-quartic lemma bounds the last pair by two. Four targets require distinct cubic witnesses, which share a vertex: write them $xy,xz$. Both old forms $A,W$ meet both pairs and therefore contain $x$, since containing the two missing neighbors instead would give a four-cycle through an old center. Safety then puts every variable of $M_1,M_2$ in $N_G(x)$; each contains its respective missing neighbor $y,z$. An old neighborhood containing that missing neighbor and a second neighbor of $x$ would again give a four-cycle or make the edge old. Thus $M_1=\{y\}$ and $M_2=\{z\}$ as full supports. These vertices lie in the path or triangle component of $N$. Every displayed term contains one of them, so none can witness the separate edge component. This final contradiction proves the claim.
\end{proof}
This classification concerns an unequal-degree hidden addition. The homogeneous hidden-addition case is not included in its hypotheses.

\begin{corollary}[The hidden unequal-cut bound for any retained count]\label{cor:hiddenunequalany}
Under the graph and original private-gate hypotheses of Proposition~\ref{prop:hiddenunequal}, the stronger bound
\[
 \vc(G)\le K+1
\]
holds for every retained product count $K$, without a cover-gain hypothesis. Consequently the hidden unequal-degree two-product, three-addition branch is excluded from the cover-gain-two, five-deletion row with an optimal quadratic remainder. The bound is attained by the extended circuit of Example~\ref{ex:hiddencutscope}.
\end{corollary}
\begin{proof}
The essentiality and original-degree normalization in Proposition~\ref{prop:hiddenunequal} precede its use of the exceptional outside graph and require no cover assumption. They give $g=UV$, hidden $h=A+g$, the four listed forms for $f$, and the two allocations of visible highs. In the $g,f$ allocation every case has at most three missing witnesses except $f=(A\vee g)W$ with zero cross. Proposition~\ref{prop:hiddendependentcover} gives $\vc(G)\le K+1$ for precisely that remaining network.

It remains to exclude four global missing edges in the $f,f$ allocation. A scalar context or positive cross makes $f$ safe; its multiples introduce no new exact edge witnesses and its listed forms have at most two. Thus both contexts $M_1,M_2$ are nonzero old linear forms and the cross is zero. The form $f=g$ gives at most two cubic witnesses. There is no visible $g$ cut in this allocation, so the other operand cannot be independently affected by $g$: the form $AW\vee g$ is unavailable. Only $A\vee g$ and $(A\vee g)W$ remain.

We use the full-support version of common-cubic rigidity. If safe old cubics $RSM_1,RSM_2$ have distinct missing witnesses, they share a vertex $x$. Both $R,S$ contain $x$: otherwise an old neighborhood would contain two distinct missing neighbors of $x$, giving a four-cycle or making the edges old. Safety then makes every variable of either context a neighbor of $x$. Since each context contains a missing neighbor, the same old-neighborhood argument forces $M_1=\{m\},M_2=\{n\}$ with $m\ne n$, and the two witnesses are $xm,xn$. All supports here are full supports in $G$.

For $f=A\vee g$, four witnesses require distinct witnesses in $AM_1,AM_2,UVM_1,UVM_2$. The safe old quadratics force $A=\{a\}$ and singleton distinct contexts $m,n$. The cubic pair gives witnesses $xm,xn$ with $x\ne m,n$. If $x=a$ these repeat the first pair; otherwise $a,m,x,n$ form a four-cycle. Thus this case is impossible.

Finally consider the four safe terms
\[
 AWM_1,\quad AWM_2,\quad UVWM_1,\quad UVWM_2.
\]
The cubics supply at most one missing witness each, and the shared quartics at most two together. Four witnesses therefore force distinct cubic edges $xm,xn$, singleton full contexts $M_1=\{m\},M_2=\{n\}$, and $x\in A\cap W$. Each quartic has at most one additional missing edge: two would force four two-element transversal factors by the old-factor equality theorem, contradicting its singleton context. Write the two additional edges $my,nz$, with $y,z\ne x$.

The old form $W$ cannot contain $y$, since it already contains $x$ and both $xm,ym$ are missing; its old center would give a four-cycle. To meet the exact witness $my$, it must therefore contain $m$. Likewise it contains $n$. But an old neighborhood containing the two distinct missing neighbors $m,n$ of $x$ again gives a four-cycle or makes those edges old. This contradiction excludes the last $f,f$ case and proves the bound. Hidden equal-degree additions and all-visible configurations are not covered by this corollary.
\end{proof}

\begin{proposition}[Two nonancestral products and three visible cuts]\label{prop:nonancestralthreecuts}
Let $G$ be finite and four-cycle-free, and let a positive variable-input circuit compute $P\in\Sup(G)$ with nonempty quadratic support $H$. Suppose at most two original products and at most three original additions are deleted. If every deleted addition is a visible unequal-degree cut and the original deleted products are pairwise nonancestral, then at most three $G$-edges are missing from $H$. Retained counts, cover gain and outside induced subgraphs are unrestricted.
\end{proposition}
\begin{proof}
Full-current-carrier erasure preserves the exact selected frame and cannot create an ancestry absent from the original DAG. A surviving private product cannot read another private product or a retained quadratic carrier. A private linear sum would be hidden, contrary to the hypothesis. Thus the surviving products are $g=UV$ and, if present, $f=WX$, on old homogeneous linear forms.

Every surviving useful cut inserts one of these raw products into a selected degree-one signal. A separate sum of them would be hidden; a retained quadratic high is absorbed; a visible-cut output has minimum degree one and cannot be a high above a nonzero low. A degree-two selected low would require a higher-degree high, which is unavailable. Introduce independent symbols $z_i$ for the at most three cut highs, keeping their extracted old lows $a_i$. The exact fixed quadratic frame expands as
\[
 P_2+\sum_i z_iL_i+\sum_{i<j}q_{ij}z_iz_j,
 \qquad q_{ij}\ge0.
\]
Each nonzero $L_i$ is old and homogeneous linear, because $a_iL_i$ is supported on $H$. A term $z_i^2$ would give a square from the nonzero old low in $P_2$, and is impossible. Nested cuts and fanout change coefficients but not this expression.

If $q_{ij}>0$, positivity in the original retained product gives
\[
 a_ia_j,\ a_iL_i,\ a_jL_j\text{ supported on }H,
 \qquad
 \supp(a_i)\subseteq\supp(L_j),\quad
 \supp(a_j)\subseteq\supp(L_i).
\]
These are the actual-low hypotheses of Lemma~\ref{lem:independentproducts}. They continue to hold when a third symbol is set to zero: its old low stays in the same fixed frame. After substituting each high $g$ or $f$, every displayed direct or cross summand is safe in the full graph. A nonzero direct term is an old cubic and has at most one missing witness.

If all crosses vanish, three direct cubics give the bound. If a positive cross joins two copies of $g$, safety of $g^2$ implies safety of $g$. Every multiple of $g$ then has no new exact edge witness beyond $g$ itself, which has at most one missing edge. Two cuts already insert $g$, leaving at most one direct cubic from $f$, so there are at most two witnesses. The symmetric case is identical.

Otherwise every positive cross joins different highs. Choose one such pair. The three terms $gL_i,fL_j,gf$ satisfy the safety-only form of Lemma~\ref{lem:independentproducts} and jointly supply at most two missing edges. There is at most one other cut, whose direct cubic adds at most one. Every remaining positive cross is the same product $gf$ up to coefficient, so adds no new support. This proves the bound. Zero contexts remove terms; a positive cross already forces both relevant contexts nonzero. Neither a reduced subsum computing the whole graph nor a fresh extraction was assumed.
\end{proof}

\begin{proposition}[The all-visible two-product cover bound]\label{prop:allvisibletwocover}
Let $G$ be finite and four-cycle-free, and let a positive variable-input circuit compute $P\in\Sup(G)$ with nonempty quadratic support $H$ and exactly four missing edges. Suppose at most two products and three additions are deleted, and all deleted additions are visible unequal-degree cuts. If $K$ products are retained, then
\[
 \vc(G)\le K+1.
\]
Neither the retained addition count nor quadratic optimality is assumed.
\end{proposition}
\begin{proof}
Lemma~\ref{lem:fivelost} makes the private counts exactly two and three and every private gate essential under exact-frame erasure. If the products are originally nonancestral, Proposition~\ref{prop:nonancestralthreecuts} gives at most three missing edges. If erasure removes an original ancestry path, apply that proposition's old-two-product, three-symbol counting argument directly to the same preserved frame; no original label or extraction is recomputed. Otherwise order the products $g,f$ by their surviving dependence. The first is $g=UV$ with old factors. Every visible high is $g$ or $f$: a selected minimum-one cut cannot be a higher input above a nonzero low, and off-frame quadratic carriers were erased. There is no hidden sum.

An operand of $f$ is raw $g$ or an actual selected linear input $A+\alpha g$. Boolean expansion therefore gives exactly the raw forms $f=gA,f=g$, or the affected forms
\[
 f=AB\vee gB,\qquad f=AB\vee g,
\]
with nonzero old $A,B$. If a visible cut keeps degree two, its scalar context makes its high safe. Only a raw higher-degree $f$ can be that high. Its at most one old-cubic witness, together with at most two remaining $g$-context cubics, gives at most three. Thus all three cuts keep degree one, and their full direct contexts are nonzero old linear forms.

Use independent symbols for these cuts. The original frame identity has direct terms $z_iL_i$ and cross terms $q_{ij}z_iz_j$, with $q_{ij}\ge0$, no diagonal square, and no higher symbol degree. Its lows satisfy $a_iL_i\subseteq H$; a positive cross gives $a_ia_j\subseteq H$, $a_i\subseteq L_j$ and $a_j\subseteq L_i$. If cut $i$ actually affects an old part $A$, its low lies in $A$ and some $k\in\supp(L_i)$ has $kA\subseteq H$. Thus the safety-only two-product bound applies to
\[
 gL_i,\quad ABM,\quad gBM
\]
and gives at most two missing witnesses. Use the highs $g,BM$, contexts $L_i,A$ and lows $a_i,k$. If a positive cross joins that same cut to context $M$, the related-cubic bound also applies to $gL_i,gB,ABM$. These uses require only safety of the indicated positive summands.

Two full-support facts will be useful. Distinct missing witnesses in the cubics $UVL$ and $UVM$ have form $xl,xm$. Their full contexts are $L=\{l\}$ and $M=\{m\}$: the common factors contain $x$, and an old context with another neighbor of $x$ would make a four-cycle or make its witnessed edge old. Also two safe quartics sharing three factors have at most two missing witnesses jointly. An old form cannot contain two distinct neighbors of a vertex when at least one of the corresponding edges is missing: its old center gives a four-cycle, and equality of that center with the common neighbor would make both edges old.

\emph{Positive cross coefficients.} A cross between two $g$ insertions makes raw $g$ safe and absorbs all its multiples. Only at most two old $AB$-context cubics remain, so there are at most three witnesses. A cross between two $f$ insertions makes $f$ safe. In the raw case it has at most one witness; in an affected case its old quadratic and cubic summands give at most two, and at most one remaining $g$ cubic adds one. For a mixed cross, raw $f$ gives $gf=f$, while $f=AB\vee g$ gives $gf=g$; the same bounds apply.

The only remaining mixed case is $f=AB\vee gB$, where $gf=gB$. Absorption leaves the direct $g$ cubics, the direct $AB$ cubics, and $gB$. There are four displayed terms, so four missing edges require one distinct witness from each. With one $g$ cut and two $f$ cuts, the actual-input and cross conditions give the related-cubic bound on $gL,gB,ABM_j$ for a crossed cut; the other $AB$ cubic adds at most one. With two $g$ cuts and one $f$ cut, the same argument applies if a crossed $g$ cut affects $A$.

Otherwise choose cut $1$ affecting $A$ and a crossed cut $2$ not affecting it. Distinct witnesses in $gL_1,gL_2,gB$ force the full contexts $L_i=\{l_i\},B=\{b\}$, with three missing star edges $xl_1,xl_2,xb$. The fourth $ABM$ witness is $by$, $y\ne x$. Actual incidence gives $l_1A\subseteq H$. The mixed coefficient at cut $2$ makes the $f$-low a singleton $l_2$, so $l_2M\subseteq H$. At least one of $A,M$ contains $y$. Whichever does supplies the old edge $l_i y$, giving the four-cycle $x-l_i-y-b-x$. All vertices are distinct: $y=x$ duplicates $xb$, $y=l_i$ would be an old loop, and safety and the distinct star edges exclude the other collisions. Consequently every $q_{ij}$ is zero in a four-witness case.

\emph{Two affected operands, zero cross.} Let $f=AB\vee g$. With one $g$ cut and two $f$ cuts, its low $a$ lies in both $A,B$. The first-pair strengthening of the related-cubic lemma says that $gL\vee gM_j$ has at most one missing witness. If $gL$ has one, all $g$ cubics together have at most one, and the two $AB$ cubics add at most two. If it has none, four witnesses would require one each from $gM_1,gM_2,ABM_1,ABM_2$. Full common-cubic rigidity gives $M_j=\{m_j\}$ and edges $xm_1,xm_2,ym_1,ym_2$, a four-cycle unless some witnesses repeat.

With two $g$ cuts and one $f$ cut, a cut affecting both $A,B$ gives the related-cubic total bound two, plus the other $g$ cubic. Otherwise choose cut $1$ affecting $A$ and cut $2$ affecting $B$. Four distinct cubic witnesses force $L_i=\{l_i\},M=\{m\}$ and missing edges $xl_1,xl_2,xm,my$, with $y\ne x$. Whichever of $A,B$ contains $y$ supplies an old edge $l_i y$ by actual incidence. The cycle $x-l_i-y-m-x$ contradicts four-cycle-freeness; coincidence $y=l_i$ would require an old loop.

\emph{One affected operand, zero cross.} Let $f=AB\vee gB$. With two $g$ cuts and one $f$ cut, choose a cut actually affecting $A$. The two-product bound gives at most two witnesses from its $g$ cubic and the two $f$ terms; the other $g$ cubic adds one. With one $g$ cut and two $f$ cuts, apply that bound to each triple
\[
 gL,\quad ABM_j,\quad gBM_j.
\]
If $gL$ has a witness, the two triples share it and have at most three jointly. Otherwise the two quartics share three factors and have at most two witnesses. Four overall would force distinct cubic witnesses from $ABM_1,ABM_2$ and two additional quartic witnesses. The cubics give a common center $z$ and full contexts $M_j=\{m_j\}$, with edges $zm_j$. Each quartic has at most one witness, since its singleton factor excludes two-edge quartic equality; its new edge is $y_jm_j$, $y_j\ne z$. The old form $B$ meets both incident missing edges $zm_j,y_jm_j$ and must therefore contain $m_j$. It would contain both $m_1,m_2$, two missing neighbors of $z$, contradicting the old-center four-cycle observation. No affected form remains.

\emph{Raw dependence and the cover.} The form $f=g$ gives at most three cubics. Let $f=gA$ with old $A$. If all three cuts insert $f$, any two-witness quartic confines the other terms through their three common full factors; if none has two, their union has at most three. With one $g$ cut and two $f$ cuts, the two shared quartics have at most two witnesses and the cubic adds one. Four therefore force two $g$ cuts and one $f$ cut, with private output
\[
 gL_1\vee gL_2\vee gAM,
\]
and exactly two witnesses from the quartic plus one additional distinct witness from each cubic. The full quartic equality gives an induced matching $xx',yy'$ and four full transversal supports $U,V,A,M$. The two additional witnesses must meet $U\cap V=\{x\}$ and hence are $xl_1,xl_2$; full common-cubic rigidity gives $L_i=\{l_i\}$. Both $l_i$ lie outside the matching's four endpoints. The set
\[
 S=\{l_1,l_2\}\cup\supp(M)
\]
has four vertices and covers all four missing edges. Its three constituent context supports are pairwise disjoint.

For each of the three original cuts mark the retained products reached by its insertion symbol along selected degree-one paths. Each marked set is nonempty, and the full opposite factor at every marked product is contained in that cut's direct context. A common marked product reached on opposite sides would give a positive cross coefficient, which was excluded. If reached on the same side, its nonempty opposite factor would lie in the intersection of two disjoint context supports. Thus the three marked sets are disjoint and contain at least three retained products. The set $S$ covers every old edge from any of them.

Every remaining retained product is a star: two nonsingleton linear factors would give a four-cycle, and intersecting factors would give a forbidden square. Choose one center for each unmarked product. These choices cover every remaining old edge, while $S$ covers all missing edges. If $m\ge3$ products were marked, this gives
\[
 \vc(G)\le (K-m)+4\le K+1.
\]
Repeated centers, repeated coefficients, overlapping stars, fanout and nested cuts do not affect this original-product incidence count.
\end{proof}



\begin{proposition}[A private linear sum with nonancestral products]\label{prop:privatelinearindependent}
Let $G$ be finite and four-cycle-free, and let a positive variable-input circuit compute $P\in\Sup(G)$ with nonempty quadratic support $H$. Suppose $H[T]$ is empty and, apart from isolates, $G[T]$ is $P_4\sqcup K_2$ or $K_3\sqcup K_2$. After exact-frame carrier erasure, it is impossible that the useful private computation consists of two nonancestral products, two visible cuts inserting them, and one private linear sum $S=U+V$ of old forms, each product reading only old forms or $S$. Sharing $S$, squaring it, positive coefficients, and arbitrary retained counts are permitted.
\end{proposition}
\begin{proof}
Each useful product must enter its own visible cut. Both highs have minimum degree two, so the full contexts are nonzero old linear forms, and the independent-symbol expansion is
\[
 P_2\vee gL\vee fM\vee qgf.
\]
For $q>0$, the old lows $a,b$ satisfy $a\subseteq M$, $b\subseteq L$ as supports, and $ab,aL,bM$ are supported on $H$. These are the original mixed-product context identities, before restriction to $T$. The target needs six positive CNF clauses, one for each of its six maximal independent sets, as in Proposition~\ref{prop:onevisibleclauses}.

If both products read $S$, write $g=SA,f=SB$, with $A,B$ old forms or $S$ itself. For $q=0$ the restricted output is
\[
 S(AL\vee BM)=S(A\vee B)(A\vee M)(L\vee B)(L\vee M),
\]
with at most five linear clauses. For $q>0$, the bracket $AL\vee BM\vee AB$ is $(A\vee B)(A\vee M)(B\vee L)$, giving at most four. Both are impossible.

Suppose only $g$ reads $S$ and $g=S^2$, with $f=BC$ old. If $q>0$, Boolean idempotence gives
\[
 SL\vee BCM\vee SBC
 =(S\vee B)(S\vee C)(S\vee M)(B\vee L)(C\vee L),
\]
again at most five clauses. If $q=0$, the safe old cubics $U^2L,V^2L,BCM$ supply all exact target witnesses and have at most three altogether.

Finally, let $g=SA$ and $f=BC$, with all three forms $A,B,C$ old. For $q=0$ the cubics $UAL$, $VAL$ and $BCM$ have at most three missing witnesses. For $q>0$, apply Lemma~\ref{lem:independentproducts} separately to the two products $WA,BC$, for $W=U,V$, with the same actual lows $a,b$ and contexts $L,M$. Every hypothesis is unchanged by splitting the high $SA$ into its positive old summands. Thus each triple
\[
 WAL,\qquad BCM,\qquad WABC
\]
has at most two missing edges. If $BCM$ witnesses a target, the triples share it, giving at most three targets in total. Otherwise it has no target witness. The two mixed quartics share $A,B,C$, so Lemma~\ref{lem:sharedquartics} bounds their union by two. Four targets therefore require two different direct witnesses from $UAL,VAL$.

Those direct witnesses share a vertex by Lemma~\ref{lem:commoncubic}, so write them $xy,xz$. They lie in the path or triangle component. Both old forms $A,L$ meet both pairs and therefore contain $x$: otherwise they contain two distinct missing neighbors of $x$, yielding a four-cycle through their old center or making the edges old. Choosing $x$ in both fixed factors shows, by safety, that every variable of $U$ and $V$ is a neighbor of $x$. Hence neither can contain an endpoint of the isolated target edge. Neither mixed quartic can witness that edge, the direct witnesses do not, and $BCM$ was assumed to have none. This contradiction exhausts the cases. Zero restrictions and repeated factors only reduce the clause or witness counts. No different extraction was used.
\end{proof}
This proposition concerns a pure private linear sum read by nonancestral products. A sum of affected retained signals or a dependent-product configuration requires a separate argument.

\begin{proposition}[A hidden addition with equal minimum degrees]\label{prop:hiddenhomogeneous}
Let $G$ be finite and four-cycle-free, and let a positive variable-input circuit compute $P\in\Sup(G)$ with nonempty quadratic support $H$. Suppose $H[T]$ is empty and, apart from isolates, $G[T]$ is $P_4\sqcup K_2$ or $K_3\sqcup K_2$. After exact-frame carrier erasure, the useful private computation cannot consist of at most two products and three additions, exactly two of the additions being visible unequal-degree cuts and the third hidden addition having operands of equal minimum degree. The full operands may contain higher terms. Retained gate counts, fanout, repeated inputs, and outside vertices are unrestricted.
\end{proposition}
\begin{proof}
Write $N=G[T]$. It has four missing edges, cover number three, and six maximal independent sets. Its edge function therefore requires at least six positive CNF clauses, by the independent-set argument in Proposition~\ref{prop:onevisibleclauses}. Lemma~\ref{lem:fivelost} forces two useful private products, which we name $g,f$ in topological order. Name the hidden addition $S$. If a frame-preserving absorption removes $S$, Proposition~\ref{prop:twosumspathtriangle} applies to the remaining two additions. We may therefore assume that $S$ is useful. The Boolean identities below analyze the output with the original selected frame and gate counts kept fixed; no gate count is obtained by extracting a new circuit.

We first state the context facts used in the cases. If one visible cut inserts each of $g,f$, their independent-symbol expansion is
\[
 P_2\vee gL\vee fM\vee qgf.                         \tag{*}
\]
The full direct contexts are nonzero old linear forms, except that a high of minimum degree greater than two may have a scalar context. A positive $q$ requires linear selected lows. If $a$ is the first low, then $aL$ is always supported on $H$; when $q>0$, additionally $a\subseteq M$. No individual insertion can be squared: its old low would then give a forbidden square in $P_2$.

An affected retained linear input has full value $C+\alpha g$, with $C$ old and $\alpha>0$. The first old low satisfies $a\subseteq C$. Moreover there is a variable $k\in\supp L$ such that $kC$ is old: for the cut output choose any $k\in\supp L$; otherwise follow the affected selected linear signal to its first retained product and choose $k$ from the opposite old operand. This proves the actual-context condition directly. Lemma~\ref{lem:independentproducts}, in its safety-only form, consequently bounds the missing witnesses in
\[
 hL,\qquad CRM,\qquad hRM                            \tag{**}
\]
by two whenever $h$ is an old quadratic product and these are safe terms of the expansion. Indeed use highs $h,RM$, contexts $L,C$, and lows $a,k$; the old products are $aL,kC,ak$. Splitting a private sum into old summands does not alter these data. The lemma is applied to safe terms, without asserting that their subsum contains every witness of $P$.

We also use the following consequence of the old-factor geometry. If safe cubics $XAB$ and $YAB$, all factors old, witness distinct missing edges $xy,xz$, then $A$ and $B$ contain $x$. Otherwise one contains $y,z$, and its old center gives a four-cycle or makes these edges old. Choosing $x$ twice shows $\supp X\subseteq N_G(x)$. Since $xy$ is missing, an old neighborhood containing $y$ and another neighbor of $x$ would again give a four-cycle; thus $X=\{y\}$ as a support, and similarly $Y=\{z\}$. In particular $X\vee Y$ is false on the separate edge component of $N$. We retain full supports in this argument, including vertices outside $T$.

\emph{The possible positions of $S$.} If its minimum degree is one, $S$ must feed a private product. A selected linear use would make its own occurrence retained; it cannot be the high branch of an unequal cut since all inputs have positive degree. Before $g$ its operands are old forms. Between $g$ and $f$ they are old forms possibly carrying positive multiples of $g$. A hidden sum of minimum degree at least two has, after full-carrier erasure, only $g,f$ available as nonlinear inputs. A repeated sum is redundant. Its only useful possibility is therefore $S=g+f$, with both products of minimum degree two. These cases exhaust the hidden addition.

\emph{Both visible cuts receive $f$, and $S$ has minimum degree one.} Neither cut precedes $f$, so all retained inputs of $S,g,f$ are unaffected. The useful multiplication network is a raw chain: Booleanly $f$ is a product of at most three linear factors, one possibly $S$. Repetitions only reduce this number. The restricted output is
\[
 fL\vee fM\vee qf^2=f(L\vee M\vee q),
\]
with at most four linear clauses, including scalar contexts. This contradicts the six-clause requirement. Hence the remaining linear cases have one insertion of each product and use~$(*)$.

\emph{$S$ precedes $g$ and is read by $g$.} Write $S=U\vee V$ with old $U,V$. Then $g=SA$ with $A$ old, or $g=S^2$. Nonancestral $f$ is excluded by Proposition~\ref{prop:privatelinearindependent}. If $f$ reads raw $g$, it is Booleanly $gR$ for an old form $R$ or $R=S$, or simply $g$. In the first case the output factors as $g[L\vee RM\vee qR]$. For $q=0$ the bracket is $(L\vee R)(L\vee M)$; for $q>0$ it is $L\vee R$. Together with the at most two clauses of $g$, there are at most four. Scalar contexts or $f=g$ only decrease this count.

If both inputs of $f$ are affected, Booleanly $f=CD\vee g$ with old $C,D$. For $q=0$ the perturbation is $gL\vee gM\vee CDM$. Split $g$ into $UA,VA$, or into $U^2,V^2$ when $g=S^2$. For each old quadratic summand $h$, Lemma~\ref{lem:relatedcubics} bounds the witnesses in $hL,hM$ by one, since $a\subseteq C\cap D$ and $aL$ is old. The two summands and the cubic $CDM$ supply at most three witnesses. For $q>0$, $g^2$ makes $g$ safe and absorbs its multiples; $g\vee CDM$ again has at most three witnesses.

It remains that $f=(C\vee g)B$ has one affected input. If $B=S$, put $g=SA$, allowing $A=1$ for the Boolean square case. Factoring $S$, the remaining bracket for $q=0$ is
\[
 AL\vee AM\vee CM=(A\vee C)(A\vee M)(L\vee M);
\]
for $q>0$ it is $A\vee CM$. At most four clauses result. Suppose instead that $B$ is old. If $g=S^2$, Booleanly $g=S$, and the two cases have respectively the CNFs
\[
\begin{split}
 SL\vee CBM\vee SBM
  &=(S\vee B)(S\vee M)(S\vee C)(B\vee L)(M\vee L),\\
 SL\vee CBM\vee SB
  &=(S\vee C)(S\vee B)(S\vee M)(L\vee B).
\end{split}
\]
They have at most five clauses. Now take $g=SA$ with $A$ old. For $q>0$ the expression is $gL\vee CBM\vee gB$. Apply Lemma~\ref{lem:relatedcubics} to each $h=UA,VA$, using $a\subseteq C\cap M$ and $aL$ old: $hL,hB$ have at most one witness together, and $CBM$ adds at most one. For $q=0$, apply~$(**)$ separately to $h=UA,VA$ in
\[
 hL,\qquad CBM,\qquad hBM.
\]
If $CBM$ witnesses a target, the two triples share it and have at most three targets in total. Otherwise the quartics $UABM,VABM$ have at most two targets by Lemma~\ref{lem:sharedquartics}. Four targets force two distinct direct witnesses from $UAL,VAL$. They are adjacent by Lemma~\ref{lem:commoncubic}. The support observation above makes $U,V$ their other endpoints, in the path or triangle component. Thus $S$, and hence $g$, is false on the separate edge component. Neither mixed term can witness that edge, and $CBM$ was assumed to supply no target. This is a contradiction.

\emph{$S$ is read only after $g$.} A pure sum computed independently before $g$ can be topologically reordered between $g$ and $f$ without changing any gate. Here $g=UV$ is old quadratic and
\[
 S=A\vee B\vee\lambda g,
\]
with old $A,B$ and $\lambda\ge0$. If $\lambda>0$, at least one operand, say the one with old part $A$, is an actual affected carrier. The other input of $f$ is $g$, $S$, or $C\vee\mu g$, with $C$ old.

If $f$ reads raw $g$, then for $\lambda>0$ it is Booleanly $g$; the output witnesses form a common-cubic star, or the star of safe $g$ if a scalar or square makes $g$ safe. For $\lambda=0$, $f=g(A\vee B)$. At zero cross coefficient the terms are $gL,gAM,gBM$. If $M$ is linear, the last two share three old factors and have at most two witnesses; if $M$ is scalar they reduce to old cubics. A scalar may also make $g$ itself safe, which only strengthens the bound. Positive cross coefficient leaves $gL,gA,gB$. In every case there are at most three targets.

If $f=S^2$, Booleanly $f=A\vee B\vee\lambda g$. All terms involving $g$ form a common-cubic missing star, unless a safe $g$ absorbs them. The remaining terms are $AM,BM$. They are safe two-factor products because the corresponding squared terms occur positively in the output. A safe product of two old forms witnessing a missing edge has both factors singleton, by the equality case of Proposition~\ref{prop:oldfactors}. Consequently their witnesses either are absent or all meet the singleton context $M$. The targets have a cover of size at most two, contradicting $\vc(N)=3$.

Consider $f=(A\vee B\vee\lambda g)(C\vee\mu g)$. If both coefficients vanish, Proposition~\ref{prop:privatelinearindependent} applies. If both are positive, $f=AC\vee BC\vee g$. At zero cross coefficient the $g$-terms form one common-cubic star and $ACM,BCM$ another. At positive cross coefficient safe $g$ absorbs its multiples, with the same two-star conclusion.

Suppose $\lambda>0$ and $\mu=0$. Then $f=AC\vee BC\vee gC$. For $q>0$, the pairs $gL,gC$ and $ACM,BCM$ each form a missing star. For $q=0$, the actual affected carrier with old part $A$ gives by~$(**)$ at most two witnesses in $gL,ACM,gCM$; the cubic $BCM$ adds at most one. Finally suppose $\lambda=0$ and $\mu>0$. For $q>0$, $gL,gA,gB$ form one star and $ACM,BCM$ another. For $q=0$ the expression is
\[
 gL\vee ACM\vee BCM\vee gAM\vee gBM.
\]
The actual affected carrier with old part $C$ permits~$(**)$ for the two triples $gL,ACM,gAM$ and $gL,BCM,gBM$. If $gL$ supplies a target they share it, giving at most three in total. Otherwise the shared quartics supply at most two, so four targets force distinct direct witnesses from $ACM,BCM$. The support observation makes $A,B$ the other endpoints of two adjacent targets in the path or triangle component. Every remaining term contains $A\vee B$, so none witnesses the separate edge. This completes the linear cases, including affected full operands.

\emph{$S$ has minimum degree at least two.} We have $S=g\vee f$, with both products of minimum degree two. If they are nonancestral they are old quadratic products. Each visible high is $g,f$, or $S$. At zero cross coefficient their expanded terms are grouped into the common-cubic star of $g$ and that of $f$, giving a two-vertex target cover. At positive cross coefficient, usefulness of $S$ means that one high is $S$ and its other high is $g,f$, or $S$. At least one square occurs, say $g^2$, so safe $g$ absorbs every multiple of $g$; the remaining $f$-terms form one star. Again the target cover number is at most two.

If $f$ depends on $g$, it cannot read raw $g$, since its minimum degree is two. It reads one or two affected retained linear signals. The first visible cut therefore inserts $g$ before $f$, and the second inserts $S$. The perturbation is
\[
 gL\vee(g\vee f)M\vee qg(g\vee f).
\]
For $q>0$, safe $g$ absorbs its multiples. A one-affected-input $f=(A\vee g)B$, or a two-affected-input $f=AB\vee g$, leaves $g\vee ABM$, with at most two witnesses. For $q=0$, the one-affected case gives $gL,gM,ABM,gBM$. The actual-context bound~$(**)$ gives at most two witnesses in $gL,ABM,gBM$, and $gM$ adds at most one. The two-affected case leaves only the three old cubics $gL,gM,ABM$. Thus the last cases are impossible as well.

All old-form arguments retain the original graph $H$, including its neighborhood centers outside $T$. Zero restrictions and repeated factors only remove targets or shorten the displayed CNFs. This proves precisely the hidden equal-minimum-degree assertion.
\end{proof}

\begin{proposition}[The surviving hidden equal-degree sum is pure linear]\label{prop:hiddenequalpure}
Let $G$ be finite and four-cycle-free, and let a positive variable-input circuit compute $P\in\Sup(G)$ with nonempty quadratic support $H$ and exactly four missing edges. Suppose extraction deletes two products and three additions, exactly two of which are visible cuts. If the hidden addition has equal original input minimum degrees, its original minimum degree is one, and after full-current-carrier erasure its full surviving value is a homogeneous sum $U+V$ of two old linear forms. There is no surviving private-product contribution to this sum. Retained counts and cover gain are unrestricted; homogeneity of the original full operands before erasure is not asserted.
\end{proposition}
\begin{proof}
Lemma~\ref{lem:fivelost} makes all five private gates essential under an exact-frame-preserving bypass. Keep their original degrees and selected occurrences during full-carrier erasure. The original-degree classification in Proposition~\ref{prop:hiddenhomogeneous} has the following possibilities. A hidden minimum-one sum before the first private product is already pure linear. Between the products $g,f$ it has full value $A+B+\lambda g$ with old $A,B$; it cannot be useful after both products, because it could neither be a high above a nonzero low nor enter a selected degree-one path without being retained. A hidden equal-degree sum of minimum at least two can only join $g,f$, both of original minimum degree two. A repeated input is Booleanly dispensable, and a retained quadratic input is an erased full carrier. We exclude that nonlinear case and the case $\lambda>0$.

The original two-insertion identity is
\[
 P_2\vee gL\vee fM\vee qgf.
\]
For a degree-one insertion of an old quadratic $g$, its full old direct context $L$ is nonzero: the selected nonzero low has a positive retained path to the quadratic output. This statement is in the full graph, with no outside restriction. If $a$ is that low, then $aL$ is old. An actually affected selected linear input with old part $A$ has $a\subseteq A$ and admits $k\in\supp(L)$ with $kA$ old, by following its selected path to the first retained product. When $q>0$, the actual cross term also gives $a\subseteq M$. These are the original-context facts used in Lemma~\ref{lem:independentproducts} and Lemma~\ref{lem:relatedcubics}. In particular, the safe terms
\[
 gL,\quad ACM,\quad gCM
\]
have at most two missing witnesses when $A$ is such an affected old part: use the two-product lemma with highs $g,CM$, contexts $L,A$ and lows $a,k$. Only safety of these displayed terms is required; their union need not supply all graph edges. We also use Lemma~\ref{lem:relatedcubics} for $UVL,UVD,ABD$ when $a\subseteq A\cap B$ and $aL$ is old.

\emph{A hidden nonlinear sum.} Suppose the hidden gate is $S=g+f$. If the two products are nonancestral after erasure, both are old quadratics, say $g=UV,f=WX$. The visible highs belong to $\{g,f,S\}$, with $S$ used at least once. With zero cross, unequal highs give at most three safe cubics. The only potentially four-term allocation inserts $S$ twice and gives
\[
 UVL,\quad UVM,\quad WXL,\quad WXM.
\]
Four distinct witnesses would force the full contexts $L=\{l\},M=\{m\}$ and the first pair of edges $xl,xm$, by the full-context common-cubic consequence. The second pair similarly gives $yl,ym$. If $x=y$ there are repeats; otherwise the four distinct vertices $x,l,y,m$ give a four-cycle. With positive cross, a square of an included high makes that high safe and absorbs its multiples. For highs $S,g$, this leaves safe $g$ and one old cubic from $f$; for two $S$ highs it makes both old quadratic summands safe. Either case has at most two missing witnesses.

If $f$ depends on $g$ while both have original minimum degree two, it cannot read raw $g$ as a factor. It reads selected linear inputs affected by the earlier visible $g$ cut, so
\[
 f=(A+\alpha g)(B+\beta g),
\]
with nonzero old parts and at least one positive coefficient. The other visible high must be $S$. Booleanly $S=g\vee f=g\vee AB$. With zero cross the private terms are therefore $gL,gM,ABM$, three old cubics. Positive cross makes $g$ safe and leaves $g\vee ABM$. Thus a nonlinear hidden sum cannot give four missing edges.

\emph{An affected minimum-one sum.} Suppose instead
\[
 S=A+B+\lambda g,\qquad \lambda>0,
\]
where $A$ is the old part of an actually affected operand. The first product $g=UV$ is old. The hidden $S$ must feed the second private product directly: an intervening selected linear use would retain it, and no further private addition is available. The other visible cut inserts $f$. Its second operand is raw $g$, another $S$, or $C+\mu g$ with old $C$ and $\mu\ge0$; a zero old part with positive coefficient belongs to the raw case.

For raw $g$, $f=gS=g$ Booleanly. Two linear contexts give at most two old cubics, while a scalar context or positive cross makes $g$ safe and absorbs its multiples. For a repeated $S$, $f=S^2=S$ Booleanly. At zero cross the terms are $gL,gM,AM,BM$. Apply the related-cubic bound to $gL,gM,A^2M$, using repeated old factors $A,A$ and $a\subseteq A$; $A^2M$ has the same Boolean support as $AM$. Those terms have at most two witnesses, and $BM$ adds at most one. At positive cross, $g^2$ makes $g$ safe and leaves $g,AM,BM$, again at most three.

If $\mu>0$ and $C\ne0$, the other input is actually affected too, so $a\subseteq C$. Booleanly
\[
 f=AC\vee BC\vee g.
\]
For zero cross apply the related-cubic bound to $gL,gM,ACM$, using $a\subseteq A\cap C$; the remaining cubic $BCM$ adds at most one. With positive cross, safe $g$ absorbs its multiples and the two old cubics add at most two to its single possible witness.

Finally let $\mu=0$, so $f=AC\vee BC\vee gC$. At zero cross the actual affected-input bound above applies to $gL,ACM,gCM$ and gives at most two witnesses; $BCM$ adds at most one. At positive cross absorption leaves
\[
 gL,\quad ACM,\quad BCM,\quad gC.
\]
Here $a\subseteq A\cap M$, the latter inclusion following specifically from $q>0$. The related-cubic bound on $gL,gC,AMC$ gives at most two, and $BCM$ again adds at most one. In these nonraw cases $f$ has original minimum degree two, so $M$ is a nonzero old linear context. The higher-minimum and scalar-context possibility was the raw case already considered. Zero coefficients remove the corresponding terms and are included in the split.

Every surviving nonlinear contribution to the hidden sum has therefore been excluded. Its remaining full value is $U+V$ with old homogeneous linear inputs, with coefficients absorbed into those forms. This conclusion concerns exactly two visible cuts; it does not exclude that pure preparation or a one-visible, two-hidden configuration.
\end{proof}

\begin{proposition}[A private linear sum with nonancestral products]\label{prop:independentlinearglobal}
Let $G$ be finite and four-cycle-free, and let a positive variable-input circuit compute $P\in\Sup(G)$ with nonempty quadratic support $H$. In the fixed full-current-carrier reduction, suppose the useful private gates are two originally nonancestral products, two visible cuts inserting them, and one hidden pure sum $S=U+V$ of old homogeneous linear forms. Each product reads old forms or $S$; sharing and squares $S^2$ are allowed. Then at most three edges are missing from $H$. No retained-count, cover-gain or outside-target hypothesis is required.
\end{proposition}
\begin{proof}
Write the highs $g,f$ and their nonzero old linear contexts $D,M$. The fixed frame gives
\[
 P_2\vee gD\vee fM\vee qgf.
\]
Both highs have original minimum degree two, so their useful direct contexts are linear, not scalar. For $q>0$ the nonzero original lows satisfy
\[
 a_0\subseteq M,\quad b_0\subseteq D,
 \qquad a_0b_0,a_0D,b_0M\subseteq H.
\]
These data are unchanged when $S$ is expanded into old summands. Every positive term is safe. All subsequent identities concern Boolean support with the original frame and counts fixed.

We first record the full-support facts used in the proof. A safe family $SXY$, with old $X,Y$, has at most two missing witnesses by expanding $S$. If it has two distinct witnesses, common-cubic rigidity gives edges $xa,xb$ and full supports $U=\{a\},V=\{b\}$, while $X,Y$ contain $x$. Indeed an old common factor omitting $x$ would contain its two missing neighbors. Safety puts the entire context $U$ in $N_G(x)$, and another vertex in $U$ would give a four-cycle with its old center or make $xa$ old; similarly for $V$. The same assertion holds if one of $X,Y$ is the Boolean unit: a safe quadratic rectangle is a star, and an old leaf factor cannot supply two missing neighbors of its singleton center. Thus the other old factor is the center. Such units only describe an original square Booleanly; they are not additional circuit inputs.

Suppose now $S=\{a,b\}$ and $xa,xb$ are missing. If an old form $X$ contains $x$, every $y\in\supp(X)\setminus\{x\}$ is nonadjacent to both $a,b$. Otherwise an old center $h$ for $X$ gives the four-cycle $h-x-a-y-h$, or its $b$ counterpart; $h=a$ would make $xa$ old, and simplicity excludes the remaining coincidences. In particular an old form cannot contain both $a,b$. If old $X,Y$ both contain $x$ and $SXY$ is safe, the occurrence of $a$ or $b$ in $X$ forces $Y=\{x\}$: any other vertex of $Y$, together with that leaf and the other $S$-leaf, would give an independent satisfying monomial.

Two more consequences will be useful. If $SXZ$ is safe, $X$ contains $x$, and it witnesses a new edge $ay$ outside $\{xa,xb\}$, then $a\in X$, $y\in Z$, and $Z$ is disjoint from $\{a,b\}$. The old form $X$ cannot contain $y$ by the preceding nonadjacency; putting either leaf in $Z$ would make an unsafe monomial with the leaf in $X$. Finally, if both $Z,T$ avoid $\{a,b\}$ and $SZT$ is safe, any exact witness has the other endpoint $y\in Z\cap T$. Safety forces both $ay,by$ to be edges. Four-cycle-freeness with $xa,xb$ forces $y=x$, so no new missing edge outside that pair is witnessed.

\emph{Both products use the sum.} Booleanly write $g=SA,f=SB$, where $A,B$ are old or a unit representing a square. For $q=0$ the two safe families are $SAD,SBM$. Four distinct witnesses would require two from each. The preceding full-support fact makes $S=\{a,b\}$ and the families two stars on those same leaves. Different centers give a four-cycle; the same center repeats the pair. Thus there are at most three.

For $q>0$ and a unit $A$, the mixed term is $SB=f$ and absorbs its multiples. The output is $SD\vee SB$, two quadratic families, and the same argument applies. The case of a unit $B$ is symmetric. With old $A,B$, the three safe families are
\[
 SAD,\quad SBM,\quad SAB.
\]
If each has at most one witness, count the terms. If one has two, it makes $S=\{a,b\}$ and witnesses $xa,xb$. Any other family with two witnesses has that same pair, since a different common center would give a four-cycle. Four witnesses therefore require a new edge from each of the other two families.

If the two-witness family is $SAB$, both $A,B$ contain $x$. A new $SAD$ edge forces a leaf, say $a$, into $A$, and then safety of $SAB$ forces $B=\{x\}$. But $SBM$ can now witness only $xa,xb$. If the two-witness family is $SAD$, a new $SAB$ edge similarly forces $a\in A$, $D=\{x\}$, and $B\cap\{a,b\}=\varnothing$. The actual cross condition gives $b_0=\{x\}$ and $xM\subseteq H$, so $M$ also avoids both missing neighbors $a,b$. The final full-support consequence then excludes a new $SBM$ witness. The case in which $SBM$ has two is symmetric. This proves the shared-sum case, including all squares.

\emph{Only one product uses the sum once.} Let $g=SA$ with old $A$, and $f=BC$ with old $B,C$. For $q=0$ the three old cubics $UAD,VAD,BCM$ give at most three witnesses. For $q>0$ apply Lemma~\ref{lem:independentproducts} separately to the old highs $UA,BC$ and $VA,BC$, with the same actual lows and contexts. Each safe triple
\[
 UAD,BCM,UABC;\qquad VAD,BCM,VABC
\]
has at most two missing witnesses. If $BCM$ has one, the shared witness bounds their union by three. Otherwise four would require two distinct direct witnesses and two additional mixed witnesses, since the two mixed quartics share three old factors and have at most two jointly.

The direct pair forces full $U=\{a\},V=\{b\}$, with $A,D$ containing $x$ and missing edges $xa,xb$. Each mixed quartic has a singleton factor and therefore cannot have two witnesses by quartic equality. Its new edge is respectively $ay$ or $bz$, with $y,z\ne x$. Since $A$ contains $x$ and meets $ay$, the old-center nonadjacency forces $a\in A$; the other mixed witness similarly forces $b\in A$. An old form cannot contain those two missing neighbors of $x$, a contradiction.

\emph{Only one product uses a square of the sum.} Let $g=S^2$, so Booleanly $g=S$, and let $f=BC$. Zero cross gives two old quadratics $UD,VD$ and one old cubic $BCM$, at most three witnesses. For positive cross, replace $g$ Booleanly by $U^2\vee V^2$ and apply the same two-product lemma to each old quadratic high with $BC$, retaining the original lows and contexts. If $BCM$ has a witness, the two triples share it and give at most three. Otherwise four witnesses require two from $SD$ and two from $SBC$. The first pair makes $S=\{a,b\}$, $D=\{x\}$ and edges $xa,xb$; the second pair gives $ya,yb$. Distinctness forces $x\ne y$, producing a four-cycle.

If neither product uses $S$, the old two-high context bound already gives at most two missing edges. These cases exhaust its occurrences in the two nonancestral products. Zero summands, repeated factors and coefficients only remove terms or repeat witnesses; no old-neighborhood hypothesis was imposed on $S$ itself. Dependent products and one-visible, two-hidden placements are outside the conclusion.
\end{proof}





\begin{lemma}[A grid of four old quartics]\label{lem:oldquarticgrid}
Let $H\subseteq G$ be finite simple graphs with $G$ four-cycle-free.
Let $E,F,A_1,A_2,B_1,B_2$ be nonzero positive homogeneous linear forms,
each supported on a singleton or in an original $H$-neighborhood.
If
\[
 Q=EF(A_1+A_2)(B_1+B_2)
\]
is safe in $G$, then at most three edges of $G\setminus H$ occur as exact
two-variable supports of its monomials. Repeated forms and arbitrary positive
coefficients are allowed. No circuit, retained-count or cover assumption is
imposed.
\end{lemma}
\begin{proof}
Identify linear forms with their supports for this argument, and put
$S=A_1\cup A_2$, $T=B_1\cup B_2$. Safety says that every choice
$e\in E,f\in F,s\in S,t\in T$ contains a $G$-edge, including repeated
choices. The sums represented by $S,T$ are not assumed old.

We first recall an old-neighborhood observation. If $pq$ is missing and an
old form $C$ contains $p$, then $q$ has no $G$-neighbor in $C\setminus\{p\}$.
Otherwise an old center $h$ of $p,r\in C$, with $qr\in E(G)$ and $r\ne p$,
gives the four-cycle $h-p-q-r-h$. Its vertices are distinct: $h=q$ would
make $pq$ old, while simplicity excludes the other coincidences. The case
$r=q$ would require a loop. It follows that an old form meeting two incident
missing edges $pq,pr$ contains their common endpoint $p$, and cannot contain
both outer endpoints $q,r$. Also, an old form cannot contain three distinct
vertices $w,u,v$ with both $wu,wv\in E(G)$: its center gives the four-cycle
$h-u-w-v-h$.

Each corner $P_{ij}=EFA_iB_j$ is safe. Suppose one corner, say $P_{11}$,
has two missing witnesses. Proposition~\ref{prop:oldfactors} makes them
an induced matching, and the four full factor supports are its four distinct
two-point transversals. Each neighboring corner shares three of these
supports, whose intersection is empty. A new missing edge would have at most
one endpoint in the matching's four vertices, and that endpoint would have
to meet all three shared supports. Hence those corners supply no new edge.
The opposite corner shares $E,F$. If these transversals are disjoint there
is no new witness; otherwise they intersect at one vertex $x$, and every
new witness has the form $xz$ outside the original quartet. The missing
witnesses of that one corner form an induced matching, so at most one such
edge occurs. The entire grid then has at most three witnesses.

It remains to exclude four distinct witnesses when every corner has at most
one. There must be one edge $e_{ij}$ in each corner, and each belongs to that
corner alone. Therefore $e_{ij}$ misses $A_{3-i}$ and $B_{3-j}$ entirely:
otherwise it already meets the other three factors of an adjacent corner
and would witness that corner too. We use this exact-support exclusivity.

\emph{Incident witnesses in a row or column.} Suppose
$e_{11}=xa$, $e_{12}=xb$, $a\ne b$. The old forms $E,F,A_1$ all contain
$x$. Safety then places every vertex of $B_j$ in $N_G(x)$. The form $B_1$
contains $a$ and cannot contain $x$, since that would give a singleton
satisfying monomial; the old-neighborhood observation forces its full support
to be $\{a\}$. Similarly $B_2=\{b\}$. The other two corner witnesses
therefore contain $a$ and $b$, respectively, and differ from $xa,xb$.
The old form $E$, meeting two incident missing edges at each of $a,b$,
must contain both. This contradicts the missing-neighbor observation at
$x$. The column case is symmetric. Thus incident witnesses can occur only
at diagonal corners, and no vertex belongs to three of them.

Here is a second consequence of safety. Suppose a missing edge $xa$ has
$x\in E\cap F$ and $a\in S\cap T$. Every vertex of
$(E\cup F)\setminus\{x\}$ is nonadjacent to $a$ by the observation.
A common vertex $z\ne x$ would therefore make the independent pair
$\{z,a\}$ satisfy $Q$. Thus $E\cap F=\{x\}$.
For $u\in E\setminus\{x\}$ and $v\in F\setminus\{x\}$, safety of
the choices $u,v,a,a$ requires $uv\in E(G)$. These two sides are disjoint;
if both have at least two vertices they give a four-cycle.

\emph{Diagonally incident witnesses.} Suppose $e_{11}=xa$, $e_{22}=xd$.
Exclusivity excludes $x$ from every $A_i,B_j$, so $a\in S\cap T$.
The common old forms $E,F$ contain $x$, and the preceding consequence
applies. The other witnesses are disjoint from the vertices of these two
edges, since any incidence would share a row or column. If those remaining
edges meet at $y$, then $y\in E\cap F$, contradicting its singleton
intersection. If they are disjoint, both $E\setminus\{x\}$ and
$F\setminus\{x\}$ meet both edges and have at least two vertices, giving
the forbidden four-cycle. All four witnesses are therefore pairwise
disjoint. No absence of other $G$-edges among their endpoints is assumed.

\emph{No witness endpoint belongs to $S\cap T$.} Suppose $w$ did, and
write its selected edge as $wx$. The common factors $E,F$ meet that edge.
They cannot both contain $w$, which would give an unsafe singleton. They
cannot both contain $x$: the preceding consequence, together with the three
other disjoint edges, would give at least three vertices on each complete
bipartite side. Thus, after exchanging the common factors,
$w\in E\setminus F$ and $x\in F\setminus E$.
In fact $E\cap F$ is empty: if $z$ were common, the old form $F$ containing
$x,z$ would make $wz$ a nonedge, whereas $z,z,w,w$ satisfies the product.

Each selected edge consequently has one endpoint $u_i$ in $E$ and its mate
$v_i$ in $F$. Index $wx$ by zero, so $u_0=w,v_0=x$. For $i\ne j$, the
old form $E$ containing $u_i,u_j$ and the missing edge $u_jv_j$ show that
$u_iv_j$ is a nonedge. Fix $i\ne0$ and choose $j$ distinct from both zero
and $i$. In the satisfying choice $u_i,v_j,w,w$, the cross pairs
$u_iv_j,wv_j$ are absent, so safety forces $wu_i\in E(G)$. This holds
for at least two different $i$. The old form $E$ would contain $w$ and two
distinct neighbors of $w$, contrary to the first observation.

Thus each witness edge has one endpoint $\alpha_i\in S\setminus T$ and
the other $\beta_i\in T\setminus S$. It must meet both unions, so these
two classes partition its endpoints.

\emph{The common factors have no common witness endpoint.} If
$\alpha_i\in E\cap F$, then for every $j\ne i$ the satisfying choice
$\alpha_i,\alpha_i,\alpha_i,\beta_j$ requires
$\alpha_i\beta_j\in E(G)$. Neither $E$ nor $F$ can then contain
$\alpha_j$, by the old-neighborhood observation for the missing edge
$\alpha_j\beta_j$. Both must instead contain $\beta_j$.
Choose distinct $i,j,k$. As $\beta_j$ is common, safety requires
$\beta_j\alpha_k\in E(G)$, but the old form $E$ contains
$\beta_j,\beta_k$ and the edge $\beta_k\alpha_k$ is missing. This is
the forbidden cross pair. The argument for a common $\beta_i$ is symmetric.

The selected edges again have opposite endpoints $u_i\in E$, $v_i\in F$,
with no common witness endpoint, and all off-diagonal pairs $u_iv_j$ are
nonedges by the old-neighborhood observation. Give every endpoint its
$S$-only or $T$-only class. The endpoints of each edge have opposite classes.
For $i\ne j$, if $u_i,v_j$ had opposite classes, those two nonadjacent
vertices would also supply the $S,T$ positions and satisfy $Q$. Hence
$u_i,v_j$ have the same class. Equivalently $u_i,u_j$ have opposite classes
for every distinct pair of indices. Three indices already contradict this
two-class requirement.

The four-single-corner possibility is impossible. Together with the
two-witness-corner case this proves the bound, retaining all original
vertices, all old cross edges, coefficients and repetitions.
\end{proof}


\begin{proposition}[A private linear sum with dependent products]\label{prop:dependentlinearglobal}
Let $G$ be finite and four-cycle-free, and let a positive variable-input
circuit compute $P\in\Sup(G)$ with nonempty quadratic support $H$.
In the fixed full-current-carrier reduction, suppose the useful private
computation has two originally ancestrally related products, two visible
insertion cuts, and one hidden pure sum $S=U+V$ of old homogeneous linear
forms. Then at most three edges are missing from $H$. Retained counts,
cover gain, fanout, squares and sharing of $S$ are unrestricted.
\end{proposition}
\begin{proof}
Suppose four or more edges were missing. Lemma~\ref{lem:fivelost}
makes all five private gates essential under the exact-frame reduction.
Keep their original degrees and selected occurrences. The pure degree-one
gate $S$ cannot enter a selected linear path, which would retain it.
It is therefore a literal private-product operand, not a newly declared
old boundary or a freely available summand in another linear input.
Name the products $g,f$ by their original dependency. If carrier erasure
removes that dependency, apply the explicit old/shared-$S$ family proof of
Proposition~\ref{prop:independentlinearglobal} to the same frame and its
unchanged low/context identities. This gives at most three witnesses,
without relabeling the original products as nonancestral. We may assume a
surviving dependency.

Every visible high is $g$ or $f$: the pure degree-one sum cannot be a high,
and off-frame full quadratic carriers have been erased. An affected
degree-one operand of $f$, other than raw $S$, has the form $B+\alpha g$
with $B$ old and $\alpha\ge0$; a zero old part is raw $g$.
Thus if $g=SA$ with $A$ old, the dependent $f$ is Booleanly a raw
$gW$ ($W$ old or $S$), $g$, or one of
\[
 BC\vee gC,\qquad BC\vee g,\qquad BS\vee g,
\]
with displayed $B,C$ old and the appropriate old parts actually affected.
Allow $A=1$ only as notation for the Boolean identity $S^2=S$, keeping
the original product and its minimum degree two. If $g$ does not use $S$,
it is an old quadratic, and useful $S$ leaves only
$f=gS$ or $f=BS\vee gS$, the latter from $(B+g)S$ with $B$ affected.
These alternatives come from the literal original operands; they include
repetitions and raw products.

If both cuts insert $f$, there is no $g$-cut to affect an old operand and
$f$ is raw. Otherwise one cut inserts each high. Their old direct contexts
$D,M$ give
\[
 P_2\vee gD\vee fM\vee qgf,\qquad q\ge0.                 \tag{*}
\]
Here $D$ is nonzero and linear. The context $M$ is linear unless the cut
keeps degree two, in which case its positive scalar makes $f$ safe.
That scalar case can occur only for a raw higher-minimum-degree $f$.
Let $a$ be the old low of the $g$-cut. Always $aD$ is old. For any old
part $B$ that this cut actually affects, $a\subseteq B$ and some
$k\in D$ has $kB$ old. For $q>0$ also $a\subseteq M$.
These are the original incidence identities in
Corollary~\ref{cor:affectedcase}; no extraction is changed below.
All displayed positive output terms are safe.

We collect the support consequences used in the cases. A safe family
$SAB$, with $A,B$ old, has at most two witnesses. If it has two, the
full supports are $U=\{u\},V=\{v\}$, its edges are $xu,xv$, and
$A,B$ contain $x$. This follows by expanding $S$ and using the full
common-cubic observation in the proof of
Proposition~\ref{prop:independentlinearglobal}. A Boolean unit among
$A,B$ represents a square; repeat the other old factor to use the same
statement. Two such safe families with the same $S$ cannot have four
distinct witnesses: two centers on the same two leaves would give a
four-cycle or duplicate an edge.

If $xu,xv$ are missing and old $A$ contains $x$, it cannot contain both
$u,v$. Moreover a new witness $uy$, $y\ne x$, in a product with factor
$A$ forces $u\in A$; otherwise $y\in A$, and the old center of $x,y$
would give a four-cycle through the missing edge $xu$ and $uy$.
Thus additional witnesses $uy,vz$ in such products force a contradiction.
Two safe old quartics sharing three factors have at most two witnesses
jointly by Lemma~\ref{lem:sharedquartics}.
Finally we use the explicit three-family part of
Proposition~\ref{prop:independentlinearglobal}: the safe families
$SAD,SBM,SAB$ have at most three witnesses when the old lows satisfy
$a\subseteq M$, $b\subseteq D$ and $aD,bM,ab\subseteq H$.
That proof uses just these supports and safety, so it is an algebraic
family statement; $A$ or $B$ may again represent a Boolean square unit.
When $g=S^2$, using the formal old quadratic highs $U^2,V^2$ makes all
applications of Lemmas~\ref{lem:independentproducts} and
\ref{lem:relatedcubics} below literal. Their Boolean supports are $U,V$;
no original gate is removed or recounted.

\emph{One cut per high, with $g=SA$.} First let $f=gW$ with $W$ old.
A positive mixed coefficient or a scalar $M$ makes $f$ safe and absorbs
$fM$. The remaining families $SAD,SAW$ cannot supply four witnesses by
the two-center observation. With $q=0$ and $M$ linear, the family $SAD$
has at most two, and the quartics $UAWM,VAWM$ share three old factors and
have at most two jointly. Four distinct witnesses force two from each.
For old $A$, the direct pair gives $S=\{u,v\}$ and $x\in A,D$.
Each mixed quartic now has a singleton factor, so cannot have two
witnesses by full-transversal equality. Its new witness is respectively
$uy$ or $vz$, distinct from $xu,xv$. Both $u,v$ would have to lie in
$A$, the contradiction just proved. With the square unit $A=1$, the two
families are $SD,SWM$, already excluded by the two-center argument.
If $W=S$ or $f$ is Booleanly $g$, zero cross gives $SAD,SAM$ and the
same argument; positive cross or a scalar context makes $g$ safe and
absorbs its multiples.

For $f=BC\vee g$ and $q=0$, the terms are $gD,gM,BCM$.
The actual low $a$ lies in $B\cap C$. Apply the first-pair conclusion
of Lemma~\ref{lem:relatedcubics} to $UAD,UAM,BCM$ and to the $V$
version. The two first pairs contribute at most two witnesses altogether,
and the old cubic $BCM$ adds one. Positive cross makes $g$ safe and leaves
$g\vee BCM$, with counts at most two and one.

For $f=BC\vee gC$ and zero cross, apply
Lemma~\ref{lem:independentproducts}, using the actual affected part $B$,
to
\[
 UAD,\ BCM,\ UACM;\qquad VAD,\ BCM,\ VACM.
\]
The highs are $UA,CM$, contexts $D,B$, and lows $a,k$; the old products
$aD,kB,ak$ give its hypotheses. Each triple has at most two witnesses.
If $BCM$ has one, it is common to both triples and the union has at most
three. Otherwise four would require two direct witnesses from $SAD$ and
two mixed witnesses from $SACM$. For old $A$, the singleton-leaf argument
in the preceding raw case forces $u,v\in A$, impossible; for the unit
case the two-center argument applies to $SD,SCM$.
With positive cross, absorption leaves $gD,BCM,gC$. Now $a\subseteq
B\cap M$. Apply the first-pair bound of Lemma~\ref{lem:relatedcubics}
to $UAD,UAC,BMC$ and the $V$ version. Those pairs give at most two, and
$BCM$ one more.

For $f=BS\vee g$ and zero cross the three families are
$SAD,SBM,SAM$. Choose $k\in D$ with $kB$ old. Use the preceding
three-family statement with its second old factor equal to the present
$M$, its second context equal to $B$, and its lows $a,k$.
The conditions are $a\subseteq B$, $k\subseteq D$ and $aD,kB,ak$
old, exactly the actual-context conditions. It gives at most three
witnesses. Positive cross makes $g$ safe and leaves $SA,SBM$, excluded
from four witnesses by the same two-center argument.

\emph{One cut per high, with $g$ old and only $f$ using $S$.}
For raw $f=gS$ and zero cross, $gD$ is one old cubic, while $gUM,gVM$
share three old factors and have at most two witnesses jointly.
A positive cross or scalar $M$ makes $f$ safe and leaves its two old
cubics, together with $gD$; again there are at most three.

Let $f=BS\vee gS$ with $B$ actually affected. At zero cross apply
Lemma~\ref{lem:independentproducts} to
\[
 gD,\ BUM,\ gUM;\qquad gD,\ BVM,\ gVM,
\]
with the same actual low/context identities. A witness in $gD$ is shared
and bounds the union by three. Otherwise four require two from $BSM$
and two additional ones from $gSM$. The first pair gives
$S=\{u,v\}$, $x\in B,M$ and missing $xu,xv$.
Each mixed quartic, having a singleton factor, supplies at most one new
edge; its witnesses $uy,vz$ force $u,v\in M$, impossible.
At positive cross, absorption leaves $gD,BSM,gS$, with
$a\subseteq B\cap M$. For $T=U,V$, the first-pair conclusion of
Lemma~\ref{lem:relatedcubics} applied to $gD,gT,BMT$ bounds
$gD\vee gT$ by one witness. If $gD$ has a witness, all three $g$-terms
jointly have one, and $BSM$ adds at most two. Otherwise four would
require two witnesses from each of $gU,gV$ and $BUM,BVM$.
The common-cubic rule makes $U,V$ singleton leaves at two distinct
centers, giving a four-cycle. This excludes the last distinct-high case.

\emph{Both cuts insert $f$.} With no $g$-cut, $f$ is raw. A scalar
direct context or positive cross makes $f$ safe and absorbs its multiples.
It contains $S$ and at most two old factors; expanding $S$ gives at most
two old cubic witnesses. A nonzero linear $S$ alone could not be safe.
Thus consider zero cross and two nonzero old linear contexts $L,M$.
The private output is $f(L+M)$. For old $g=XY$ and $f=gS$ this is the
grid $XY(U+V)(L+M)$ of Lemma~\ref{lem:oldquarticgrid}.
For $g=SA$ and $f=gW$ with $W$ old, use the same grid with
$X=A,Y=W$. The possibilities $f=gS$ or $f=g^2$ reduce to $X=Y=A$;
a square unit is handled by repeating the remaining old factor.
All have at most three missing witnesses.
The sole linear degeneration is $f=S$, giving the safe quadratic
rectangle $S(L+M)$. Its two supports are disjoint, since a common vertex
would give a square monomial. Four-cycle-freeness makes one side singleton.
The opposite side is a union of two old forms, each of which supplies at
most one missing neighbor of that singleton by the old-center observation.
There are at most two witnesses.

Two cuts inserting only $g$ cannot make dependent $f$ useful. The cases
therefore exhaust the two high allocations, all literal operand forms,
raw squares, shared $S$, actual affected inputs, zero or positive cross
coefficients and scalar degree-two contexts. Every missing edge must be
supplied by a displayed positive term, since $P_2$ supplies only $H$.
This proves the global bound. One-visible, two-hidden profiles are outside
this proposition.
\end{proof}


\begin{lemma}[Seven old factor supports cannot supply four witnesses]
\label{lem:sevenoldfactors}
Let $H\subseteq G$ be finite simple graphs, with $G$ four-cycle-free.
Let $P,Q$ be safe products of nonzero old homogeneous linear forms.
Count equal factor supports only once, including repetitions within either
product, and let $b$ be the number of distinct supports in their union.
If the union of their exact missing-edge witnesses contains four distinct
edges whose graph has vertex-cover number at least three, then $b\ge8$.
There is no actual-context, circuit or retained-count assumption.
\end{lemma}
\begin{proof}
Boolean repetition changes neither safety nor exact witnesses, so first
collapse repeated supports. Suppose $b\le7$. By
Proposition~\ref{prop:oldfactors}, each product has at most two missing
witnesses: three would require eight distinct factor supports.
The four stipulated edges therefore split into two disjoint pairs,
one pair from each product. Each pair is an induced matching in $G$.

Color the pairs by their products. Their union has maximum degree two
and cannot contain a three-edge path: its colors alternate, and its first
and third edges would belong to the same induced matching while the middle
edge joins their endpoints. A cycle would contain such a path, and a
triangle cannot be colored by two matchings. Thus every nontrivial
component is $K_2$ or $P_3$. Four edges with cover number at least three
give precisely $4K_2$ or $P_3\sqcup2K_2$.
In either case each product's matching has an edge entirely outside the
four-vertex endpoint set of the other's matching. Equivalently those two
endpoint sets meet in at most one vertex.

Consider the pair witnessed by $P$, with endpoint set $W$.
Its four choices of one endpoint from each edge are independent sets in
$G$. Every factor meets both edges. A factor can avoid at most one such
transversal, and every transversal must be avoided by some factor, since
otherwise it would satisfy the safe product on an independent assignment.

Suppose a transversal $T$ has exactly one rejecting distinct factor
support $F$. Then $F\cap W=W\setminus T$. In fact this is its full support.
If $z\in F\setminus W$, every other factor meets $T$. For any $t\in T$,
let $s\in F$ be its matching mate and let $h$ be an old center of $s,z$.
An edge $tz$ would give the four-cycle $h-s-t-z-h$.
The vertices are distinct: $h=t$ would make the missing edge $st$ old,
and simplicity excludes $h=s,z$. Hence $z$ is nonadjacent to every
vertex of $T$. The independent assignment $T\cup\{z\}$ would satisfy
$P$, a contradiction. Thus $F=W\setminus T$ exactly.

Such a unique rejector cannot also be a factor of $Q$: it is a two-point
subset of $W$, whereas every $Q$-factor must meet the $Q$-witness edge
entirely outside $W$. Unique rejecting supports are therefore private to
their product. Different transversals have different rejectors.

Let $p,q$ be the numbers of distinct supports in $P,Q$, respectively,
and let $u_P$ count the $P$-transversals with exactly one rejector.
Every remaining transversal has at least two, and one factor rejects at
most one transversal. Consequently
\[
 p\ge u_P+2(4-u_P)=8-u_P.
\]
There are at least $u_P\ge8-p$ factor supports private to $P$, so
$b-q\ge8-p$. Interchanging the products gives $b-p\ge8-q$.
Adding yields $2b\ge16$, contrary to $b\le7$.
Supports rejecting no transversal only increase the counts.
All coefficients, outside vertices and repeated neighborhood centers have
been retained in this argument.
\end{proof}

\begin{proposition}[Two boundary profiles with four private products]
\label{prop:sevenfactorprofiles}
Let $G$ be finite and four-cycle-free, and let a positive variable-input
circuit compute $P\in\Sup(G)$ with nonempty quadratic support $H$.
Suppose exactly four edges are missing, their graph $J$ satisfies
$\vc(J)\ge3$, and the private budget is four products and two additions.
Neither of the following fixed-frame configurations is possible:
\begin{enumerate}
\item exactly one addition is a visible insertion cut;
\item both additions are visible cuts, and all four private products
are ancestral to the first visible high.
\end{enumerate}
Retained gate counts and cover gain are unrestricted. In particular the
conclusion applies to these configurations whenever
$\vc(G)-\vc(H)\ge3$.
\end{proposition}
\begin{proof}
Use the full-current-carrier reduction while preserving the original
minimum-degree labels, selected occurrences and gate charges. A retained
quadratic carrier contributes its full value on a protected addition-only
output path, so its further private uses are erased as already proved.
Before the first visible cut, all remaining retained boundary signals of
a useful private network are old homogeneous linear forms.
Every useful private gate feeds a visible high. A removed use or a vanished
branch can only decrease the estimates below; it does not create a new
selected gate or boundary form.

We use a connected binary-DAG count. If there are $n$ internal gates and
$d$ distinct boundary sources, the underlying connected graph has at least
$n+d-1$ operand incidences, and the gates have at most $2n$ in total.
Hence $d\le n+1$. Fanout and repeated operands cannot increase this bound.
Distinct factor supports are no more numerous than distinct boundary
sources. A scalar context adds no factor support, and an old linear context
adds at most one.

With one visible cut, all useful private gates lie in its high $F$.
There are at most four products and one hidden addition, so the high has
at most six distinct old boundary sources. Below the hidden addition
there are only products of those forms; write its two inputs as $A,B$.
Above it there are again only products. Repeated uses and shared ancestors
therefore give, at Boolean support level,
\[
 F=C(A\vee B),
\]
where $A,B,C$ are products of the original boundary forms.
A positive power of $A\vee B$ is Boolean-equivalent to one copy, so DAG
sharing does not create further terms. If the addition becomes redundant,
there is only one product and the same bound applies.

The unique insertion has a scalar or old linear direct context $D$.
There is no squared insertion symbol: its nonzero old low input would
otherwise produce a square in the squarefree quadratic remainder.
All fanout is included in this one full context. The private output thus
has the two safe products $CAD$ and $CBD$, using at most seven distinct
old factor supports altogether. Every missing edge must be witnessed by
one of them, since the quadratic remainder supplies only $H$.
Lemma~\ref{lem:sevenoldfactors} gives a contradiction. A zero term,
scalar context or repeated factor only lowers the count.

Now suppose both additions are visible, and all products precede and feed
the first high $g$. Its multiplication-only network has at most five
distinct old boundary sources. The second high $f$ is $g$ or one of its
ancestral products. Thus $f$ divides $g$ at Boolean support level.
The exact two-insertion frame has private contribution
\[
 gL\vee fM\vee qgf,
\]
where $q\ge0$ and the direct contexts are old homogeneous linear forms
or scalars. If $q=0$, the two displayed products have at most the five
boundary supports together with those of $L,M$, hence at most seven.
If $q>0$, then $gf=g$ Booleanly. The safe summand $g$ absorbs $gL$,
leaving $g\vee fM$, with at most six supports. Zero contexts or coincident
highs give smaller counts. Positivity makes the retained displayed products
safe; they supply all four missing edges. The same lemma again contradicts
$\vc(J)\ge3$.

These are identities used to analyze the fixed original frame, not a
recomputation of its extraction. Finally
$\vc(G)\le\vc(H)+\vc(J)$, because the union of an $H$-cover and a
$J$-cover covers $G$. Thus cover gain at least three implies the required
hypothesis on $J$. The other two-visible ancestry allocations and the
three-product, three-addition configurations are not excluded here.
\end{proof}


\begin{proposition}[One visible cut and two hidden additions]
\label{prop:onevisibletwocover}
Let $G$ be finite and four-cycle-free, and let a positive variable-input
circuit compute $P\in\Sup(G)$ with nonempty quadratic support $H$ and
exactly four missing edges. Suppose canonical extraction deletes two
products and three additions, exactly one of which is a visible cut.
The other additions may have arbitrary original input degrees, full
inhomogeneous values, dependencies and fanout. If $K$ products are retained,
then $\vc(G)\le K+1$. Retained additions and cover gain are unrestricted.
\end{proposition}
\begin{proof}
Erase off-frame full quadratic carriers with the original selected frame
protected. Every useful private node must feed the unique visible high
$F$: otherwise its first entry into the selected subgraph would be another
visible cut. The high therefore has at most two products and two additions.
Its boundary signals are old homogeneous linear forms. A retained quadratic
boundary is an erased full carrier. A retained linear boundary cannot
already contain a private contribution, since its first selected entry
would be the same cut, creating a directed cycle back into that cut's high.
This argument concerns full carriers and original degrees.

The exact one-symbol output is Booleanly $P_2\vee FL$. If the cut
retains degree two, $L$ is a positive scalar, written $1$ below.
If it retains degree one, $L$ is a nonzero full old linear context:
the old low $a$ satisfies $aL\subseteq H$. There is no squared insertion
term, since the nonzero old low on both sides of a retained product would
put a square in $P_2$. Fanout is included in $L$. All four missing edges
are witnessed by the safe term $FL$. We will show that $L$ must be linear
with a full two-vertex support, then charge that support to the actual frame.

\emph{A four-gate sharing reduction.} Any connected binary AND/OR DAG
with at most four gates, at most two of each type, is Boolean-equivalent
to a formula with no larger count of either type, over its original atoms.
Bypass equal-input gates. If an internal signal $s$ is shared, take two
paths to their first reconvergence. The at most four available nodes form
a triangle, a four-node shortcut, or a four-node diamond. A triangle has
value $s\circ(s\mathbin{*}a)$: equal operations give the already computed
$s\circ a$, while different operations give $s$ by absorption.
Reduce any such earlier triangle first, irrespective of other fanout.

For a shortcut $s\circ((s\mathbin{*}a)\mathbin{\#}b)$, the node
$s$ itself is the remaining fourth gate, so $a,b$ are boundary atoms.
If the outer operation is AND, the case with both inner operations AND
would require three AND gates. The other cases satisfy
\[
\begin{split}
 s\wedge((s\wedge a)\vee b)&=s\wedge(a\vee b),\\
 s\wedge((s\vee a)\wedge b)&=s\wedge b,\\
 s\wedge((s\vee a)\vee b)&=s.
\end{split}
\]
The OR cases are dual. A diamond is
$(s\mathbin{*}a)\circ(s\mathbin{\#}b)$. With outer AND, two inner
ANDs again exceed the budget; two inner ORs give
$(s\vee a)\wedge(s\vee b)=s\vee(a\wedge b)$; different inner
operations leave the inner AND by absorption. The dual cases cover outer
OR. Each replacement removes sharing without increasing either gate count.
The shared node and its first reconvergence already use two nodes, so at
most two path nodes remain: the list is exhaustive. Repeating these
reductions proves the formula assertion; shared input atoms are permitted.

Apply this assertion to $F$. Distributive expansion of the formula has
at most four terms, each a product of at most three old atoms. If there
are four, the two OR gates must be incomparable, with an AND as their
least common ancestor. The other possible AND is either outside the two
arms or inside one arm. Up to names, the only alternatives are
\[
 F=E(A\vee B)(C\vee D),\qquad
 F=(A\vee B)(CD\vee E).                         \tag{*}
\]
All letters are old nonzero forms; only the common $E$ in the first form
may be $1$ when that extra AND is absent. Nested ORs give at most three
terms. A common factor along one branch is pulled outside by associativity.
Coincident atoms or absorbed terms only reduce the number. These are
semantic identities, with the original extraction and its charges fixed.

We recall two full-support observations. If $uv$ is missing and an old
form contains $u$, then $v$ has no neighbor in its support other than $u$.
A contrary neighbor $w$ and the old center $h$ give $h-u-v-w-h$;
$h=v$ would make $uv$ old and the other coincidences are excluded by
simplicity. Thus an old form meeting two incident missing edges $xu,xv$
contains $x$ and cannot contain both outer endpoints.
If safe cubics $CDA,CDB$ have distinct missing witnesses, the common-cubic
lemma and this observation give edges $xa,xb$, with $x\in C\cap D$ and
full supports $A=\{a\}$, $B=\{b\}$. Indeed both common forms contain
$x$; choosing $x$ twice makes every variable of $A$ a neighbor of $x$.
The old-center argument forbids any second variable alongside its missing
neighbor $a$. The same holds for $B$. These observations need safety only.

If the expansion of $F$ has at most three terms, four missing witnesses
force one term of $FL$ to witness at least two. A safe product of at most
three old forms cannot do so. Hence $L$ is linear, the term is an old
quartic, and the full equality case of Proposition~\ref{prop:oldfactors}
forces $|\supp(L)|=2$.

Consider the first form in~$(*)$. If $E,L$ are old, $FL$ is the grid of
Lemma~\ref{lem:oldquarticgrid}, with at most three missing witnesses.
If exactly one is $1$, repeat the other old common factor Booleanly and
apply the same lemma. If both are $1$, put
$S=\supp(A)\cup\supp(B)$ and $T=\supp(C)\cup\supp(D)$.
Safety makes $S,T$ disjoint and every pair in $S\times T$ an edge.
Four-cycle-freeness makes one side singleton. Each of the two old forms
on the other side contains at most one missing neighbor of that singleton,
by the preceding observation. There are at most two witnesses.
Thus this pure-grid form cannot supply all four.

For the second form with linear $L$, the safe terms are
\[
 ACDL,\quad BCDL,\quad AEL,\quad BEL.
\]
If either quartic has two witnesses, its equality case again forces
$|\supp(L)|=2$. Otherwise all four terms must have one different witness.
The two cubics give full $A=\{a\},B=\{b\}$, missing $xa,xb$, and
$x\in E\cap L$. The quartic with $A$ has a different edge through $a$;
the old form $L$ meets it and also contains $x$, so the missing-neighbor
observation forces $a\in L$. Similarly $b\in L$. This is impossible
for the old form $L$ containing the two missing neighbors of $x$.

For scalar $L=1$, the mixed-grid terms are $ACD,BCD,AE,BE$.
Each has at most one witness. Four distinct witnesses would force safe
quadratic equality $E=\{x\},A=\{a\},B=\{b\}$, with missing
$xa,xb$. The other common cubics then have new witnesses $ya,yb$.
Distinctness gives a four-cycle $x-a-y-b-x$. Thus scalar context is
impossible here too. We have proved in all cases that a four-witness
circuit has linear $L$ with a full two-vertex support.

Let $S=\supp(L)$. It covers every missing edge, since each witness occurs
in $FL$. The unique cut retains degree one, so follow a selected route
to its first retained product. Its entire opposite old operand lies in
the full direct context $L$ by positivity, and therefore $S$ covers
all old edges generated by that product. The insertion cannot reach both
operands, by the forbidden old square already noted. Every other retained
quadratic product is a star: its two disjoint linear supports span a
complete bipartite graph, so one side is singleton in a four-cycle-free
graph. Choose one center for each of these $K-1$ products and add $S$.
This covers the full graph, including all missing edges, with at most
$K-1+2=K+1$ vertices. Repeated centers, coefficients and fanout do not
increase the bound. No changed formula has been re-extracted.
\end{proof}


\begin{proposition}[Two multiplicative highs with zero cross coefficient]
\label{prop:unaffectedzerocross}
Let a positive variable-input circuit compute $P\in\Sup(G)$, where $G$
is finite and four-cycle-free, the quadratic support $H$ is nonempty,
and exactly four edges are missing. Suppose there are at most four
private products and two private additions, both visible cuts.
After full-current-carrier erasure, assume their highs $g,f$ are
multiplication-only networks on old homogeneous linear boundaries, using
the original private products. Raw ancestor sharing and repeated inputs
are permitted, but neither network reads a surviving perturbed selected
linear input. Suppose also that the exact independent-symbol frame has
zero mixed coefficient, so its Boolean output is
\[
 P_2\vee gL\vee fM,
\]
with the actual scalar or old linear direct contexts $L,M$.
If $K$ products are retained, then $\vc(G)\le K+2$.
Retained addition counts and cover gain are unrestricted.
\end{proposition}
\begin{proof}
Let $J=E(G)\setminus E(H)$. If $\vc(J)\le2$, choose one center
for each retained quadratic product to cover $H$, and add a two-vertex
$J$-cover. This gives the result. Suppose $\vc(J)\ge3$.
Each positive term $gL,fM$ is safe, and their union supplies every missing
witness. Scalar factors may be omitted; a zero term only lowers the count.
Keep the original frame and count supports without changing its gates.

Let $p\le4$ be the number of useful private multiplication nodes in the
union of the two high ancestor networks. Compress each old boundary
signal to an atom, identifying equal full supports only for this count.
The resulting undirected multigraph has $p$ operation vertices, $b_0$
boundary vertices and $2p$ operand edges with multiplicity. Every vertex
feeds one of two outputs, so it has at most two connected components.
The elementary inequality edges $\ge$ vertices minus components gives
\[
 b_0\le p+2\le6.
\]
This includes repeated inputs, shared private nodes, shared old supports,
and a high ancestral to the other. If an erased high reduces to a boundary
atom, it can also be counted as a one-vertex component.

The two direct contexts add at most two distinct supports. Thus the safe
products $gL,fM$ use at most eight distinct old supports together.
Lemma~\ref{lem:sevenoldfactors} forces exactly eight. Equality gives
$p=4$, $b_0=6$, two components, and two nonzero linear contexts whose
supports are new and distinct. Moreover the multigraph has exactly
vertices minus components edges, so both components are trees.
No private operation or old boundary support is shared between the
outputs, and neither component has a repeated operand or internal
fanout. The highs are therefore two multiplication trees.
Write their operation counts as $p_g,p_f$, with $p_g+p_f=4$.
The two direct products have respectively $p_g+2,p_f+2$ old factors.

Allowing $p_g=0$ causes no problem: the distributions $0+4$ and $1+3$
give factor counts $2+6$ and $3+5$, respectively, and the old-factor
bound permits at most three missing witnesses in their union.
Four witnesses consequently force $p_g=p_f=2$.
Both direct products are quartics and each has exactly two missing
witnesses, with disjoint edge sets. The full equality statement in
Proposition~\ref{prop:oldfactors} makes each of its four supports a
two-point transversal of its witnessed induced matching. In particular
\[
 |\supp(L)|=|\supp(M)|=2.
\]
This is the full actual context support, including any outside vertices;
the union of the two witness pairs need not be an induced matching.

Put $S=\supp(L)\cup\supp(M)$. It covers all missing edges.
If $|S|=2$, the $K$ retained-star centers together with $S$ prove the claim.
If $|S|=3$, follow either insertion's selected degree-one path to its first
retained product. Its entire nonzero opposite operand lies in that cut's
full direct context, by positivity of the addition-only paths to the
quadratic output. Thus $S$ covers every old edge supplied by this product.
One center for every other retained product and $S$ give at most
$K-1+3=K+2$ vertices.

If $|S|=4$, the contexts are disjoint. Choose a retained product reached
by each cut. These products are distinct. If both symbols reached the
same operand of one product, its nonzero opposite operand would lie in
both context supports, impossible. If they reached opposite operands,
that retained product and its positive output path would contribute a
positive mixed coefficient, contrary to the independent-symbol hypothesis.
Hence $S$ covers the old edges of two distinct retained products as well
as every missing edge. The other $K-2$ centers and $S$ form a cover of size
at most $K+2$.

Each retained quadratic product is indeed a star: its disjoint selected
linear supports form a complete bipartite subgraph of $H$, and two
nonsingleton sides would make a four-cycle. Repeated centers, coefficients,
shared addition paths and additional marked products can only reduce the
size of the covers constructed above. This proves the bound on the full
graph, without any optimality hypothesis. A positive mixed coefficient
or a surviving affected old input is outside the stated scope.
\end{proof}


\begin{proposition}[Two direct witnesses from a cubic high]
\label{prop:twowitnesscover}
Let $G$ be finite and four-cycle-free, and let a positive variable-input
circuit compute $P\in\Sup(G)$ with nonempty quadratic support $H$ and
exactly four missing edges. Keep its original extracted frame, with $K$
retained products. Suppose two visible insertion cuts have exact
independent-symbol output
\[
 P_2+z_1L+z_2M+qz_1z_2,\qquad q\ge0,
\]
where $L,M$ are their actual old direct contexts, scalar or homogeneous
linear according to the selected degrees. After full-current-carrier
erasure, substitute $z_1=g,z_2=f$. Suppose $g$ is a product of three old
nonzero homogeneous linear forms, $gL$ witnesses two distinct missing
edges, and Booleanly
\[
 f=R\vee gT,
\]
where $R$ is a product of at most three old homogeneous linear forms and
$T$ is any nonnegative polynomial; zero terms are permitted. Then
\[
 \vc(G)\le K+2.
\]
There is no private-product budget, retained optimality, nonancestry or
zero-cross requirement. The displayed decomposition is a hypothesis about
the actual high, not a normalization of every two-cut circuit.
\end{proposition}
\begin{proof}
The safe term $gL$ has two missing witnesses, so $L$ is nonzero and linear.
The full equality case of Proposition~\ref{prop:oldfactors} makes its four
factor supports the four two-element transversals of an induced matching
of two missing edges. In particular $\supp(L)$ has two elements. The three
full factor supports of $g$ have empty common intersection. Any two-variable
support in any multiple of $g$ meets all three supports. It cannot use an
outside vertex, because its other endpoint would then lie in their empty
intersection; among the four matching endpoints, the only $G$-edges are
the original two matching edges. Thus every missing witness of any
multiple of $g$ belongs to that pair, regardless of the multiplier's
degree, support or dependence on $g$.

Substitute $f=R\vee gT$ into the exact frame. All contributions other than
$P_2$ and $RM$ are divisible by $g$, including the positive-cross terms
when $q>0$. Consequently $RM$ must witness both other missing edges.
A scalar $M$ would leave at most three old factors, which is impossible.
The old-factor bound forces $R$ to have exactly three nonzero factors and
$M$ to be a nonzero old linear form. Full quartic equality gives
$|\supp(M)|=2$. Therefore
\[
 S=\supp(L)\cup\supp(M)
\]
has at most four vertices and covers every missing edge.

Mark every original retained quadratic product whose selected input is
reached from either cut. Both contexts are linear, so both selected lows
have degree one and each cut marks at least one product. A cut cannot
reach both inputs of one product, because its nonzero old low would
contribute a square to $P_2$. The entire opposite old operand of every
marked product occurs, with positive coefficient, in that cut's full
direct context. This follows by expanding the actual selected linear
paths with independent insertion symbols; all paths from a retained
product to the quadratic output use additions. Hence $S$ covers every
$H$-edge supplied by every marked product.

If at least two products are marked, take $S$ and one star center for each
unmarked product. This covers $H$ and all missing edges with at most
$K-2+4=K+2$ vertices. Such a center exists because a squarefree quadratic
product is complete bipartite and four-cycle-freeness makes one operand
singleton.

Suppose exactly one product is marked. Its old operands are $U,V$.
If both cuts reach the same operand, their full direct contexts are
positive scalar multiples of the same opposite operand: no other marked
product contributes a direct term. Thus their supports coincide and
$|S|=2$, giving a cover with at most $K-1+2$ vertices. If the cuts reach
opposite operands, their full direct supports are $\supp(V)$ and
$\supp(U)$. Both have two elements by the two quartic equalities, and they
are disjoint because $P_2$ is squarefree. The original retained product
would contain a four-cycle. This is impossible, even when the mixed
coefficient is positive. Arbitrary fanout, repeated centers and positive
coefficients do not change these support and incidence arguments.
\end{proof}


\begin{proposition}[Two affected factors after a cubic high]
\label{prop:twoaffectedcover}
Keep the graph, four-missing-edge and original-frame hypotheses of
Proposition~\ref{prop:twowitnesscover}, with exactly two deleted additions,
both visible. Suppose the first high is an old cubic $g=UVW$ and the
second high is Booleanly
\[
 f=C(A\vee g)(B\vee g),
\]
with $U,V,W,C,A,B$ nonzero old homogeneous linear forms. The first
selected low $a_0$ is nonzero and linear, with
$\supp(a_0)\subseteq\supp(A)\cap\supp(B)$ and $a_0L$ supported on $H$.
If $q>0$, also require the actual mixed-context inclusion
$\supp(a_0)\subseteq\supp(M)$. With the full direct contexts and fixed
frame as above, $\vc(G)\le K+2$. Retained counts, coefficients, fanout and
repetitions are unrestricted. These are hypotheses on the actual two
affected inputs, not a general normalization assertion.
\end{proposition}
\begin{proof}
We use the following full-support observation. If $xy$ is missing and an
old form contains $y$, then $x$ has no other $G$-neighbor in that support:
an additional neighbor and the old neighborhood center would give a
four-cycle. The center cannot be $x$, since $xy$ is missing.

Suppose a safe old quintic $UV WCM$ witnesses an induced missing pair
with endpoint set $E$. If $M$ is not a full two-element transversal of
that pair, every $G$-edge witnessed by any multiple of $g=UVW$ belongs
to the pair. Here is the required equality argument, including possible
extra support vertices. The four independent endpoint transversals each
require a rejecting factor, and one factor can reject only one of them.
A factor that uniquely rejects a transversal $T$ has full support
$E\setminus T$: if it had an extra vertex $z$, all other factors meet
$T$, while the preceding old-neighbor observation makes $z$ nonadjacent
to every vertex of $T$. This would give an unsafe monomial.

If $M$ rejects no transversal, the other four factors are distinct full
transversals. If $M$ rejects one but is not full, another factor shares
that rejecting class, and the remaining three factors are distinct full
transversals. When all three $g$-factors are full, their empty common
intersection gives the confinement proved in
Proposition~\ref{prop:twowitnesscover}. Otherwise two $g$-factors, say
$U,V$, and $C$ are the three full unique rejectors, while $W,M$ have the
same fourth endpoint restriction. If $U,V$ are disjoint, an edge hitting
both has both endpoints in $E$, and is one of the induced pair. If they
meet at $x$, the fourth endpoint restriction contains the mate $x'$ of
$x$ and omits $x$. A new edge hitting $U,V$ would have to be $xz$ with
$z\notin E$. To meet $W$ it requires $z\in\supp(W)$; that support also
contains $x'$, so the old-neighbor observation forbids $xz$. This proves
the quintic confinement in every case.

Boolean distribution gives $f=ABC\vee gC$. First let $q=0$. The missing
witnesses are supplied by the safe terms
\[
 gL,\qquad X=ABCM,\qquad Y=gCM.
\]
If $gL$ has two witnesses, Proposition~\ref{prop:twowitnesscover} applies.
Assume it has at most one. A zero $M$ is impossible. If $M$ is scalar,
$X$ has at most one witness; if the quartic $Y$ has two, its full equality
confines the $gL$ witness to the same pair, and otherwise all three terms
have at most one. Thus four witnesses force $M$ to be old and linear.

If $X$ has two witnesses, full quartic equality gives
$|\supp(M)|=2$. Otherwise $X$ has at most one, so $Y$ must have two and
$gL$ must have a witness outside that pair. The quintic confinement just
proved again forces $M$ to be a full two-element transversal. In both
cases $\supp(M)$, supplemented by one endpoint of the possible $gL$
witness, is a set $S$ of at most three vertices covering all missing
edges. The second actual cut marks a retained product whose entire
opposite old factor is contained in $M$. Thus $S$ covers all old edges of
that product; one center for each other retained product gives
\[
 \vc(G)\le K-1+|S|\le K+2.
\]
The charge uses the original full direct context and permits fanout.

Now let $q>0$. The actual low inclusion supplies a vertex in
$\supp(A)\cap\supp(B)\cap\supp(M)$. Hence the safe old quartic $X=ABCM$
cannot have two missing witnesses: the four full transversals at equality
have empty intersection in every three distinct supports. Absorption of
the positive mixed term gives
\[
 P=P_2\vee gL\vee X\vee gC
\]
at Boolean-support level. Both $gL$ and $gC$ are safe old quartics.
If either has two missing witnesses, its full equality confines every
missing witness of the other $g$-multiple to that same pair. Otherwise
each has at most one. In either case these two terms together supply at
most two missing edges, and $X$ supplies at most one. Four missing edges
are impossible. This proves the asserted cover bound for all $q\ge0$,
while keeping the original frame and gate counts fixed.
\end{proof}


\begin{proposition}[Three old or affected boundaries after a cubic high]
\label{prop:cubicthreeboundary}
Keep the original-frame, graph and four-missing-edge assumptions of
Proposition~\ref{prop:twoaffectedcover}. Let $g$ be an old cubic and suppose
the actual second high has Boolean-support form
\[
 f=\prod_{i\in I}(A_i\vee g)\prod_{j\notin I}A_j,
 \qquad I\subseteq\{1,2,3\},
\]
with three nonzero old homogeneous linear forms $A_1,A_2,A_3$.
For $I\ne\varnothing$, the first nonzero old linear low $a_0$ is contained
in every affected $A_i$, and $a_0L$ is supported on $H$. When $q>0$,
also require the actual inclusion $\supp(a_0)\subseteq\supp(M)$.
For $I=\varnothing$, retain the ordinary two-symbol low/context identities.
Then $\vc(G)\le K+2$, for arbitrary $K$ and $q\ge0$.
No claim is made that every four-product circuit has these three boundary
factors; a direct input from a proper private-product ancestor is separate.
\end{proposition}
\begin{proof}
Put $J=E(G)\setminus E(H)$. If $\vc(J)\le2$, add a two-vertex
$J$-cover to one center for every retained star. Otherwise
Theorem~\ref{thm:twosumscomponents} makes $J$ either $4K_2$ or
$P_3\sqcup2K_2$. In particular, given two disjoint $J$-edges, at least one
of the other two edges avoids their four endpoints: two attachments would
give cover number at most two, or a component other than $K_2,P_3$.

We first prove a support observation. Let $X,Y$ each be a product of at
most three old forms. Safe expressions $XL,YD,XY$ cannot together witness
all four $J$-edges if the first two each witness at most one. Indeed $XY$
has at most six factors, so it must witness exactly two edges, an induced
matching with endpoint set $E$, while the other terms supply distinct
extra edges $e_X,e_Y$. Each of the four independent endpoint transversals
must have a rejecting factor. A support rejects at most one, and a
uniquely rejecting support is the full complementary transversal: any
extra vertex, together with the uniquely rejected independent assignment,
would satisfy the product, and the old-center four-cycle argument makes
that extra vertex nonadjacent to the assignment. This is the same
full-support argument as in Lemma~\ref{lem:sevenoldfactors}.

If both extra edges avoid $E$, no factor can be a unique rejector, since
every $X$-factor meets $e_X$ and every $Y$-factor meets $e_Y$. Four classes
would need eight distinct supports, impossible. Otherwise suppose $e_X$
meets $E$ at $v$ and $e_Y$ avoids $E$. Every $X$-factor meets two incident
missing edges and hence contains $v$: omitting it would put their two
other endpoints in one old neighborhood and create a four-cycle. The
two endpoint transversals containing $v$ must therefore be rejected by
$Y$-factors. None can reject uniquely, because it also meets $e_Y$.
They would require four distinct $Y$-supports, although $Y$ has at most
three factors. The symmetric case is identical. This proves the observation
for arbitrary positive multipliers $L,D$; only the old factors of $X,Y$
were counted.

Write $R=A_1A_2A_3$ and $k=|I|$. If $gL$ has two missing witnesses,
Proposition~\ref{prop:twowitnesscover} applies immediately, since the
second high is $R$ plus a multiple of $g$. Assume henceforth that $gL$
has at most one. For $k>0$, put $T=\prod_{j\notin I}A_j$, with the
empty product used only as the Boolean unit in this analysis. Then
\[
 f=R\vee gT.
\]

Let $q=0$. For $k=0$, the only terms $gL,RM$ supply at most one and two
witnesses, respectively. For $k>0$ the safe terms are
\[
 gL,\qquad RM,\qquad gTM.
\]
If $RM$ has two witnesses, $M$ is linear and full quartic equality gives
$|\supp(M)|=2$. This support meets all witnesses of $RM$ and $gTM$;
add one endpoint of the possible $gL$ witness to obtain a set $S$ of at
most three vertices covering $J$. The actual second cut marks a retained
product whose entire opposite old factor lies in $M$. Thus $S$ covers
its old edges, and the other $K-1$ centers give $\vc(G)\le K+2$.
If $RM$ has at most one witness, use the preceding observation with
$X=g,Y=TM$: the latter has at most $3-k+1\le3$ old factors, and every
one occurs in $RM$. A scalar context shortens the count; a zero term
cannot supply four witnesses. This settles $q=0$.

Now let $q>0,k=0$. The three terms are $gL,RM,gR$. If $RM$ has two
witnesses, its three full $R$-supports are distinct transversals of that
matching and have empty common intersection. Every $G$-edge witnessed by
a multiple of $R$ is confined to the same pair, as in
Proposition~\ref{prop:twowitnesscover}. Thus $gR$ adds no edge and $gL$
adds at most one. Otherwise the observation applies with $X=g,Y=R$.

Finally let $q>0,k>0$. Absorption leaves
\[
 gL,\qquad RM,\qquad gT.
\]
If $RM$ has at most one witness, apply the observation with $X=g,Y=T$;
when $T=1$, safe $g$ itself has at most one witness and the conclusion
is immediate. Suppose $RM$ has two. Its four full supports
$A_1,A_2,A_3,M$ are distinct two-point transversals. If $k\ge2$, the
nonzero $a_0$ lies in $M$ and two affected supports, contradicting the
empty intersection of any three distinct transversals. Hence $k=1$.
Rename the affected form $A$ and the other two $B,C$. The nonempty
intersection of $A,M$ permits the following relabeling:
\[
 \begin{gathered}
 RM\text{ witnesses }ab,cd,\\
 A=\{a,d\},\quad B=\{a,c\},\quad C=\{b,c\},
 \quad M=\{b,d\},\quad\supp(a_0)=\{d\}.
 \end{gathered}
\]
These are full supports. Booleanly $BC=c\vee ab$. Any new edge of
$gBC$ must therefore be $cv$ with $v$ outside the quartet. There is at
most one, since $cd\in J$ and $J$ has maximum degree two. Four witnesses
force this edge and a further edge $xy$ from $gL$, disjoint from the five
vertices $a,b,c,d,v$. Thus $J$ consists of the path $d-c-v$ and the two
separate edges $ab,xy$. We prove that the full first context is singleton.

Choose the names so $y\in\supp(L)$. The old identity gives
$\supp(L)\subseteq N_H(d)$. The safe old product $gc$ witnesses $cv$.
A missing edge is isolated in the graph of all exact edge supports of a
safe old product, by the first part of Proposition~\ref{prop:oldfactors}.
Consequently $g$ is false on $cx$ and $cy$. The same isolation for the
$xy$ witness of $gL$, using $y\in L$, makes $g$ false on $vy$.
Each of its three old factors meets $cv$ and $xy$ and can reject at most
one of these three cross pairs. Their restrictions to $\{c,v,x,y\}$ are
therefore, in some order,
\[
 \{v,y\},\qquad\{v,x\},\qquad\{c,x\}.
\]
In particular $g$ is true on $vx$. The last old factor contains $c,x$;
an edge $vx$ would give a four-cycle through its old center and the
missing edge $cv$. That center cannot equal $v$, since $cv$ is not old.
Thus $vx$ is a $G$-nonedge. No assertion about the factors' additional
full support is needed here.

Take $z\in\supp(L)\setminus\{y\}$. The old edge $dz$ excludes $d,c$.
The choice $z=v$ would give $gL$ the additional missing witness $cv$,
and $z=x$ would make it true on the nonedge $vx$. Otherwise all relevant
vertices are distinct. An edge $xz$ would make the four-cycle
$d-y-x-z-d$, and an edge $vz$ would make $d-c-v-z-d$. Both are forbidden.
The independent set $\{v,x,z\}$ would then satisfy $gL$, a contradiction.
Hence the full context is $L=\{y\}$.

Choose $y$ as the center covering the old edges of one retained product
marked by the first cut, and one center for each of the other products.
This gives an $H$-cover with at most $K$ vertices that already contains
$y$. Add $a,c$ to cover the remaining missing edges $ab,cd,cv$; the edge
$xy$ is already covered. The total is at most $K+2$. All counts concern
the original selected frame and full direct contexts. The cases exhaust
every affected subset, scalar or zero context and mixed coefficient.
\end{proof}


\begin{proposition}[Nonancestral products with only visible cuts]\label{prop:nonancestralexceptional}
Let $G$ be finite and four-cycle-free, and let a positive variable-input circuit compute $P\in\Sup(G)$ with nonempty quadratic support $H$, exactly four missing edges, and $\vc(G)-\vc(H)\ge3$. Suppose the extracted frame has exactly $\vc(H)$ products, and let $C$ be its distinct minimum cover of assigned centers. Assume $G[V(G)\setminus C]$, with isolates retained, is $P_4\sqcup K_2$ or $K_3\sqcup K_2$. If every deleted addition is a visible unequal-degree cut and the original deleted products are pairwise nonancestral, at least four products are deleted. No bound on retained additions or visible cuts is imposed.
\end{proposition}
\begin{proof}
Use full-current-quadratic-carrier erasure, preserving the exact frame. It does not create an ancestry relation absent from the original DAG, since each shortcut follows an original directed path. A surviving private product cannot read another private product, directly or through retained gates. It cannot read a retained quadratic carrier, whose off-frame multiples have been absorbed. Its two surviving operands are therefore degree-one forms. No private degree-one addition can occur in these operands: it would be a hidden equal-degree addition, contrary to the all-visible hypothesis. The high of a visible cut in such an operand either contains an earlier private product, forbidden by nonancestrality, or is an absorbed carrier. Thus every surviving private product has the form $g_i=U_iV_i$ with two old linear forms. In particular it has minimum degree two, and every useful cut inserting it retains degree one.

The minimum-product frame has one product for each $c\in C$, with its selected center factor supported only at $c$. Its full center input is consequently
\[
 A_c=\alpha_c c+\sum_i\alpha_{ci}g_i,
 \qquad \alpha_c>0,\quad\alpha_{ci}\ge0.
\]
All intervening operations are additions; another private product on such a path would violate nonancestrality, and a full quadratic carrier was erased. This allows arbitrary fanout, nested cuts, repeated coefficients and center-input addition networks.

Put $T=V(G)\setminus C$ and $N=G[T]$. The safe-center argument in Proposition~\ref{prop:exceptionaloutside} applies: $N$ is one of the two well-covered exceptional graphs, all four missing edges lie outside $C$, and the cover gain is at least three. After $C=0$, each full center input is safe for $N$, and replacing every retained product by that full center input preserves $f_N$. For clarity, if a center input were true on an independent set, extend it to a maximal independent set of $N$. Its complement is a three-vertex cover containing the entire outside old neighborhood of that center; replacing the center by this cover would cover all of $G$ with at most $\vc(H)+2$ vertices. This contradicts the assumed gain. Thus the safe-center replacement gives
\[
 f_N=\bigvee_{i\in I}\overline{g_i},
 \qquad
 \overline{g_i}=(U_i|_T)(V_i|_T),
\]
for a set $I$ of surviving private products. Each included nonzero product is individually safe, since it is a positive summand of a safe full center input. Products appearing only in leaf inputs need not belong to $I$; their effect is accounted for by the safe-center replacement, rather than discarded through a new minimum-part computation.

The factors remain old relative to the original pair $H\subseteq G$, even when their old neighborhood centers lie outside $T$. A safe product of two old forms witnesses at most one missing edge by Proposition~\ref{prop:oldfactors}. Directly, two disjoint cross edges would give a four-cycle, while two adjacent missing edges would put their two outer endpoints in one old neighborhood and again give a four-cycle or make the edges old. Every one of the four minimal true assignments of $N$ must be witnessed by a disjunct $\overline{g_i}$. No such disjunct witnesses two. Hence $|I|\ge4$, and the original deleted-product count is at least $|I|$. Absorbed products only increase the original count. The argument retains every outside isolate and requires neither coefficient one nor a minimal retained addition count.
\end{proof}

\begin{lemma}[Independent old products with five combined factors]\label{lem:independentfive}
Let $H\subseteq G$ be finite simple graphs with $G$ four-cycle-free. Suppose $J=G\setminus H$ has exactly four edges and every nontrivial component of $J$ is $K_2$ or $P_3$. Let $A,B,C,D,E,L,M$ be nonzero old homogeneous linear forms, with repetitions permitted. If the three products
\[
 X=ABCL,\qquad Y=DEM,\qquad Z=ABCDE
\]
are safe in $G$, their union cannot witness all four $J$-edges as exact two-variable supports. No actual-context, cover-gain or retained-count assumption is imposed.
\end{lemma}
\begin{proof}
First consider any safe product of five old forms witnessing two missing edges $ab,cd$. These form an induced matching by Proposition~\ref{prop:oldfactors}; put $W=\{a,b,c,d\}$. Each factor meets both pairs. Each of the four independent transversals $T$ of this matching must be missed by some factor, whose restriction to $W$ is then exactly $W\setminus T$. A factor can miss at most one transversal.

If $F$ is the unique factor missing $T$, its full support is $W\setminus T$. Indeed suppose an extra vertex $z\notin W$ belongs to $F$. Every other factor meets $T$. For $t\in T$, its matching mate $s$ lies in $F$. Let $h$ be an old $H$-neighborhood center for $F$. An edge $zt$ would give the four-cycle $h,s,t,z$; the only possible center coincidence $h=t$ would make $st$ old. Thus $z$ has no neighbor in $T$. Choosing $z$ in $F$ and a $T$-variable in every other factor produces a monomial supported in the independent set $T\cup\{z\}$, contradicting safety.

Four transversal classes must be missed by five factors. At most one class can have more than one rejecting factor, so at least three factors have full supports equal to distinct two-element transversals. A factor meeting all four transversals does not affect this count.

Now $X,Y,Z$ have at most two, one and two missing witnesses respectively. If $X$ has two, its four full factor supports are the four transversals of their matching quartet. In particular $A,B,C$ are three distinct transversals, with empty common intersection. Any exact edge witnessed by $Z$ must meet all three. It therefore cannot have an endpoint outside that quartet, and the only $G$-edges inside it are the two already witnessed by $X$. Adding $Y$ gives at most three witnesses.

Hence a hypothetical union witnessing all four edges has the disjoint allocation: one edge $e$ from $X$, one edge $f$ from $Y$, and an induced pair $ab,cd$ from $Z$. Use its endpoint set $W$. At least three of $A,B,C,D,E$ have full two-element transversal supports, so at least one of $A,B,C$ is contained in $W$. The edge $e$ must meet it and cannot lie wholly in $W$. Write $e=xw$ with $x\in W$ and $w\notin W$, and let $x'$ be the matching mate of $x$; call the other matching edge $yy'$.

Each of $A,B,C$ meets both missing edges $xx'$ and $xw$, so each contains $x$. Otherwise its old neighborhood would contain $x',w$, yielding a four-cycle through $x$ or making the missing edges old. The independent transversals $\{x,y\}$ and $\{x,y'\}$ must consequently be rejected by $D,E$. One factor cannot reject both while meeting both matching edges. They are therefore the unique rejectors, and the full-support result gives, after exchanging $D,E$,
\[
 \supp(D)=\{x',y'\},\qquad \supp(E)=\{x',y\}.
\]
The different witness $f$ meets both these supports and cannot lie wholly in $W$. Thus $f=x'v$ for some $v\notin W$. The three distinct missing edges $xw,xx',x'v$ lie in one connected component. They form a triangle if $v=w$, and a three-edge path otherwise. Both contradict the stipulated $K_2/P_3$ component structure.
\end{proof}

\begin{corollary}[The independent-high five-factor template]\label{cor:independentfive}
Under the same graph hypotheses, suppose $g$ is a product of at most three old forms, $f$ a product of at most two, $L,M$ are nonnegative affine old contexts, and $q\ge0$. The safe terms
\[
 gL,\qquad fM,\qquad qgf
\]
cannot witness all four missing edges. Here an affine old context means a nonnegative scalar plus an old homogeneous linear form, including zero. This is an algebraic template statement, without a claim that all circuits of a given gate budget have this form.
\end{corollary}
\begin{proof}
A positive constant in a context makes its corresponding high safe. That high absorbs the mixed term Booleanly, and the two direct terms supply at most three missing witnesses. Thus it suffices to use homogeneous contexts. Shorter products can be padded by repeating old factors, without changing Boolean supports; this is an identity for support analysis, not a new charged circuit.

If $q=0$ the direct bounds are two and one. If $L=0$, the bounds for $fM$ and $gf$ are one and two. If $M=0$ and $gL$ has at most one witness, the same total bound holds; if $gL$ has two, its four-transversal equality confines every $gf$ witness to those same two edges, exactly as in the preceding proof. Zero highs simplify these cases further. The remaining nonzero homogeneous-context case is Lemma~\ref{lem:independentfive}.
\end{proof}

\begin{proposition}[Three products and two additions cannot supply four missing edges]\label{prop:threeproductstwosums}
Let $G$ be finite and four-cycle-free, and let a positive variable-input circuit compute $P\in\Sup(G)$ with nonempty quadratic support $H$. If exactly four edges of $G$ are absent from $H$, canonical extraction cannot delete at most three multiplications and at most two additions. Retained counts, cover gain and quadratic optimality are unrestricted.
\end{proposition}
\begin{proof}
Put $J=E(G)\setminus E(H)$. Lemma~\ref{lem:fivelost} forces exactly three private products and two private additions under the assumed budgets. Any Boolean-preserving removal that preserves the exact extracted frame would leave at most four private gates and is therefore impossible. Apply full-current-carrier erasure, keeping the original gate identities, minimum degrees and selected frame. All three private products remain essential. Theorem~\ref{thm:twosumscomponents} says that every nontrivial component of the full missing graph $J$ is $K_2$ or $P_3$. Thus no three distinct missing edges form a connected subgraph.

All identities below analyze Boolean support with that frame and those gate counts fixed. Every positive output summand is safe for $G$. A safe summand witnesses a missing edge precisely when it is true on that edge's two-variable assignment. Boolean absorption and removal of repeated factors therefore preserve the relevant witness sets. An old factor means a nonzero homogeneous linear form supported in an original $H$-neighborhood, with singleton supports allowed. Scalars and zero contexts are dealt with explicitly below, rather than treated as old factors.

\emph{Four support bounds.} Proposition~\ref{prop:oldfactors} implies that the missing witnesses of a safe product of $k$ old factors form an induced matching and number at most $\lfloor\log_2 k\rfloor$. A quartic with two witnesses has as its full factor supports the four distinct transversals of their matching. A quintic with two witnesses has at least three distinct full transversal factors: each of the four independent transversals must have a rejecting factor; a uniquely rejecting factor is full by the old-center four-cycle argument in Lemma~\ref{lem:independentfive}; among five factors at least three of the four rejecting classes are singletons.

Two distinct full transversals of an induced matching have intersection empty or a singleton $x$. Any other missing edge meeting both must then have one endpoint $x$ and the other outside the matching's four endpoints. Indeed an edge wholly within those endpoints is already one of the matching edges. This gives the following four bounds.

\emph{(A)} Two safe old products using at most six distinct boundary factors together have at most three missing witnesses. Remove repetitions, and write them $RA,RB$, with $m$ common factors and $a,b$ noncommon factors. If an arm is empty, absorption leaves one product with at most six factors and hence at most two witnesses. Otherwise $a,b\ge1$ and $m+a+b\le6$. Four witnesses would require two from each product, so $m+a,m+b\ge4$. For $m\le1$ this is impossible. For $m=2$ both products are quartics, with two common full transversals. For $m=3$ at least one product is a quartic, with three common full transversals. For $m=4$ both are quintics; at least two of the three full transversals of the first are common. Larger $m$ violates the factor budget. In every case two common full transversals confine any new witness of the second product to their singleton intersection. It cannot have two such new edges, because its own missing witnesses form an induced matching. Thus four distinct witnesses are impossible. This particular bound does not need the component restriction on $J$.

\emph{(B)} At most three safe products, each with at most four old factors and all sharing two factors $U,V$, together have at most three missing witnesses. If no term has two, count the terms. Otherwise its quartic equality makes $U,V$ distinct full transversals. Every additional missing edge must attach to their intersection $x$. Two additional edges, together with the original matching edge incident to $x$, would make a connected three-edge subgraph of $J$.

\emph{(C)} Let $g=UV$, let $L,M$ be old linear contexts or scalars, and let $Q$ be a product of at most three old factors. If $gL,gM,QM$ are safe, their union has at most three missing witnesses. Zero contexts give the separate factor bounds. A nonzero scalar $L$ makes $g$ absorb $gM$, giving at most $1+2$ witnesses; a scalar $M$ makes $g$ absorb $gL$, giving a smaller bound. Otherwise the two cubics have at most one witness each, and their witnesses form a star by Lemma~\ref{lem:commoncubic}. Four total witnesses would require two distinct incident edges from those cubics and two different edges from a quartic $QM$. Its full $M$-support is a transversal of its matching. The new edge from $gM$ meets that support and attaches to one of the matching edges. Its incident companion from $gL$ would put three distinct missing edges in one component, contrary to the shape of $J$.

\emph{(D)} Let $g=UV,h=XY$, let $Z$ be an old product of at most two factors, and let $L$ be an old linear context or scalar. Safe terms $gL,hZ,gh$ together have at most three missing witnesses. Zero terms cause no difficulty; scalar $L$ lets $g$ absorb $gh$, giving $1+2$. Suppose $L$ is linear. If $gh$ has two witnesses, its four factors are the four full transversals. If $U,V$ intersect at $x$, the complementary pair $X,Y$ intersects at the matching mate $x'$; if the first pair is disjoint, so is the second. Any new edge from $gL$ attaches at $x$, and one from $hZ$ attaches at $x'$. There is at most one new edge on either side, and one on each side would join three missing edges through $xx'$. Thus there are at most three altogether. If $gh$ has at most one witness, four in total would require two from $hZ$, one new edge $e$ from $gh$, and one further edge from $gL$. The full $X,Y$ supports must intersect at an endpoint $x$ of the $hZ$ matching; write them $\{x,y\},\{x,y'\}$. Booleanly $h=x\vee yy'$. Hence $gx$ is safe and witnesses $e$: the summand $gyy'$ can witness no new edge other than the already present $yy'$. The common-cubic lemma makes the new witnesses of $gx$ and $gL$ incident. Together with the matching edge at $x$, they again give three connected missing edges. A shorter nonzero $h$ may be padded by repeating a factor; equality demanding distinct full transversals then excludes that degeneration.

We also use Lemma~\ref{lem:sharedquartics}: two safe products sharing three old factors and having at most four factors each jointly have at most two missing witnesses. If either has two, its three common full transversals have empty intersection and confine every witness of the other to the same matching. Otherwise each has at most one. Repeated-factor padding includes shorter products.

\emph{The fixed-frame cases.} The classification in the proof of Lemma~\ref{lem:fourmatchinglost} uses original private products, visible cut occurrences and full carrier erasure, independently of its later matching restriction. There must be a visible cut, because a private signal first enters a selected path through the discarded high input of a selected unequal-degree sum. With one visible cut, its high uses the three products and the other addition, which may be an equal-degree sum or a hidden cut. This connected binary four-gate DAG has at most five distinct old linear boundaries. Its Boolean value is $C_0(A_0\vee B_0)$, with monomial factors; extract common factors and remove repetitions. Its scalar or old linear output context gives at most two safe old products using at most six distinct boundaries. Bound (A) excludes four witnesses, including zero contexts or absorbed arms.

Both additions are consequently visible. Give their highs separate symbols before substituting their actual dependence. The significant output has the form
\[
 P_2\vee gL\vee fM\vee qgf,\qquad q\ge0,
\]
where $L,M$ are full direct contexts, each zero, scalar or old homogeneous linear. A diagonal insertion term would force a square of an old low in the squarefree quadratic frame, so is absent. Let $k$ be the number of private products ancestral to the first high $g$, including that product. Before its visible cut there are only products of old boundaries, and every remaining private product feeds $f$. Thus $k=1,2,3$ exhausts the cases.

If $k=3$, then $f$ is $g$ or a product ancestor, and $g=fA$ Booleanly, with at most four distinct old factors in their union. For $q=0$ the two terms $gL,fM$ use at most six boundaries, so (A) applies. For $q>0$, $gf=g$ and absorption leaves $g\vee fM$, also covered by (A).

If $k=2$, write $u=UV$ and $g=uW$, allowing $W=1$ for a repeated-input degeneration. When the remaining product $f$ reads the raw uninserted $u$, its possibilities are
\[
 f=uA,\quad uA\vee g,\quad u,\quad g.
\]
The first gives $uWL,uAM,q\,uWA$; the second, when $q=0$, gives $uWL,uAM,uWM$. These are at most three quartics sharing $U,V$, so (B) applies. When $q>0$ in the second case, absorption leaves $g\vee uAM$, and (A) applies. The last two possibilities similarly leave at most two products using six boundaries or fewer.

Without a raw $u$ input, the possibilities are $f=AB$, $AB\vee gB$, or $AB\vee g$, where old parts may vanish. For $f=AB$, $q=0$ gives bounds two and one; for $q>0$ the terms $gL,ABM,gAB$ are Corollary~\ref{cor:independentfive}. For $f=AB\vee gB$ and $q=0$, the terms are
\[
 gL,\quad ABM,\quad gBM.
\]
With linear $M$, these are that same independent-high template, with three-factor high $g$, two-factor high $BM$, and contexts $L,A$. With scalar $M$, the two quartics $gL,gB$ share three factors and give at most two witnesses, while $AB$ gives at most one. For $q>0$, absorption leaves $gL,ABM,gB$, with the same shared-three-factor bound plus one cubic. Finally, $f=AB\vee g$ gives $gL,ABM,gM$ when $q=0$, again at most two plus one; when $q>0$ it reduces to $g\vee ABM$.

It remains that $k=1$, so $g=UV$. If $f$ is unaffected, it is an old product of at most three factors. The $q=0$ bounds are one and two; $q>0$ is Corollary~\ref{cor:independentfive} with the highs exchanged. Otherwise the two remaining products multiply at most three factors consisting of old parts perturbed by $g$, with raw $g$ permitted. Boolean distribution and absorption put
\[
 f=DE\vee gD,
\]
where $D$ is a product of $d\le2$ unaffected old factors and $DE$ has at most three old factors. Here $D$ is an actual unaffected subproduct, not an independent formal parameter. Repeated operands and more than one affected factor satisfy the same identity. The two output forms are
\[
 q=0:\quad gL\vee DEM\vee gDM,
 \qquad
 q>0:\quad gL\vee DEM\vee gD.
\]
For the moment take $L$ nonzero linear. If $d=2$, then $E$ has at most one old factor. For $q=0$ and linear $M$, use the independent-high template with highs $DM,g$ and contexts $E,L$. For scalar $M$ use (D) with $h=D,Z=E$; for $q>0$ use (D) with $h=D,Z=EM$. If $d=1$ and $q=0$, use (D) with $h=DM,Z=E$ when $M$ is linear; scalar $M$ leaves three cubics. For $q>0$ use (C) on $gL,gD,D(EM)$, since $EM$ has at most three factors. If $d=0$ and $q=0$, use (C) on $gL,gM,EM$; for $q>0$ the raw $g$ absorbs $gL$, leaving $g\vee EM$ with at most one plus two witnesses. An empty $E$ makes $f=D$ by absorption and belongs to the unaffected case.

\emph{Zero and scalar boundaries.} The preceding uses of the independent-high template never need an unproved extension to a scalar old factor. If $L$ is a nonzero scalar, the safe $g$ absorbs $qgf$ and every affected multiple of $g$, leaving $g\vee DEM$. These use at most six old boundaries: two for $g$, at most three for $DE$, and one for $M$ when $k=1$; three, two and one respectively when $k=2$. Ancestor and raw-$u$ cases have fewer. Bound (A) applies. If $L=0$, an affected high leaves $DEM\vee gDM$ for $q=0$, or $DEM\vee gD$ for $q>0$; if $M=0$, it leaves $gL\vee qgD$. Their boundary unions have size at most six, so (A) applies. An unaffected high leaves respectively $fM\vee qgf$ or $gL\vee qgf$, with the same bound. The ancestor and raw-$u$ cases only shorten these expressions.

For nonzero scalar $M$, the $k=2$ cases above reduce to the shared-three-factor pair plus $AB$, or to $g\vee AB$. In the affected $k=1$ cases, $d=2$ gives (D), $d=1$ gives three cubics, and $d=0$ gives $g\vee E$. An unaffected $f$ has at most three factors; it absorbs $qgf$, and $f\vee gL$ has at most one plus two witnesses. For $k=3$, $f$ absorbs $gL$ since $g=fA$. A vanishing old part either removes a term, gives $f=0$, or gives $f=gD$; the latter leaves two safe old products using at most six boundaries and is (A). Repetitions or empty Boolean subproducts do not increase any bound. A nonzero scalar output summand would itself violate safety.

Every original placement and ancestry case therefore has at most three missing witnesses, contrary to $|J|=4$. No retained-count or cover hypothesis was used.
\end{proof}


\begin{lemma}[Three recursive products with safe output disjuncts]\label{lem:recursivesafeoutputs}
Let $H\subseteq G$ be finite simple graphs, with $G$ four-cycle-free. Suppose $H[T]$ is empty and $N=G[T]$, apart from isolates, is $P_4\sqcup K_2$ or $K_3\sqcup K_2$. Let $U,V,A,B,C,D$ be zero or old homogeneous linear forms supported on $T$, where old neighborhoods refer to the original $H$ and may have centers outside $T$. For nonnegative coefficients, put
\[
\begin{split}
 g_1&=UV,\\
 g_2&=(A+\alpha g_1)(B+\beta g_1),\\
 g_3&=(C+\gamma g_1+\delta g_2)(D+\eta g_1+\theta g_2).
\end{split}
\]
The Boolean disjunction of any subset of $g_1,g_2,g_3$ differs from $f_N$. Nonzero scalar constants are not old forms here. Repeated operands, zero old parts and unsafe unselected intermediate products are permitted.
\end{lemma}
\begin{proof}
Call a product safe if every monomial contains an $N$-edge. A selected disjunct is safe. Any additional safe $g_i$ can be adjoined to the asserted output without changing $f_N$, so we may take all safe products as output disjuncts. Zero is safe and supplies no witness. A multiple of a safe product cannot supply a new exact edge support: its safe factor already contains the same edge.

A safe product of two old forms has at most one missing witness, and at equality both factors are singleton. Safe old cubics have at most one witness, and quartics at most two, with the full transversal equality case of Proposition~\ref{prop:oldfactors}. The union of two missing-edge stars cannot equal $N$, whose cover number is three. Also $f_N$ needs six positive CNF clauses. Booleanly the second product, after exchanging $A,B$ if necessary, has one of the forms
\[
 g_2=AB,\qquad AB\vee g_1,\qquad B(A\vee g_1).
\]
Zero old parts are included. The first two products alone cannot suffice: if $g_1$ is safe it supplies at most one edge and a safe $g_2$ adds at most the old edge of $AB$; if $g_1$ is unsafe, a safe $g_2$ is an old quadratic or $AB\vee UVB$, with at most two witnesses. We can therefore assume $g_3$ is selected and safe.

\emph{$g_1$ is safe.} Put $h=AB$. All terms divisible by $g_1$ supply at most its single witness. If $h$ is safe, its multiples supply at most its single witness, and the only other possible part of $g_3$ is the safe old quadratic $CD$. There are at most three edges. If $h$ is unsafe, $g_2$ is unsafe and is not selected. It can occur in at most one operand of safe $g_3$, since its square would retain an unsafe monomial. Apart from multiples of safe $g_1$, the remaining parts are $CD$ and $h$ times the old linear part of the opposite operand. They are one safe old quadratic and one safe old cubic. Including the witness of $g_1$ still gives at most three.

\emph{$g_1$ is unsafe and $g_2$ is safe.} Both coefficients in $g_2$ cannot be positive, since they would insert $g_1^2$. Thus $g_2$ is an old quadratic, or $AB\vee UVB$. The latter two terms are safe. If $AB$ has an edge, its old factors are singleton, and all these witnesses meet the vertex in $B$; if $AB=0$, only the one cubic edge remains. Thus the witnesses of $g_2$ lie in one star. Unsafe $g_1$ can occur in at most one operand of $g_3$. After absorbing multiples of safe $g_2$, the remaining terms are $CD$ and one old cubic $g_1D$, or symmetrically $g_1C$. These form another star: a nonzero safe $CD$ has singleton factors, while $CD=0$ leaves just one cubic edge. Two stars cannot supply $N$.

\emph{Both $g_1,g_2$ are unsafe.} Neither is selected, and neither can occur in both operands of safe $g_3$. If they occur on the same side, or one is absent, write
\[
 g_3=D(C\vee\epsilon g_1\vee\zeta g_2).
\]
The second-product forms $AB$ and $AB\vee g_1$ give at most $CD,UVD,ABD$, three old products with at most one witness each. The same holds for $g_2=B(A\vee g_1)$ when $\epsilon>0$, since $UVBD$ is absorbed by $UVD$. When $\epsilon=0$, only
\[
 CD,\qquad ABD,\qquad UVBD
\]
remain. Four witnesses would require two from the quartic and one from each other term. Quartic equality forces $D$ to have two vertices, whereas the witnessed safe quadratic $CD$ forces $D$ singleton. If a coefficient or term vanishes, the count only decreases.

It remains that the two unsafe products occur on opposite sides:
\[
 g_3=(C\vee g_1)(D\vee g_2).
\]
For $g_2=AB$, this has the four-clause CNF
\[
 (C\vee U)(C\vee V)(D\vee A)(D\vee B),
\]
contrary to the six-clause requirement. For $g_2=AB\vee g_1$, the mixed term contains a Boolean copy of unsafe $g_1$, impossible. Finally, if $g_2=B(A\vee g_1)$, distribution and absorption give exactly
\[
 g_3=CD\vee CAB\vee UVD\vee UVB.
\]
All four terms are safe. If $CD=0$, the three old cubics supply at most three edges. Otherwise $C,D$ are singleton; $CD,CAB$ form a star at the vertex of $C$, and $UVD,UVB$ form a common-cubic missing star by Lemma~\ref{lem:commoncubic}. Again the targets have a two-vertex cover, impossible. These cases cover all safety choices, zero restrictions and repeated inputs.
\end{proof}

\begin{corollary}[All visible cuts on the exceptional outside graphs]\label{cor:allvisibleexceptional}
Under the graph, cover-gain, minimum-retained-product and exceptional-outside hypotheses of Proposition~\ref{prop:nonancestralexceptional}, suppose every deleted addition is a visible unequal-degree cut. Then at least four products are deleted, without an ancestry restriction. Retained and deleted addition counts are otherwise unrestricted.
\end{corollary}
\begin{proof}
Suppose at most three products are deleted and apply exact full-current-carrier erasure. A surviving private product cannot read a retained quadratic carrier. Every remaining operand is an old linear part plus a nonnegative linear combination of earlier private-product signals: retained linear paths contain only additions, and no extra private homogeneous or hidden sum is available. A raw private product corresponds to a zero old part. Topologically ordering the at most three products therefore gives exactly the recursive forms of Lemma~\ref{lem:recursivesafeoutputs}, with zero unused slots permitted as algebraic padding rather than charged circuit gates.

Set the assigned minimum center cover to zero and keep all outside isolates. The full selected center factor at a center $c$ has the form $\lambda_c c+\sum_i\lambda_{ci}g_i$, with $\lambda_c>0$. As proved in Proposition~\ref{prop:nonancestralexceptional}, the exceptional-outside safe-center argument makes every full center input safe and permits replacing each retained product by that input. This argument uses only the original selected center and outside old neighborhood, not private-product nonancestry.

A visible cut that retains degree two also needs to be included. Its higher-degree high enters the protected quadratic output through additions and is a complete output disjunct. After carrier erasure such a useful high is a private product; a separate private sum in that high would be hidden, contrary to the hypothesis. Include all such direct high disjuncts as well. Replacing retained products by their center inputs increases the restricted Boolean output, but every resulting disjunct is safe: it is either a positive summand of a safe center input or already an additive output summand. The resulting disjunction is therefore both at least and at most $f_N$.

We have obtained $f_N$ as a disjunction of a subset of at most three recursive private products. Restricted old factors still use their original $H$-neighborhoods, including centers outside the restricted vertex set. Lemma~\ref{lem:recursivesafeoutputs} gives the contradiction. Erased products only increase the original deleted count. Hidden additions and other outside shapes remain outside this conclusion.
\end{proof}

\begin{proposition}[Three private products require a hidden sum of two degree-one inputs]\label{prop:hiddenlinearnecessary}
Under the graph, cover-gain, minimum-retained-product and exceptional-outside hypotheses of Proposition~\ref{prop:nonancestralexceptional}, at most three deleted products force an original deleted addition whose two inputs both have minimum degree one. The number, visibility and degrees of all other deleted additions and all product dependencies are unrestricted. This equal-minimum-degree addition is hidden; its full operands need not be homogeneous.
\end{proposition}
\begin{proof}
Suppose no such original deleted addition exists. For an original node $v$, let $\ell_v$ be its homogeneous degree-one part. If $v$ is retained and $\ell_v\ne0$, this is its selected old linear value. Every nonsingleton selected linear support lies in an original $H$-neighborhood. A variable is an old singleton.

If $v$ is deleted with $\ell_v\ne0$, it is an addition, since variable-input products have minimum degree at least two. By the contrary assumption exactly one child has minimum degree one. Thus $\ell_v$ equals the linear part of that child. Repeating this descent reaches a retained old form or variable. Every original degree-one part is therefore old or zero. The argument concerns original degree labels; a minimum-degree-one operand may contain higher terms.

Apply full-current-quadratic-carrier erasure, retaining the exact selected frame. It removes contributions containing retained quadratic carriers from private-product operands and preserves their degree-one parts. Between variables, earlier private products and any private-product operand, only additions remain. Expanding these paths as an identity gives an old linear part plus a nonnegative linear combination of earlier private products. No gate is duplicated or recounted. Topologically ordering the at most three products yields exactly the recursive forms in Lemma~\ref{lem:recursivesafeoutputs}, including zero old parts, repeated inputs and both affected operands. If necessary, algebraic zero slots pad the list; no new original gate is asserted.

Set the assigned minimum center cover to zero, retaining all outside isolates. The old parts remain old for the original pair $H\subseteq G$, even if their neighborhood centers lie outside the restricted set. Exceptional-outside safe-center replacement turns each retained product into its full center input while preserving $f_N$. The old part of that input is supported on its assigned center and vanishes. All remaining operations outside the fixed private products are additions, so the entire outside output is a nonnegative linear combination of those products, possibly together with a nonnegative linear form in outside variables. The latter must vanish: any such variable would make the output true on a singleton, including an isolated singleton.

Direct higher-degree insertions into retained degree-two output paths are included in this same additive output representation. Hence $f_N$ is a disjunction of a subset of at most three recursive private products. Included products are safe; unselected intermediate ones need not be. This is exactly the possibility excluded by Lemma~\ref{lem:recursivesafeoutputs}. The contradiction proves the claim. An original equal-minimum-one addition would be retained if its occurrence lay on a selected path, so the required gate is hidden.
\end{proof}

\begin{lemma}[One arbitrary linear factor on an exceptional target]\label{lem:onearbitraryfactor}
Let $H\subseteq G$ be finite simple graphs, with $G$ four-cycle-free. Suppose $H[T]$ is empty and $N=G[T]$, with isolates retained, is $P_4\sqcup K_2$ or $K_3\sqcup K_2$. Let $S$ be any nonnegative homogeneous linear form on $T$. If $F=S\prod_iL_i$ is safe in $N$, where there is at least one nonzero old linear factor $L_i$, then $F$ witnesses at most two of the four $N$-edges. Old neighborhoods are those of the original $H$, with centers allowed outside $T$. The number and repetitions of old factors are unrestricted.
\end{lemma}
\begin{proof}
If $S=0$ there is nothing to prove. An old factor cannot contain two vertices of $T$ with a common $N$-neighbor: its original $H$-center is outside $T$, and together these vertices and their common neighbor would give a four-cycle. Every witnessed edge meets every factor.

For a triangle $abc$ and separate edge $xy$, an old factor contains at most one triangle vertex, so cannot meet all three triangle edges. Three witnessed edges would therefore include $xy$ and two triangle edges, say $ab,bc$. Every old factor must contain $b$; otherwise it contains the forbidden pair $a,c$. Choosing $b$ from each old factor forces $\supp(S)\subseteq N_N(b)=\{a,c\}$ by safety. This prevents the required $xy$ witness.

For a path $a-b-c-d$ and separate edge $xy$, an old factor cannot contain $a,c$ or $b,d$. If all three path edges were witnessed, every old factor would contain both $b,c$. Safety would force $S$ into the disjoint neighborhoods $N_N(b)=\{a,c\}$ and $N_N(c)=\{b,d\}$, a contradiction. Two adjacent path edges together with $xy$ similarly force every old factor to contain their common center, again preventing the isolated-edge witness.

The remaining three-edge subset is $ab,cd,xy$. Every old factor meeting $ab$ and $cd$ meets each independent pair $I=\{a,c\}$ and $J=\{b,d\}$: omitting one would force the other forbidden pair into that old factor. Their product is therefore true on both $I,J$. Safety forces every variable of $S$ into both open neighborhoods $N_N(I)=\{b,d\}$ and $N_N(J)=\{a,c\}$, which are disjoint. This final contradiction exhausts all three-edge subsets. Isolates belong to none of the displayed neighborhoods and cannot evade the argument.
\end{proof}

\begin{proposition}[One visible cut in the exceptional six-gate profile]\label{prop:onevisibleexceptional}
Under the graph, cover-gain, minimum-retained-product and exceptional-outside hypotheses of Proposition~\ref{prop:nonancestralexceptional}, it is impossible that at most three products and at most three additions are deleted and exactly one deleted addition is a visible cut. The other additions may have arbitrary original degrees and inhomogeneous full operands; retained additions, fanout and repeated inputs are unrestricted.
\end{proposition}
\begin{proof}
Proposition~\ref{prop:gainthreesix} forces all six private gates to remain useful: an exact-frame-preserving removal of any one would give at most five deletions. Thus the budgets are three products and three additions. By Proposition~\ref{prop:hiddenlinearnecessary}, at least one of the two hidden additions has both original input minimum degrees one. Original degrees and selected occurrences are kept throughout.

\emph{The unique high.} After full-current-carrier erasure every useful private gate is an ancestor of the sole visible cut's high $F$. Indeed its first entry into the selected subgraph must be a discarded high of a selected unequal-degree sum; products and equal-degree sums select both inputs. A private gate following the unique cut cannot feed its earlier high without a cycle. Retained nonlinear boundaries into $F$ have been erased. Retained linear boundaries are unchanged old forms: a prior perturbation would already require a visible entry, and a selected linear path through a hidden minimum-one sum would select that sum. Hence $F$ is computed from old linear boundaries by three products and two hidden additions.

The whole output is Booleanly $P_2\vee FM$. A selected low of degree two gives a positive scalar context $M$; a degree-one low gives a nonzero old homogeneous linear context, since its product with that old low is supported on $H$. Repeated insertion within a product would force a square of the nonzero old low in $P_2$, and is forbidden. On $T=V\setminus C$, retaining isolates, we therefore have
\[
 f_N=FM.                                                        \tag{*}
\]
A linear $M$ meets every $N$-edge, so has at least three support vertices. Either exceptional target has six distinct minimal positive CNF clauses.

\emph{Positions of a degree-one hidden sum.} Call one such addition $S$ and the other hidden addition $h$. If both are originally sums of two degree-one inputs, they are homogeneous linear computations from old boundaries: there is no other unequal minimum-one sum through which a private product could enter. The high is then a monomial in linear forms, computed by three products and therefore having at most four distinct factors. Including $M$ gives at most five CNF clauses, impossible.

If only $S$ has this original degree type, it is a pure linear sum of old forms unless it reads $h$. Such a dependency is direct: a product has minimum degree at least two, and a retained selected linear path through $h$ would select it. In the dependent case $h=A+g$ with old $A$ and monomial $g$ on old forms; then $S=A+B+g$ with old $B$, or a repeated $h$. All other operations are products, and $h$ implies $S$ Booleanly. Thus the high simplifies to $R(A\vee g)$ or $R(Q\vee g)$ with $Q=A\vee B$, while raw $g$ and repeated inputs only simplify further. Replacing $h,S$ by $g$ in the original product DAG computes the monomial $Rg$ with at most three products, hence at most four distinct old factors. Distributing the displayed disjunction gives at most four linear clauses; $M$ adds at most one. This also contradicts six. Zero restricted old parts do not increase the count.

It remains that $S$ is one homogeneous linear boundary form, not required to be old, and $h$ is the only other addition in the high circuit. Regard $S$ as one atom. A circuit of products and a single disjunction has Boolean value
\[
 F=R(P\vee Q),
\]
where $P,Q,R$ are monomials in linear boundary atoms. Repeated occurrences of $h$ are absorbed; common atoms of $P,Q$ are extracted into $R$, and atoms already in $R$ are removed from the other factors. The three atom sets may therefore be taken disjoint, with at most one nonold atom $S$ in their union. A connected DAG with three products and one addition has at most five distinct boundary sources. Writing $p,q,r$ for these disjoint atom counts gives
\[
 p+q+r\le5,\qquad F\text{ has at most }r+pq\text{ clauses}.       \tag{**}
\]
If a branch is absorbed, $F$ is a monomial with at most four factors and the previous bound applies. Otherwise $p,q\ge1$.

For scalar $M$, the only way~(**) reaches six is $r=0$, $\{p,q\}=\{2,3\}$. For linear $M$, which adds one clause, the only possibilities are that case and $r=1,p=q=2$.

\emph{Two-versus-three branches.} When $r=0,p=2,q=3$ and $M$ is scalar, the two safe terms are products of two and three linear atoms. At most one term contains $S$. The other is old and has at most one missing witness, while the term containing $S$ has at least one old factor and at most two witnesses by Lemma~\ref{lem:onearbitraryfactor}. Their union has at most three.

For old linear $M$, the two safe terms have three and four factors. An all-old three-factor term has at most one witness. An all-old four-factor term cannot have two: the full four-transversal equality would force $M$ to have two support vertices, whereas it covers $N$ and has at least three. The term containing $S$ again has at most two witnesses by Lemma~\ref{lem:onearbitraryfactor}. Thus four target edges remain impossible.

\emph{The six-clause grid.} Finally suppose $M$ is linear and $r=1,p=q=2$. Write the single atom in $R$ as $R$ and $P=P_1P_2,Q=Q_1Q_2$. Formula~(*) has six clauses
\[
 M,\quad R,\quad P_1\vee Q_1,\quad P_1\vee Q_2,
       \quad P_2\vee Q_1,\quad P_2\vee Q_2.                      \tag{***}
\]
Each must be a distinct minimal vertex cover of $N$. To justify this exactly, every maximal independent assignment falsifies a clause; that clause is a vertex cover because it is implied by $f_N$, and minimality forces it to equal the complementary minimal cover. Six assignments and six clauses leave no repetitions, extra isolated vertices or redundant clauses.

On $K_3\sqcup K_2$, an old form cannot cover the triangle, contradicting $M$. On the path $a-b-c-d$ and separate edge $xy$, the six covers are
\[
 acx,\ acy,\ bcx,\ bcy,\ bdx,\ bdy.
\]
An old form cannot contain $a,c$ or $b,d$, so $M$ is $bcx$ or $bcy$. In the four union clauses in~(***), a variable cannot occur in just one cell: each occurrence comes from a whole row factor $P_i$ or column factor $Q_j$. Consequently the number of $ac$-family covers in this grid is zero or two, and the same holds for the $bd$ family. Four distinct cells must consist of two whole families. If one family were $bc$, the omitted pair $M,R$ would be an $ac$ or $bd$ family, contradicting old $M$. Hence the grid is precisely
\[
 acx,\ acy,\ bdx,\ bdy,
\]
and the two separate clauses are $bcx,bcy$.

The two $ac$ cells must be one row or column; a diagonal pair cannot be supplied by unions of row and column factors without additional occurrences of $a$. Both $a,c$ must belong to that same row or column factor, since either in a perpendicular factor would enter a $bd$ cell. The complementary row or column similarly forces a different factor containing $b,d$. Both factors are nonold, contradicting that at most one of the four factors can be $S$. This exhausts~(**) and proves the proposition. The two-visible configuration is not excluded by this proof.
\end{proof}



\begin{proposition}[Two visible cuts in the exceptional six-gate profile]\label{prop:twovisibleexceptional}
Under the graph, cover-gain, minimum-retained-product and exceptional-outside
hypotheses of Proposition~\ref{prop:nonancestralexceptional}, it is impossible
that at most three products and three additions are deleted, exactly two
deleted additions are visible cuts, and the third addition has both original
input minimum degrees equal to one. The full operands may be inhomogeneous;
product dependencies, repeated inputs, fanout and retained additions are
unrestricted.
\end{proposition}
\begin{proof}
Put $T=V(G)\setminus C$ and $N=G[T]$, retaining isolates. Its four edges
are all missing, its cover number is three, and its edge function needs six
distinct positive linear CNF clauses. Proposition~\ref{prop:gainthreesix}
makes all six private gates essential under exact-frame-preserving carrier
erasure. Keep their original identities and minimum degrees. Call the two
visible highs $g,f$ in cut order and the hidden addition $S$.

The original $1+1$ gate $S$ cannot lie on a selected degree-one path, which
would retain it. It cannot be a visible high: both of its input minima are
one. After full-current-carrier erasure, each visible high is therefore a
private-product output. Every useful private gate is ancestral to one of
these highs, since its first entry into a selected path must be through a
visible cut. The highs may coincide or have a product-ancestry relation.

The exceptional-outside safe-center argument used in
Proposition~\ref{prop:nonancestralexceptional} does not require nonancestral
products. If a full center input were true on an $N$-independent assignment,
extend that assignment to a maximal independent set. Its complementary
three-vertex cover contains every outside old neighbor of the center.
Replacing that center by this cover would cover $G$ with at most
$\vc(H)+2$ vertices, contrary to the gain. Thus every full center input is
safe, and replacing retained products by these inputs preserves the outside
output. Only $g,f$ can enter such an input as private signals. A cut retaining
degree two contributes its high directly along a protected quadratic-output
addition path; that high is again $g$ or $f$. Consequently
\[
 f_N=\bigvee_{h\in I}h,\qquad\varnothing\ne I\subseteq\{g,f\},       \tag{*}
\]
with restrictions to $T$ understood. Only included highs, and their positive
summands, are asserted safe. An uninserted private product is not an extra
output disjunct. No new extraction is performed.

All old forms below remain supported in original $H$-neighborhoods, including
neighborhood centers outside $T$. A safe old product of at most three linear
forms has at most one missing witness. Lemma~\ref{lem:onearbitraryfactor}
bounds a safe product with one arbitrary homogeneous linear factor and at
least one old factor by two $N$-edge witnesses. Repetitions of that arbitrary
factor collapse Booleanly.

We first record a useful \emph{pair bound}. A safe union of two monomial
products, of lengths at most two and three respectively, cannot equal $f_N$
if at most one distinct linear factor $S_0$ across both products is nonold.
If $S_0$ occurs in just one term, that term has at most two witnesses by the
preceding lemma and the old term has at most one. A term without any old
factor is a linear Boolean function and, unless zero, cannot be safe. If
$S_0$ occurs in both terms, factor it out after removing repetitions: the
remaining lengths are at most one and two, giving at most
$1+1\cdot2=3$ linear CNF clauses. An empty arm gives absorption and fewer
clauses. If $S_0$ is absent, both terms have at most one witness. These
possibilities exclude $f_N$. Zero terms only reduce the counts.

\emph{A homogeneous hidden sum.} Suppose no visible insertion reaches an
original degree-one operand of $S$. Then $S=U+V$ is genuinely homogeneous,
with old $U,V$ and positive coefficients absorbed into the forms. There is
no other hidden addition through which a private product could enter such
an operand. A retained selected linear path cannot read $S$ without
retaining it. Thus $S$ is a literal boundary factor of private products,
rather than a freely available summand in another linear boundary. Denote
this one possibly nonold atom by $S_0$; all other linear boundaries are old.
A pure sum computed later can be placed earlier in this dependency analysis
without changing the original circuit or its counts.

Let $k$ count the private products ancestral to the first high $g$, including
$g$. Before that cut they compute monomials in old forms and possibly $S_0$.
If $k=3$, the other high is $g$ or a product ancestor and Booleanly divides
$g$. Any selected subset in (*) therefore reduces to a monomial with at
most four distinct linear factors, giving fewer than six clauses.

If $k=2$, write $u=AB$, $g=uD$, allowing repeated-factor degenerations.
The last product forming $f$ can read old or $S_0$ boundaries, old selected
linear forms perturbed by $g$, or a raw private ancestor $u$ or $g$.
There is no extra private addition to introduce an additive uninserted
$u$. A literal $S_0$ is not itself perturbed by the visible cut: that would
require another hidden addition after it.
If $f$ reads raw $u$, its Boolean possibilities are
$uE$, $uE\vee g$, $u$, and $g$, using $ug=g$.
A selected union is a monomial or $u(D\vee E)$, with at most three clauses.
Without raw $u$, the possibilities are
\[
 f=EF,\qquad EF\vee gF,\qquad EF\vee g.
\]
Here unaffected factors may be $S_0$, and zero old parts include raw-$g$
inputs. If $g$ is selected, absorption gives $g\vee EF$, excluded by the
pair bound. If only $f$ is selected, the first form has two clauses; the
second is $F(E\vee ABD)$, with at most four; and the third is again covered
by the pair bound. Zero parts and repeated factors only shorten them.

If $k=1$, then $g=AB$ has at most two linear factors. The two remaining
private products forming $f$ multiply at most three boundary factors,
unaffected old/$S_0$ forms or old parts perturbed by $g$; raw $g$ is allowed.
Without an affected operand, $f$ is a monomial of length at most three;
alone it has at most three clauses, and together with $g$ the pair bound
applies. Otherwise distribution and absorption give
\[
 f=HK\vee gH,
\]
where $H$ has $h$ unaffected factors, $K$ has $j$ old parts of affected
factors, and $h+j\le3$. A raw $g$ makes $K=0$.
If $g$ is selected, absorption gives $g\vee HK$, again excluded by the
pair bound. If only $f$ is selected and $h\ge1$, the CNF
$H(K\vee AB)$ has at most $h+2j\le5$ clauses.
For $h=0$ it is $K\vee g$, so the pair bound applies. For $K=0$ it is
$gH$, a monomial with at most four distinct factors. This exhausts the
homogeneous case, without declaring an unselected earlier product safe.

\emph{An affected hidden sum.} If $S$ has higher terms, a visible cut must
reach one of its original degree-one operands. It must be the first cut:
an effect from the last cut would have to return from $S$ to a later visible
high, creating a directed cycle. Thus $S$ lies after the first cut and
before a product feeding $f$, and the first cut retains degree one.
Before it, $g$ is a monomial in old forms. Only additions occur along the
degree-one paths to $S$, so its restricted Boolean value is
\[
 S=T_0\vee g,
\]
where $T_0$ is one arbitrary homogeneous linear form, possibly zero.
Positive coefficients, or both operands affected, give the same identity.
If the higher contribution vanishes upon restriction, the homogeneous case
already applies. A selected linear path cannot read $S$, so its later use
is as a literal factor. Other affected boundaries have the form $A\vee g$
with $A$ old, or are raw $g$. The only possibly nonold atom after expansion
is $T_0$.

The value $k=3$ leaves no product after $S$ to turn its minimum-one output
into a visible high; $S$ would be useless, contrary to essentiality.
For $k=1$, write $g=AB$ old. The two remaining products again give
$f=HK\vee gH$, with $h+j\le3$, where $H$ consists of unaffected old
factors and $K$ contains at most the one nonold atom $T_0$ after repetitions.
If $g$ is selected, the output $g\vee HK$ has at most one old-quadratic
witness and two witnesses from the other term. If that term has no old
factor it is linear and cannot be nonzero and safe. If only $f$ is selected
and $h\ge1$, there are at most $h+2j\le5$ clauses.
For $h=0$ the same witness argument applies to $K\vee g$; for $K=0$ the
monomial $gH$ has at most four clauses.

Finally let $k=2$, so $u=AB$ and $g=uD$ are old. Only one product remains
to form $f$. If it reads raw $u$, usefulness of affected $S$ forces the
other operand to be $S$, and
\[
 f=uS=uT_0\vee g=u(T_0\vee D).
\]
This absorbs $g$ whether or not that high is selected, and has at most three
clauses. A raw-$g$ input gives $gS=g$, and is simpler.
Without raw $u$, the last product again gives $f=HK\vee gH$, now with
$h+j\le2$, with $H$ old and at most the atom $T_0$ nonold in $K$.
If $g$ is selected, $g\vee HK$ has at most one old-cubic witness and two
other witnesses. If only $f$ is selected and $h\ge1$, the CNF
$H(K\vee ABD)$ has at most $h+3j\le4$ clauses. For $h=0$ the preceding
witness count applies; a raw-$g$ factor gives a monomial of length at most
four. Repeated $S$ inputs and zero restrictions reduce these bounds.

An original hidden $1+1$ addition is either homogeneous or affected as
described, since there is no other hidden addition and a minimum-one operand
cannot be a product. The cases include every private-product ancestry and
every selected subset in (*), with original frame, labels and counts kept
fixed. Each contradicts the required four witnesses or six clauses.
\end{proof}

\begin{proposition}[Nonancestral private products require four additions]\label{prop:nonancestraladditions}
Under the graph, cover-gain, minimum-retained-product and exceptional-outside hypotheses of Proposition~\ref{prop:nonancestralexceptional}, suppose the original private products are pairwise nonancestral. Then at least four original additions are deleted. The number of deleted products, the number of retained additions, and the types and placements of deleted additions are unrestricted.
\end{proposition}
\begin{proof}
Put $T=V(G)\setminus C$ and $N=G[T]$, keeping isolates. All four $N$-edges are missing. Apply exact full-current-carrier erasure, without introducing gates or new ancestry. A surviving private product cannot read another private product, and its off-frame retained quadratic inputs have been absorbed. Thus it is $g_i=U_iV_i$ with homogeneous linear factors. These may initially use private linear additions and need not yet be old.

Every full selected center input is $\alpha_c c+\sum_i\alpha_{ci}g_i$. Its unchanged degree-one part has support only at $c$, by positivity, and intervening nonlinear contributions use additions only. The exceptional-outside safe-center argument therefore gives, after setting $C=0$,
\[
 f_N=\bigvee_{i\in I}\overline{g_i},
 \qquad\overline{g_i}=(U_i|_T)(V_i|_T),
\]
where each included nonzero product is individually safe. This safety follows after center replacement. Let $s=|I|$ count distinct included original products. Products omitted from this representation will not be charged.

A safe quadratic product is a star, since two nonsingleton sides would give a four-cycle. Choose its singleton center $u_i$ and let $R_i$ be its leaf support; either orientation is allowed for one edge. Group the products by these chosen centers. For center $u$, let $k_u$ be its number of products and $m_u$ the size of the union of their leaf supports. Then
\[
 \sum_u k_u=s,\qquad \sum_u m_u\ge |E(N)|=4.          \tag{a}
\]
Every edge is supplied, while multiple appearances only increase the second sum.

Consider the linear preparation of the $k_u$ leaf factors. An old boundary form contributing to one has at most one restricted variable. Indeed its support lies in $N_G(u)$ by positivity. Two distinct leaves $v,w$ in the same old neighborhood, with center $z$, would give the four-cycle $u,v,z,w$; $z=u$ would make the two edges old, also impossible. Original neighborhood centers outside $T$ are included in this argument. Selected linear boundary values cannot contain a private linear sum that is missing from the extraction, nor can they contain a private-product contribution when feeding another private product, by nonancestrality. Thus all further leaf-preparation nodes are original private linear additions.

Erase empty restricted branches for counting, replace singleton-support nodes by their variable, identify repeated variable sources, and prune to these leaf outputs. If $d_u$ nonsingleton addition nodes remain, their operand graph has at least $m_u+d_u$ vertices, at most $2d_u$ distinct edges, and at most $k_u$ components. Hence
\[
 d_u\ge m_u-k_u.
\]
The charged nodes for different centers are disjoint. A node with two distinct restricted leaves used at centers $u,u'$ would give two common neighbors of those centers and a four-cycle; singleton nodes were not charged. If $a_{\mathrm{in}}$ counts the original private additions available for these preparations,~(a) therefore gives
\[
 a_{\mathrm{in}}\ge\sum_u d_u\ge4-s.                \tag{b}
\]
Sharing, repeated centers and positive coefficients are allowed throughout this graph count.

There is a separate charge for introducing the $s$ quadratic sources. For each included $g_i$, choose an addition-only route to a full selected center input and stop at its first visible entry cut. Such a cut retains degree one and has a quadratic high. Before entry the route uses only private quadratic additions: a retained quadratic carrier is an erased off-frame use, while a minimum-degree-one signal would need a further multiplication to become a high of larger degree. Such a private multiplication would create forbidden ancestry, and a retained multiplication is an erased carrier. A hidden unequal sum with linear low cannot stay hidden on an addition-only route into a selected degree-one input, since its own occurrence would then be selected.

Take the union of the private quadratic-addition ancestors of these entry highs, including both operands of every sum. Let $s'\ge s$ be its product sources, $d$ its private quadratic addition nodes, and $t$ its distinct visible-cut terminals. Every undirected component contains a terminal; if there are $c$ components, $c\le t$. Counting operand occurrences on $s'+d+t$ vertices gives
\[
 2d+t\ge s'+d+t-c,\qquad d+t\ge s'+t-c\ge s.        \tag{c}
\]
Direct source-to-cut edges, shared subgraphs and repeated operands all satisfy this count. No cut is made free merely because its old low vanishes after restriction. Its original gate is one of the terminals, even when that low already carries another insertion.

The additions in~(c) are disjoint from the preparation additions in~(b). Internal nodes in~(c) have degree two, and an entry cut contains a private-product source; it could not then feed a private-product operand without violating nonancestry. Preparation nodes lie before products and contain no such source. Thus the original deleted-addition count is at least
\[
 a\ge a_{\mathrm{in}}+(d+t)\ge(4-s)+s=4.
\]
Removed gates only increase the original count. The proof uses the minimum-product, gain and exceptional-outside assumptions through safe-center replacement; it does not extend that replacement to other outside graphs or to dependent products.
\end{proof}

\begin{proposition}[Four outside matching edges with two actual contexts]\label{prop:fouroutsidegain}
Suppose exactly four edges of a four-cycle-free graph $G$ are missing from the nonempty quadratic support $H$. Suppose the extracted frame has $|E(H)|+\vc(H)-1$ gates, and let $C$ be its minimum cover of assigned star centers. Assume the four missing edges form an induced matching with all eight endpoints outside $C$. If at most four products and at most two additions are deleted, then
\[
 \vc(G)\le\vc(H)+2.
\]
In particular the cover-gain-three, six-deletion profile with four deleted products and two additions is impossible when the outside graph is this matching and isolates.
\end{proposition}
\begin{proof}
Write $T$ for the eight matching endpoints. Full-carrier erasure preserves the exact frame. Lemma~\ref{lem:fourmatchinglost} rules out at most five remaining deletions, so the assumed budgets leave exactly four useful products and two additions. The algebraic classification in the proof of Lemma~\ref{lem:fourproductssix}, before its maximal-cover-gain argument, applies. After restriction to $T$, one visible cut gives at most ten positive linear clauses; the $k=4$ case at most six; $k=3$ at most fifteen, or nine with a direct private ancestor; $k=1$ at most fifteen. For $k=2$ a direct-ancestor case gives at most eight, and the other cases at most fifteen except two independent cubic highs with zero cross coefficient. These bounds include zero, scalar and repeated-input degenerations. The four-edge matching requires sixteen clauses. Thus only
\[
 P_2\vee gL\vee fM
\]
remains, where $g,f$ are independent products of three old homogeneous linear boundary forms and $L,M$ are their actual scalar or homogeneous-linear direct contexts.

Every term is safe on the full graph. A nonconstant context is an old form: substituting any variable of its selected low linear form gives an $H$-neighbor of every context variable. By Proposition~\ref{prop:oldfactors}, a scalar context permits at most one missing witness, and a linear context at most two. The four missing witnesses must therefore split into two disjoint pairs, each witnessed by a term with four factors. Equality in that proposition applies on the full graph: each term's four full factor supports are exactly the four transversals of its two missing edges. Consequently the full supports $U=\supp(L)$ and $V=\supp(M)$ are disjoint two-element transversals contained in $T$. In particular no context vertex in $C$, or elsewhere outside $T$, has been silently discarded.

Let $A$ be the first cut's selected low form. It is linear and nonzero, and $AL$ is supported on $H$. Since $U$ has two vertices, four-cycle-freeness permits at most one common $H$-neighbor of $U$, so $A$ is supported on a single vertex $h\in C$. Similarly the other low is supported on $k\in C$. A leaf-only use of the first cut in a retained product would put the opposite star center in its actual context, contrary to $U\cap C=\varnothing$. Hence that cut reaches the literal center input of the unique assigned $h$-star. Its full assigned leaf set $S_h$ is contained in $U$. Likewise $S_k\subseteq V$. Every $H$-edge from $h$ to outside $C$ is assigned to its own star, and similarly for $k$.

Both assigned leaf sets are nonempty. The centers are distinct, since otherwise their common nonempty assigned leaf set would lie in $U\cap V$. An old edge $hk$ is impossible: whichever endpoint owns that edge would have the other center in its assigned leaf set, contradicting its containment in $T$. Thus
\[
 C'=(C\setminus\{h,k\})\cup U\cup V
\]
is a cover of the full graph. Indeed old edges not incident with the removed centers remain covered; their outside neighbors lie in the added leaf sets; edges to other centers remain covered there; and the possible edge between the two removed centers was excluded. Finally $U,V$ meet all four missing edges. Hence $|C'|=|C|+2$, proving the claim. This exchange includes old edges outside the matching and all edges among the original centers.
\end{proof}

\begin{corollary}[The three-product, two-addition profile at cover gain three]\label{cor:gainthreeprofile}
Let $G$ be finite and four-cycle-free, and let a positive variable-input circuit compute $P\in\Sup(G)$. Suppose exactly four $G$-edges are missing from the nonempty quadratic support $H$ and $\vc(G)-\vc(H)\geq3$. With arbitrary retained counts, extraction cannot delete at most three products and at most two additions. If a below-star circuit had $d=3,D=5,e=1$, its deleted counts would therefore be two products and three additions.
\end{corollary}
\begin{proof}
Choose any minimum $H$-cover $C$, independently of the extracted product centers, and put $N=G[V\setminus C]$. Every outside edge is missing, so $|E(N)|\leq4$, while adjoining an $N$-cover to $C$ gives $\vc(N)\geq3$. The component enumeration in Proposition~\ref{prop:gainthreesix} leaves either an exceptional graph $P_4\sqcup K_2$ or $K_3\sqcup K_2$, or an induced three-edge matching. Proposition~\ref{prop:exceptionalarbitrary} excludes the former cases without a retained-count hypothesis. In the latter, choose the six matching endpoints $T$. A fourth missing edge exists and cannot have both endpoints in $T$. Propositions~\ref{prop:adjacentfourth} and~\ref{prop:disjointfourth} exclude its meeting $T$ once or avoiding it. This proves the first assertion.

For $D=5$ and cover gain three, Proposition~\ref{prop:oneproductcover} excludes at most one useful private product. With none, full-carrier erasure leaves no source of a higher missing witness; with one, the gain is at most two. No deleted addition would leave no unequal-degree cut, while exactly one requires at least $2^4-2=14$ private products by Proposition~\ref{prop:oneadditionall}. Thus the only possible budgets were $(p,a)=(2,3)$ or $(3,2)$, and the first assertion removes the latter. The retained-count requirement in Corollary~\ref{cor:fourresidualfinal} is unchanged.
\end{proof}

\begin{proposition}[Five deletions at cover gain three]\label{prop:gainthreefiveprivate}
Let $G$ be finite and four-cycle-free, and let a positive variable-input circuit compute $P\in\Sup(G)$ with nonempty quadratic support $H$ and exactly four missing quadratic edges. If
\[
 K\le\vc(H)+1,\qquad \vc(G)-\vc(H)\ge3,
\]
where $K$ counts retained products, canonical extraction cannot delete at most two products and at most three additions. Retained addition counts are unrestricted.
\end{proposition}
\begin{proof}
Let $p,a$ count the deleted products and additions. Lemma~\ref{lem:fivelost} gives $p+a\ge5$, so the assumed budgets force $p=2,a=3$. In particular any exact-frame carrier erasure or Boolean absorption that removed a private gate would contradict that lemma. Apply the full-carrier reduction, retaining the exact selected frame, its $K$ products, and all five private gates. The resulting circuit and its preserved extraction are considered together; visibility and operand minimum degrees refer to that circuit.

Choose any minimum $H$-cover $C$ and put $T=V(G)\setminus C$, retaining isolates. Every edge of $N=G[T]$ is missing, so $|E(N)|\le4$. Adjoining an $N$-cover to $C$ shows $\vc(N)\ge3$. The small-graph classification in Proposition~\ref{prop:gainthreesix} gives an induced three-edge matching, or $P_4\sqcup K_2$ or $K_3\sqcup K_2$. Corollary~\ref{cor:threedeletedproducts} rules out the matching with only two private products. Thus $N$ is one of the two exceptional graphs. Its four edges are all the global missing edges, and it has six maximal independent sets.

At least one private addition is visible: a path carrying a missing witness must first enter the selected subgraph through the high branch of a selected unequal-degree sum, since selected products and equal-degree additions retain both inputs. With exactly one visible cut, Proposition~\ref{prop:onevisibleclauses} bounds the number of maximal independent sets of $N$ by $2^2+1=5$, a contradiction. With three visible cuts, the two private products are either nonancestral or ancestrally related in this same preserved-frame circuit. Propositions~\ref{prop:twooldthreecuts} and~\ref{prop:twodependentthreecuts} exclude those two possibilities.

There are consequently two visible cuts and one hidden addition. If its operand minimum degrees are unequal, Proposition~\ref{prop:hiddenunequal} gives cover gain at most two, contrary to the hypothesis. If they are equal, Proposition~\ref{prop:hiddenhomogeneous} excludes the exceptional target without a retained-count hypothesis. Equal minimum degrees allow affected operands with higher terms, so no case has been discarded by calling their full values homogeneous. The two alternatives exhaust the hidden addition and prove the claim.
\end{proof}

\begin{lemma}[Three conjunctions and two disjunctions]\label{lem:threeandtwoor}
Let $F$ be a single-output monotone Boolean circuit with at most three binary conjunctions and two binary disjunctions. Its boundary inputs are positive linear clauses, with constants zero and one allowed. Internal fanout and repeated inputs are unrestricted. Then $F$ has a positive linear CNF with at most six clauses, or it is a disjunction of at most three products, each using at most two of the original boundary clauses. The alternatives need not be exclusive.
\end{lemma}
\begin{proof}
We first regard the boundary inputs as independent formal atoms. All identities below are lattice identities, so the result remains valid after substituting the actual clauses, even if their supports overlap. Propagate constants and discard unused nodes.

\emph{One disjunction.} A circuit with at most one disjunction and $t$ conjunctions has Boolean form $C(U\vee V)$, where $C,U,V$ are products of boundary atoms. Remove repetitions, remove from the two arms every atom already in $C$, and move their common atoms into $C$. The remaining three atom sets are disjoint. If both arms are nonempty and the union of all three atom sets has $b$ atoms, the corresponding formula uses $b-2$ conjunctions. A connected binary DAG with $n$ gates has at most $n+1$ distinct boundary vertices: its underlying graph has at least $n+b-1$ edges and at most $2n$ operand edges. Here $n\le t+1$, so $b-2\le t$. Thus the formula uses no more of either gate type, and each atom occurs at most once.

If an arm is empty, absorption leaves a product of distinct atoms. A product of $b$ independent atoms requires at least $b-1$ conjunctions in the original circuit: choose an input at each disjunction to obtain a nonzero conjunctive parse; its monomial must contain every required atom, and a connected conjunction DAG with $t$ gates has at most $t+1$ boundary atoms. This also covers the case of no disjunction. Consequently the same formula conclusion holds in every degeneration.

\emph{Four-gate sharing.} A circuit with at most two gates of each type can be made a formula without increasing either count. Bypass identical inputs. If an internal node $s$ is shared, follow two paths to their first reconvergence. With at most four gates, an immediate triangle has the form $s\mathbin{\circ}(s\mathbin{*}a)$ and is removed by idempotence or absorption. Without a triangle, the remaining possibilities are a shortcut
\[
s\mathbin{\circ}\bigl((s\mathbin{*}a)\mathbin{\#}b\bigr)
\quad\hbox{or a diamond}\quad
(s\mathbin{*}a)\mathbin{\circ}(s\mathbin{\#}b).
\]
For an outer conjunction, the shortcut reductions are
\[
s\wedge((s\wedge a)\vee b)=s\wedge(a\vee b),\qquad
s\wedge((s\vee a)\wedge b)=s\wedge b,
\]
and $s\wedge((s\vee a)\vee b)=s$. The all-conjunction shortcut exceeds the stipulated conjunction budget. In the diamond, two inner disjunctions give $s\vee(a\wedge b)$, while mixed inner operations leave the inner conjunction term. The outer-disjunction cases are the lattice duals. No replacement increases either gate count. In a genuine four-node shortcut or diamond, $s$ is the only other internal node and $a,b$ are boundary atoms. A larger ancestor cone leaves fewer path nodes and therefore gives a triangle. Repeating these reductions removes all internal sharing.

We also need to control repeated occurrences of a \emph{designated boundary atom} $s$ in the resulting four-gate formula. They can be reduced to at most one without increasing either count. First reduce a proper repeated subtree; it has at most three gates. An arm equal to $s$ is handled by
\[
s\wedge H(s)=s\wedge H(1),\qquad
s\vee H(s)=s\vee H(0),
\]
followed by constant propagation. The remaining four-leaf shape is a diamond and is covered by the preceding identities. After these reductions, a repeated $s$ occurs once in each root arm. There are at most five leaves. Unless an arm is just $s$, the only split is two leaves versus three. If the three-leaf arm is $s\mathbin{*}(b\mathbin{\#}c)$, apply the diamond identities with $b\mathbin{\#}c$ as a single subexpression. Otherwise the shape is
\[
(s\mathbin{\circ}a)\mathbin{*}
       ((s\mathbin{\#}b)\mathbin{@}c).
\]
If $@=*$, associate the root operation and use the diamond reduction. If they differ, duality permits $*=\wedge$ and $@=\vee$. The remaining operations must be one conjunction and one disjunction. One choice absorbs to $s\wedge a$; the other is exactly
\[
(s\vee a)\wedge((s\wedge b)\vee c)
   =(s\wedge(b\vee c))\vee(a\wedge c).
\]
Both sides use two gates of each type, and the right side uses $s$ once. This exhausts the five-leaf case. Repetitions among the other atoms are harmless.

\emph{Five gates.} In the original circuit choose the earliest shared internal node $s$ in a topological order. Its proper ancestor cone has no external use, for that would be an earlier shared node; the cone is therefore already a formula. Contract it to the atom $s$. At most four gates remain. If the gate producing $s$ is a disjunction, the context has at most one disjunction, so the one-disjunction argument produces a formula using $s$ at most once. If it is a conjunction, the context has at most two gates of each type; the two four-gate arguments give the same conclusion. Substitute the original ancestor formula once. The entire circuit is now a formula with no increase in either gate count. No internal subcircuit has been duplicated.

\emph{Formula counts.} With at most one disjunction, write the formula as $C(U\vee V)$. If $c,u,v$ count its factor occurrences, then $c+u+v\le5$, and its CNF has at most $c+uv\le6$ clauses when both arms are nonempty. Empty arms or no disjunction give fewer clauses.

With two disjunctions there are at most six leaf occurrences. If the disjunctions are incomparable, the formula has form
\[
C(U\vee V)(A\vee B).
\]
All four arms are nonempty after constant propagation. Its CNF has $c+uv+ab$ clauses, with $c+u+v+a+b\le6$. This is at most five: when $c=0$ the largest count comes from arm lengths $2,2,1,1$; for $c=1$ or $2$ the count is at most four.

If the disjunctions are nested, the form is
\[
C\bigl(A(U\vee V)\vee B\bigr),\qquad
c+a+u+v+b\le6,\quad u,v,b\ge1,\quad a,c\ge0.
\]
It has at most $c+b(a+uv)$ clauses. If $a\ge1$, the possible values $u+v=2,3,4$ give maxima $6,6,5$, respectively; larger values are impossible. If $a=0$ and $c\ge1$, the maximum is five, since at $c=1$ the three positive arm lengths sum to at most five and their product is at most four. The cases $c\ge2$ are smaller. Finally, if $a=c=0$, the formula is $U\vee V\vee B$. With $u+v+b\le6$, the product $uvb$ exceeds six only when $u=v=b=2$. That is precisely the asserted disjunction of three products of two original boundary atoms. Substituting the original positive linear clauses preserves every CNF bound and the original-boundary description of this exception. Constants and repeated leaves only simplify the result.
\end{proof}

\begin{proposition}[One visible cut on a path and two separate edges]\label{prop:onevisiblepath}
Let $G$ be finite and four-cycle-free, and let a positive variable-input circuit compute $P\in\Sup(G)$ with nonempty quadratic support $H$. Suppose $H[T]$ is empty and $G[T]$, apart from isolated vertices, is $P_3\sqcup2K_2$. Canonical extraction cannot delete at most three products and three additions while exactly one deleted addition is a visible cut. Retained counts, cover gain, the original degrees of hidden additions, squares and fanout are unrestricted.
\end{proposition}
\begin{proof}
Apply exact full-current-carrier erasure with the original selected frame and minimum-degree labels fixed. Every useful private gate feeds the unique visible cut's high $F$: its first entry into the selected output must be through that cut, and a return from that cut to its high would be a cycle. Off-frame retained quadratic carriers have been erased. A retained linear boundary containing a private contribution would already follow the visible cut and therefore could not feed its high. Thus $F$ is computed from old homogeneous linear boundaries by at most three products and two hidden additions, regardless of the hidden additions' original degree types.

The full Boolean output is $P_2\vee FD$. If the cut keeps degree two, $D$ is a positive scalar; if it keeps degree one, $D$ is one nonzero old homogeneous linear form, including all fanout contexts. A squared insertion is impossible because the nonzero selected old low would enter both sides of a retained product and create a square in $P_2$. These are the original actual-context identities, not an extraction of a simplified Boolean expression.

By Lemma~\ref{lem:threeandtwoor}, either $F$ has at most six positive linear CNF clauses, or it is a union of at most three products of at most two old boundary forms. In the first case $FD$ has at most seven clauses. Set every variable outside $T$ to zero. The quadratic part vanishes and the result is exactly $f_{G[T]}$, whose eight maximal independent sets arise from the two choices on the path and two on each separate edge, with every isolate included. Each maximal independent assignment must falsify a different valid clause: one clause false on two such sets would also be false on their union, which contains an edge. Eight clauses are therefore necessary, a contradiction.

In the second case, multiply each of the at most three products by $D$ before restricting variables. These are globally safe products of at most three old homogeneous linear forms; a scalar context contributes no factor. Keep the original ambient pair $H\subseteq G$ throughout, including old-neighborhood centers outside $T$. Proposition~\ref{prop:oldfactors} gives at most one missing-edge witness per nonzero product. A zero term contributes none, and a nonzero constant or linear unsafe term cannot appear in the positive output. The three products therefore cannot supply all four edges of $G[T]$, each absent from $H$. This proves the assertion.
\end{proof}


\begin{proposition}[Three nonancestral products on a path and two edges]\label{prop:paththreevisible}
Let $G$ be finite and four-cycle-free, and let a positive variable-input circuit compute $P\in\Sup(G)$ with nonempty quadratic support $H$ and exactly four missing edges. Suppose $H[T]$ is empty and $G[T]$, apart from isolates, is $P_3\sqcup2K_2$. If extraction deletes at most three products and three additions, every deleted addition is visible, and the deleted products are pairwise nonancestral in the original circuit, then
\[
\vc(G)\le K+2,
\]
where $K$ is the retained product count. Retained additions, fanout, repeated centers and positive coefficients are unrestricted. Thus the minimum-product frame cannot have cover gain at least three in this profile.
\end{proposition}
\begin{proof}
All four edges of the induced target are the global missing edges; write their graph as $J$. A minimum $J$-cover consists of the path center and one endpoint of each separate edge. Every target vertex except the two path leaves therefore belongs to a three-vertex $J$-cover. We call these five vertices \emph{cover vertices} within this proof.

\emph{The original frame.} Apply exact full-current-carrier erasure. It cannot create product ancestry. Every surviving private product has two old homogeneous linear operands: an off-frame retained quadratic operand was erased; an affected linear operand would carry an earlier private product and create forbidden ancestry; and a private linear preparation would be hidden. The induced three-edge matching requires three useful private products by Corollary~\ref{cor:threedeletedproducts}. Proposition~\ref{prop:threeproductstwosums} then forces all three additions to remain useful. A hidden sum cannot combine highs, and no product can feed another private product. Each product must enter a visible cut. Thus there are exactly three products
\[
g_i=U_iV_i\quad(1\le i\le3),
\]
and the three cuts insert them once each. Each high has minimum degree two and each selected old low $a_i$ is nonzero and linear.

Introduce independent insertion symbols before substituting the highs. The full Boolean output has form
\[
P_2\vee\bigvee_i g_iL_i\vee
       \bigvee_{i<j:q_{ij}>0}g_ig_j,
\]
where every $L_i$ is a full nonzero old homogeneous linear direct context. An individual insertion cannot enter both operands of a retained product, since its selected old low would create a square in $P_2$. The actual old identities give
\[
a_iL_i\text{ supported on }H,
\qquad q_{ij}>0\Longrightarrow
\left\{
\begin{aligned}
&a_ia_j\text{ supported on }H,\\
&\supp(a_i)\subseteq\supp(L_j),\quad
  \supp(a_j)\subseteq\supp(L_i).
\end{aligned}\right.
\tag{*}
\]
Every displayed term is safe. Each direct cubic $D_i=g_iL_i$ has at most one missing witness. Each positive mixed quartic $Q_{ij}=g_ig_j$ has at most two; equality gives the full four-transversal description of Proposition~\ref{prop:oldfactors}.

\emph{A productive index.} Call $i$ productive if $D_i$ witnesses an edge $d$ and some positive $Q_{ij}$ witnesses a different edge $e$. Then the full support of $L_i$ is a singleton endpoint of $d$. If neither $U_i$ nor $V_i$ is singleton, this follows directly from the safe-old-cubic singleton-factor argument. Otherwise a singleton operand forces $d,e$ to share its vertex. The only adjacent missing edges are the two path arms, say $ba,bc$, with center $b$. Every old form meeting both must contain $b$: otherwise it contains $a,c$, and its old neighborhood center either makes the arms old or completes a four-cycle through $b$. Hence both factors of $g_i$ contain $b$, so $g_i$ is true on the singleton $b$. Safety makes every member of $L_i$ a $G$-neighbor of $b$. An old neighborhood containing a missing neighbor of $b$ cannot contain any other $G$-neighbor of $b$, by the same four-cycle argument. Thus $L_i$ is precisely the singleton claimed.

If $L_i=\{v\}$ and $v$ is a cover vertex, choose $v$ to cover the old edges of one original retained product marked by this cut: its entire opposite old operand is the nonempty singleton $v$. One center for each other retained product gives an $H$-cover with at most $K$ vertices containing $v$. Complete $v$ to a three-vertex $J$-cover by adding two vertices. This proves the desired bound. It remains to rule out four witnesses when every productive index has its singleton context at a path leaf.

\emph{Mixed terms with at most one witness.} Suppose every $Q_{ij}$ has at most one missing witness. If no index is productive and $m$ direct terms have witnesses, a mixed term incident with such an index cannot add a different edge. Only pairs between the other $3-m$ indices can contribute new edges. The total is at most
\[
m+\binom{3-m}{2}\le3\qquad(0\le m\le3).
\]
Therefore some index $i$ is productive. Its singleton context is a path leaf, say $L_i=\{v\}$, unless the cover bound is already proved. Choose $j$ with $q_{ij}>0$. Formula~(*) gives $a_j=\{v\}$ as a support. Since $a_jL_j$ is old and $H[T]$ is empty, $L_j$ contains no vertex of $T$. Thus $D_j$ witnesses no missing edge.

Let $k$ be the third index. If $q_{ik}>0$, the same reasoning eliminates the direct witness of $D_k$. Moreover $q_{jk}=0$, for otherwise $a_ja_k$ contains the square $v^2$ in the old quadratic output. At most $D_i,Q_{ij},Q_{ik}$ contribute, so there are at most three edges. Consequently four witnesses require $q_{ik}=0$. The only possible contributions are then $D_i,D_k,Q_{ij},Q_{jk}$, each with one witness. Four distinct edges force all four contributions to be distinct, so $k$ is also productive. Its context must again be a path-leaf singleton $L_k=\{w\}$ unless it gives the cover bound. But $q_{jk}>0$ and (*) give $a_j=\{v\}\subseteq L_k$, hence $w=v$. Both direct terms would witness the unique missing edge at that same path leaf, a contradiction.

\emph{A mixed term with two witnesses.} Suppose $Q_{12}$ witnesses $ab,cd$. Its four full old factor supports are the four transversals of this induced matching. If the two factors of $g_1$ are complementary transversals, those of $g_2$ are also complementary. Neither signal can be true on an edge involving a vertex outside the quartet, and within it the only $G$-edges are $ab,cd$. Thus every term involving $g_1$ or $g_2$ is confined to those two edges; $D_3$ adds at most one, insufficient.

Otherwise each pair consists of adjacent transversals. Relabeling gives the global support identities
\[
g_1=(a\vee c)(a\vee d)=a\vee cd,
\qquad
g_2=(b\vee c)(b\vee d)=b\vee cd.
\]
An additional edge available to $g_1$ must be $ax$ with $x$ outside the quartet; for $g_2$ it must be $bx$. At least one of the other two $J$-edges lies wholly outside the quartet. If the matching chose the two separate edges, both remaining path arms do; if it chose one path arm and a separate edge, the other separate edge does. Every wholly outside edge must be supplied by $D_3$, which has at most one witness. Therefore exactly one such edge $e$ is possible, and $D_3$ witnesses it. The last edge is the other path arm, say $ax$ after exchanging indices. It must come from $D_1$ or $Q_{13}$.

If $Q_{13}$ witnesses $ax$, index $3$ is productive with direct witness $e$. Its singleton context is an endpoint of the separate edge $e$, hence a cover vertex, giving the required cover. Otherwise $D_1$ witnesses $ax$. Since $g_1$ is true on singleton $a$, safety and the old-neighborhood argument give the full equality $L_1=\{x\}$. The positive $q_{12}$ forces $a_2=\{x\}$, and therefore $xL_2$ is old. The nonzero $L_2$ supplies a vertex $z$. Since $g_2$ is true on singleton $b$, safety gives $bz\in E(G)$. The old edge $xz$ excludes $z=x$ and $z=a$, because $ax$ is missing; safety excludes $z=b$. Thus $x,a,b,z$ are distinct and $x-a-b-z-x$ is a forbidden four-cycle.

The cases exhaust all mixed terms. In every realizable case the original product count supplies an $H$-cover containing a vertex of a three-vertex $J$-cover, and adding its other two vertices gives $\vc(G)\le K+2$. Full supports, all original products and all old edges were retained in the charge; no safe-center replacement for this outside graph was assumed.
\end{proof}


\begin{proposition}[An old quartic high and two boundary factors]\label{prop:quartictwoboundary}
Keep the original two-visible-cut frame and the graph hypotheses of Proposition~\ref{prop:twoaffectedcover}, with exactly four missing edges and $K$ retained products. There are exactly two deleted additions, both visible. Suppose $g$ is Boolean-equivalent to a product of at most four nonzero old homogeneous linear supports, and the second high has Boolean form
\[
f=(A\vee\epsilon g)(B\vee\delta g),\qquad
\epsilon,\delta\in\{0,1\},
\]
where $A,B$ are nonzero old homogeneous linear forms. Each affected input is an actual boundary carrying the first cut. In particular, for positive mixed coefficient $q$, the nonzero old linear first low $a_0$ satisfies $\supp(a_0)\subseteq\supp(M)$ and $a_0L$ is supported on $H$. Then $\vc(G)\le K+2$, without retained optimality or a retained-addition restriction. Repetition may pad the Boolean representation of $g$; its original arithmetic degree and the actual context types remain fixed. A direct proper-private-ancestor input of the second high is outside this statement.
\end{proposition}
\begin{proof}
A nonzero variable-input high is false at the all-zero assignment, so its old-support representation has at least one factor. Repeat factors to write it Booleanly as a product of four old forms. This changes neither the original frame nor any arithmetic degree. A safe product is considered through its supports throughout the proof.

We recall two support observations, including their full-support qualifications. If $xy$ is missing and an old form contains $y$, then $x$ is adjacent to no other member of that form: a second neighbor $z$ and the old center would give a four-cycle, while a center equal to $x$ would make $xy$ old. Now let a safe product of $k$ old factor occurrences witness two missing edges. They are an induced matching with endpoint set $E$. Each of the four independent endpoint transversals must be rejected by a factor, and a factor meeting both edges can reject at most one. A uniquely rejecting occurrence has full support equal to the complementary two-point transversal. Indeed, an extra vertex $z$ can be selected from that factor and a vertex of the rejected transversal from every other factor. The old-neighbor observation makes $z$ nonadjacent to both selected transversal vertices, producing an unsafe monomial. If $u$ rejection classes have a unique occurrence, then
\[
k\ge u+2(4-u)=8-u.
\]
Thus a quintic has at least three unique full transversals and a sextic at least two. Repeated factors are counted as occurrences; an idle factor only strengthens the bound.

In particular, if safe $gD$, with $D$ old linear or a positive scalar, witnesses two missing edges, every $G$-edge satisfying $g$ is confined to that pair. For scalar $D$, use quartic equality. For linear $D$, at least three unique full transversals occur in the quintic. If three belong to $g$, their common intersection is empty and confinement follows. Otherwise $D$ and exactly two $g$-factors are the three unique classes; the other two $g$-factors reject the fourth class. If the two full $g$-supports are disjoint, no edge with an outside endpoint meets both. If they intersect at $x$, an additional edge would be $xz$ outside the quartet. The two remaining $g$-factors contain the matching mate $x'$ and omit $x$. To meet $xz$ each would contain $z$, contrary to the old-neighbor observation for the missing edge $xx'$. This proves confinement for full ambient supports and repeated factors.

Write $J=E(G)\setminus E(H)$. If $\vc(J)\le2$, add a two-vertex $J$-cover to one center for every retained product. Hence assume $\vc(J)\ge3$. Theorem~\ref{thm:twosumscomponents} says that every nontrivial component of $J$ is $K_2$ or $P_3$. We will use the following consequence: given two disjoint matching edges of $J$, its other two distinct edges cannot both meet their endpoint quartet. Two attachments on one component, or attachments joining the components, would violate the component restriction. Extending separate components gives two copies of $P_3$, whose cover number is only two.

Put $P'=gL$ and $Q=ABM$; $Q$ has at most one missing witness. Expand the high and absorb positive multiples Booleanly, with the actual frame fixed. With no affected input the terms are $P',Q$, and also $gAB$ when $q>0$. With one affected input, say $f=(A\vee g)B$, they are $P',Q,gBM$ when $q=0$, and $P',Q,gB$ when $q>0$. With both inputs affected they are $P',Q,gM$ when $q=0$, and $Q,g$ when $q>0$.

If $P'$ has two witnesses, the quintic confinement just proved confines every other $g$-multiple to that pair, leaving at most one further witness from $Q$. The same argument applies to any other $g$-multiple with at most one additional old factor. Consequently the only possible four-witness cases are
\[
\begin{array}{lll}
\mathrm{I}:& q=0,\ \hbox{one affected input},
   & P'=gL,\quad Q=ABM,\quad W=gBM;\\
\mathrm{II}:& q>0,\ \hbox{no affected input},
   & P'=gL,\quad Q=ABM,\quad W=gAB.
\end{array}
\]
Here $W$ has two witnesses, and $P',Q$ each supply one distinct new edge, denoted $e_P,e_Q$. A zero $L$ removes its term; a positive scalar $L$ makes $g$ itself safe and absorbing, leaving at most three witnesses. Thus $L$ is nonzero old linear. In case I a scalar or zero $M$ leaves at most a quintic $g$-multiple and is already excluded. In case II the actual positive-cross inclusion of the nonzero linear low forces $M$ to be nonzero old linear. These conclusions concern original context types, not padded degrees.

Let $E$ be the endpoint quartet of $W$'s two missing witnesses. At least two of its six factor occurrences are unique full transversals.

\emph{Case I.} If $M$ is a full two-point transversal, its support together with one endpoint of $e_P$ is a set $S$ of at most three vertices covering all missing edges. The actual second cut marks an original retained product whose entire opposite old operand lies in $M$. Thus $S$ also covers all old edges of that product. One center for each of the other $K-1$ products gives $\vc(G)\le K-1+3=K+2$.

Suppose instead that $M$ is not a full two-point transversal. It cannot be a unique rejector. If $B$ is unique, at least one $g$-factor is unique too. Then $e_P$ meets $E$ through the full $g$-factor and $e_Q$ meets $E$ through $B$, contradicting the preceding missing-component consequence. If $B$ is not unique, at least two $g$-factors are unique. Three full transversals confine $e_P$ to the old pair. Two disjoint ones also prevent a new edge. The only possibility is two intersecting full supports, with common endpoint $x$, and $e_P=xz$ for $z\notin E$. Each of the two other rejecting classes has a complementary support containing $x'$ and omitting $x$. No $g$-factor can occupy either class, since it would need to contain $z$ to meet $e_P$, forbidden by old-neighbor exclusion. Those classes must therefore be rejected by $B$ and $M$ alone. Both would be unique full transversals, contradicting the assumption on $M$. This completes case I.

\emph{Case II.} If a unique full factor lies in $g$ and another in $A$ or $B$, then both $e_P,e_Q$ meet $E$, again impossible. If at least two unique classes lie in $g$ and neither $A$ nor $B$ is unique, the preceding argument applies: three classes or two disjoint classes confine $e_P$; two intersecting classes force the remaining two rejecting classes outside $g$, making $A,B$ unique after all. Exactly one unique $g$-class with neither $A$ nor $B$ unique is incompatible with the requirement of at least two unique classes.

It remains that $A,B$ are the two unique full transversals and no $g$-factor is unique. The four $g$-occurrences must occupy the remaining two classes twice each, leaving no idle factor. If $A,B$ are disjoint, $Q$ cannot have a new witness. Otherwise they intersect at $x$, and the two other complementary transversals both contain its matching mate $x'$ and omit $x$. Every $g$-factor therefore contains $x'$, so $g$ is true on singleton $x'$. The new edge of $Q$ is $xz$ with $z\notin E$. Since $AB$ is true on singleton $x$, safety makes every vertex of $M$ a $G$-neighbor of $x$. The missing neighbor $z$ belongs to $M$, so old-neighbor exclusion gives the full equality $M=\{z\}$.

The actual positive-cross inclusion now gives $\supp(a_0)=\{z\}$, hence $zL$ is supported on $H$. Safety of $gL$ makes every vertex $l$ of the nonzero $L$ a $G$-neighbor of $x'$. The old edge $zl$ excludes $l=z$ and $l=x$, since $xz$ is missing; safety excludes $l=x'$. The four distinct vertices $x,z,l,x'$ form the forbidden cycle $x-z-l-x'-x$. Case II is impossible.

The sole realizable alternative supplied the set $S$ covering every missing edge and all old edges of a marked retained product. Every other retained product is a star, independently of coefficients or addition sharing. The resulting original-product charge proves the asserted bound without an optimality assumption. Proper private-ancestor operands have not been normalized by this proof.
\end{proof}


\begin{proposition}[Raw divisors of an old four-factor high]\label{prop:rawdivisor}
Let $H\subseteq G$ be finite simple graphs with $G$ four-cycle-free. Let $g$ be a product of at most four distinct old homogeneous linear supports, and let $u$ be a nonempty product of at most three of those same supports, so that $g$ is a Boolean multiple of $u$. Let $A$ be old nonzero linear, let $L,M$ be old linear forms or nonnegative scalars, and let $q\ge0$. For either
\[
 f=uA\quad\hbox{or}\quad f=uA\vee g,
\]
assume $gL\vee fM\vee qgf$ is safe. Any four distinct missing-edge witnesses of this expression have a two-vertex cover. The same conclusion holds for $f=gA$, $f=g$, or $f$ a product of two Boolean old-product divisors of $g$. Consequently an actual two-cut frame of this form with four missing edges satisfies $\vc(G)\le K+2$, for arbitrary retained product count $K$. These are support templates, without a classification assertion for arbitrary polynomial ancestors.
\end{proposition}
\begin{proof}
We first record a confinement fact. If safe $XD$ has two missing witnesses, $X$ uses at most four distinct old supports and $D$ is old linear or a positive scalar, every $G$-edge meeting every factor of $X$ is one of those two edges. The witnesses form an induced matching. Among its four independent transversals, a factor rejects at most one; a unique rejector is its full complementary two-point transversal, by the old-center four-cycle argument. With at most five supports at least three rejection classes are unique. Removing $D$ leaves three full transversals, which confine every $G$-edge, except possibly when the multiplicities are $1,1,1,2$ and $D$ is a unique factor. Then $X$ contains two full unique supports and both factors in the duplicated class. Disjoint full supports already confine. If they intersect at $x$, an outside edge must be $xz$. Both duplicated factors contain the matching mate $x'$ and omit $x$; to meet $xz$ they must contain $z$. An old factor containing $x',z$ cannot coexist with missing $xx'$ and edge $xz$: its old center gives a four-cycle, or makes $xx'$ old. This proves confinement. Repetitions, idle factors and a scalar $D$ only strengthen the count or give quartic equality.

For $f=uA$ and $q=0$, the two terms $P'=gL$ and $Q=uAM$ use at most seven distinct old supports, so Lemma~\ref{lem:sevenoldfactors} applies. For $q>0$, the third term is $R=gA$. Each term has at most five old supports and at most two missing witnesses. Two witnesses of $P'$ confine $R$ by the preceding fact with prefix $g$; two of $Q$ confine $R$ with prefix $uA$. Thus either four witnesses come from $P'\vee Q$, or both direct terms have at most one. In the latter case two witnesses of $R$ confine $P'$, leaving at most three altogether. Hence every four-witness set is supplied by the seven-support union $P'\vee Q$.

For $f=uA\vee g$ and $q>0$, absorption gives $g\vee uAM$, with at most six supports. For $q=0$, the terms are $P'=gL$, $Q=uAM$, $R=gM$. Two witnesses of $P'$ confine $R$ with prefix $g$; two of $Q$ confine $R$ with prefix $uM$ and final factor $A$. If both $P',Q$ have at most one, two of $R$ confine $P'$ and again leave at most three. The seven-support lemma therefore applies in every four-witness case.

For $f=gA$, the two terms at $q=0$ use at most seven supports, while positive $q$ absorbs to $gL\vee gA$, with at most six. The case $f=g$ is shorter. If $f$ is a product of two old-product divisors of $g$, then $gf=g$ Booleanly; the remaining two products use at most six supports when $q=0$ and at most five when $q>0$. Zero terms disappear and scalar factors are omitted throughout. Finally, one star center per original retained product covers $H$; adding the two missing-cover vertices proves the circuit consequence without changing any original gate or degree.
\end{proof}

\begin{proposition}[A raw quadratic ancestor and two boundaries]\label{prop:rawquadratic}
Let $G$ be finite, simple and four-cycle-free, let $P\in\Sup(G)$ have nonempty quadratic support $H$, and suppose exactly four edges are missing. Suppose extraction deletes exactly two additions, both visible, and keep their actual frame $P_2\vee gL\vee fM\vee qgf$. Suppose Booleanly $u=UV$, $g=uW$ and
\[
 f=u(A\vee\epsilon g)(B\vee\delta g),\qquad \epsilon,\delta\in\{0,1\},
\]
where $U,V,W,A,B$ are nonzero old homogeneous linear forms. Each affected boundary carries the actual first cut and contains its nonzero old linear low $a_0$; $a_0L$ is old, and positive $q$ also gives $\supp(a_0)\subseteq\supp(M)$. Contexts, occurrences and arithmetic degrees are the original ones. Then the missing graph $J$ has $\vc(J)\le2$, and $\vc(G)\le K+2$. No optimality or retained-addition hypothesis is required.
\end{proposition}
\begin{proof}
Suppose $\vc(J)\ge3$. By Theorem~\ref{thm:twosumscomponents}, $J$ is $4K_2$ or $P_3\sqcup2K_2$. Every missing witness meets both $U,V$, since all off-frame terms have factor $UV$. Thus both full supports have size at least three; otherwise either is a two-cover.

A safe product containing $U,V$ and at most two other old factors has at most one missing witness: two would force full quartic equality, contrary to their support sizes. If $UVFGH$ has two, the unique-rejector count from the preceding proof forces $F,G,H$ to be three distinct full two-point transversals of their induced matching. Neither $U$ nor $V$ can be unique, and both reject the fourth transversal $D$; their restrictions to the quartet are exactly $D$. Any $G$-edge meeting $UV$ and two of $F,G,H$ belongs to that matching. Indeed disjoint full factors confine immediately; if they meet at $x$, an outside edge $xz$ forces $U,V$ to contain $z$ together with $x$'s missing mate, which the old-neighbor four-cycle rule forbids.

We need one collision argument. Suppose $UVFGH$ has the two matching witnesses, $gL$ and $gF$ are safe, and the actual nonzero low satisfies $a_0\subseteq G\cap H$ and $a_0L$ is old. These three terms cannot supply four edges of $J$. Both latter terms have at most one witness, so they would supply distinct new edges. The $gF$ edge meets the matching quartet; the $gL$ edge is disjoint from it, since two extra edges meeting a matching quartet would contradict the stated shapes and cover number. A new $gF$ edge requires the full transversal $F$ to intersect $D$. Relabel
\[
 F=\{x,y\},\quad D=\{x,y'\},\quad
 \{G,H\}=\{\{x',y\},\{x',y'\}\},
\]
with matching $xx',yy'$. The new edge is $xz$, with $z$ outside the quartet: an edge through $y$ would force $U$ to contain $y'$ and its forbidden new neighbor. The low is supported on $x'$, so $x'L$ is old. Safety of $gF$ forbids $x\in W$ and its witness forces $z\in W$.

Write the other new edge $lt$, disjoint from the quartet and from $z$. The relation $x'L$ old and the old-neighbor rule in $U,V$, which contain $x$, imply $L\cap U=L\cap V=\varnothing$. Hence name its endpoints so that $l\in L$ and $t\in U\cap V$. The factor $W$ contains $l$ or $t$. If it contains $l$, choosing $x$ from $U,V$ and $l$ from $W,L$ gives the nonedge $xl$, since the old form $W$ contains $z,l$ and $xz$ is missing. Thus $t\in W$. Choosing $y'$ from $U,V$, $z$ from $W$ and $l$ from $L$ now gives the independent set $\{y',z,l\}$: $y'z$ is forbidden by the old rule using $xz$ in $U$, $y'l$ by using missing $tl$ in $U$, and $zl$ by using $tl$ in $W$. Its vertices are distinct by the endpoint restrictions. This unsafe term proves the collision assertion.

Put $P'=gL$ and $Q=uABM$. With no affected boundary, $q=0$ gives at most $1+2$ witnesses. For positive $q$, the third term is $gAB$; if either this term or $Q$ has two, its full $A,B$ transversals confine the other, while $P'$ adds at most one. With exactly one affected boundary, say $A$, we have $f=uAB\vee gB$. At $q=0$ the third term is $gBM$; two witnesses in it or $Q$ force full $B,M$ transversals and confine the other. At $q>0$ absorption leaves $P',Q,gB$. Four witnesses require two in $Q$ and one new witness in each other term. Apply the collision assertion with distinguished $F=B$ and common-low factors $A,M$.

With both boundaries affected, $f=uAB\vee g$. For $q=0$, the terms are $P',Q,gM$; the same collision assertion uses $F=M$ and common-low factors $A,B$. For $q>0$, safe $g$ absorbs all its multiples and $g\vee Q$ has at most three witnesses. In the collision cases, two witnesses of $Q$ force $M$ linear; $L=0$ removes a required witness and a scalar $L$ makes $g$ absorb the other term, so $L$ is genuinely linear. All repetitions and compatible scalar cases are therefore included. The contradiction proves $\vc(J)\le2$, and the original $K$ stars give the asserted cover of $G$.
\end{proof}

\begin{proposition}[Four boundaries after a quadratic high]\label{prop:quadraticfourboundary}
Let $G$ be finite, simple and four-cycle-free, let $P\in\Sup(G)$ have nonempty quadratic support $H$ and exactly four missing edges $J$. Suppose extraction deletes exactly two additions, both visible, and keep their original two-symbol frame. Suppose Booleanly $g=UV$ and
\[
 f=\prod_{i=1}^4(A_i\vee\epsilon_i g),\qquad \epsilon_i\in\{0,1\},
\]
with all displayed forms nonzero, old and homogeneous linear. Every affected boundary contains the actual first low $a_0$, with $a_0L$ old; positive $q$ gives $a_0\subseteq M$. All selected occurrences and original degrees remain fixed, including repetitions used only in a Boolean representation. Then $\vc(J)\le2$ and $\vc(G)\le K+2$, without a retained-count or private-product-count assumption. This does not assert an original-operand normalization.
\end{proposition}
\begin{proof}
Assume $\vc(J)\ge3$, so again $J$ is $4K_2$ or $P_3\sqcup2K_2$. We prove a common three-term obstruction. Let $X$ use at most two old factors, $Y$ at most four, and $D$ be an old product, with $YD$ using at most five factors. Suppose
\[
 XL,\qquad YD,\qquad XY
\]
are safe and supply all four edges, where $L$ is old linear and a nonzero old low $a$ belongs to every factor of $D$, with $aL$ old. The witness bounds are respectively one, two and two. Unique rejectors are full transversals, and a product with $n$ distinct supports witnessing an induced pair has at least $8-n$ unique rejection classes.

First suppose $YD$ has at most one witness. Then $XY$ supplies a matching pair and the two direct terms supply distinct extra edges. If both extras avoid the matching quartet, no factor can be a unique full rejector, requiring at least eight supports, whereas $XY$ has at most six. If the $XL$ extra avoids it and the $YD$ extra attaches at $v$, every $Y$-factor contains $v$. The two independent transversals containing $v$ must each have at least two rejecting $X$-factors: a unique one would be full and could not meet the outside $XL$ edge. This needs four $X$-supports, impossible.

In the remaining case the $XL$ edge is $vt$ attached to the quartet and the $YD$ edge avoids it. Every $X$-factor contains $v$. The two transversals containing $v$ require four distinct $Y$-factors, since none can be a unique full rejector while meeting the outside edge. These exhaust $Y$; each contains the matching mate $v'$. Safety of $XL$, with $X$ true on singleton $v$, and oldness of $L$ force the full equality $L=\{t\}$. Choose $w\in a$. Then $wt$ is old, and safety of $YD$, which is true on $\{v',w\}$, gives $v'w\in E(G)$. Safety excludes $w=v'$; the old edge excludes $w=t$ and $w=v$, the latter because $vt$ is missing. The four distinct vertices form the forbidden cycle $v-t-w-v'-v$.

Now suppose $YD$ has two witnesses. Its union with $XY$ uses at most seven supports, so Lemma~\ref{lem:sevenoldfactors} implies it supplies exactly three of the four edges; $XL$ supplies the fourth. At least three unique full transversals occur in $YD$. They cannot all lie in $D$, whose factors share the nonzero low $a$, since three distinct full transversals have empty intersection. One is therefore a $Y$-factor. The new $XY$ edge meets the matching quartet through this factor; write it $vt$, with original matching edge $vv'$. Every $Y$-factor meets both edges and hence contains $v$. The remaining $XL$ edge is disjoint from $vt$ and the quartet by the shapes of $J$. But $vX$ is a safe Boolean subterm of $XY$ and witnesses $vt$. The two safe cubics $vX$ and $XL$ share their two old factors, so their missing edges form a star by the common-cubic lemma, contradicting disjointness. Fewer factors can be repeated Booleanly; a unit $Y$ would already contradict the unique-support conclusion. This completes the common obstruction.

Let $k$ be the number of affected boundaries. Write $R=\prod A_i$, and for $k>0$ write $R=TB$, where $T$ contains the unaffected and $B$ the affected factors. Then $f=TB\vee gT$. If $q=0,k=0$, the two terms $gL,RM$ supply at most three witnesses. If $q=0,k>0$, use the common obstruction with
\[
 X=g,\quad Y=TM,\quad D=B,\quad a=a_0.
\]
Here $Y$ uses at most $5-k\le4$ factors and $YD$ at most five. If $q>0,k=0$, use $X=g,Y=R,D=M,a=a_0$; the actual mixed identity makes $M$ linear and contains $a_0$. If $q>0,k>0$, absorption leaves $gL,TBM,gT$; use $X=g,Y=T,D=BM,a=a_0$. Every factor of $D$ contains the actual low, including $M$. These cases exhaust all choices.

A zero term only removes witnesses. A positive scalar $L$ makes safe $g$ absorb all its multiples, leaving at most three witnesses with the quintic other term. Scalar factors in the other applications are omitted; the actual positive mixed identity retains the required linear context. Thus no scalar equality case is lost. All supports and old edges are ambient, and no Boolean factor padding changes an original gate. The contradiction gives the two-cover of $J$ and hence $\vc(G)\le K+2$.
\end{proof}

\begin{proposition}[One linear preparation and two unaffected highs]\label{prop:pathpurelinear}
Let $G$ be finite, simple and four-cycle-free, and let $P\in\Sup(G)$ have nonempty quadratic support $H$ and exactly four global missing edges. Let $T$ be their endpoint set, and assume $G[T]=P_3\sqcup2K_2$ and $H[T]$ is empty. Keep an actual two-visible-cut frame
\[
 P_2\vee gL\vee fM
\]
with zero mixed coefficient and full original contexts $L,M$. Suppose the joint networks computing $g,f$ use at most three multiplication gates, no additions, and old homogeneous linear boundaries together with at most one further homogeneous linear form $S$. Internal sharing, raw product ancestry and repeated factors are allowed. Then $\vc(G)\le K+2$ for arbitrary retained product count $K$. No normalization of an affected nonlinear operand is asserted.
\end{proposition}
\begin{proof}
Write the path as $uxv$ and the separate edges as $ab,cd$. On their endpoint set $T$, the quadratic part vanishes and the output is the target edge function. Its eight maximal independent sets require eight distinct positive CNF clauses: one clause false on two such sets would also be false on their union, which contains an edge.

If one high is ancestral to the other, write $f=gR$ after interchanging highs and contexts. With at most three joint products, $g$ and $R$ together use at most four distinct boundary atoms. The output $g(L\vee RM)$ has at most five linear clauses. This boundary count follows from the operand count of the connected multiplication DAG, including repeated inputs and fanout. If neither high is ancestral but they share a private product, the three-product budget forces $h=AB$, $g=hC$, $f=hD$; repeated-$h$ inputs only simplify. The output $AB(CL\vee DM)$ has at most six clauses. Both cases contradict eight.

The private cones are therefore disjoint. Relabel them so that the smaller uses at most one product. Its contribution has at most three linear factors, and the other at most five, or at most four when both highs are product roots. Every old factor meeting both missing path arms contains $x$: otherwise it contains $u,v$ and its old center gives a four-cycle. Consequently a safe product of old factors and repetitions of the single $S$ cannot witness both arms and a separate edge. All old factors would contain $x$; safety forces every variable of $S$ to be a $G$-neighbor of $x$, so none is an endpoint of a separate induced edge. A product without old factors or without $S$ would be true on a singleton and unsafe. Any other three-edge target subset is an induced three-edge matching, requiring eight clauses; at most five factors cannot witness it. Thus each contribution has at most two witnesses. The smaller, with at most three factors, cannot witness two disjoint edges, which require four clauses. Four total witnesses therefore force it to supply the path arms, and the other to supply $ab,cd$. The form $S$ is present in the path term and has no separate-edge endpoint, so it is absent from the other term.

When the smaller high uses one product, the other contribution has at most four old factors. Its two missing witnesses force full quartic equality, including its actual context $M$, whose full support is a two-point transversal of $ab,cd$. The set $\supp(M)\cup\{x\}$ covers all missing edges and every old edge of an original retained product marked by the second cut: that product's entire opposite operand lies in the full context $M$. One star center for each other retained product gives $(K-1)+3=K+2$ vertices.

For completeness, if a high is only a boundary output, the smaller contribution must be $SL$ and supply the path arms. The old context $L$ contains $x$. Safety makes every $l\in L$ adjacent to both $u,v$; another such vertex gives the four-cycle $x-u-l-v-x$, while $l=u,v$ gives a forbidden singleton. Hence the full context is $L=\{x\}$. It covers a marked original retained product and both path arms. Adding one endpoint of each separate edge and centers for the other $K-1$ products again costs at most $K+2$. Zero or scalar contexts only shorten clause counts or eliminate a required witness. All cover charges retain the original products and their ambient old edges.
\end{proof}

\begin{lemma}[A product followed by a quadratic union]\label{lem:pathproductunion}
Let $H\subseteq G$ be spanning, with $G$ finite, simple and four-cycle-free. Let $A,B,C,D$ be nonzero old linear forms, let $L,M$ be old linear forms or nonnegative scalars, and put $u=AB$, $v=CD$, $g=uv$, $f=u\vee v$. For $q\ge0$, a safe expression
\[
                 gL\vee fM\vee qgf
\]
cannot witness four missing edges forming a subgraph $P_3\sqcup2K_2$. The selected four edges need not exhaust the global missing edges or induce the entire graph on their endpoints. Consequently the same exclusion holds in an exact two-visible-cut frame with these high identities. This is a support-template statement, without an assertion that other hidden additions have this form.
\end{lemma}
\begin{proof}
Booleanly $gf=g$. For $q>0$ the expression is $g\vee uM\vee vM$, and for $q=0$ it is $gL\vee uM\vee vM$. Each displayed term is safe. If $M=0$, only the first term remains and witnesses at most two missing edges. If $M$ is a positive scalar, $f$ absorbs the other terms; its two safe old quadratic terms witness at most two in total. Hence $M$ is nonzero linear. The two old cubics $uM,vM$ each witness at most one missing edge, while the remaining term, with at most five old factors, witnesses at most two. Four selected witnesses therefore require exactly two from that remaining term and one new, distinct witness from each cubic.

Let $E$ be the endpoint quartet of the two witnesses in the remaining term. They are an induced matching in $G$ by the old-factor theorem. Of the other two edges in the selected $P_3\sqcup2K_2$, at least one avoids $E$: a disjoint pair consists either of the two isolated edges, leaving both path edges outside, or of one path edge and one isolated edge, leaving the other isolated edge outside.

For $q>0$, or for $q=0$ and positive scalar $L$, four-factor equality makes $A,B,C,D$ full two-point transversals of $E$. Both extra cubic witnesses then meet $E$, a contradiction. A zero $L$ cannot supply the required two witnesses when $q=0$.

It remains that $gL$ is an old quintic with two missing witnesses. At least three factor occurrences are unique rejectors of the four independent transversals of $E$: if $t$ classes have unique rejectors, then $5\ge t+2(4-t)$ gives $t\ge3$. Every uniquely rejecting old factor has full support equal to its complementary two-point transversal. At least two such factors lie among $A,B,C,D$. If one lies in each of the pairs $\{A,B\}$ and $\{C,D\}$, both extra witnesses meet $E$, again impossible. After interchanging the pairs, the remaining case has $A,B$ distinct full transversals. If they are disjoint, any edge meeting both lies entirely in $E$, so $uM$ has no new witness. Otherwise $A\cap B=\{x\}$, and its new witness is $xz$ with $z\notin E$. The singleton $x$ satisfies $u$, so safety forces every member of $M$ to be adjacent to $x$. Since $xz$ is missing and $M$ is old, the old-center four-cycle argument forces the full equality $M=\{z\}$. The new witness from $vM$ also contains $z$. Together with the matching edge at $x$, these are three distinct selected edges in one connected component, contrary to $P_3\sqcup2K_2$. All arguments concern full supports; repeated factors and scalar degenerations were included.
\end{proof}

\begin{proposition}[A linear preparation with arbitrary cross coefficient]\label{prop:pathpurelineargeneral}
Let $G$ be finite, simple and four-cycle-free, and let a positive variable-input circuit compute $P\in\Sup(G)$ with nonempty quadratic support $H$ and exactly four global missing edges. On their endpoint set $T$, assume $G[T]=P_3\sqcup2K_2$ and $H[T]=\varnothing$. Keep an exact original two-visible-occurrence frame
\[
                  P_2\vee gL\vee fM\vee qgf,\qquad q\ge0.
\]
Suppose the original private product roots $g,f$ jointly use at most three multiplication gates in an unaffected multiplication-only network whose boundaries are old homogeneous linear forms and at most one arbitrary nonnegative homogeneous linear form $S$. Sharing, original raw-product ancestry and repeated factors are allowed. The contexts are the actual old linear or scalar contexts, retaining all low-part identities. If $K$ counts original retained products, then $\vc(G)\le K+2$. No claim is made that an arbitrary hidden addition gives such an unaffected network.
\end{proposition}
\begin{proof}
For $q=0$, apply Proposition~\ref{prop:pathpurelinear}. Suppose $q>0$. The actual cross identities make both $L,M$ nonzero old homogeneous linear forms: each contains the nonzero linear low part of the other occurrence. Write the target as $a-b-c$, $de$, $ij$. Its edge function has eight prime positive clauses, because its minimal covers choose either $\{b\}$ or $\{a,c\}$ and one endpoint from each isolated edge. Any safe expression witnessing all four targets restricts to this edge function on $T$.

If one high is ancestral to the other, write $f=gR$ Booleanly. The mixed term absorbs $fM$, leaving $g(L\vee R)$. A connected multiplication network with at most three nodes has at most four distinct external boundary signals. Distributing the displayed disjunction therefore gives at most four linear clauses, contradicting eight. If the roots are nonancestral but share a private product, the budget forces $h=AB$, $g=hC$, $f=hD$, with a repeated raw input represented by the constant~1 when necessary. The high part is
\[
 h(CL\vee DM\vee CD)=h(C\vee M)(D\vee L)(C\vee D),
\]
with at most five clauses, also impossible. These are output identities, with original degrees and gate counts fixed.

The remaining root cones are disjoint. Relabel so $g=UV$ has one private product, while $f$ has at most three boundary linear factors. We establish three facts on $T$.

First, a safe product of at most four linear factors, all old except possibly $S$, has at most two target witnesses. Three disjoint witnesses require eight clauses. Any other three consist of both path arms and an isolated edge. An old factor meeting both arms must contain $b$, since containing $a,c$ instead creates a four-cycle through its old center. If every factor is old, selecting $b$ gives an unsafe singleton. Otherwise selecting $b$ from all old factors forces $S\cap T\subseteq N_G(b)$, excluding both endpoints of the isolated edge. In particular $fM$ has at most two target witnesses.

Second, $g$ cannot be true on a singleton $z\in T$. Safety of $gL$ makes $L\cap T\subseteq N_G(z)$, and safety of $gf$ makes $zf$ safe. If $z$ is on the path, neither isolated edge can be supplied by $gL$, so $f$ must witness both. Its restriction to those four endpoints is safe, because $z$ has no neighbor there, and hence is the two-edge-matching function. That requires four clauses, but $f$ has at most three. If $z$ is on an isolated edge, $gL$ can supply only that edge; the restriction of $f$ to the path and the other isolated edge is safe and must supply all three, again requiring four clauses. This proves the second fact.

Third, if $gL$ supplies any target witness, $g$ cannot be true on two disjoint target edges. Its two supports on $T$ are disjoint by the preceding fact. Orient two supposed edges as $xx'$ and $yy'$, with $x,y$ in one factor and $x',y'$ in the other. Both cross pairs $xy',yx'$ are nonedges. Choose a variable $t\in L\cap T$ from a direct target witness. No endpoint among these four can make both triples $txy'$ and $tyx'$ safe: one remains an unsafe cross pair. Every other target vertex is adjacent to at most one of those four endpoints, so it also cannot make both disjoint cross pairs safe. This contradicts safety of $gL$.

The term $gL$ must supply a target, since otherwise $f(M\vee g)$ alone supplies all four and has at most five clauses. Let $W$ be the set of target edges on which $g$ is true. Both $gL$ and $gf$ supply only edges in $W$, while $fM$ supplies at most two. Thus $W$ contains two edges. The second and third facts force these to be the two path arms. Consequently both isolated edges are witnessed by $fM$.

This safe product has at most four linear factors. On the four isolated-edge endpoints, the four matching clauses force exactly four uniquely rejecting factors. In particular $M$ restricts to a two-point transversal. Its full support is that same pair, because $M$ is old. Indeed an extra variable $w\in M$ is nonadjacent to both members of the opposite independent transversal by the old-center four-cycle argument. Each other factor meets that opposite transversal; selecting $w$ and those members would give an unsafe monomial. This argument requires oldness only for $M$, so it remains valid when another factor is $S$.

The full support of $M$, together with $b$, is a three-vertex cover of all missing edges. Since $M$ is an actual nonzero linear context, at least one original retained product has its entire opposite support contained in $M$; positivity includes every such support in the context. Keep one center from each of the other at most $K-1$ original products and replace the marked product centers by this three-vertex cover. It covers all old and missing edges, giving $\vc(G)\le(K-1)+3=K+2$. All outside vertices, original incidences, repeated centers and fanout were retained in this charge.
\end{proof}

\begin{lemma}[Three occurrences with a common quadratic ancestor]\label{lem:commonancestorpath}
Let $u=UV$ with old linear $U,V$, and let $A,B$ be old linear forms or~1. Choose at most three high occurrences, with repetitions allowed, either from the chain $\{u,uA,uAB\}$ or the fork $\{u,uA,uB\}$. Form an output from their direct old linear or scalar contexts and arbitrary nonnegative pairwise occurrence cross terms. A safe output cannot witness an induced $H$-empty $P_3\sqcup2K_2$.

In an original all-visible frame this excludes the following precise profile: there are at most three useful private products, every private operand after exact carrier erasure is old linear or a raw earlier private product, and the first quadratic product is ancestral to every useful private product. It does not include affected old-plus-high operands, a join of two independent products, or a dependent pair with an independent third product.
\end{lemma}
\begin{proof}
Aggregate contexts of repeated high values by disjunction; each aggregate is still a positive linear clause or a constant. Cross products in a divisibility chain equal the deeper high and can be absorbed by adding~1 to that high's aggregate context. The chain output therefore has form
\[
 u[C_0\vee AC_1\vee ABC_2]
 =UV(C_0\vee A)(C_0\vee C_1\vee B)(C_0\vee C_1\vee C_2),
\]
with at most five linear clauses. For the fork, all cross terms except $uAB$ are similarly absorbed, giving
\[
 u[C_0\vee AC_A\vee BC_B\vee\rho AB],\qquad\rho\in\{0,1\}.
\]
When $\rho=0$, the bracket is the conjunction of
\[
 C_0\vee A\vee B,\quad C_0\vee A\vee C_B,\quad
 C_0\vee C_A\vee B,\quad C_0\vee C_A\vee C_B.
\]
When $\rho=1$, the first three clauses suffice. Including $U,V$, these are at most six clauses. The identities hold with zero contexts, constant-one factors, repeated highs, and equal supports. A safe output with all four target witnesses would restrict to the edge function of $P_3\sqcup2K_2$, which needs eight clauses by its eight maximal independent sets. Thus neither family can supply the target.

For the stated circuit consequence, keep original nodes after exact carrier erasure. The second private product is $uA$ or $u^2$, represented Booleanly by $A=1$ in the latter case. A third product multiplying it by an old form gives the chain; multiplying $u$ by another old form gives the fork. Multiplying or squaring raw products gives a repeated level, also represented by a constant-one degeneration. A product of two old forms with no $u$ ancestor is excluded by hypothesis, and there is no other private product available. Visible highs are private outputs, since bare old linear values cannot be the high input of an unequal-degree cut and off-frame full quadratic carriers have been erased. The exact selected frame has precisely the direct and pairwise terms above, including crosses of distinct occurrences carrying the same high. Its old quadratic part vanishes on the target endpoints. Applying the algebraic argument proves the bounded consequence without changing an original degree or gate count.
\end{proof}

\begin{proposition}[A literal join with three visible occurrences]\label{prop:literaljoinpath}
Let $G$ be finite, simple and four-cycle-free, and let a positive variable-input circuit compute $P\in\Sup(G)$ with nonempty quadratic support $H$. Suppose its four global missing edges have endpoint set $T$ with $G[T]=P_3\sqcup2K_2$ and $H[T]=\varnothing$. Keep the original quadratic frame with $K$ products. Suppose three visible occurrences have highs $u=AB$, $v=CD$, $w=uv$, one of each, where $A,B,C,D$ are old linear forms and these are literal original private product outputs. Allow all actual old linear or scalar direct contexts and pairwise occurrence cross terms. Then $\vc(G)\le K+2$. This assertion concerns this literal join only, with no affected nonlinear boundary or additional insertion multiplicity assumed.
\end{proposition}
\begin{proof}
Write the direct contexts as $L,M,N$. Every pairwise high product is Booleanly $w$, so the missing part is $uL\vee vM\vee w$ if any cross coefficient is positive, and $uL\vee vM\vee wN$ otherwise. Original degrees and incidences remain fixed. The direct terms each witness at most one missing edge; the last term witnesses at most two. All four targets therefore require two from the last term and one distinct new edge from each direct term. The last two form an induced matching with endpoint quartet $E$. At least one of the other target edges avoids $E$.

A positive cross coefficient or positive scalar $N$ gives four-factor equality: all $A,B,C,D$ are full two-point transversals, making both extra edges meet $E$, impossible. Zero $N$ with zero crosses leaves only the two direct terms. Hence only an old quintic $ABCDN$ remains. Its two witnesses give at least three uniquely rejecting full transversal factors. If one is in each of $\{A,B\}$ and $\{C,D\}$, both extra edges again meet $E$. Otherwise interchange the pairs so $A,B$ are the two unique factors in the first pair and $N$ the third. The remaining occurrences $C,D$ must both reject the fourth independent transversal; there is no idle factor.

If $A,B$ are disjoint, $uL$ has no new edge, since every edge meeting both is contained in $E$. Thus $A\cap B=\{x\}$, and the direct new edge is $xz$ with $z\notin E$. Let $x'$ be the matching mate of $x$. The two transversal types containing $x$ are $A,B$; both remaining types contain $x'$. As $C,D$ reject the same remaining class, both contain $x'$, so $v$ is true on singleton $x'$. The fourth target, from $vM$, is the other isolated edge $rs$, disjoint from the path $x'-x-z$ and the other matching edge. A scalar $M$ violates safety at $x'$; otherwise $M$ contains some $y\in\{r,s\}$. Then $vM$ is true on $\{x',y\}$, an independent pair in the explicitly induced graph $G[T]$, a contradiction. This proves the stated bound, in fact by excluding its exact support template. No broader normalization or noninduced-target conclusion is used.
\end{proof}

\begin{proposition}[No original linear sum in a matching attachment]\label{prop:matchinglinearsum}
Let a positive variable-input circuit compute $P\in\Sup(G)$ for finite simple four-cycle-free $G$, with nonempty quadratic support $H$ and four missing edges. Suppose $\vc(G)-\vc(H)=3$, the original frame has $K=\vc(H)$ retained products, and extraction deletes exactly three products and three additions. Let $C$ be its distinct assigned center cover and suppose $G[V\setminus C]$ is a three-edge matching with arbitrary isolates. Write $e_0$ for the fourth missing edge. Assume either $e_0\subseteq C$, or $e_0=hz$ with $h\in C$, $z$ a matching endpoint, and $hz'\notin H$ for its mate $z'$. Then no original deleted addition has both input wires of original minimum degree one. The inputs may be inhomogeneous. This does not cover an isolated attachment endpoint or an endpoint whose mate is an old neighbor.
\end{proposition}
\begin{proof}
Put $T=V\setminus C$ and let $M=G[T]$. For $h\in C$, the full outside old leaf support of its retained star is $N_h=N_H(h)\cap T$: an edge from $h$ to a noncenter can be assigned only to this product. This support is contained in no matching transversal. Otherwise complete it to a three-vertex cover $R$ of $M$. The set $(C\setminus\{h\})\cup R$ covers $H$ and $M$. An internal $e_0$ retains its other center endpoint; an attachment $h_0z$ retains $h_0$ unless $h=h_0$. In that last case neither $z$ nor $z'$ belongs to $N_h$, so choose $z$ in the completion $R$. In every case this covers $G$ with $K+2$ vertices, contradicting the assumed gain.

Apply exact full-carrier erasure and set $C=0$, retaining outside isolates and every original deleted-addition operand occurrence. The output is the matching function. Each full center-side signal $A_h$ is safe for that function. Indeed, if it were true on an independent assignment, extend to a maximal independent set $S$ of $M$, containing all isolates. The retained product has an addition-only path to the output, so its old leaf form must vanish on $S$. This puts $N_h$ inside $T\setminus S$, a matching transversal, contrary to the preceding paragraph. Replace every retained product by its full safe center signal. The output remains exactly the matching function; carrier erasure has already removed feedback through retained quadratic products. The matching conjunction bound forces all three original private products to remain useful.

Use the admitted matching equality rigidity and its original marked-addition operand count. Each private product first becomes true on its own distinct matching edge, on which all preceding private products vanish. It cannot read a predecessor product, or a disjunction of predecessor products alone, as a required input. Keep all original marked additions and zero boundary anchors. If $a$ useful marked additions, $b$ boundary vertices and $c$ connected components occur, the operand multigraph has three product vertices, $a$ addition vertices, $b$ other immediate operands, and $2a$ operand incidences, including repetitions. Every product has a useful exit through these additions and every component has a boundary. For completeness, a private product cannot enter a retained degree-one position by its original degree; equal-degree export to a retained quadratic carrier would retain that private product, while higher-degree export crosses a marked cut. An unanchored component feeding another product violates its own-edge condition. A retained singleton-center network supplies an original degree-one boundary even when restriction makes its value zero. Thus the original boundary argument applies without deleting zero operands or inventing a newly homogeneous input.

Connectedness gives
\[
                 2a\ge3+a+b-c,\qquad b\ge c,\qquad a\le3.
\]
Equality follows throughout: $a=3$, every component has exactly one boundary, and the actual operand multigraph is a forest. In particular all three original deleted additions survive, and parallel incidences would constitute a forbidden cycle.

Suppose one original deleted addition has two original minimum-degree-one inputs. Trace backward from each of its two operand incidences, always following a degree-one operand of any encountered marked degree-one addition. An original multiplication has minimum degree at least two, so cannot lie on such a path. Acyclicity makes each path terminate at a boundary vertex of the counted graph. These are original paths: full-carrier erasure and retained-product replacement cannot change their degree-one provenance, and original zero boundary incidences have been kept. Both paths lie in the same one-boundary component and have different first incidences. They therefore give two paths to the sole boundary, hence an undirected cycle, including the parallel-edge case when the immediate operands coincide. This contradicts the forest. The argument does not assert that other retained boundaries are unaffected or that all private products are nonancestral.
\end{proof}

\begin{proposition}[A preparation and multiple affected old boundaries]\label{prop:pathmultipleaffected}
Let $G$ be finite, simple and four-cycle-free, and let a positive variable-input circuit compute $P\in\Sup(G)$ with nonempty quadratic support $H$ and exactly four global missing edges. Suppose their endpoints $T$ induce $P_3\sqcup2K_2$ in $G$, with $H[T]=\varnothing$. Keep an exact original two-visible-cut frame $P_2\vee gL\vee fM\vee qgf$, $q\ge0$, with its actual old linear or scalar contexts and $K$ original retained products. Suppose
\[
 g=Z(U\vee V),\qquad f=\prod_{i=1}^3(A_i\vee\epsilon_i g),
 \qquad\epsilon_i\in\{0,1\},\quad\sum_i\epsilon_i\ge2,
\]
where all $Z,U,V,A_i$ are nonzero old homogeneous linear forms. Every affected boundary is an actual retained degree-one signal carrying the first cut, with nonzero old low $a_0$ contained in each affected $A_i$, and $a_0L$ old; when $q>0$, also $a_0\subseteq M$. Then $\vc(G)\le K+2$. Original degrees, occurrences, repetitions and fanout are retained. The claim does not allow the raw nonold form $U\vee V$ as one of the $A_i$, or assert an exhaustive hidden-addition normalization.
\end{proposition}
\begin{proof}
Write the target as $a-b-c$, $de$, $ij$, and put $S=U\vee V$. An old form cannot contain both $a,c$, by the four-cycle through its old center. A safe product with only one possibly nonold linear factor cannot witness all edges of an induced $P_3\sqcup K_2$: every positive CNF for that graph must contain its four exact minimal-cover clauses, two of which contain both path leaves. Only its one nonold factor could supply such a clause. Three disjoint target edges require eight clauses. These observations apply after restriction, with the old neighborhoods still in the ambient $H$.

We first record the support facts used below. The direct term $P'=gL$ has at most three linear clauses, so its target witnesses have a one-vertex cover; if there are two, they are both path arms. In that case old $Z$ and nonconstant old $L$ contain $b$, safety excludes $b$ from $S$, and $U,V$ supply different leaves. Selecting $b$ from the old common factors and applying old-neighbor exclusion to each missing arm forces the full supports $U=\{a\}$, $V=\{c\}$, up to exchange. The same argument omits a scalar $L$. Thus no multiple of $g$ can witness an isolated target edge in this case.

For any old linear or scalar $E,F$, a safe grid $g(E\vee F)$ has at most two target witnesses. Zero or positive scalar contexts reduce to at most three clauses. Otherwise its two rows are the safe old cubics $ZUE,ZUF$ and $ZVE,ZVF$. The common-cubic lemma makes each row's missing witnesses a star. Three independent target edges cannot be covered by these two stars. For both path arms and one isolated edge, one row must supply both arms while the other supplies the isolated edge. In the first row, its two common old factors contain $b$; old-cubic rigidity forces the full contexts $E=\{a\}$ and $F=\{c\}$. Then every grid witness meets a path leaf, excluding the isolated edge. This proves the grid bound.

A safe $R=gCM=ZSCM$, with $C$ old, likewise has at most two target witnesses. Three disjoint edges exceed its four-clause count, and the other three-edge shape requires two leaf-containing clauses although only $S$ can be nonold. If $R$ has two disjoint witnesses, its four factors uniquely reject the four independent transversals of that quartet. Every old factor, in particular $M$, then has full two-point complementary-transversal support. For an extra outside variable in such an old factor is nonadjacent to the rejected transversal by old-neighbor exclusion; selecting it and members of that transversal from the other factors would violate safety. No full-support assertion about the nonold factor $S$ is needed.

Whenever the full actual $M$ has size two and $q=0$, it covers every missing witness of $fM$, while $P'$ has a one-vertex witness cover. Mark an original retained product whose entire opposite support lies in $M$. These two vertices cover all its old edges. Keep centers for the other $K-1$ products and add the one $P'$-cover vertex, giving at most $K+2$ vertices covering $G$.

If all three boundaries are affected, write their old parts $A,B,C$. Booleanly $f=ABC\vee g$. The safe old quartic $Q=ABCM$ has at most one target witness: two would make its four factors full distinct transversals, whereas $A,B,C$ have the nonzero common low $a_0$. For $q=0$ the output is $Q\vee g(L\vee M)$, so the grid bound gives at most three targets. For $q>0$, $g$ absorbs its other multiples and has at most two target witnesses, while $Q$ has at most one. Both cases are impossible.

Suppose exactly $A,B$ are affected and $C$ is old. Then $f=ABC\vee gC$. If $q>0$, the output is $Q\vee g(L\vee C)$. Two witnesses in $Q$ would force four distinct full transversals, contradicting $a_0\subseteq A\cap B\cap M$. Thus $Q$ contributes at most one and the grid at most two, again impossible.

Finally let $q=0$. The three terms are $P'=gL$, $Q=ABCM$, and $R=gCM$. If $Q$ supplies two targets, full quartic equality gives full two-point $M$ and the preceding original-product cover charge applies. Otherwise $Q$ supplies at most one. If $P'$ supplies both path arms, the full-leaf conclusion prevents $R$ from supplying either isolated edge, so four are impossible. Hence four targets require exactly two from $R$ and one distinct extra from each of $P',Q$. A disjoint pair from $R$ again gives full two-point $M$ and the cover charge. The remaining case has $R$ supplying both path arms. Old $Z,C$ and nonconstant old $M$ contain $b$, while $S$ excludes $b$ and supplies both leaves. If $P'$ witnesses $de$, one safe summand $ZUL$ does so, where $Z$ contains $b$ and $U$ a path leaf. Neither is a singleton endpoint of $de$. A scalar $L$ is impossible by old-quadratic rigidity; old-cubic rigidity forces, say, the full $L=\{e\}$ and both $Z,U$ to contain $d$. Choosing $d$ from $Z,S$ and $b$ from $C,M$ (omitting scalar $M$) gives an $R$ monomial on the independent pair $\{b,d\}$, a contradiction. This exhausts all compatible zero/scalar and witness cases, proving the bound with full original contexts.
\end{proof}

\begin{proposition}[One affected old boundary with zero cross]\label{prop:pathsingleaffectedzero}
Keep the graph, four global missing edges, induced target, original frame and retained count $K$ of Proposition~\ref{prop:pathmultipleaffected}, but suppose the mixed coefficient is zero and
\[
 g=Z(U\vee V),\qquad f=(A\vee g)BC,
\]
with $Z,U,V,A,B,C$ nonzero old homogeneous linear forms. The affected boundary is an actual retained degree-one signal with nonzero low $a_0\subseteq A$ and $a_0L$ old. Then $\vc(G)\le K+2$, without retained optimality. This statement keeps the old-base and actual-context hypotheses and does not classify other hidden placements.
\end{proposition}
\begin{proof}
Write the path as $y-x-z$ and the isolated edges as $ab,rs$. Set $S=U\vee V$. Beyond $P_2$ the safe output is
\[
 gL\vee Q,\qquad Q=BCM(A\vee ZS)
   =BCM(A\vee Z)(A\vee S).
\]
The affected occurrence has original selected degree one, so $L$ is homogeneous linear or zero; a positive scalar is not its actual degree. The second context $M$ may have any compatible scalar, zero or linear type. The displayed CNF gives $Q$ at most five clauses. It cannot supply all four targets or any three forming an induced matching, each of which requires eight clauses.

The common-cubic lemma applied to $ZUL,ZVL$ makes the target witnesses of $gL$ a star. Thus $L$ is nonzero, and $gL$ supplies either one target or both path arms. Suppose first it supplies both arms. If its full support is the singleton $L=\{x\}$, mark a retained product with complete opposite support in $L$. The three-vertex cover consisting of $x$ and one endpoint of each isolated edge covers that product and all missing edges; centers for the other $K-1$ products prove $K+2$.

Assume $L$ is nonsingleton. If $Z$ is singleton, it is $\{x\}$, so no $g$-multiple witnesses an isolated target. If $Z$ is also nonsingleton, old-cubic singleton rigidity forces $U,V$ to be the full singletons $y,z$: their two cubes supply different arms, and their repeated endpoint in $Z\cap L$ is $x$. Again no $g$-multiple witnesses an isolated target. Therefore $ABCM$ witnesses both $ab$ and $rs$. Full quartic equality makes $A$ a two-point transversal contained in $T$. The nonzero $a_0\subseteq A$ then has a target vertex, as does $L$. This contradicts $a_0L$ old and $H[T]=\varnothing$, including the forbidden-square case. Only the already charged singleton-context case can remain.

If $gL$ has only one target, that target cannot be a path arm: $Q$ would then need the other arm and both isolated edges, an induced three-edge matching. Rename its unique target $rs$. Thus $Q$ supplies $xy,xz,ab$. Every nonconstant old factor among $B,C,M$ meets these three edges and contains $x$, since an old factor containing both $y,z$ would give a four-cycle. Hence $BCM$ is true on singleton $x$, with a positive scalar $M$ omitted.

The two remaining clauses $A\vee Z$ and $A\vee S$ must reject the independent pairs $\{x,a\}$ and $\{x,b\}$, one each: each clause meets $ab$, so cannot reject both pairs. Both clauses therefore omit $x$. Since both path arms are witnesses of $Q$, each clause contains both $y,z$. The old $A$ cannot contain both leaves. A leaf $w$ omitted by $A$ consequently belongs to both $Z,S$, so $g$ is true on singleton $w$. But $gL$ witnesses $rs$, and $L$ contains an endpoint $t$ of that edge. This makes $gL$ true on the independent pair $\{w,t\}$ in the induced target, a contradiction. The cases exhaust the first contribution. All cover charges used full actual contexts and original product identities, proving the asserted scope.
\end{proof}

\begin{proposition}[One affected old boundary with positive cross]\label{prop:pathsingleaffectedpositive}
Let $H\subseteq G$ with $G$ finite, simple and four-cycle-free, and let $T$ induce $P_3\sqcup2K_2$ in $G$, with $H[T]=\varnothing$. Keep a safe exact original two-visible-cut frame $P_2\vee gL\vee fM\vee qgf$ with $q>0$, actual old contexts and original low identities. Suppose
\[
 g=Z(U\vee V),\qquad f=(A\vee g)BC,
\]
where all six displayed base forms are nonzero old homogeneous linear forms, the affected boundary genuinely carries the first cut, and its nonzero low satisfies $a_0\subseteq A\cap M$ and $a_0L$ old. Then this frame cannot witness all four target edges. Additional global missing edges are permitted; no arbitrary hidden-addition normalization is asserted.
\end{proposition}
\begin{proof}
Put $S=U\vee V$. Since $gf=gBC$ absorbs the extra $M$ multiple, the safe missing part is
\[
                 Y=gL\vee ABCM\vee gBC,\qquad g=ZS.
\]
On $T$, a safe expression with all four witnesses equals the target edge function, which needs eight positive clauses. If $gL$ supplied no target, $BC(AM\vee ZS)$ alone would supply them with at most six clauses. Thus $gL$ has a target witness. Its context is nonzero linear: a positive scalar $L$ contradicts the old quadratic relation $a_0L$, and zero cannot supply a witness.

The high $g$ has no true singleton $z\in T$. Otherwise safety gives $L\cap T\subseteq N_G(z)$ and makes $zBC$ safe. If $z$ is on the path, $gL$ supplies neither isolated edge, so $BC$ must witness both. Its restriction to those endpoints is safe because $z$ has no neighbor there. It would compute a two-edge matching with only two clauses, although four are required. If $z$ is on an isolated edge, $gL$ supplies only that edge; $BC$ must safely compute the other path and isolated edge, again requiring four clauses. Both are impossible.

Therefore $Z\cap S\cap T$ is empty. If $g$ were true on disjoint target edges $xx',yy'$, orient them with $x,y$ in one factor and $x',y'$ in the other. Both cross pairs $xy',yx'$ are nonedges. Choose $t\in L\cap T$ from the direct target witness. No endpoint among these four makes both triples $txy',tyx'$ safe; one remains a cross nonedge. Every other target vertex is adjacent to at most one of the four, and also cannot make both cross pairs safe. This contradicts safety of $gL$. Thus the target edges on which $g$ is true contain no disjoint pair.

The safe old quartic $ABCM$ supplies at most two targets. All others, from $gL$ or $gBC$, are true for $g$. To supply four, $g$ must therefore be true on exactly the two path arms, and $ABCM$ must witness both isolated edges. Full quartic equality makes $A,B,C,M$ the four full distinct transversal supports of that matching. Since $a_0\ne0$ lies in $A\cap M$, this intersection is a single target endpoint $x$, and $a_0$ has full support $\{x\}$. The old relation $a_0L$ and $H[T]=\varnothing$ then force $L\cap T=\varnothing$, contradicting the direct target witness. This proof keeps every full support and original low relation, and uses no assumption about other global missing edges.
\end{proof}

\begin{lemma}[A shared raw linear preparation]\label{lem:rawlinearpreparation}
Let $S,Z,A,B,L,M$ be positive linear Boolean clauses, with constants zero and one allowed. Put
\[
 g=SZ,\qquad f=S(A\vee\epsilon g)(B\vee\delta g),\qquad
 Y=gL\vee fM\vee qgf,
 \qquad \epsilon,\delta,q\in\{0,1\}.
\]
Then $Y$ has a positive linear CNF with at most seven clauses. Consequently these exact high identities in a safe two-visible-cut frame cannot supply all four edges of an induced $H$-empty $P_3\sqcup2K_2$. The clauses need not be old. This statement permits repetitions and constants, but asserts no general normalization of hidden additions or arithmetic degrees.
\end{lemma}
\begin{proof}
Factor out $S$ and use idempotence only at the Boolean output:
\[
 Y=S[ZL\vee TM\vee qZT],\qquad
 T=(A\vee\epsilon Z)(B\vee\delta Z).
\]
For $q=0$, the bracket is a two-clause product disjoined with a three-clause product. Distributing gives at most six linear clauses, hence seven after adjoining $S$.

Let $q=1$. If neither boundary is affected, the bracket is
\[
 ZL\vee ABM\vee ZAB=(Z\vee M)(AB\vee L)(Z\vee AB).
\]
The last two factors each distribute into two linear clauses, giving five in the bracket. If exactly one boundary is affected, exchange $A,B$ to make it $A$. Absorption gives
\[
 ZL\vee ABM\vee ZBM\vee ZB=Z(L\vee B)\vee ABM.
\]
This again has at most six linear clauses by the two-by-three distribution. If both are affected, $T=AB\vee Z$ and the bracket is $Z\vee ABM$, with three clauses. These identities include zero, one and coinciding forms, so all cases give the stated bound.

After restriction to the induced target, safety and its four witnesses make $Y$ exactly its edge function. That graph has eight maximal independent sets. A clause false on two different maximal independent sets is false on their union, which contains an edge; thus a valid CNF needs a different clause for each of the eight sets. Restriction cannot increase the seven-clause count, a contradiction.
\end{proof}

\begin{lemma}[A linear first high and an affected boundary]\label{lem:linearfirstaffected}
Let $g,A_1,A_2,A_3,L,M$ be arbitrary positive linear Boolean clauses, with constants zero and one allowed. Suppose at least one $\epsilon_i$ is one, the others are zero or one, and put
\[
 f=\prod_{i=1}^3(A_i\vee\epsilon_i g),\qquad
 Y=gL\vee fM\vee qgf,\qquad q\in\{0,1\}.
\]
Then $Y$ has a positive linear CNF with at most seven clauses. Hence it cannot be a safe contribution witnessing all four edges of an induced $H$-empty $P_3\sqcup2K_2$ in an exact graph frame. This is a Boolean support statement, with no old-neighborhood assumption or general hidden-addition normalization. Representing an original square of a raw linear preparation by its Boolean linear support does not change its arithmetic degree or gate identity.
\end{lemma}
\begin{proof}
Let $I=\{i:\epsilon_i=1\}$, $J=\{1,2,3\}\setminus I$, $k=|J|\le2$, and write
\[
 T=\prod_{j\in J}A_j,\qquad A=A_1A_2A_3.
\]
Empty products are one. The common affected disjunct gives $f=A\vee gT$ and $gf=gT$. If $q=0$, distribution and absorption give
\[
 Y=gL\vee AM\vee gTM=(g\vee AM)(L\vee TM).
\]
Indeed $A$ contains $T$, so the right side's term $AMTM$ is $AM$, which absorbs $AML$. The first factor has four linear clauses $g\vee A_i$ and $g\vee M$. The second has the $k$ clauses $L\vee A_j$ for $j\in J$, together with $L\vee M$. Thus there are at most $5+k\le7$ clauses.

If $q=1$, the term $gT$ absorbs $gTM$, and
\[
 Y=gL\vee AM\vee gT=(g\vee AM)(L\vee T).
\]
The same four clauses occur in the first factor, while the second contributes at most $k$, none when $T=1$. This gives at most $4+k\le6$. Both identities hold with constants and repeated clauses. Finally the induced target has eight maximal independent sets, hence requires eight positive linear clauses by the rejecting-clause argument in Lemma~\ref{lem:rawlinearpreparation}. Safety and four witnesses would make the restricted contribution exactly that edge function, contradicting the bound. No arithmetic degree or original gate count was inferred from these identities.
\end{proof}

\begin{proposition}[Two occurrences of a common first high]\label{prop:doublefirsthigh}
Let $G$ be finite, simple and four-cycle-free, with nonempty old quadratic support $H$ and four global missing edges whose endpoints $T$ induce $P_3\sqcup2K_2$, with $H[T]=\varnothing$. Let $U,V,A,B,C$ be nonzero old homogeneous linear forms, put $g=UV$, and suppose
\[
 f=ABC\vee gC\quad\hbox{or}\quad f=ABC\vee g.
\]
Keep the exact original three-occurrence frame with first two highs $g$ and third high $f$:
\[
 Y=gL_1\vee gL_2\vee fM\vee qgf\vee rg,
 \qquad q,r\in\{0,1\}.
\]
Here $L_1,L_2$ are actual old homogeneous linear contexts or zero, $M$ is an actual old linear or scalar context, and the two nonzero original lows satisfy
$a_1\subseteq A$, $a_2\subseteq B$, with $a_iL_i$ supported on $H$. Retain the other original frame identities. Then a safe $Y$ cannot witness all four target edges.

In particular this excludes the bounded realization with two distinct insertions of $g$ affecting both inputs of $h=(A\vee g)(B\vee g)$, followed by $hC$ or $h(C\vee g)$. It does not classify other insertion orders or dependency patterns.
\end{proposition}
\begin{proof}
Write the target as $xy,xz,rs,tu$. No old form contains both $y,z$: its old center and $x$ would give a four-cycle, with coincident centers either giving a loop or making a target edge old. Hence an old form meeting both path arms contains $x$. The missing witnesses of safe $UV D_i$, for old linear contexts $D_i$, form a star by the common-cubic lemma; on this target its size is at most two. A safe old quartic has at most two missing witnesses, and two force its four full supports to be the distinct transversals of that missing matching. In particular, if $Q=ABCM$ has two targets, full $A,B$ lie in $T$. The original nonzero lows and $a_iL_i$ old then force $L_i\cap T=\varnothing$: otherwise they produce an old edge in $T$ or a forbidden square.

Whenever $Y$ contains $g$ itself, that safe old quadratic and all its multiples supply at most one missing edge. The only remaining term is $Q$, with at most two. This covers $r=1$, and also $q=1$ when $f=ABC\vee g$.

For $f=ABC\vee g$ and $q=r=0$, the remaining terms are
\[
 Q\vee UVL_1\vee UVL_2\vee UVM.
\]
A scalar $M$ gives the already excluded direct $g$ term; zero removes $Q,UVM$. Otherwise the last three terms form a target star of size at most two. Four targets force that star to be the path and $Q$ to supply both isolated edges. Full quartic equality puts $A,B,C,M$ inside the isolated endpoints. The low relations remove $gL_1,gL_2$, and $M$ has no path vertex, a contradiction.

For $f=ABC\vee gC$ and $q=1,r=0$, absorption gives
\[
 Q\vee UVL_1\vee UVL_2\vee UVC.
\]
The same star and quartic argument applies, now using that $C$ has no path vertex. A scalar $M$ leaves $Q$ with only three genuine old factors and at most one target, so does not create an exception.

It remains to consider $f=ABC\vee gC$ and $q=r=0$. Set
\[
 P'=UVL_1\vee UVL_2,\qquad Q=ABCM,\qquad R=UVCM,\qquad L=L_1\vee L_2.
\]
The targets of $P'$ form a star. The union
\[
 Q\vee R=CM(AB\vee UV)
\]
has at most six positive linear clauses. If $P'$ supplies no target, this union would compute the full target after restriction and would need eight clauses. If $P'$ supplies just one path arm, the same union must supply the other arm and both isolated edges, an induced three-edge matching, again requiring eight clauses.

Suppose $P'$ supplies just one isolated edge, say $rs$. At least one of $Q,R$ supplies two of the remaining three targets. Its pair must consist of one path arm and $tu$, since the old-factor theorem makes a pair in one product an induced matching. If it is $Q$, full $A,B\subseteq T$ and the low relations eliminate $P'$. If it is $R$, full quartic equality puts $U,V$ inside the endpoints of that path arm and $tu$, so neither meets $rs$; again $P'$ cannot supply its edge. With scalar $M$, neither quartic has four genuine factors, so the necessary pair is impossible directly.

Finally suppose $P'$ supplies both path arms. Then $x\in U\cap V$. The exact identity
\[
 Y=UVL\vee ABCM\vee UVCM
   =(UV\vee AB)(UV\vee CM)(L\vee CM)
\]
gives eight clauses from its first two factors, each containing $U$ or $V$, and at most two more from the last. All first eight clauses are true on every assignment containing $x$. But there are four maximal independent sets of the target containing $x$, one choice on each isolated edge. The last two valid clauses cannot reject all four: each rejects at most one maximal independent set. This contradiction exhausts the star possibilities. Constants and repetitions only reduce the displayed clause counts. Throughout, the two original lows remain separate, and no Boolean identity changes the original degree or occurrence ledger.
\end{proof}

\begin{proposition}[A shared first cut and a raw third multiplier]\label{prop:sharedcutrawchain}
Let $G$ be finite, simple and four-cycle-free, with nonempty old quadratic support $H$ and four global missing edges whose endpoints $T$ induce $P_3\sqcup2K_2$, with $H[T]=\varnothing$. Let $U,V,A,B,C$ be nonzero old homogeneous linear forms, and set
\[
 g=UV,\qquad h=(A\vee g)(B\vee g),\qquad f=hC.
\]
Keep the exact original three-visible-occurrence frame with consecutive highs $g,h,f$ and contribution
\[
 Y=gL\vee hN\vee fM\vee q_{12}gh\vee q_{13}gf\vee q_{23}hf,
 \qquad q_{ij}\in\{0,1\}.
\]
Here $L,N$ are actual old homogeneous linear contexts or zero, $M$ is an actual compatible old linear or scalar context, and the first insertion is reused in both affected inputs of $h$. Its nonzero old low satisfies $a_0\subseteq A\cap B$ and $a_0L$ old. Retain every original cross identity, in particular $q_{12}>0\Rightarrow a_0\subseteq N$ and $q_{13}>0\Rightarrow a_0\subseteq M$. Then safe $Y$ cannot supply all four target witnesses. This excludes exactly the displayed shared-first-cut and literal third-multiplier template, without classifying other operands or insertion orders.
\end{proposition}
\begin{proof}
Booleanly $h=AB\vee g$, $gh=g$, $gf=gC$, and $hf=hC$. Write $q=q_{12}$, $r=q_{13}$, $s=q_{23}$. Thus
\[
 Y=gL\vee hN\vee hCM\vee qg\vee rgC\vee shC.
\]
Every positive summand is safe. If $q>0$, any variable of $a_0$ lies in $A,B,N$, making the summand $ABN$ true on a singleton. Thus $q=0$.

If $s>0$, absorption gives $Y=gL\vee h(N\vee C)=h(N\vee C\vee gL)$, since $gL$ implies $h$. The first factor has four linear clauses, and the second has the three clauses $N\vee C\vee U$, $N\vee C\vee V$, $N\vee C\vee L$. This gives at most seven clauses, fewer than the target's eight. Hence $s=0$ as well.

Every variable of $a_0$ lies outside $T$. Otherwise the old relation $a_0L$ and $H[T]=\varnothing$ force $L\cap T=\varnothing$, so $gL$ vanishes on the target. If $r=0$, the remaining expression $h(N\vee CM)$ has at most six clauses. If $r>0$, it is
\[
 h[N\vee C(M\vee g)],
\]
with four clauses for $h$ and the three clauses $N\vee C$, $N\vee M\vee U$, $N\vee M\vee V$ for the bracket. Both contradict the eight-clause bound. These identities include scalar and zero $M$.

The safe old quartic $Q=ABCM$ consequently has at most one target: two would put the full $A,B,C,M$ inside their target quartet, contradicting $a_0\subseteq A$ outside $T$. Also $ABN$ has no target. A putative target $xy$ would make $A,B$ nonsingleton, so cubic rigidity forces full $N=\{y\}$ and $x\in A\cap B$, after exchange. Choose $z\in a_0$ outside $T$. Safety would make $zy$ an edge, but the old factor $A$ contains $x,z$ and $xy$ is missing, so old-neighbor exclusion forbids $zy$. This is a contradiction.

If $r>0$, $gC$ absorbs $gCM$. The only remaining targets come from $Q$ and the common-cubic star $gL\vee gN\vee gC$, with at most one plus two targets. Thus $r=0$. Put $R=gCM$ and $G_0=gL\vee gN$. Their witness bounds are respectively two and a star of size two, while $Q$ has at most one. If $R$ has two targets, full quartic equality confines $U,V,C,M$ to their matching quartet. Every $G_0$ target meets it through $U$, and every $Q$ target through $C$. But any matching pair chosen from $P_3\sqcup2K_2$ leaves a target disjoint from that quartet. Hence $R$ has at most one target.

Four targets now force $G_0$ to supply both path arms and $Q,R$ to supply the two isolated edges. If $x$ is the path center, each old factor $U,V$ meeting both arms contains $x$, since it cannot contain both leaves. Thus $g$ is true on singleton $x$. Both $C,M$ are true on both isolated edges. Their restriction to the four isolated vertices is safe: an independent assignment making $CM$ true, together with $x$, would make $R$ true on an independent assignment of the induced target. Therefore $CM$ computes the two-edge matching there, with at most two linear clauses although four are required. This final contradiction exhausts the coefficient cases. All original lows, arithmetic degrees and occurrences have been kept fixed.
\end{proof}

\begin{proposition}[One visible cut and independent sets]\label{prop:onevisibleindependentsets}
Let $G$ be finite, simple and four-cycle-free, and let a positive variable-input circuit compute $P\in\Sup(G)$ with nonempty quadratic support $H$. Suppose canonical extraction and pruning delete at most three products and at most three additions, exactly one of the latter being a visible cut, and suppose at least four global edges are missing from $H$. For every vertex set $T$ with $H[T]=\varnothing$, the induced graph $G[T]$ has at most seven maximal independent sets.

Consequently the missing edges cannot include an induced three-edge matching. In particular the one-visible profile is impossible if $G[V\setminus C]=3K_2$ with arbitrary isolates for any old center cover $C$, irrespective of where further missing edges meet $C$. Retained counts, retained optimality and cover gain are unrestricted.
\end{proposition}
\begin{proof}
Use exact full-carrier erasure, preserving the original quadratic frame and degree labels. Every useful private gate feeds the high $F$ of the unique visible cut: its first entry to the selected output passes through that cut, and a return to its own high would be a directed cycle. Off-frame retained quadratic carriers have been erased; a retained linear signal carrying a private contribution lies after the cut and cannot feed its high. Thus $F$ has old homogeneous linear boundaries and at most three private products and two hidden additions, with their original degrees, repetitions and sharing allowed.

The exact Boolean output is $P_2\vee FD$. The actual full context $D$ is one old homogeneous linear form when the selected low has degree one, and a positive scalar when it has degree two. In the former case oldness follows by multiplying $D$ by any variable of the nonzero selected low in $P_2$. A square of the insertion symbol would put that low in both operands of a retained product and create a forbidden square. All global missing edges must therefore be witnessed by the safe contribution $FD$.

Lemma~\ref{lem:threeandtwoor} gives either at most six positive linear CNF clauses for $F$, or a union of at most three products using at most two original old boundary forms each. In the latter alternative multiply each term by the full $D$ before restricting any variable. The resulting positive terms are globally safe products of at most three old forms. Each supplies at most one missing edge by the old-factor theorem. A scalar context adds no factor, zero terms add no witness, and a nonzero constant or linear output term would itself be unsafe. Thus that alternative supplies at most three global missing edges, contrary to the hypothesis. The old neighborhoods here remain in the full ambient $H$, including centers outside a later restriction.

Hence $FD$ has at most seven positive linear clauses. Set variables outside any $T$ with $H[T]=\varnothing$ to zero. The entire $P_2$ vanishes and $FD$ becomes exactly $f_{G[T]}$, with at most seven clauses. Every maximal independent set must falsify a clause, and a clause cannot be false on two different maximal independent sets: it would be false on their union, which contains an edge. Constants cause no exception; an edge-free induced graph has exactly one maximal independent set, including when $T$ is empty. This proves the bound. A three-edge matching has eight maximal independent sets, and adjoining isolates does not change that count, proving both consequences.
\end{proof}

\begin{proposition}[One literal dependent pair and one independent product]\label{prop:mixedrawthreevisible}
Let $G$ be finite, simple and four-cycle-free, and let a positive variable-input circuit compute $P\in\Sup(G)$ with nonempty quadratic support $H$. Suppose the four global missing edges have endpoint set $T$ inducing $P_3\sqcup2K_2$ in $G$, with $H[T]=\varnothing$. Keep an exact original three-visible-occurrence frame, one occurrence for each of
\[
 u=AB,\qquad t=uC=ABC,\qquad v=DE,
\]
where all five factors are nonzero old homogeneous linear forms. The dependence $t=uC$ is literal raw reuse, without an affected old-plus-high operand. Retain all actual old linear or scalar contexts, original lows and pairwise occurrence coefficients. If $K$ is the original retained multiplication count, then $\vc(G)\le K+2$. No retained optimality, further insertion multiplicity or general dependency normalization is asserted.
\end{proposition}
\begin{proof}
Write the actual contexts as $L,M,N$, and the coefficients of the pairs $(u,t),(u,v),(t,v)$ as $\alpha,\beta,\gamma$. Booleanly, beyond $P_2$ the safe output is
\[
 Y=uL\vee tM\vee vN\vee\alpha t\vee\beta uv\vee\gamma tv.
\]
Put $P'=uL$, $Q=tM$, $R=vN$, $W=uv$, $X=tv$. The direct cubics $P',R$ have at most one missing witness each; $Q,W,X$ have at most two. A positive $\alpha t$ is a safe old cubic with at most one. Scalar and zero contexts omit factors or terms. A direct cubic whose first two old factors are nonsingleton and which supplies a missing edge has a full singleton context at one endpoint; its two high factors both contain the other endpoint. Otherwise selecting the context endpoint from one high factor would force the other high factor to be singleton as well.

We first record two original-product cover charges. If a full direct context is a singleton $z$ at an isolated target endpoint or at the path center, choose a three-vertex cover of all missing edges containing $z$. The actual context marks an original retained product whose complete opposite support lies in $\{z\}$. Those three vertices cover that product and the missing edges; centers of the other $K-1$ retained stars give $K+2$.

For the second charge, suppose $\alpha=0$, the path is $x'-x-z$, full $L=\{z\}$, and full $M=\{x',a\}$ where $a$ is on one isolated target edge. The two contexts mark distinct original products. If they entered the same side of one product, its nonzero opposite support would lie in their empty intersection; if they entered opposite sides, positivity would make the original $\alpha$ positive. Choose an endpoint $r$ of the other isolated edge. The four vertices $\{z,x',a,r\}$ cover the missing edges and both marked products, so the other $K-2$ centers again give $K+2$. Repeated centers only improve either charge.

Here is a uniform two-witness case. Suppose $\alpha=0$ and $Q=ABCM$ supplies two targets, with endpoint quartet $T_0$. Full quartic equality makes $A,B,C,M$ the four distinct full transversals. In particular $A\cap B\cap C=\varnothing$, so an edge satisfying $t=ABC$ must lie inside $T_0$; the induced quartet has only the two matching edges. Thus $X$ adds no target. Every extra edge from $P'$ or $W$ meets $T_0$, while $R$ supplies at most one. If the pair in $T_0$ were both isolated edges, the two remaining path arms would be uncovered. Thus it consists of a path edge $xx'$ and one isolated edge $ab$; the other targets are $xz$ and $rs$, and $R$ supplies $rs$. If a present $W$ supplies $xz$, both factors $D,E$ meet the disjoint edges $xz,rs$ and are nonsingleton. Cubic rigidity makes full $N$ an endpoint of $rs$, so the first charge applies. Otherwise $P'$ supplies $xz$. The full transversals $A,B$ must intersect at $x$, so $u$ is true on singleton $x$. Safety and old-neighbor exclusion force full $L=\{z\}$. The remaining transversal $M$ has form $\{x',a\}$ after renaming the isolated endpoints. The second charge applies. This closes every $Q$-two case with $\alpha=0$.

If $\alpha>0$, $t$ absorbs $Q$ and $X$, leaving $P'\vee t\vee R\vee\beta W$. The common-cubic lemma makes the targets of $P'\vee t$ a star of size at most two. If $W$ has two targets, full equality puts $A,B,D,E$ inside their quartet. Every other term then contains one of those full factors, so every other target meets that quartet. But every matching pair selected from $P_3\sqcup2K_2$ leaves another target disjoint from it, a contradiction. Thus $W$ has at most one target. Four targets force the common star to be both path arms and $R,W$ to be the two isolated edges. Both $D,E$ are nonsingleton, and full $N$ is an isolated endpoint by cubic rigidity. The first charge proves the bound.

Let $\alpha=0$ and $\beta>0$. Now $W$ absorbs $X$. The preceding quartet argument excludes two targets in $W$, and the uniform case handles two in $Q$. Otherwise four targets require one distinct target from each of $P',Q,R,W$. Both $A,B$ meet three distinct target edges, from $P',Q,W$; since no target vertex has degree three, both are nonsingleton. Hence full $L$ is a target singleton. An isolated $P'$ target gives the first charge. If $P'$ is a path arm and $R$ is isolated, the $R,W$ targets are disjoint, so $D,E$ are nonsingleton and $N$ is an isolated singleton, again giving the charge. The remaining assignment has $P',R$ as the two path arms and $Q,W$ isolated. Here $N$ is singleton too. If either context is the path center, charge it. Otherwise the two contexts are the distinct path leaves; cubic rigidity makes both $u,v$ true on the path-center singleton. This contradicts safety of $W=uv$.

It remains that $\alpha=\beta=0$. The uniform case already handles $Q$ with two targets. Suppose $X$ is present and has two. Its five old factor occurrences $A,B,C,D,E$ must reject the four independent transversals of the matching quartet $T_1$. Each rejects at most one. If $s$ rejection classes occur only once, then $5\ge s+2(4-s)$, so $s\ge3$. Each unique old rejector has its full complementary two-point support. At least one of $A,B,C$ is unique; consequently every $Q$ edge meets $T_1$. If $T_1$ consists of the two isolated edges, both path arms lie outside it and must be supplied by $P',R$. Then none of $A,B,D,E$ could be a full unique factor, leaving only $C$, a contradiction. Otherwise the pair is one path arm and one isolated edge. The other isolated target, disjoint from $T_1$, must come from $P'$ or $R$. Its quadratic high has two nonsingleton factors, because each also meets both edges of $T_1$. Its full direct context is therefore an isolated singleton, proving the bound by the first charge.

Finally $Q$ and $X$ have at most one target each. If $\gamma=0$, the three direct terms cannot supply four; otherwise all $P',Q,R,X$ supply distinct targets. Again $A,B$ meet three target edges, so full $L$ is singleton. An isolated $P'$ target is charged. If $P'$ is a path arm and $R$ isolated, the disjoint $R,X$ witnesses force full $N$ to be an isolated singleton and are charged. The only remaining assignment has $P',R$ as the two path arms and $Q,X$ as the isolated edges. Full $N$ is singleton. If it is the path center, charge it; otherwise write $N=\{z\}$ for its path leaf. The positive original cross $\gamma$ puts the nonzero linear $t$-low inside full $N$, so that low has support $\{z\}$. Its own old product with $M$ is supported on $H$. But $Q=tM$ witnesses an isolated target, so $M$ contains an endpoint $w\in T$ of that edge. This produces the forbidden old edge $zw$ in $H[T]$. A scalar $M$ is already incompatible with the positive-cross original degree types.

The coefficient cases and target-witness cases are exhaustive. Each yields a contradiction or one of the full-context charges. These charges use all four global missing edges, original product identities and symmetric original low/context facts; no extraction or degree count has been recomputed after a Boolean identity.
\end{proof}

\begin{proposition}[The original operands with four products and two additions]\label{prop:fourproductreduction}
Let $G$ be finite, simple and four-cycle-free, let $P\in\Sup(G)$ have nonempty quadratic support $H$ and exactly four missing edges, and suppose extraction deletes at most four products and at most two additions. Let $K$ count original retained products. If $\vc(G)>K+2$, exact-frame carrier erasure keeps all four original private products and both additions useful, and both additions are visible cuts. In their chronological cut order the first high $g$ is Boolean-equivalent to $UV$, with $U,V$ nonzero old homogeneous linear, and the second high has form
\[
 f=\prod_{i=1}^4(A_i\vee\epsilon_i g),\qquad \epsilon_i\in\{0,1\},
\]
with nonzero old homogeneous linear $A_i$. Every affected factor is an actual degree-one signal carrying the first cut, with its original low and full context identities. Repetition only pads the Boolean representation. A repeated-private-square first high may still have two original product ancestors and original degree four; neither its degree nor the original gate ledger is changed. No raw private-ancestor operand remains unclassified in this residual.
\end{proposition}
\begin{proof}
One center of each retained quadratic star covers $H$, so the hypothesis gives $\vc(J)\ge3$ for the four missing edges $J$. Apply the full-carrier erasure, preserving the exact selected frame and $P_2$. If a private product became dispensable, pruning would leave at most three products and two additions, contrary to Proposition~\ref{prop:threeproductstwosums}. If an addition became dispensable, exactly one remaining addition would require at least $2^4-2=14$ private products by Proposition~\ref{prop:oneadditionall}; with none, no higher witness could first enter the selected quadratic output through a cut. Thus the original bounds are exact and all six private gates remain useful. At least one cut is visible, and Proposition~\ref{prop:sevenfactorprofiles} excludes exactly one. Both are therefore visible unequal-degree cuts.

Before substituting their highs, the exact independent-symbol output is
\[
 P_2\vee gL\vee fM\vee qgf.
\]
An inserted symbol cannot occupy both factors of a selected quadratic product: its nonzero old low would create a square in $P_2$. Thus the direct contexts are actual old homogeneous linear forms for selected degree one, and scalars for selected degree two. If $a_0$ is the first linear low, $a_0L$ is old. A positive mixed coefficient forces both lows linear and gives $a_0\subseteq M$ and the symmetric relation for the other low. These are original selected-occurrence identities.

Every full retained quadratic carrier has an addition-only selected output path; its off-frame uses into a private high have been erased. The remaining retained boundaries are degree-one signals. Before the first cut they are old homogeneous linear; afterward each is old $A$ or $A+\alpha g$, with $A$ nonzero old linear and $\alpha>0$. There is no third private addition to insert another private value into such a signal. A raw earlier private product may instead feed another product directly. Consequently the only boundaries of the later private network are old linear signals, actual first-cut-affected linear signals, and raw private ancestors of $g$, including $g$. Neither high can itself be a retained degree-one boundary, whose positive minimum degree cannot be the high side of an unequal cut, or an erased off-frame retained quadratic carrier. Both highs are private product outputs.

All useful private products reach $g$ or $f$, the only remaining entries into the selected frame. Let $k$ count original private product nodes ancestral to $g$ in the erased circuit, including its root. This is not an assertion about ancestry through erased carrier uses. Its multiplication-only network has at most $k+1$ distinct old boundary signals: $2k$ operand occurrences include at least $k-1$ internal uses. Sharing and repetitions only lower the count. The remaining $4-k$ products form a connected ancestor network of $f$ on the boundary types just listed, with at most $5-k$ external signals.

First remove a uniform raw-$g$ case. If $g$ itself is an external boundary of that later network, then Booleanly $f=gR$. Raw ancestor divisors of $g$ and affected factors $A\vee g$ are absorbed. At most $4-k$ ordinary old factors remain in $R$. For $q=0$, the two safe products $gL,gRM$ use at most
\[
 (k+1)+(4-k)+2=7
\]
distinct old supports. For $q>0$, $gf=f$ absorbs $fM$ and leaves at most six supports in $gL\vee f$. Lemma~\ref{lem:sevenoldfactors} contradicts $\vc(J)\ge3$. This is an output identity, not a deletion or recount of original gates.

The case $k=4$ is excluded by Proposition~\ref{prop:sevenfactorprofiles}. If $k=3$, only one product remains for $f$. With two old or affected linear operands it is exactly Proposition~\ref{prop:quartictwoboundary}, since $g$ uses at most four old supports. A proper raw ancestor $u$ uses at most three such supports. Paired with old $A$ or affected $A\vee g$, it gives $uA$ or $uA\vee g$, covered by Proposition~\ref{prop:rawdivisor}. Two raw ancestor operands are its explicit two-divisor alternative. Raw $g$ was already excluded. Every $k=3$ allocation therefore gives the forbidden cover bound.

Suppose $k=2$. Let $h=UV$ be the proper private ancestor. The original second product is either $g=hW$ for old linear $W$, or the repeated-input square $g=h^2$. The remaining two-product network for $f$ has at most three external boundaries. If it has neither raw $h$ nor raw $g$, these are old or actual affected linear signals. In the first case $g=UVW$ is an actual old cubic, so Proposition~\ref{prop:cubicthreeboundary} applies after Boolean repetition of fewer later factors. In the square case retain the residual: $g$ is Boolean-equivalent to $UV$, and $f$ uses at most three old-or-affected boundaries. Calling this original degree-four square an old cubic would not justify the preceding proposition, and is not done here.

If raw $h$ occurs in the later network, factoring it out leaves at most two old-or-affected boundaries. Thus
\[
 f=h(A\vee\epsilon g)(B\vee\delta g)
\]
Booleanly. This includes a raw ancestor at an internal product, an intermediate repeated square, or a direct $h$ input at the high root. With one other boundary repeat it; with none use $A=U,B=V$ and no affected factor. Proposition~\ref{prop:rawquadratic} covers these cases. For the original square $g=h^2$, its permitted Boolean identity $g=hU$ supplies the displayed three-support notation while retaining the original degree and two nodes. Alternatively raw $h$ is then Booleanly raw $g$, and the uniform seven-support count applies. Thus the sole $k=2$ residual is the no-raw-ancestor repeated-private-square case, already of the asserted Boolean-quadratic form.

Finally $k=1$ gives $g=UV$ directly. It has no proper private ancestor; raw $g$ has been excluded. The remaining three-product network for $f$ therefore has at most four old-or-affected linear boundaries, whose product is its Boolean value. Repeat a boundary if fewer occur. Each affected copy is the same actual retained degree-one signal, so its old part contains the original first low and its context identities are unchanged. The original $k=2$ square residual has exactly the same representation with at most three actual boundaries. If that square's selected first low had degree two, no affected degree-one boundary or mixed retained product could occur; the first context would be a nonzero scalar. Then safe Boolean-old-quadratic $g$ has at most one missing witness and the other old product $fM$ at most two, already impossible. Thus no low-degree qualification needed in the later template is lost. All cases have been accounted for, with original node identities and arithmetic degrees fixed.
\end{proof}

\begin{proposition}[Four products and two additions]\label{prop:fourproductstwosumscover}
Let $G$ be finite, simple and four-cycle-free. Let a positive variable-input circuit compute $P\in\Sup(G)$ with nonempty quadratic support $H$ and exactly four missing quadratic edges. If canonical extraction and pruning delete at most four multiplication gates and at most two addition gates, then
\[
 \vc(G)\le K+2,
\]
where $K$ is the multiplication count in the original extracted quadratic frame. Retained optimality, retained addition count and cover gain are unrestricted.
\end{proposition}
\begin{proof}
Suppose $\vc(G)>K+2$. Proposition~\ref{prop:fourproductreduction} applies to the original at-most gate counts and preserves the entire quadratic frame, hence $H$ and $K$. It produces two visible cuts with their exact two-symbol frame, a first high Boolean-equivalent to $UV$, and a second high represented by four actual old-or-first-cut-affected linear boundaries. All original low/context identities survive. In particular the repeated-private-square possibility retains its original degree four and original ancestor count; Boolean repetition creates no arithmetic gate.

These are exactly the hypotheses of Proposition~\ref{prop:quadraticfourboundary}, which gives a two-vertex cover of the four missing edges. One chosen center for each of the $K$ original retained stars covers $H$. Adding those two vertices covers $G$ with at most $K+2$ vertices, a contradiction. No replacement of $K$ by $\vc(H)$ was made.
\end{proof}

\begin{corollary}[The remaining four-witness possibility]\label{cor:fourresidualfinal}
Let $G$ be finite and four-cycle-free. For a circuit below the star count with exactly four missing quadratic witnesses and nonempty quadratic support $H$, write $d=\vc(G)-\vc(H)$, let $D$ be the deleted-gate count, and let $e$ be the extracted circuit's excess over $|E(H)|+\vc(H)-1$. The only remaining necessary numerical possibility is
\[
 d=3,\qquad D=6,\qquad e=0,
\]
and the six deleted gates must be three products and three additions.

The quadratic remainder is optimal. The outside graph of a minimum $H$-cover cannot contain an induced four-edge matching. The exceptional outside graphs $P_4\sqcup K_2$ and $K_3\sqcup K_2$ are excluded. On either remaining outside shape, at least two cuts are visible. On $P_3\sqcup2K_2$, if all three are visible, the private products must have an original ancestry relation. With two visible cuts, the unaffected linear-preparation template of Proposition~\ref{prop:pathpurelineargeneral} is excluded for every cross coefficient. Lemmas~\ref{lem:pathproductunion} and~\ref{lem:commonancestorpath} exclude their precise product--union and common-ancestor profiles. These are necessary conditions, not an asserted saving circuit or a complete four-witness lower bound.
\end{corollary}
\begin{proof}
Theorem~\ref{thm:maxfour} removes the maximal-gain row from Corollary~\ref{cor:fourmissingledger}. For retained product and addition counts $K,A$, write
\[
 e=(K-\vc(H))+(A-(|E(H)|-1)).
\]
Both summands are nonnegative by the separate quadratic bounds. Proposition~\ref{prop:gainthreesix} removes $d=3,D=5,e=0$ and the $e=1$ allocation $K=\vc(H)$, $A=|E(H)|$. The other $e=1$ allocation has $K=\vc(H)+1$, $A=|E(H)|-1$. Corollary~\ref{cor:gainthreeprofile} reduces that allocation to two private products and three additions, and Proposition~\ref{prop:gainthreefiveprivate} excludes it. This removes the whole $d=3,(D,e)=(5,1)$ row.

For cover gain two with $D=5,e=0$, Proposition~\ref{prop:oneproductcover} excludes at most one private product, because $K=\vc(H)$; Proposition~\ref{prop:oneadditionall} excludes at most one private addition. Proposition~\ref{prop:threeproductstwosums} excludes the three-product/two-addition allocation. In the remaining two-product/three-addition case, Proposition~\ref{prop:allvisibletwocover} excludes three visible cuts. With two visible cuts, Corollary~\ref{cor:hiddenunequalany} excludes a hidden unequal-degree sum; Proposition~\ref{prop:hiddenequalpure} makes the equal-degree alternative pure old-linear, and Propositions~\ref{prop:independentlinearglobal} and~\ref{prop:dependentlinearglobal} exclude its two ancestry alternatives. With one visible cut, Proposition~\ref{prop:onevisibletwocover} gives $\vc(G)\le K+1=\vc(H)+1$. Thus the entire cover-gain-two row is excluded.

In the remaining $d=3,D=6,e=0$ row, $K=\vc(H)$. Proposition~\ref{prop:gainthreeproducts} requires at least three deleted products without an addition-count restriction. At most one deleted addition is excluded by Proposition~\ref{prop:oneadditionall}. The only allocations were therefore $(p,a)=(3,3)$ and $(4,2)$. Proposition~\ref{prop:fourproductstwosumscover} gives $\vc(G)\le K+2$ for the latter, contrary to $\vc(G)=K+3$. Hence only three products and three additions remain.

Corollary~\ref{cor:fourdeletedproducts} excludes an outside induced four-edge matching in this three-product allocation. On either exceptional outside graph, Proposition~\ref{prop:hiddenlinearnecessary} requires a hidden original $1+1$ addition. At least one cut is visible; Proposition~\ref{prop:onevisibleexceptional}, Corollary~\ref{cor:allvisibleexceptional}, and Proposition~\ref{prop:twovisibleexceptional} exclude respectively one, three and two visible cuts. Proposition~\ref{prop:onevisibleindependentsets} excludes one visible cut on either remaining outside shape: $P_3\sqcup2K_2$ and $3K_2$ with isolates each have eight maximal independent sets, and their outside vertex sets are $H$-independent. Proposition~\ref{prop:paththreevisible} excludes nonancestral products when all cuts are visible, and Proposition~\ref{prop:pathpurelineargeneral} excludes its exact two-visible unaffected linear-preparation template for arbitrary cross coefficient, each by the contradiction $\vc(G)\le K+2$. Lemmas~\ref{lem:pathproductunion} and~\ref{lem:commonancestorpath} exclude their stated product--union and common-ancestor support templates. Proposition~\ref{prop:literaljoinpath} excludes its literal-join profile. Proposition~\ref{prop:matchinglinearsum} excludes original $1+1$ deleted additions in its specified three-matching attachment cases. Propositions~\ref{prop:pathmultipleaffected}--\ref{prop:pathsingleaffectedpositive} exclude their prepared-high templates with affected old base boundaries. Lemmas~\ref{lem:rawlinearpreparation} and~\ref{lem:linearfirstaffected} exclude their exact shared raw-preparation and affected linear-first-high templates. Propositions~\ref{prop:doublefirsthigh} and~\ref{prop:sharedcutrawchain} exclude their bounded double-insertion and shared-cut/raw-chain templates, and Proposition~\ref{prop:mixedrawthreevisible} excludes the literal dependent pair plus independent product with one occurrence each. Other raw-preparation uses, hidden and ancestral configurations, and the remaining attachment cases are not settled by this corollary.
\end{proof}



\section{A disjunction count permitting conjunction gates}
The following lemma is the counting ingredient needed after circuit transposition. It is stated for ordinary Boolean circuits and is independent of inverse-semiring terminology.

\begin{lemma}\label{lem:orsums}
Let $S_1,\ldots,S_r$ be nonempty subsets of a finite input set $Y$, with $|S_i\cap S_j|\le1$ for $i\ne j$. A monotone Boolean circuit with these $r$ designated outputs
\[
 h_i(y)=\bigvee_{u\in S_i}y_u
\]
has at least $\sum_{i=1}^r(|S_i|-1)$ disjunction gates. Conjunction gates, shared wires and constants zero and one are permitted.
\end{lemma}
\begin{proof}
For a wire $g$, let
\[
 A(g)=\{u\in Y:g(1_{\{u\}})=1\}.
\]
Fix $S=S_i$, and let $b_S$ be the input equal to zero on $S$ and one outside $S$. Call a wire \emph{false} when $g(b_S)=0$. Monotonicity implies that a false wire has $A(g)\subseteq S$, since $1_{\{u\}}\le b_S$ for $u\notin S$.

Process the gates in their original order, maintaining a partition of $S$, initially into singletons. We maintain that for every processed false wire, its set $A(g)$ is empty or is contained in one part. False variable inputs have singleton sets and a false constant has the empty set. At a false conjunction gate, at least one operand is false, and $A(g)$ is contained in that operand's set. No change to the partition is needed. At a false disjunction gate both operands are false, and $A(g)$ is their union. If their nonempty sets lie in different parts, merge these two parts; otherwise make no change. This preserves the invariant and makes at most one merge at that gate. A merge gate is false on $b_S$ and is true on at least two distinct singleton inputs from $S$.

The designated output $h_i$ is false on $b_S$ and has $A(h_i)=S$. Thus all elements of $S$ must belong to one part at the end, requiring at least $|S|-1$ merges. Finally a merge gate for $S_i$ cannot also be a merge gate for $S_j$: being false on both complementary assignments forces its singleton set into $S_i\cap S_j$, whereas a merge requires that set to have at least two elements. The gates counted for distinct outputs are disjoint. Summing proves the assertion.
\end{proof}

\section{The one-sided Boolean count}
This section proves Theorem~\ref{thm:onesided}, including the role of degree in the transposition.

Trim the circuit to the gates feeding its output and delete isolated graph vertices. Since syntactic degree never decreases along a wire, every remaining marked-dependent wire has degree one, and every conjunction has at most one such operand. A wire of degree zero depends on $Y$ alone and takes value one on the all-one $Y$-assignment.

\begin{lemma}\label{lem:normal}
A useful disjunction cannot have one marked-dependent operand and one degree-zero operand. Consequently the dependent part consists of $|X|$ marked inputs, $M$ disjunctions of two dependent wires, and $K$ conjunctions with one dependent and one degree-zero operand.
\end{lemma}
\begin{proof}
Such a mixed disjunction would be true at $x=0,y=1$. Along a path to the output, every later conjunction has a degree-zero other operand, true at this assignment, and every disjunction preserves truth. The output would be true there, contradicting $f_G(0,1)=0$. A conjunction of two dependent operands would have degree two. These observations exhaust the possibilities.
\end{proof}

For a dependent wire $w$, distributivity gives the expansion
\[
 F_w(x,y)=\bigvee_{v\in X}\bigl(x_v\wedge\alpha_{w,v}(y)\bigr),
\]
where $\alpha_{w,v}$ is the disjunction, over directed paths from $x_v$ to $w$ in the dependent part, of the conjunction of the degree-zero multipliers along that path. A missing path contributes zero and an empty multiplier list contributes one. This is a finite path expansion, used as an identity rather than as a circuit construction. At the output $o$,
\[
 \alpha_{o,v}(y)=\bigvee_{u\in N(v)}y_u,                 \tag{\thesection.1}\label{eq:neighborhood}
\]
as follows by setting the marked inputs to the singleton $v$.

\begin{lemma}\label{lem:centers}
The circuit has $K\ge\vc(G)$ mixed conjunction gates.
\end{lemma}
\begin{proof}
For such a gate $g$, let $\operatorname{Src}(g)$ be the set of marked inputs connected to its dependent operand by paths containing no conjunction. Some gates may have an empty source set. Call $g$ relevant if it is the first conjunction on a path witnessing some edge $vu$: that is, a path from $x_v$ to the output whose multipliers are all true at the singleton $y=1_{\{u\}}$.

Every edge has such a witnessing path by \eqref{eq:neighborhood}, and it has a first conjunction: a path with none would make the output true when $y=0$. For a relevant $g$, if $\operatorname{Src}(g)=\{v\}$, choose $v$ as its center. Otherwise, following the same suffix path from any $v'\in\operatorname{Src}(g)$ shows that every neighbor $u$ witnessed through $g$ is a common neighbor of all these sources. Two distinct sources have at most one common neighbor, by the absence of a four-cycle. Since relevance gives at least one, that neighbor is unique; choose it as the center of $g$.

For every edge $vu$, the first conjunction of its witnessing path has center either $v$ or $u$. The centers of the relevant gates form a vertex cover, and there are at most $K$ of them.
\end{proof}

The wire reversal below is an instance of the classical transposition principle. \citet[p.~46]{JuknaBook} gives the exact OR-network count and its historical attributions. We include the read accounting because the degree-zero multipliers here are themselves Boolean circuits, and conjunctions remain in the transposed network. Isolated marked inputs have already been removed.

\begin{lemma}[Transposition with an exact read count]\label{lem:transpose}
Let $P_\vee$ be the number of degree-zero disjunction gates. There is a monotone circuit with outputs $\alpha_{o,v}$ for $v\in X$ and with at most
\[
 P_\vee+M-|X|+1
\]
disjunction gates.
\end{lemma}
\begin{proof}
Keep all degree-zero gates. Process the dependent wires in reverse order. Put $\beta_o=1$. For a wire $a\ne o$, form the disjunction of one term for each occurrence in which a later dependent gate reads it. If that reader $r$ is a disjunction, the term is $\beta_r$. If it is a conjunction with degree-zero multiplier $m_r$, the term is $\beta_r\wedge m_r$. The latter term needs at most one conjunction, which is irrelevant to the present count. Reverse induction shows that $\beta_a$ is the disjunction of the products of multipliers on paths from $a$ to the output. Therefore $\beta_{x_v}=\alpha_{o,v}$.

Let $\operatorname{reads}(a)$ count operand occurrences, so using the same wire twice counts twice. Every useful dependent wire other than the output has at least one read. Joining its terms requires at most $\operatorname{reads}(a)-1$ disjunctions. There are $2M+K$ read occurrences in total, and $|X|+M+K-1$ dependent wires other than the output. Thus the new disjunction count is
\[
 \sum_{a\ne o}\bigl(\operatorname{reads}(a)-1\bigr)
 =2M+K-(|X|+M+K-1)=M-|X|+1.
\]
Adding the $P_\vee$ degree-zero disjunctions proves the bound. Constants or duplicate terms may reduce it further.
\end{proof}

\begin{proof}[Proof of Theorem~\ref{thm:onesided}]
The neighborhoods $N(v)$ for distinct $v\in X$ have pairwise intersections of size at most one. Lemma~\ref{lem:orsums} applied to the transposed circuit gives
\[
 P_\vee+M-|X|+1\ \ge\ \sum_{v\in X}(|N(v)|-1)
 =|E|-|X|.
\]
Thus the original number $P_\vee+M$ of disjunctions is at least $|E|-1$. Lemma~\ref{lem:centers} gives at least $\vc(G)$ conjunctions, even before counting any degree-zero conjunctions. The lower bound follows. The star circuit from the proof of Theorem~\ref{thm:main} has marked degree one: every star multiplies a marked linear sum by an unmarked variable, or a marked variable by an unmarked sum. It attains the bound.
\end{proof}

\begin{corollary}\label{cor:degreewitness}
If a monotone circuit beats $|E|+\vc(G)-1$ for a four-cycle-free bipartite graph, then after deleting unused gates its syntactic degree is at least two with respect to each side of the bipartition. In each orientation, some conjunction reads two positive-degree wires.
\end{corollary}
\begin{proof}
Apply Theorem~\ref{thm:onesided} to each choice of marked side. If every conjunction had a degree-zero operand, induction would keep every useful wire's degree at most one.
\end{proof}

\section{What supported zeros say about Boolean complexity}
Here we distinguish an output that is a supported zero on an input from a supported zero coefficient of a formal polynomial. The distinction is necessary even for a linear polynomial: $x-1$ has no zero coefficient, but evaluates to the old zero at $x=1$.

The map
\[
 R_0\hookrightarrow R\times\B,
 \quad\tau\longmapsto(0,0),\quad a\longmapsto(a,1)
\]
is an injective semiring homomorphism. Its two coordinates will be denoted $\pi$ and $\chi$. For a circuit over $R_0$, the first coordinate is a ring circuit and the second is its monotone shadow. Unary additive inverse, if present, acts trivially on the second coordinate.

For the following proposition, inputs to $(\F_2)_0$ are $\tau$ and $1$, representing false and true, and constants $\tau,e,1$ are allowed. Boolean circuits in this section may also use free constants zero and one. For a nonconstant graph function these constants can be removed by evaluating gates with constant operands, without increasing size. Read the answer by $\pi$; define $C_0(f)$ to be the minimum size when the \emph{final output} is never $e$ on Boolean inputs, and $C_*(f)$ without that condition. Intermediate outputs may be $e$ in both definitions. All sizes count the relevant binary gates; Boolean NOT gates count one. The unpublished Fable consultation already states this output qualifier and the factor-three comparison explicitly \citep[Q2, corollary]{FableConsult}; we restate and prove them to fix the model.

\begin{proposition}\label{prop:complexity}
For every monotone Boolean function $f$,
\[
 \cmon(f)\le C_0(f)\le3\cmon(f),
 \qquad
 C_*(f)=\Theta(\cbool(f)).
\]
Reading the answer by $\chi$, instead of by $\pi$, always computes a monotone Boolean function.
\end{proposition}
\begin{proof}
The shadow of a final-output-zero-free circuit agrees with its ring value, since its only possible final values are $\tau$ and $1$. Taking the shadow therefore gives the first inequality. Conversely conjunction lifts as multiplication. Disjunction lifts on $\{\tau,1\}$ by
\[
 a\mathbin{\operatorname{OR}}b=a+b+ab,
\]
using one multiplication and two additions. Checking the four pairs shows its result is $\tau$ for two false inputs and $1$ otherwise. Replacing every monotone gate by these formulas gives the second inequality. The intermediate sum $1+1=e$ is allowed.

For unrestricted circuits, the same formula implements OR, multiplication implements AND, and $1+a$ implements NOT in the ring coordinate. Thus $C_*(f)\le3\cbool(f)$. Conversely addition in $\F_2$ is XOR and multiplication is AND. XOR has a constant-size AND/OR/NOT implementation, while both $e$ and $\tau$ map to the Boolean constant zero. Projecting and replacing gates therefore gives $\cbool(f)\le cC_*(f)$ for an absolute constant $c$, with the harmless constant-function cases treated directly. The last assertion follows because $\chi$ sends addition to OR and multiplication to AND.
\end{proof}

\begin{remark}[The output condition over a general ring]\label{rem:generalring}
Let $R$ be any nonzero commutative unital ring. Encode each Boolean input by $\tau,1$, and allow free constants $\tau,1$. Require the final output to be $\tau$ for false and $1$ for true, and allow addition, multiplication and unary additive inverse, counting each gate once. The minimum size $C_{0,R}(f)$ for a monotone function then satisfies
\[
 \cmon(f)\le C_{0,R}(f)\le4\cmon(f).
\]
For the lower bound, take the shadow and delete inverse gates, whose shadows are identities. For the upper bound, AND is multiplication and OR is $a+b-ab$, using at most four gates. Checking the four Boolean input pairs gives precisely $\tau,1$ as the final results. Intermediate old zeros are allowed: for example, $1+1$ is the old zero in characteristic two. Thus this extends the output-level comparison, while the unrestricted $\Theta(\cbool)$ statement above remains specific to $(\F_2)_0$.
\end{remark}

\begin{remark}[Why the output qualifier is essential]
If $e$ is forbidden at \emph{every} useful gate and every Boolean input, and constants are $\tau$ and $1$, the remaining nonzero shadows are monotone and true on the all-one input. Adding two such wires would produce $1+1=e$ there. Consequently every useful addition has a zero operand and can be deleted. The only resulting nonconstant functions are conjunctions of variables. The all-gates formulation of the monotone-equivalence assertion is therefore false.
\end{remark}

\begin{corollary}[The precise consequence for clique]\label{cor:clique}
For every clique-function family to which the monotone lower bound of \citet[Theorem~3.9, p.~12]{AB} applies, final-output-zero-free $(\F_2)_0$ circuits obey the same lower bound up to constant factors. The family has polynomial-size unrestricted $(\F_2)_0$ circuits if and only if it has polynomial-size Boolean circuits.
\end{corollary}
\begin{proof}
Apply Proposition~\ref{prop:complexity}. The cited monotone lower bound is the only external input.
\end{proof}
This identifies the required comparison without proving a new nonuniform Boolean lower bound. A statement merely that the ratio of the two complexities is exponentially large does not imply that the smaller complexity is polynomial. Thus ``exponential savings'' must be specified as a reduction to polynomial size if it is used in an equivalence with membership in $\mathrm P/\mathrm{poly}$.

\subsection{Slice functions}
A $k$-slice function is zero at weights below $k$, one at weights above $k$, and arbitrary at weight $k$. It is monotone. The polynomial-overhead simulation is classical, attributed to \citet{Berkowitz} by \citet[Theorem~2.6 and Remark~2.7, pp.~57--58]{WegenerSlice}. Wegener proves the stronger overhead $O(n\min\{k,n+1-k,\log^2n\})$. We give a simpler cubic-overhead construction because it specifies exactly the output-zero condition used here.

\begin{proposition}\label{prop:slice}
If an $n$-variable $k$-slice function has a Boolean circuit of size $s$, it has a monotone circuit of size $O(s+n^3)$. Its final supported zeros in the preceding circuit model can therefore be eliminated with polynomial overhead.
\end{proposition}
\begin{proof}
First compute each wire and its complement together. By De Morgan's laws, an AND gate produces the positive output by AND and its complement by OR of the complementary operands; OR is dual, and NOT swaps the pair. This yields a circuit with at most $2s$ binary gates whose only negations are input literals.

On the exact weight-$k$ slice, $\neg x_i$ equals the monotone threshold $T_i$ asserting that at least $k$ of the other $n-1$ variables are true. Replace every negative input literal by $T_i$. Let the resulting monotone circuit be $g$. It agrees with the original function at weight $k$.

A threshold on $m$ inputs can be computed by the recurrence
\[
 T_{j,r}=T_{j-1,r}\vee(x_j\wedge T_{j-1,r-1}),
\]
with $T_{j,0}=1$, $T_{0,r}=0$ for $r>0$, omitting impossible indices. This uses $O(m^2)$ binary gates. Compute all $n$ functions $T_i$ and the thresholds $T_{n,k}$ and $T_{n,k+1}$ in $O(n^3)$ gates. Then
\[
 T_{n,k+1}\vee\bigl(T_{n,k}\wedge g\bigr)
\]
is zero below the slice, agrees with $g$ on it and is one above it. Boundary values $k=0,n$ use the same constant conventions. Proposition~\ref{prop:complexity} converts this monotone circuit into a final-output-zero-free one with constant-factor overhead.
\end{proof}

\subsection{Parity tests}
\begin{proposition}\label{prop:parity}
For bits $b_1,\ldots,b_m$, lift $0$ to $\tau$ and $1$ to $1\in(\F_2)_0$. Their sum is $\tau$ if all bits vanish, $e$ if a positive even number are one, and $1$ if an odd number are one. Thus an even-parity test passes exactly when the sum belongs to $\{\tau,e\}$, a Boolean semiring under the inherited operations, with its own multiplicative identity $e$. Its inclusion is not unital.
\end{proposition}
\begin{proof}
The ring coordinate of the sum is parity, while the support coordinate is the OR of all bits. These two values distinguish the three cases. The pair $\{\tau,e\}$ has $e+e=e$, $e^2=e$, additive identity $\tau$ and multiplicative identity $e$, so it is Boolean.
\end{proof}

\section{Two lifting problems and two independent losses of information}
This section answers which comparison follows solely from the coefficient and support homomorphisms.
There are homomorphisms
\[
 \N[M]\longrightarrow\Z_0[M]\longrightarrow\B[M],
\]
where a zero coefficient maps to $\tau$ and a positive coefficient maps to that positive old integer. The first map is injective and its image has no supported coefficient zeros. For $M=\N^n$, append the homomorphism $\B[M]\to Q_n$ that remembers only the monotone predicate generated by the monomials. Theorem~\ref{thm:main} is exactly the minimization of gate count over lifts of $f_G$ along this composite.

For comparison, effective sign coefficients form $\N[C_2\times M]$, with $C_2=\{1,\varepsilon\}$ and $\varepsilon^2=1$. A coefficient $(a^+,a^-)$ maps to the old integer $a^+-a^-$ when $a^++a^->0$ and to $\tau$ otherwise. This gives a surjective homomorphism to $\Z_0[M]$: any integer is a difference of two nonnegative integers, and an old zero is represented by $(1,1)$. Addition and multiplication preserve presence because dimensions add and multiply without cancellation; traces are additive and multiplicative. An associative effective lift of structure constants must in addition satisfy all associativity equations in this larger semiring.

\begin{proposition}\label{prop:independentlosses}
Supported coefficient zeros and repeated-variable identification in $\B[\N^n]\to Q_n$ are independent phenomena. No nontrivial supported coefficient zero occurs in a lift from $\N[\N^n]$, even when an edge is represented only by higher powers. Conversely a supported coefficient zero can occur without repeated variables.
\end{proposition}
\begin{proof}
Every occurring coefficient in the image of $\N[\N^n]$ is a positive integer. The monomial $x^2y$ therefore has no supported coefficient zero, although it and $xy$ have the same Boolean minimal support and different total degrees. Conversely $x-x$ in $\Z_0[x]$ has the supported coefficient zero at $x$, whose exponent is one. For effective coefficients, $1+\varepsilon$ gives a supported zero already at monoid degree zero. These examples establish both separations.
\end{proof}

What transfers is evaluation through homomorphisms, including the same gate-list image and the preservation of equations already satisfied by a lift. A solution to associativity equations need not have a small circuit, and the existence of a small circuit does not impose associativity on an unrelated multiplication table. The homomorphisms alone therefore give no equality between a dimension-lift defect and the cost of edges represented through Boolean idempotence.

\subsection{The canonical formal determinant}
The formal target also matters for the determinant claim in the unpublished translation notes. The following is a consequence of classical monotone permanent support bounds, not a new separation. \citet[Section~4.3, pp.~887--889]{JS} prove a stronger multiplication bound by their content method. We give an elementary exponential bound with a complete proof to specify exactly what is computed.

For a nonzero commutative unital ring $R$, let $\Delta_n^{\mathrm{can}}\in R_0[x_{ij}:1\le i,j\le n]$ have precisely the permutation monomials as its occupied exponents, with coefficient $\operatorname{sgn}(\sigma)$ on $\prod_i x_{i,\sigma(i)}$. The coefficient is interpreted in $R$.

\begin{proposition}\label{prop:canonicaldet}
For $n\ge3$, a circuit computing $\Delta_n^{\mathrm{can}}$ from variables and free scalar constants in $R_0$, using binary addition and multiplication and unary additive inverse, has at least
\[
 \binom{n}{\lfloor n/3\rfloor}
\]
binary gates. The bound holds over every such $R$, including characteristic two.
\end{proposition}
\begin{proof}
Apply the occupancy map $R_0[x]\to\B[x]$. An old scalar, even the old zero, maps to $1$, while $\tau$ maps to $0$; inverse gates become identities. The target is the formal Boolean permanent polynomial, whose exponents have not been identified by $x^2=x$. Delete inverse gates and zero branches and read the remaining circuit over $\N$, with each nonzero scalar replaced by $1$. Its output $Q$ has positive coefficients and exactly the permanent's monomial support. Its binary size has not increased.

Every useful gate polynomial $Q_g$ is homogeneous. Indeed, fix one parse occurrence of $g$ and all monomial choices outside it, giving a monomial context $m$ such that every monomial of $Q_gm$ occurs in the output. Positivity and the output degree $n$ force all monomials of $Q_g$ to have the same degree. They also force each such monomial to be a partial permutation. Degree-zero gates can be eliminated supportwise, since positive scalar multiplication changes no monomial support.

Take a parse tree for an output permutation monomial. Bypass addition steps and descend a product child whenever it has degree greater than $2n/3$. At the last product with degree greater than $2n/3$, both children have degree at most $2n/3$, and one has degree $d>n/3$. Choose this child occurrence. For $n\ge3$ it is a computation gate, and
\[
 n/3<d\le2n/3.
\]

For a fixed chosen gate $g$, one monomial context $m_g$ forces every monomial of $Q_g$ to use exactly the rows $I_g$ and columns $J_g$ absent from $m_g$. Both sets have size $d$. Every other possible context of $g$ must use their complements, since multiplying that context by any fixed monomial of $Q_g$ gives an output permutation. Hence the total number of permutations whose parses can be assigned to $g$ is at most $d!(n-d)!$. This bound counts all permutations respecting these two sets, and so remains valid for shared gates and repeated descriptions of a context.

Put $k=\lfloor n/3\rfloor$. Since $k\le d\le n-k$, we have $d!(n-d)!\le k!(n-k)!$. Assigning every output permutation one of its chosen gates, and writing $s$ for the number of binary gates, gives
\[
 n!\le s\,k!(n-k)!,
\]
which proves the claim.
\end{proof}

Projection to $R[x]$ can erase additional occupied monomials whose coefficients are old zeros. A small ring circuit for the determinant lifts gate by gate to $R_0$, but its lift may retain exactly such additional exponents and therefore fail to compute $\Delta_n^{\mathrm{can}}$. Proposition~\ref{prop:canonicaldet} concerns this specified formal element. It gives no lower bound for unrestricted Boolean circuits or for an arbitrary formal lift whose ring projection is the determinant.

\section{The remaining exact question}
The unrestricted question is the following: if $G$ is bipartite and has no four-cycle, must every $P\in\Sup(G)$ require at least $|E|+\vc(G)-1$ variable-input nonnegative arithmetic gates? The analogous question for every graph of girth at least five is the formulation attributed to L-meta-608; removing the triangle restriction gives the formulation attributed to L-meta-612. These questions remain undecided in this paper.

The obstruction can be stated without metaphor. A monomial whose support is an edge may have degree greater than two. Taking the degree-two component can delete the only witness for that edge. The polynomial $x_u^2x_v$ for a one-edge graph is the smallest example. Proposition~\ref{prop:missing} gives a lower bound in terms of the number of these missing quadratic witnesses, and Theorem~\ref{thm:twomissing} charges their full cost when at most two are missing. Theorem~\ref{thm:threemissing} extends the full count to at most three missing witnesses, including empty quadratic support. Proposition~\ref{prop:fouredgeboundary} settles the empty-quadratic case with four missing witnesses. For nonempty quadratic support and four missing witnesses, Theorem~\ref{thm:maxfour} settles maximal cover gain and Corollary~\ref{cor:fourresidualfinal} records the remaining lower-gain cases; the general four-or-more-missing-witness bound remains open. For marked degree at least two, the path expansion also has products of two dependent factors; replacing repeated marked variables by one can create new minimal edge witnesses. The transposed neighborhood calculation no longer has the linear form required by Lemma~\ref{lem:transpose}. Corollary~\ref{cor:degreewitness} proves that any counterexample in the bipartite case must encounter this obstruction on both sides.

\begin{example}[A discarded-gate charge that is false]\label{ex:chargefailure}
It would suffice to charge at least $r+\vc(G)-\vc(G_2)$ discarded gates, but that proposed inequality is false. Take four vertices $u,v,w,t$, four centers $a,b,c,d$, and four private leaves $\ell_a,\ell_b,\ell_c,\ell_d$. In $G_2$, join the centers respectively to $\{u,w\},\{u,t\},\{v,w\},\{v,t\}$, and join each center to its private leaf. Add $uv$ and $wt$ to obtain $G$. The old core is an eight-cycle. A cycle using one new edge has length at least five; one using both has length at least six. Thus $G$ has girth five. The four pendant edges certify $\vc(G_2)=4$, and together with the two new edges form a six-edge matching. The four centers together with $u,w$ cover $G$, so $\vc(G)=6$. There are twelve old edges and fourteen edges in all.

Put $A=u+w$, $B=u+t$, $C=v+w$, $D=v+t$. Compute
\[
 a(A+\ell_a)+b(B+\ell_b)+c(C+\ell_c)
 +(d+ABC)D+d\ell_d.
\]
The four clause sums cost four gates, the three pendant additions cost three, forming $ABC$ and adding $d$ costs three, the five remaining products cost five, and joining their outputs costs four. The total is nineteen, the star bound. Boolean distributivity gives $ABCD=uv\vee wt$, so the circuit computes $f_G$. Its formal quadratic part is exactly $H_{G_2}$. Extraction deletes only the two products forming $ABC$ and the addition $d+ABC$, leaving sixteen gates, one more than the optimal fifteen. Thus $r=2$ and the cover number drops by two, but only three gates disappear.

Folding the last pendant into $(d+ABC)(D+\ell_d)$ would save a gate, but accepts the independent input $\{u,w,\ell_d\}$. The missing charge in this example is therefore supplied by inefficiency in the retained quadratic circuit. Counting discarded gates alone cannot establish the desired universal bound.
\end{example}

The public project documented L-meta-608, L-meta-612 and L-meta-647 as conjectural and L-meta-649 as checked on 23 November 2026. These are historical statuses. Repeated public retrieval attempts on 24--25 November did not recover the live node records. The latest direct readings, observed on 25 November 2026 at 12:06:03 UTC, returned navigation-only shells at nodes 606, 608, 612, 647 and 649, without individual node content or status. Earlier requests to the data endpoint were unsuccessful, and exact recency-filtered searches for node 608 returned general indexed homepage activity rather than a current node-specific status. The detailed dated retrieval chronology accompanies the workbench records. Neither that indexed activity nor the unsuccessful direct readings establish the nodes' current statuses. The mathematical statements and proofs here stand independently of those statuses.

\section*{Method and provenance}
The construction was explored in unpublished project notes beginning in November 2025, independently of \citet{FGJM}, whose preprint appeared on 19 November 2024. The first texts were produced by Gemini 2.5 Deep Think in the Gemini app's Deep Think mode from a conversation with Gemini 2.5 Pro. Versions were deposited on Zenodo under HI-AI from 6 November 2025 and later withdrawn. A GPT-5-series Pro model in ChatGPT Pro mode at reasoning effort ultra extracted the algebra in March 2026. GPT-5-series Pro models in the same mode and effort produced the April--July 2026 corpus, and Codex runs at reasoning effort ultra produced the August--September workbench. The circuit translation notes of 21--23 November 2026 were produced by Claude Opus 5.5 at effort max, with consultations by Claude Fable 5.2 at effort max and Claude Opus 5.5 at effort high. The circuit map was produced by Claude Opus 5.5 at effort high; its degree-one argument by Claude Opus 5.5 at effort max. These are cited only as unpublished project work.

The present mathematics and text were produced by GPT-6 Astra Pro in ChatGPT Pro mode. The checks comprise independent derivations from statements, targeted proof review and finite exact computations, with their separate scopes recorded through the workbench; the reasoning-effort setting was not exposed and is not inferred. No human has checked this text. The text and original verification programs are dedicated to the public domain under CC0~1.0. The programs \emph{Circuit support checks} and \emph{Exact edge-circuit synthesis}, and the detailed check records, accompany this paper. The workbench record address is \url{\workbenchurl}; the permanent record address is \url{\recordurl}.

\begin{thebibliography}{99}
\bibitem[Alon and Boppana(1987)]{AB}
Alon, N., \& Boppana, R.~B. (1987). The monotone circuit complexity of Boolean functions. \emph{Combinatorica, 7}, 1--22. \url{https://doi.org/10.1007/BF02579196}.
\bibitem[Amano and Maruoka(2004)]{AM}
Amano, K., \& Maruoka, A. (2004). On the monotone circuit complexity of quadratic Boolean functions. In \emph{Algorithms and computation: ISAAC 2004}, Lecture Notes in Computer Science, 3341, 28--40. Springer. \url{https://doi.org/10.1007/978-3-540-30551-4_5}. The $K_{2,2}$-free single-level bound is cited through Jukna's account; the original full text was not retrieved.
\bibitem[Anthropic(2026)]{Anthropic}
Anthropic. (2026). \emph{VP $\ne$ VNP over every field}. Hexagon 2611.00954v1, Theorem~1.1 and Section~1.3. \url{https://hexagonmath.org/pdf/2611.00954v1}.
\bibitem[Berkowitz(1982)]{Berkowitz}
Berkowitz, S.~J. (1982). \emph{On some relationships between monotone and non-monotone circuit complexity}. Technical report, University of Toronto. Attribution verified through the contemporary slice-function literature; the report itself was not retrieved.
\bibitem[Caviglia et~al.(2026)]{CFMM}
Caviglia, E., Faul, P.~F., Manuell, G., \& Mesiti, L. (2026). \emph{Clifford semigroups and the monoidal Grothendieck construction} (Version~2). arXiv:2607.12944. \url{https://arxiv.org/abs/2607.12944}.
\bibitem[Faul et~al.(2024)]{FGJM}
Faul, P.~F., Goswami, A., Joubert, G., \& Manuell, G. (2024). \emph{The case for inverse semirings}. arXiv:2411.12318. \url{https://arxiv.org/abs/2411.12318}.
\bibitem[Jerrum and Snir(1982)]{JS}
Jerrum, M., \& Snir, M. (1982). Some exact complexity results for straight-line computations over semirings. \emph{Journal of the ACM, 29}(3), 874--897. \url{https://snir.cs.illinois.edu/listed/J7.pdf}.
\bibitem[Jukna(2004)]{JuknaGraphs}
Jukna, S. (2004). \emph{On graph complexity}. ECCC TR04-005, revision~1, Theorem~7.3, p.~18. \url{https://eccc.weizmann.ac.il/report/2004/005/revision/1/download/}.
\bibitem[Jukna(2006)]{JuknaSingle}
Jukna, S. (2006). Disproving the single level conjecture. \emph{SIAM Journal on Computing, 36}(1), 83--98. \url{https://doi.org/10.1137/S0097539705447001}.
\bibitem[Jukna(2012)]{JuknaBook}
Jukna, S. (2012). \emph{Boolean function complexity: Advances and frontiers}. Algorithms and Combinatorics, 27. Springer. Transposition principle, p.~46. \url{https://doi.org/10.1007/978-3-642-24508-4}; author-hosted chapter: \url{https://web.vu.lt/mif/s.jukna/boolean/chapter1.pdf}.
\bibitem[Jukna(2015)]{JuknaTropical}
Jukna, S. (2015). Lower bounds for tropical circuits and dynamic programs. \emph{Theory of Computing Systems, 57}, 160--194. ECCC TR14-080, revision~1. \url{https://eccc.weizmann.ac.il/report/2014/080/revision/1/}.
\bibitem[Jukna(2022)]{JuknaRead}
Jukna, S. (2022). \emph{Notes on Boolean read-$k$ and multilinear circuits}. ECCC TR22-094; author manuscript, Lemmas~1--2 and Theorem~1, pp.~5--7. \url{https://web.vu.lt/mif/s.jukna/ftp/readk-submit.pdf}; \url{https://eccc.weizmann.ac.il/report/2022/094/}.
\bibitem[Krapivin et~al.(2026)]{KPS}
Krapivin, A., Przybocki, B., \& Subercaseaux, B. (2026). Near-optimal encodings of cardinality constraints. \emph{29th International Conference on Theory and Applications of Satisfiability Testing (SAT 2026)}, Article~23. \url{https://doi.org/10.4230/LIPIcs.SAT.2026.23}. Full version: arXiv:2603.28954v2, Section~3.4, Theorem~6, pp.~7--8, Appendix~A.2, p.~19. \url{https://arxiv.org/abs/2603.28954v2}.
\bibitem[OpenAI(2026)]{OpenAICubic}
OpenAI. (2026). \emph{A cubic lower bound for border determinantal complexity of the permanent}. September~24. Paper and formalization scope. \url{https://github.com/openai/math/blob/main/lean/docs/108.md}.
\bibitem[P = NP? project(2026)]{Fleet}
P = NP? project. (2026). \emph{Public proof graph and proof repository}. Nodes L-meta-608, L-meta-612, L-meta-647 and L-meta-649. Historical statements dated 23 November 2026; access attempted 24 November 2026. \url{https://pequalnp.com/}; \url{https://github.com/open-goat/pequalnp-proofs}.
\bibitem[Project consultation(2026)]{FableConsult}
Project consultation. (2026, November~21). \emph{The shape statement and two regimes} [Unpublished consultation by Claude Fable~5.2, effort max, for Claude Opus~5.5, effort max]. Q2, Theorem~6 and corollary.
\bibitem[Project notes(2026)]{ProjectNotes}
Project notes. (2026). \emph{Zero-adjoined rings and monotone circuits}. Unpublished project notes, November~21--23. Provenance is given in the method note; the workbench record address is \url{\workbenchurl}.
\bibitem[Schnorr(1976)]{Schnorr}
Schnorr, C.~P. (1976). A lower bound on the number of additions in monotone computations. \emph{Theoretical Computer Science, 2}(3), 305--315. \url{https://doi.org/10.1016/0304-3975(76)90083-9}.
\bibitem[Valiant(1980)]{Valiant}
Valiant, L.~G. (1980). Negation can be exponentially powerful. \emph{Theoretical Computer Science, 12}, 303--314. \url{https://doi.org/10.1016/0304-3975(80)90060-2}.
\bibitem[Wegener(1985)]{WegenerSlice}
Wegener, I. (1985). On the complexity of slice functions. \emph{Theoretical Computer Science, 38}, 55--68. \url{https://doi.org/10.1016/0304-3975(85)90209-9}.
\end{thebibliography}
\end{document}
